% !TeX program = xelatex
% Standalone English volume. Use XeLaTeX until the references stabilize.
\documentclass[a4paper,11pt]{article}
\usepackage[T1]{fontenc}
\usepackage{lmodern}
\usepackage[english]{babel}
\usepackage[left=26mm,right=26mm,top=25mm,bottom=25mm,headheight=15pt,headsep=7mm,footskip=11mm]{geometry}
\usepackage{amsmath,amssymb,mathtools,amsthm}
\newtheorem{proposition}{Proposition}[section]
\newtheorem{lemma}[proposition]{Lemma}
\newtheorem{theorem}[proposition]{Theorem}
\newtheorem{corollary}[proposition]{Corollary}
\newtheorem{maintheorem}[proposition]{Main Theorem}
\theoremstyle{definition}
\newtheorem{definition}[proposition]{Definition}
\theoremstyle{remark}
\newtheorem{remark}[proposition]{Remark}
\makeatletter
\renewcommand*\l@section{\@dottedtocline{1}{0em}{2.4em}}
\renewcommand*\l@subsection{\@dottedtocline{2}{1.45em}{3.45em}}
\renewcommand*\l@subsubsection{\@dottedtocline{3}{3.8em}{4.55em}}
\makeatother
\usepackage{xcolor,fancyhdr,titlesec,needspace,xurl}
\usepackage{longtable,multicol}
\usepackage{hyperref}
\definecolor{ink}{HTML}{1A3542}
\definecolor{muted}{HTML}{4F606A}
\hypersetup{unicode=true,colorlinks=true,linkcolor=ink,filecolor=ink,urlcolor=ink,citecolor=ink,bookmarksnumbered=true,
 pdfdisplaydoctitle=true,pdflang={en},
 pdftitle={The generalized Riemann hypothesis for Dirichlet L-functions. Edition 930, Volume VI},
 pdfsubject={A natural-language reconstruction of the proposed proof},
 pdfauthor={A mathematical reconstruction}}
\titleformat{\section}{\Large\bfseries\color{ink}}{\thesection}{.65em}{}
\titlespacing*{\section}{0pt}{.8em}{.65em}
\setlength{\parindent}{1.2em}
\linespread{1.045}
\setlength{\parskip}{.22em plus .05em minus .03em}
\setlength{\emergencystretch}{2em}
\setcounter{topnumber}{2}
\setcounter{tocdepth}{3}\setcounter{secnumdepth}{3}
\widowpenalty=4500
\clubpenalty=4500
\pagestyle{fancy}\fancyhf{}
\fancyhead[L]{\small\color{muted}Dirichlet GRH \textperiodcentered\ Volume I}
\fancyhead[R]{\small\color{muted}Dirichlet zeros and power intervals}
\fancyfoot[C]{\small\thepage}
\renewcommand{\headrulewidth}{.3pt}
\newcommand{\lead}[1]{\par\smallskip\noindent{\bfseries\color{ink}#1.}\enspace}
\newcommand{\sour}[2]{\unskip\penalty150\hspace{.35em}{\small\hyperlink{ref#1}{[#2]}}}
\titleformat{\subsection}{\large\bfseries}{\thesubsection}{.6em}{}
\titlespacing*{\subsection}{0pt}{1.15em plus .15em}{.4em}
\setlength{\abovedisplayskip}{9pt plus 2pt minus 3pt}
\setlength{\belowdisplayskip}{9pt plus 2pt minus 3pt}
\newcommand{\RR}{\mathbb R}\newcommand{\CC}{\mathbb C}\newcommand{\QQ}{\mathbb Q}
\newcommand{\ZZ}{\mathbb Z}\newcommand{\ord}{\operatorname{ord}}
\newcommand{\Frob}{\operatorname{Frob}}\newcommand{\Gal}{\operatorname{Gal}}
\newcommand{\len}{\operatorname{length}}\newcommand{\Res}{\operatorname*{Res}}
\newcommand{\wmean}{\operatorname{av}}\newcommand{\Ker}{\operatorname{ker}}
\newcommand{\Lh}{\mathcal L}\newcommand{\Gg}{\mathcal G}
\newcommand{\real}{\operatorname{Re}}\newcommand{\imag}{\operatorname{Im}}
\newcommand{\abs}[1]{\left|#1\right|}\newcommand{\eps}{\varepsilon}
\newcommand{\anchoritem}[2]{\par\Needspace{2.8\baselineskip}\hypertarget{ref#1}{}\noindent{\bfseries\color{ink}#2.}\enspace}
\newcommand{\srcloc}[2]{{\small\nolinkurl{#1.lean}}, line~#2}
\numberwithin{equation}{section}
\raggedbottom

\newcommand{\Sel}{\mathcal S}
\newcommand{\Cl}{\operatorname{Cl}}
\newcommand{\Aut}{\operatorname{Aut}}
\newcommand{\GL}{\operatorname{GL}}
\newcommand{\OO}{\mathcal O}
\hypersetup{bookmarksdepth=3}


\addto\captionsenglish{\renewcommand{\partname}{Volume}}


\makeatletter
\expandafter\def\csname r@cor:coc145\endcsname{{25.15}{36}{}{proposition.25.15}{Dirichlet-GRH-v930-Volume-V.pdf}}
\expandafter\def\csname r@cor:minor67\endcsname{{66.19}{60}{}{proposition.66.19}{Dirichlet-GRH-v930-Volume-XIII.pdf}}
\expandafter\def\csname r@cor:res28\endcsname{{24.10}{10}{}{proposition.24.10}{Dirichlet-GRH-v930-Volume-V.pdf}}
\expandafter\def\csname r@cor:tail114\endcsname{{19.20}{17}{}{proposition.19.20}{Dirichlet-GRH-v930-Volume-IV.pdf}}
\expandafter\def\csname r@eq:level146\endcsname{{25.13}{38}{}{equation.25.13}{Dirichlet-GRH-v930-Volume-V.pdf}}
\expandafter\def\csname r@eq:tail114\endcsname{{19.15}{17}{}{equation.19.15}{Dirichlet-GRH-v930-Volume-IV.pdf}}
\expandafter\def\csname r@eq:window\endcsname{{{W}}{2}{}{AMS.2}{Dirichlet-GRH-v930-Volume-I.pdf}}
\expandafter\def\csname r@lem:bern99\endcsname{{25.4}{26}{}{proposition.25.4}{Dirichlet-GRH-v930-Volume-V.pdf}}
\expandafter\def\csname r@lem:des28\endcsname{{24.8}{8}{}{proposition.24.8}{Dirichlet-GRH-v930-Volume-V.pdf}}
\expandafter\def\csname r@lem:fit13\endcsname{{11.2}{3}{}{proposition.11.2}{Dirichlet-GRH-v930-Volume-II.pdf}}
\expandafter\def\csname r@lem:int13\endcsname{{11.3}{4}{}{proposition.11.3}{Dirichlet-GRH-v930-Volume-II.pdf}}
\expandafter\def\csname r@lem:top\endcsname{{19.2}{3}{}{proposition.19.2}{Dirichlet-GRH-v930-Volume-IV.pdf}}
\expandafter\def\csname r@nav:I-28\endcsname{{6.6}{32}{}{subsection.6.6}{Dirichlet-GRH-v930-Volume-I.pdf}}
\expandafter\def\csname r@nav:I-37\endcsname{{70.1}{3}{}{subsection.70.1}{Dirichlet-GRH-v930-Volume-XIV.pdf}}
\expandafter\def\csname r@nav:volume-IV\endcsname{{19}{3}{}{section.19}{Dirichlet-GRH-v930-Volume-IV.pdf}}
\expandafter\def\csname r@prop:example30\endcsname{{24.15}{14}{}{proposition.24.15}{Dirichlet-GRH-v930-Volume-V.pdf}}
\expandafter\def\csname r@prop:finite15\endcsname{{11.4}{6}{}{proposition.11.4}{Dirichlet-GRH-v930-Volume-II.pdf}}
\expandafter\def\csname r@prop:minors33\endcsname{{6.14}{46}{}{proposition.6.14}{Dirichlet-GRH-v930-Volume-I.pdf}}
\expandafter\def\csname r@prop:short107\endcsname{{9.9}{114}{}{proposition.9.9}{Dirichlet-GRH-v930-Volume-I.pdf}}
\expandafter\def\csname r@sec:comparison\endcsname{{6}{24}{}{section.6}{Dirichlet-GRH-v930-Volume-I.pdf}}
\expandafter\def\csname r@sec:missing\endcsname{{6.6}{32}{}{subsection.6.6}{Dirichlet-GRH-v930-Volume-I.pdf}}
\expandafter\def\csname r@seriesguide\endcsname{{1.1}{4}{}{subsection.1.1}{Dirichlet-GRH-v930-Volume-I.pdf}}
\expandafter\def\csname r@thm:char\endcsname{{24.1}{3}{}{proposition.24.1}{Dirichlet-GRH-v930-Volume-V.pdf}}
\expandafter\def\csname r@thm:common16\endcsname{{11.9}{10}{}{proposition.11.9}{Dirichlet-GRH-v930-Volume-II.pdf}}
\expandafter\def\csname r@thm:conductor232\endcsname{{19.15}{12}{}{proposition.19.15}{Dirichlet-GRH-v930-Volume-IV.pdf}}
\expandafter\def\csname r@thm:critical13\endcsname{{11.1}{3}{}{proposition.11.1}{Dirichlet-GRH-v930-Volume-II.pdf}}
\expandafter\def\csname r@thm:effect146\endcsname{{25.17}{38}{}{proposition.25.17}{Dirichlet-GRH-v930-Volume-V.pdf}}
\expandafter\def\csname r@thm:eigen240\endcsname{{21.20}{57}{}{proposition.21.20}{Dirichlet-GRH-v930-Volume-IV.pdf}}
\expandafter\def\csname r@thm:fourier145\endcsname{{25.14}{35}{}{proposition.25.14}{Dirichlet-GRH-v930-Volume-V.pdf}}
\expandafter\def\csname r@thm:integerlayers333\endcsname{{22.9}{69}{}{proposition.22.9}{Dirichlet-GRH-v930-Volume-IV.pdf}}
\expandafter\def\csname r@thm:invar148\endcsname{{20.11}{31}{}{proposition.20.11}{Dirichlet-GRH-v930-Volume-IV.pdf}}
\expandafter\def\csname r@thm:positive306\endcsname{{28.1}{84}{}{proposition.28.1}{Dirichlet-GRH-v930-Volume-V.pdf}}
\expandafter\def\csname r@thm:primitive307\endcsname{{28.2}{85}{}{proposition.28.2}{Dirichlet-GRH-v930-Volume-V.pdf}}
\expandafter\def\csname r@thm:twists\endcsname{{5.1}{19}{}{proposition.5.1}{Dirichlet-GRH-v930-Volume-I.pdf}}
\expandafter\def\csname r@thm:wild98\endcsname{{25.1}{24}{}{proposition.25.1}{Dirichlet-GRH-v930-Volume-V.pdf}}
\makeatother
\begin{document}
\setcounter{part}{5}
\fancyhead[L]{\small\color{muted}Dirichlet GRH \textperiodcentered\ Volume VI}
\fancyhead[R]{\small\color{muted}Marked deformations and traces}
\thispagestyle{plain}\begin{center}
{\LARGE\bfseries The generalized Riemann hypothesis\\[2pt]for Dirichlet \(L\)-functions}\\[6pt]
{\small 27 November 2026 \quad\textperiodcentered\quad Edition 930}\end{center}
\part{Marked Galois deformations and trace congruences}\label{part:VI}
\begin{abstract}\noindent We relate strict extension classes to rank-three Galois representations. We calculate tangent spaces of ordinary deformations, distinguish orthogonal components, and compare their trace congruences on the common interval.\end{abstract}
\setcounter{section}{29}
\noindent The \hyperref[seriesguide]{guide in Volume~I} lists all twenty-seven volumes.
\begingroup\small\setlength{\parskip}{0pt}\tableofcontents\endgroup
\clearpage

\section{Extension classes and trace families}\label{nav:g-IV-0}\label{nav:volume-VI}

\subsection{Passing the strict condition to finite coefficients}\label{nav:IV-00}
\label{sec:strict31}

Fix an odd prime \(\ell\). Let \(E/\QQ\) be totally imaginary abelian, let \(P\) be the \(\ell\)-Sylow subgroup of \(\Gal(E/\QQ)\) and let \(H\) be its prime-to-\(\ell\) factor. Put \(K=E^H\). Write \(\OO\) for a finite unramified extension of \(\ZZ_\ell\) containing the values of all characters of \(H\), \(J=\operatorname{Frac}(\OO)\), \(v_\ell(\ell)=1\). Extend a character \(\alpha\) of \(H\) trivially across \(P\). For a positive twist \(r\) set
\[
 \vartheta=\kappa_\ell^{-r}\alpha^{-1}|_{G_K},\qquad
 D=(J/\OO)(\vartheta),\qquad D_m=D[\ell^m].
\]
When referring to the complementary sign take \(\alpha(c)=(-1)^r\), with \(c\) complex conjugation. The strict group \(S_\vartheta\subset H^1(K,D)\) consists of classes trivial on inertia at every finite place away from \(\ell\) and trivial on the full local group at every place above \(\ell\). It is the character component of the strict group over \(E\) by Lemma \ref{lem:des28}. That lemma also covers \(\ell\mid [E:\QQ]\). For the fields of the prime count, the components for which \(\bar\vartheta=1\) have already been bounded in Corollary \ref{cor:res28}. We consider every other component here.

Define \(S_{\vartheta,m}^{\rm f}\subset H^1(K,D_m)\) by the same local conditions with coefficient \(D_m\), including at \(w\mid\ell\). For such a place set \(s_w=0\) if \(\bar\vartheta|_{G_{K_w}}\ne1\). Otherwise put
\[
 s_w=\min_{g\in G_{K_w},\ \vartheta(g)\ne1}v_\ell(\vartheta(g)-1).
\]
The minimum exists since the cyclotomic character has infinite image there and \(r\ne0\).

\begin{theorem}[The cost of imposing finite local conditions]\label{thm:finite31}
Suppose \(\bar\vartheta\ne1\) on \(G_K\). For every \(m\ge1\), finite-to-divisible coefficients give an injection with
\begin{equation}
 0\le\len_{\OO}S_\vartheta[\ell^m]-\len_{\OO}S_{\vartheta,m}^{\rm f}
       \le\sum_{w\mid\ell}\min(m,s_w).                        \label{eq:loss31}
\end{equation}
If the residual character is locally trivial at \(w\), then
\[
  s_w= a_w+v_\ell(r),\qquad
  1\le a_w\le 1+v_\ell([K_w:\QQ_\ell]).
\]
In a representative of a strict class killed by \(\ell^m\), the local coboundary at \(w\) can be taken with values of its defining element in \(D_{m+s_w}\).
\end{theorem}
\begin{proof}
The exact sequence for multiplication by \(\ell^m\) on \(D\) gives an isomorphism \(H^1(K,D_m)\simeq H^1(K,D)[\ell^m]\), since some global element has \(\vartheta(g)-1\) a unit and thus \(D^{G_K}=0\). On inertia away from \(\ell\), the character has finite image of order prime to \(\ell\). If this image is trivial its invariant module is divisible, and if nontrivial its invariant module is zero, since some Teichm\"uller value differs from one by a unit. In either case the kernel of the local coefficient map, the quotient of the invariant module by \(\ell^m\), is zero. The finite lift already satisfies the required inertia conditions.

At \(w\mid\ell\) the possible local images of that lift lie in the connecting image of \(D^{G_{K_w}}/\ell^mD^{G_{K_w}}\). Its length is \(\min(m,s_w)\), since the invariant module is precisely \(D_{s_w}\) (with \(D_0=0\)). Simultaneous localization gives an injection of the quotient in \eqref{eq:loss31} into the product of these local kernels. All groups in the inequality are finite, for instance by the positive-twist finiteness in Proposition \ref{prop:finite15} and character descent. If a local finite cocycle represents a coboundary in \(D\), its values have the form \((\vartheta(g)-1)y\). An element attaining the minimum shows \(\ell^{m+s_w}y=0\); use a unit difference if \(s_w=0\). This proves the assertion about the local element. More precisely, put \(s_* = \max_{w\mid\ell}s_w\). Under the global coefficient isomorphism, \(S_\vartheta[\ell^m]=(S_{\vartheta,m+s_*}^{\rm f})[\ell^m]\). Indeed the chosen local elements place every lifted class inside this finite strict group after inclusion of coefficients. Conversely all its classes map to strict classes in divisible coefficients, and the global coefficient map is injective.

Let \(\langle u\rangle\) be the principal-unit part of \(u\in\ZZ_\ell^\times\). The local cyclotomic character on \(G_{\QQ_\ell}\) is onto \(\ZZ_\ell^\times\), as follows from the local cyclotomic extension of degree \((\ell-1)\ell^{n-1}\) at level \(n\). Hence the image of \(\langle\kappa_\ell\rangle\) on the subgroup \(G_{K_w}\) has index dividing \([K_w:\QQ_\ell]\) in \(1+\ell\ZZ_\ell\). By the \(\ell\)-adic logarithm it is \(1+\ell^{a_w}\ZZ_\ell\) with the asserted bound on \(a_w\). The finite character \(\alpha\) has only Teichm\"uller values. If \(\bar\vartheta\) is locally trivial, its entire Teichm\"uller part cancels, so that its image is the \((-r)\)-th power of this principal-unit group. The logarithm shows that its first nonzero valuation is \(a_w+v_\ell(r)\).
\end{proof}

\begin{corollary}[A bound at all odd primes on the power interval]\label{cor:sum31}
Use the fields \(E=E_x\subset\QQ(\zeta_N)\) and \(t=\lfloor x^{401}\rfloor\) of \hyperref[nav:volume-IV]{Volume~IV}. Let \(d_E=[E:\QQ]\). Independently choose a positive \(m_{\ell,r,\alpha}\) for each nontrivial residual component. The sum of the left-hand length differences in \eqref{eq:loss31} over all odd primes, all \(1\le r<t\) and all such complementary components, weighted by \(\log\ell\), is at most
\begin{equation}
 C d_E t\bigl(t+\log(2N)\bigr)\log(2td_E)
                  \le \exp(x^{\eps/3})                         \label{eq:family31}
\end{equation}
on every sufficiently distant common interval. In particular one can take \(m_{\ell,r,\alpha}\) large enough to contain each entire strict group.
\end{corollary}
\begin{proof}
If \(\ell\nmid N\), both \(K/\QQ\) and \(\alpha\) are unramified at \(\ell\). Local inertia over \(K_w\) then has full mod-\(\ell\) cyclotomic image, since the local cyclotomic extension is totally ramified. Local residual triviality forces \(\ell-1\mid r\), hence \(\ell\le t\). At any contributing prime and twist, Theorem~\ref{thm:finite31} bounds \(s_w\log\ell\) by \(\log\ell+\log r+\log d_E\). The product of the number of characters \(|H|\) and the number of places over \(\ell\) is at most \(|H|[K:\QQ]=d_E\). Sum first over \(\ell\le t\). For the remaining primes dividing \(N\) use \(\sum_{\ell\mid N}\log\ell\le\log N\) and \(\#\{\ell:\ell\mid N\}\le\log(2N)/\log2\). The displayed bound follows after summing over \(r<t\). The estimates \(d_E\ll_\psi x^{2+8\eps}\) and \(\log(2N)\ll_{\psi,\eps}x^{\eps/4}\) from Volume~IV give the second inequality. The last assertion follows from finiteness of each strict group.
\end{proof}

\subsection{An integral presentation of the global extension classes}\label{nav:IV-01}
\label{sec:present31}

Keep \(\bar\vartheta\ne1\). Take a finite set \(S\) of finite places of \(K\) containing those over \(\ell\) and the ramification of \(\alpha\), and write \(G=G_{K,S}\). Fix \(g_*\in G\) with \(\vartheta(g_*)-1\) a unit. At \(w\mid\ell\) let \(H_w\) be the kernel of \(\vartheta\) on the full local group. At \(w\in S\) away from \(\ell\) let it be the kernel on inertia. Local groups are mapped into \(G\) using fixed embeddings. Outside \(S\) the finite modules are unramified, and classes vanishing on inertia descend to \(G\).

\begin{lemma}[Normalized finite cocycles]\label{lem:norm31}
The group \(S_\vartheta[\ell^m]\) is isomorphic to the module of continuous functions \(z:G\to D_m\) satisfying
\[
 z(gh)=z(g)+\vartheta(g)z(h),\qquad z(g_*)=0,\qquad
 z(H_w)=0\quad(w\in S).
\]
\end{lemma}
\begin{proof}
The global lift is unique as a cohomology class by Theorem~\ref{thm:finite31}. Every finite cocycle has a unique cohomologous representative vanishing at \(g_*\), obtained by subtracting \((\vartheta(g)-1)z(g_*)/(\vartheta(g_*)-1)\). Its restriction to \(H_w\) is a homomorphism. The map from finite to divisible coefficients on its values is injective there since the action is trivial. A local class in divisible coefficients trivial on \(H_w\) comes by inflation from the local character image. Away from \(\ell\) the latter is finite of order prime to \(\ell\), so its degree-one cohomology vanishes. At \(w\mid\ell\) its principal-unit subgroup is open procyclic and its quotient has order prime to \(\ell\). A topological generator acts by a scalar different from one, whose difference from one is onto on \(J/\OO\). Matching a cocycle at this generator by a coboundary and using averaging on the quotient gives the same vanishing. Thus killing the specified kernels is equivalent to the divisible strict condition at all places.
\end{proof}

\begin{theorem}[The relation module and its determinants]\label{thm:present31}
There are finitely many topological generators in a common quotient of \(G\) sufficient for all the cocycles in Lemma~\ref{lem:norm31}. If their number is \(n\), there is a full-rank relation submodule \(B\subset\OO^n\) such that, for all \(m\),
\[
    S_\vartheta[\ell^m]\simeq\operatorname{Hom}_{\OO}
                (\OO^n/B,\ell^{-m}\OO/\OO),\qquad
    \len_{\OO}S_\vartheta=\min v_\ell\det(b_1,\ldots,b_n).
\]
The minimum runs over \(n\) relation vectors described in the proof, with valuation \(+\infty\) for zero determinants.
\end{theorem}
\begin{proof}
Let \(K'/K\) be the field cut out by \(\bar\vartheta\). In \(G\), factor its kernel by the kernel of its maximal pro-\(\ell\) quotient. This latter kernel is characteristic in the open normal subgroup, so we obtain a quotient \(\Gamma\) of \(G\). The character \(\vartheta\) factors through it: on \(G_{K'}\) its values are principal units. For a finite cocycle the homomorphism \(g\mapsto(z(g),\vartheta(g)\bmod\ell^m)\) takes this subgroup into a finite \(\ell\)-group, a semidirect product of the additive finite module by a principal-unit image. Hence every such cocycle also factors through \(\Gamma\).

The normal pro-\(\ell\) subgroup has finite-dimensional continuous homomorphisms to \(\mathbf F_\ell\). Indeed these homomorphisms are unramified away from the primes over \(S\). Adjoin \(\mu_\ell\) to \(K'\), an extension of degree dividing \(\ell-1\), and use injectivity of restriction by corestriction. After identifying \(\mathbf F_\ell\) with \(\mu_\ell\), Kummer theory bounds the resulting space by the dimension of the units of the ring of \(S\)-integers modulo \(\ell\)-th powers and the dimension of its \(\ell\)-torsion class group. Both are finite. Thus the maximal elementary abelian quotient of the pro-\(\ell\) subgroup is finite. Lifts of a basis topologically generate: if their closed subgroup were proper, its image would be proper in a finite \(\ell\)-group quotient, hence contained in a subgroup of index \(\ell\). That would give a nonzero character killing the basis. Adjoin lifts of generators of the finite residual quotient to obtain generators of \(\Gamma\).~\cite[Ch.~4]{VI-IV-VI-wilkes}

Choose these \(n\) generators, let \(F_n\) be the free profinite group on their lifts \(u_i\), and pull \(\vartheta\) back to \(F_n\). The free profinite group is the profinite completion of the ordinary free group on this finite basis. There is a unique continuous crossed homomorphism \(v:F_n\to\OO^n\) with \(v(u_i)\) the \(i\)-th basis vector: map \(u_i\) to \((v(u_i),\vartheta(u_i))\) in the compact profinite group \(\OO^n\rtimes\vartheta(G)\) and use the defining property of the completion on each finite quotient. Give the vectors \(v(h)\) as relations for all \(h\) in the kernel of \(F_n\to\Gamma\), for a lift of \(g_*\), and for the preimages of the images of all \(H_w\). Their \(\OO\)-span is \(B\). This submodule is finitely generated because \(\OO\) is Noetherian, and is closed since a finite sum of compact cyclic submodules is compact. A linear map from \(\OO^n\) to a finite coefficient group sends \(v\) to a cocycle. Conversely every finite cocycle on the free group is given by this recipe by its generator values and continuity. It descends and satisfies the conditions of Lemma~\ref{lem:norm31} exactly when the linear map kills \(B\). This proves the first formula.

Passing to the union in \(m\) identifies \(S_\vartheta\) with \(\operatorname{Hom}_{\OO}(\OO^n/B,J/\OO)\), since the source has finitely many generators. Finiteness of \(S_\vartheta\) excludes a free summand in \(\OO^n/B\). The elementary-divisor theorem over the discrete valuation ring now gives equality of its length with the length of its dual and with the valuation of the ideal generated by the maximal minors of a relation matrix. The originally specified vectors contain a finite generating set for their span, by finite generation. The determinant ideal is unchanged when all of them are allowed, by multilinearity. Its valuation is the claimed minimum.
\end{proof}

The theorem determines precisely which global determinants would suffice. It supplies finite generation but no uniform upper bound for their valuations. The following calculation connects these linear equations to the deformation argument of the supplied source.

\subsection{A rank-three deformation and the trace comparison}\label{nav:IV-02}
\label{sec:deform31}

Put \(\Psi=\vartheta^{-1}\), and let \(\xi:G\to\{1,-1\}\) be a continuous quadratic character, allowing the trivial character. Write \(J_3\) for the matrix with antidiagonal entries one and all other entries zero. The nontrivial element in \(\{1,-1\}\) acts on \(\GL_3\) by \(h\mapsto J_3(h^{\mathsf T})^{-1}J_3\). The diagonal matrix \(d(g)=\operatorname{diag}(\Psi(g),1,\Psi(g)^{-1})\) is fixed by this involution. Thus \(\rho_0(g)=(d(g),\xi(g))\) is a homomorphism to the semidirect product. For a finite module \(M=\OO/\ell^m\OO\), denote by \(\OO\oplus\epsilon M\) the ring in which \(\epsilon M\) is a square-zero ideal.

\begin{proposition}[An integral root direction for the marked classes]\label{prop:root31}
Let \(N_0=E_{21}-E_{32}\), with \(E_{ij}\) the standard matrix units. Identify \(D_m\) with \(M\) in the chosen coefficient coordinate. For a normalized cocycle of Lemma~\ref{lem:norm31} put
\[
       \rho_z(g)=(1+\epsilon z(g)N_0)\rho_0(g).
\]
These are distinct homomorphisms lifting \(\rho_0\), with \(\rho_z(g_*)=\rho_0(g_*)\). For every \(g\in\ker\xi\) their characteristic polynomials agree with that of \(d(g)\). On local groups subject to strictness they are conjugate to the base representation by elements of the same root direction after the coefficient enlargement of Theorem~\ref{thm:finite31}.
\end{proposition}
\begin{proof}
On the square-zero ideal the involution acts by \(X\mapsto-J_3X^{\mathsf T}J_3\); it fixes \(N_0\). Conjugation by \(d(g)\) scales each of \(E_{21},E_{32}\) by \(\Psi(g)^{-1}=\vartheta(g)\). Multiplying two displayed matrices therefore gives the cocycle law, since products of two elements of \(\epsilon M\) vanish. The lower triangular matrix part recovers \(z(g)\) from its \((2,1)\)-entry since \(\Psi(g)\) is a unit. Its diagonal equals that of \(d(g)\), giving the characteristic polynomial. If locally \(z(g)=(\vartheta(g)-1)y\), conjugate \(\rho_0(g)\) by \(1-\epsilon yN_0\) to obtain \(\rho_z(g)\). At primes above \(\ell\) the element \(y\) has the bound just proved. Away from \(\ell\) any inertia coboundary can be taken in \(D_m\), by the inertia calculation in Theorem~\ref{thm:finite31}.
\end{proof}

For the quadratic character of discriminant \(-1208\), \(\ell=3\) and odd \(r\), take \(E=\QQ(\sqrt{-302})\). The complementary character is its nontrivial quadratic character, with nontrivial residual \(\vartheta\) on inertia at 151. Proposition \ref{prop:example30}, with its proof over this quadratic field, gives a nonzero class already at \(m=1\). Thus the equality of all traces on \(\ker\xi\) in Proposition~\ref{prop:root31} actually occurs for different marked deformations arising from the arithmetic problem.

\begin{proposition}[What an integral cotangent bound detects]\label{prop:sat31}
Let \(A\) be a graded \(\OO\)-algebra with an augmentation \(x:A\to\OO\) killing all nonzero degrees. Write \(I=\ker x\), and
\[
 I^{(2)}=\{z\in I:\text{some }0\ne u\in\OO\text{ satisfies }uz\in I^2\}.
\]
Fix a nonzero degree \(j\) and homogeneous equations \(F_i\in A_j\). Suppose a finite \(\OO\)-module \(V\) injects additively into the augmentation derivations \(A\to M\) that vanish outside degree \(j\) and kill each \(F_i\). Suppose every derivation in the image lifts to an \(\OO\)-valued augmentation derivation of degree support \(j\). If the module
\[
 C_j=A_j/\bigl((A_j\cap I^{(2)})+\textstyle\sum_i\OO F_i\bigr)
\]
has finite length, then \(\len_{\OO}V\le\len_{\OO}C_j\).
\end{proposition}
\begin{proof}
An augmentation derivation obeys \(D(ab)=x(a)D(b)+x(b)D(a)\), so kills \(I^2\). An integral lift takes values in the torsion-free module \(\OO\); it therefore kills \(I^{(2)}\). Reducing the lift shows that the original derivation kills this saturation. Its values are determined by its restriction to \(A_j\), which factors through \(C_j\). The resulting additive injection into \(\operatorname{Hom}_{\OO}(C_j,M)\) bounds cardinalities. A cyclic summand \(\OO/\ell^b\) contributes length \(\min(m,b)\) to this Hom module. Taking logarithms to base \(|\OO/\ell|\) gives the inequality even for an injection only known to be additive. The linear presentation in Theorem~\ref{thm:present31} realizes this calculation in a polynomial algebra on \(n\) variables of degree \(-1\): the equations are linear forms with coefficient vectors in \(B\). Lifts of derivations on the ambient polynomial ring are obtained by lifting the values on its variables.
\end{proof}

\begin{theorem}[A trace estimate for a finite order]\label{thm:order31}
Let \(R\) be a commutative \(\OO\)-algebra finite as a module, \(x:R\to\OO\) an augmentation, \(I=\ker x\), and \(T\) its \(\OO\)-torsion ideal. Suppose the generic fiber \(R\otimes J\) is reduced. Write \(\Sigma\) for its geometric points other than \(x\). Let \(a\in I\) have \(a(P)\ne0\) at every \(P\in\Sigma\). Then
\begin{equation}
 \len_{\OO}(I/I^2)\le
 \len_{\OO}\bigl(T/(T\cap I^2)\bigr)
       +\sum_{P\in\Sigma}v_\ell(a(P)).                         \label{eq:order31}
\end{equation}
In particular if \(T=0\) and for every \(m\) the classes of \(S_\vartheta[\ell^m]\) inject into augmentation derivations \(R\to M\), their full length is bounded by this sum.
\end{theorem}
\begin{proof}
Torsion maps to zero under \(x\). Set \(R'=R/T\), \(I'=I/T\). The kernel of \(I/I^2\to I'/I'^2\) is \(T/(T\cap I^2)\). The ideal \(I'\) is finite free over the discrete valuation ring. Over an algebraic closure of \(J\) it has one coordinate line for each point of \(\Sigma\): a finite reduced algebra in characteristic zero is a product of separable fields, and augmentation removes precisely its own coordinate. Multiplication by \(a\) on \(I'\) is injective and has determinant \(\prod_{P\in\Sigma} a(P)\), with the empty product one. The elementary-divisor theorem gives \(\len(I'/aI')=\sum_P v_\ell(a(P))\). Since \(aI'\subset I'^2\), additivity of length gives \eqref{eq:order31}. Augmentation derivations to \(M\) are precisely \(\operatorname{Hom}_{\OO}(I/I^2,M)\), by the decomposition \(R=\OO\oplus I\) and the product rule. Its length is at most \(\len(I/I^2)\). Taking \(m\) above the exponent of the finite strict group proves the consequence (also for additive injections).
\end{proof}

The hypotheses on an integral model matter. In a finite torsion-free \(\ZZ\)-graded \(\OO\)-algebra with reduced generic fiber every nonzero-degree piece vanishes. Indeed, after extending scalars to an algebraic closure of \(J\), scaling homogeneous elements of degree \(j\) by \(t^j\) permutes the finite set of geometric points. A nonzero value in nonzero degree would consequently take the infinitely many values \(t^j a(P)\) in a finite set. Reducedness then forces each such element to vanish, and torsion-freeness gives vanishing integrally. A comparison preserving the root direction must therefore account for the scaling orbits or supply a different integral model on which the required derivations survive. Traces alone take identical values on the deformations of Proposition~\ref{prop:root31}.

\lead{Connection with the source} In the supplied archive the normalization and the reflected root occur in
\srcloc{Development/AutomorphicArithmetic/AutomorphicVarieties/UnitaryThreeGroups/ShaInjGen}{30} and line 101 of that file. The integral lifting condition of Proposition~\ref{prop:sat31} occurs in \srcloc{Development/AutomorphicArithmetic/AutomorphicVarieties/CommutativeAlgebra/OrbitVanishingLocal}{120}. Lines 103--160 of \nolinkurl{Development/AutomorphicArithmetic/AutomorphicVarieties/UnitaryThreeGroups/ModelUpper.lean} feature a graded model with a finite weight-zero algebra and a sum over scaling orbits. That setting differs from the finite-order hypothesis in Theorem~\ref{thm:order31}. To complete the length bound by this route one must construct the relevant integral model for the global local conditions, control its orbits and tangents, and show that its points yield the bounded automorphic systems of Section~\ref{nav:III-01} with the required multiplicities. The local bound \eqref{eq:family31} applies in advance of that comparison. The integral theta-period estimate of Volume~I remains a separate requirement.

\subsection{The two families of traces}\label{nav:III-00}
Fix an imaginary quadratic field \(K\) with discriminant of absolute value \(d_K\), a positive integer \(N\), and \(r\ge1\). Write \(c\) for complex conjugation. Choose a rational prime \(p\ge7\) split in \(K\) and prime to \(3Nd_K\), and fix \(w\mid p\). We use the split condition so that the residue cardinality is \(p\). At the end two such primes will be fixed for the entire interval.

Here are the automorphic representations we count. The congruence level in this definition is on the automorphic side. Write \(U_K(N)=\ker(\GL_3(\widehat{\OO}_K)\to\GL_3(\OO_K/N\OO_K))\).
\begin{enumerate}
\item Take the cuspidal \(\Pi\) of \(\GL_3(\mathbb A_K)\) with \(\Pi^\vee\simeq\Pi^c\), with \(U_K(N)\)-fixed vectors, and cohomological for a coefficient system of highest weights \(\lambda_{\sigma,1}\ge\lambda_{\sigma,2}\ge\lambda_{\sigma,3}\) at the two embeddings of \(K\). Take all such systems with \(|\lambda_{\sigma,i}|\le h=4(r+2)\) and trivial action of the positive real split center. Cohomological means that relative Lie algebra cohomology of \(\Pi_\infty\) tensored with the coefficient system is nonzero, modulo that split center. At \(w\), let \(\alpha_1,\alpha_2,\alpha_3\) be the unitarily normalized Satake parameters of \(\Pi_w\). Record two labels, with values \(b_{\Pi,\pm}=\sum_{j=1}^3\alpha_j^{\pm1}\).
\item Take all normalized primitive holomorphic cusp eigenforms \(f\) of weight \(2r+1\) and level dividing \(M=Nd_K\), with arbitrary nebentypus \(\omega_f\). At \(p\) write \(a_p(f)\) for the Fourier coefficient. Record
\[
        b_{f,+}=1+p^{-r}a_p(f),\qquad
        b_{f,-}=1+p^{-r}\omega_f(p)^{-1}a_p(f).
\]
\end{enumerate}
A label records its representation and sign. Repeated values on different labels are retained. In the second case the characteristic polynomial at \(p\) has roots \(\beta_1,\beta_2\) with sum \(a_p(f)\) and product \(\omega_f(p)p^{2r}\). The two numbers are the traces on \(1\oplus(\beta_1/p^r,\beta_2/p^r)\) and its inverse. Since \(p\) splits in \(K\), restriction from \(\QQ\) to \(K_w\) leaves the residue degree one. Thus these are also the traces for the corresponding restricted, normalized Deligne representation with a trivial line adjoined. In rank three the two signs likewise accommodate either choice of Frobenius convention.

\begin{theorem}[Size, purity and integrality]\label{thm:traces}
Let \(\mathcal T=\mathcal T(N,r,p)\) be the labelled multiset above. There are constants \(C_K,A\) such that
\begin{equation}
 |\mathcal T|\le C_K[N(r+2)]^A,\qquad
 [\QQ(b):\QQ]\le C_K[N(r+2)]^A\quad(b\in\mathcal T).
 \label{eq:size}
\end{equation}
Under every complex embedding \(|b|\le3\). With \(u=18h+1\), every \(p^u b\) is an algebraic integer. The constants and \(u\) are independent of the split auxiliary prime.
\end{theorem}

\begin{lemma}\label{lem:coh}
For one coefficient system of the specified weights, the total dimension of cohomology of the \(\GL_3/K\) locally symmetric space of level \(U_K(3N)\) with these coefficients is \(O_K(N^{36}(h+1)^6)\). The number of distinct cuspidal representations from this coefficient system at level \(U_K(N)\) has the same upper bound.
\end{lemma}
\begin{proof}
The locally symmetric space is the adelic double quotient by the maximal compact at infinity, by the positive real split center and by the indicated finite level. The quotient at level \(U_K(3)\) has finitely many components and the homotopy type of a finite cell complex by the Borel--Serre compactification. Here stabilizers at the archimedean symmetric space are trivial. An arithmetic stabilizer preserves a rank-six integral lattice, has determinant of absolute norm one, and is compact at infinity after removing the positive real center. An element fixing a point up to real scaling has scaling one by this determinant condition. The stabilizer is therefore finite and embedded, in a basis of that lattice, into matrices congruent to the identity modulo three. This principal congruence group is torsion free by Lemma \ref{lem:top}.~\cite{VI-IV-V-bs}

The space at level \(U_K(3N)\) is a finite cover of this fixed space of degree at most \([U_K(3):U_K(3N)]\). This index is at most \((3N)^{36}\), by reduction of the action on \((\OO_K/3N)^3\), a free module of rank six over \(\ZZ/3N\). Pulling back cells, allowing separate components, gives \(O_K(N^{36})\) cells in total. For one embedding the Weyl dimension formula is
\[
 \dim W_\lambda
  =\frac{(\lambda_1-\lambda_2+1)(\lambda_1-\lambda_3+2)
                             (\lambda_2-\lambda_3+1)}{2}.
\]
The rank of the product coefficient system for both embeddings is therefore \(O((h+1)^6)\). Cellular cochains with local coefficients have one copy of the fiber per cell. This proves the dimension bound.

The cuspidal part of automorphic cohomology injects into ordinary cohomology with the corresponding coefficients. In this form one uses Franke's automorphic cohomology theorem and the cuspidal direct summand: in the complex of automorphic forms the projection to the cuspidal support for the full group splits off the cuspidal cohomology. At finite level its summands are the tensor products of relative Lie algebra cohomology at infinity with finite-level fixed vectors of the cuspidal representations, counted with automorphic multiplicity. Our cohomological assumption makes the first factor nonzero and \(U_K(3N)\subset U_K(N)\) makes the second nonzero. Each distinct \(\Pi\) therefore contributes at least one dimension. There are \(O((h+1)^6)\) possible ordered weights over the two embeddings, which will allow us to sum this bound.~\cite{VI-IV-V-franke}
\end{proof}

\begin{lemma}[An integral action of the spherical operator]\label{lem:int}
In the first family \(p^{18h+1}b_{\Pi,\pm}\) is an algebraic integer of degree at most \(O_K(N^{36}(h+1)^6)\).
\end{lemma}
\begin{proof}
Let \(T^+\) be the double coset at \(w\) of \(\operatorname{diag}(p,1,1)\), and let \(T^-\) be its inverse double coset. Use hyperspecial volume one. The unramified Satake formula gives as their eigenvalues on a spherical vector \(p\sum_i\alpha_i\) and \(p\sum_i\alpha_i^{-1}\), in either order if the opposite normalization for right translation is used. The factor \(p\) is \(p^{(3-1)/2}\). For example, in the normalized unramified principal series, the index-\(p\) sublattices in the double coset divide into the three triangular orbits of sizes \(p^2,p,1\). The factors of the square root of the Borel modulus for their diagonal parts are \(p^{-1},1,p\). This gives \(p\) times the sum of the three unramified character values. The inverse calculation gives the inverse values. The \(\Pi\)-summand of Lemma~\ref{lem:coh} contains the spherical action at \(w\) since the level there is hyperspecial.

We explain integrality on the entire cohomology lattice. The coefficient system is defined over \(K\): the two factors of \((\operatorname{Res}_{K/\QQ}\GL_3)_K\) correspond to the two embeddings of \(K\) into itself. At a finite-adelic representative \(g\) let \(\Lambda_g=K^3\cap g\widehat{\OO}_K^3\). Using its images by the two embeddings and the Schur and determinant construction of Section \ref{nav:I-37}, gives an integral coefficient lattice at that representative. It transforms equivariantly under \(\GL_3(K)\) and is fixed by the finite level on the right. These lattices thus form an integral local system.

On the two legs of either Hecke correspondence, the underlying lattices, compared in the common rational space, obey \(p\Lambda\subseteq\Lambda'\subseteq p^{-1}\Lambda\). At \(w\) this follows from the elementary divisors of the double coset, and away from \(w\) the local lattices agree before the harmless multiplication by the rational integer \(p\). The same two-sided inclusion holds for dual lattices, in the reversed direction. At one embedding \(W_\lambda\) uses tensor degree \(|\lambda-(\lambda_3,\lambda_3,\lambda_3)|\le4h\) and the determinant to the power \(\lambda_3\), costing at most \(3h\) powers of \(p\) in either direction. To see the first bound, sum the two differences \(\lambda_1-\lambda_3,\lambda_2-\lambda_3\). Images under an integral Young symmetrizer preserve the resulting inclusions, since the same symmetrizer is used on the common rational tensor power. The determinant assertion follows from the third exterior power, with its dual for negative exponents. Both embeddings together cost at most \(14h\), so \(p^{18h}\) suffices.

The cohomological action of a double coset is pullback, change on coefficients, and transfer (the sum over the sheets). Pullback and transfer are integral. The preceding inclusions make the coefficient map integral after multiplication by \(p^{18h}\). Thus \(p^{18h}T^\pm\) preserves integral cohomology modulo torsion, an \(\OO_K\)-lattice in the cohomology over \(K\). It has a monic characteristic polynomial over \(\OO_K\), by localization at all finite primes, and its degree is bounded by Lemma~\ref{lem:coh}. The eigenvalue \(p^{18h+1}b_{\Pi,\pm}\) is therefore integral and has degree over \(\QQ\) at most twice that dimension.
\end{proof}

\begin{proof}[Proof of Theorem~\ref{thm:traces}]
Lemmas~\ref{lem:coh} and \ref{lem:int} give the rank-three count and degree bounds, after summing over the weights. For purity use Shin's theorem for regular algebraic conjugate-self-dual cuspidal representations of a CM field. Its infinitesimal-character condition is met by cohomology for our algebraic coefficient system. The extra weight restriction in even rank is absent in rank three. At \(w\) the representation is unramified. Shin's Theorem 1.2(iii) says that, in his motivic normalization, all its Frobenius roots have absolute value \(p^{(3-1)/2}=p\) under every complex embedding. Dividing by \(p\) gives the unitarily normalized roots. Both these numbers and their inverses have modulus one at all embeddings, so \(|b_{\Pi,\pm}|\le3\).~\cite[Thm.~1.2]{VI-IV-V-shin}

For the second family, put \(k=2r+1\). The newforms embed as independent Hecke eigensystems in \(S_k(\Gamma_1(M))\) by the classical newform theorem. Independence also follows by using their Hecke eigenvalues away from \(M\), which distinguish primitive eigensystems by multiplicity one. We recall a dimension bound. If \(d=[\operatorname{SL}_2(\ZZ):\Gamma_1(M)]\), multiply the weight-\(k\) translates of a nonzero form over a set of representatives of \(\Gamma_1(M)\backslash\operatorname{SL}_2(\ZZ)\). The product is a nonzero level-one modular form of weight \(kd\). Its order at infinity is at least that of the original form because each other translate is holomorphic there in a suitable cusp parameter. The level-one valence formula bounds this order by \(kd/12\). Hence the first \(\lfloor kd/12\rfloor+1\) Fourier coefficients determine a cusp form. Reduction modulo \(M\) gives \(d\le M^4\). This proves the needed polynomial count.

The Fourier coefficients of a normalized newform are algebraic integers, and conjugating its coefficients by a field embedding produces a normalized newform of the same weight and level with conjugate nebentypus. These facts follow from the arithmetic theory of newforms and their Hecke algebra. Consequently \(p^r b_{f,\pm}\) is integral, and the number of conjugates of either trace is bounded by the number of labels just counted. Deligne's purity theorem for holomorphic newforms of weight \(k\ge2\) gives \(|\beta_i|=p^{(k-1)/2}=p^r\) at every embedding. Thus their normalized traces with the extra unit have absolute value at most three. Here \(p\nmid M\), as required. Multiplying by \(p^{u-r}\) gives the stated common normalization.~\cite{VI-IV-V-newforms,VI-IV-V-deligne,VI-IV-V-weil}
\end{proof}

\subsection{A simultaneous bound for the Hecke congruence ideals}\label{nav:III-01}
Let \(G_N=(\ZZ/N\ZZ)^\times\). For a subgroup \(H\) with \(G_N/H\) cyclic, write \(\mathcal A_H=\{\alpha:\ker\alpha=H\}\), the full orbit of complex characters with this kernel. For the chosen \(p\) define
\[
 z_\alpha=\alpha(p)^{-1},\qquad
 d_\alpha=1+p^{-r}z_\alpha+p^r z_\alpha^{-1}.
\]
The \(\alpha\)'s index the factors even if their values at \(p\) coincide. To each label \(s\in\mathcal T\) attach the set of all conjugates of its trace \(b_s\), once per embedding of \(\QQ(b_s)\). Keep the sets for different labels disjoint. Denote the resulting labelled multiset by \(\mathcal T^*\).

\begin{theorem}[The rational integer for an orbit]\leavevmode\label{thm:integer}
There is a nonzero integer \(I_{p,H,r}\), the same for every prime \(\ell\ne p\), with
\[
     \log |I_{p,H,r}|\le C_K|\mathcal A_H|[N(r+2)]^{A'}\log p .
\]
For any common embedding of the algebraic numbers into \(\overline{\QQ}_\ell\), normalize \(v_\ell(\ell)=1\). For each \(\alpha\) choose any submultiset \(T_\alpha\) of \(\mathcal T^*\). Then
\[
  \sum_{\alpha\in\mathcal A_H}\ \sum_{b\in T_\alpha}
                 v_\ell(b-d_\alpha)\le v_\ell(I_{p,H,r}).
\]
\end{theorem}
\begin{proof}
With \(u=18h+1\) from Theorem~\ref{thm:traces} put
\[
 I_{p,H,r}=\prod_{\alpha\in\mathcal A_H}\prod_{b\in\mathcal T^*}
                    p^u(b-d_\alpha).
\]
The product is an algebraic integer since \(u\ge r\). Both multisets of values are invariant under field automorphisms and the factors are taken over their full product. Hence the integer is rational. As \(|d_\alpha|\ge p^r-1-p^{-r}>3\) and \(|b|\le3\), every factor is nonzero. Its absolute value is at most \(6p^{u+r}\). By \eqref{eq:size}, adjoining all conjugates increases cardinality by at most another polynomial in \(N(r+2)\). Taking logarithms proves the size bound. At \(\ell\ne p\) all factors have nonnegative valuation and \(p^u\) is a unit. Summing over fewer factors proves the local assertion, including duplicate traces carried by different labels.
\end{proof}

\begin{proposition}[Congruence length at a Hecke operator]\label{prop:order}
Let \(\OO\) be the valuation ring of a finite extension of \(\QQ_\ell\) containing \(d=d_\alpha\) and the trace values under consideration, where \(\ell\ne p\). Write \(e_{\OO}\) for its ramification index. For any list \(b_1,\ldots,b_s\) of such values, set \(t=(d,b_1,\ldots,b_s)\) in \(\OO^{1+s}\), and let \(\mathbb T\) be an \(\OO\)-subalgebra of that product containing \(t\) and the diagonal scalars. Define the congruence ideal
\[
   \eta(\mathbb T)=\{a\in\OO:(a,0,\ldots,0)\in\mathbb T\}.
\]
It measures the finite index along the distinguished eigensystem in the joint algebra of eigenvalues. In the algebra generated by the one operator one has
\begin{equation}
 \eta(\OO[t])=\left(\prod_{\beta\in\{b_1,\ldots,b_s\}_{\rm distinct}}
                                      (d-\beta)\right)\OO .
 \label{eq:ideal}
\end{equation}
In particular \(e_{\OO}^{-1}\len_{\OO}\OO/\eta(\mathbb T)\le\sum_{j=1}^s v_\ell(d-b_j)\).
\end{proposition}
\begin{proof}
A polynomial \(P\in\OO[X]\) represents an element supported in the distinguished coordinate precisely when it vanishes at all the \(b_j\). Since these values lie in \(\OO\), successive division by the distinct monic factors \(X-\beta\) shows that their product divides \(P\) in \(\OO[X]\). At each division any other root remains a root because \(\OO\) has no zero divisors. Evaluating at \(d\) gives \eqref{eq:ideal}, including the reverse inclusion by taking the product itself. This ideal is nonzero by Theorem~\ref{thm:integer}. Enlarging the algebra enlarges the ideal, and \(\len_{\OO}\OO/(a)=e_{\OO}v_\ell(a)\). The bound follows.
\end{proof}

This construction applies to the algebra generated by the indicated spherical trace on the two automorphic lists and the diagonal comparison parameters. The normalized trace is integral at \(\ell\ne p\). On the newform block it is a sum of the identity and the normalized \(T_p\), with the diamond operator for inverse trace. On the rank-three block it is the normalized \(T^\pm\). Equivalently one can take either orientation in separate lists. The algebra of joint eigenvalues of additional good Hecke operators, when acting on the same lists and integral at \(\ell\), contains this single-operator algebra. The proposition applies to it as well.

\begin{corollary}[The two auxiliary primes on the power interval]\label{cor:interval}
Take \(N=M_x=2^{l+2}Fqd\), \(t_x=\lfloor x^{401}\rfloor\), on the intervals of Volume~IV. Fix distinct primes \(p_1,p_2\ge7\) split in \(K\) and prime to \(3Fd_K\). For sufficiently large intervals they are prime to \(N\) everywhere. Choose the operator at \(p_2\) for \(\ell=p_1\) and at \(p_1\) otherwise. For every \(r<t_x\) and every complex character \(\alpha\) of \(G_N\), form any algebra in Proposition~\ref{prop:order} from a submultiset of the prescribed list at the chosen prime. Denote its ideal by \(\eta_{\ell,r,\alpha}\), allowing the local coefficient ring \(\OO\) to vary. Then
\begin{equation}
 \sum_{1\le r<t_x}\ \sum_{\alpha\in\widehat G_N}\ \sum_\ell
   (\log\ell)\,\frac{\len_\OO\OO/\eta_{\ell,r,\alpha}}{e_\OO}
 \ \le\ C_{K,p_1,p_2}[N(t_x+2)]^{A''}
 \ \le\ \exp(x^{\eps/3})
 \label{eq:budget}
\end{equation}
for every sufficiently large \(x\) on these intervals.
\end{corollary}
\begin{proof}
Such fixed primes exist, for instance by Dirichlet's theorem in the progression \(1\bmod 3Fd_K\) after excluding its prime divisors. In Theorem \ref{thm:twists} the auxiliary primes \(q,d\) tend to infinity with the initial scale. Thus they eventually avoid \(p_1,p_2\). For a fixed orbit and twist, Propositions~\ref{prop:order} and Theorem~\ref{thm:integer} bound the weighted sum over all \(\ell\) by \(\log|I_{p_1,H,r}|+\log|I_{p_2,H,r}|\), since \(\sum_\ell v_\ell(I)\log\ell=\log|I|\). This is uniform over all the local choices. The orbits partition the characters, so \(\sum_H|\mathcal A_H|=\varphi(N)\le N\). Sum next over \(r<t_x\) and enlarge the polynomial exponent. Finally \(\log N\ll_{\psi,\eps} x^{\eps/4}\) and \(\log(2t_x)\ll\log(2x)\) on the common family. The logarithm of the middle bound is eventually at most \(x^{\eps/3}\).
\end{proof}

\section{Orbits in the marked deformation space}\label{nav:g-IV-1}

\subsection{The contributions of the scaling orbits}\label{nav:IV-03}
\label{sec:orbit34}

We calculate the terms on the characteristic-zero side of the proposed cotangent comparison for a graded algebra. Throughout this section \(B=\bigoplus_{i\in\ZZ}B_i\) is a commutative graded \(\OO\)-algebra whose degree-zero algebra \(B_0\) is finite as an \(\OO\)-module. Assume that \(B\) is generated over \(B_0\) by finitely many homogeneous elements of degrees \(-2,-1,1,2\). Let \(k\) be an algebraic closure of the fraction field \(J\). Scaling by \(t\in k^\times\) multiplies the value of a homogeneous element of degree \(j\) at a geometric point by \(t^j\). Call an orbit \emph{mixed} if its points have nonzero values on some positive-degree and some negative-degree elements. We assume the Zariski tangent space of \(\operatorname{Spec}(B\otimes k)\) has dimension at most one at each mixed point. The Zariski tangent space at a \(k\)-point with maximal ideal \(\mathfrak m\) is the dual of \(\mathfrak m/\mathfrak m^2\).

\begin{theorem}[The components at mixed points]\leavevmode\label{thm:orbits34}
There are finitely many mixed orbits, at most \(\dim_J(B_0\otimes J)\). Each is an open and closed reduced subscheme of \(\operatorname{Spec}(B\otimes k)\) with graded coordinate ring \(k[T^g,T^{-g}]\), where \(g\in\{1,2\}\) is the greatest common divisor of the absolute degrees of the generators nonzero at any of its points. Distinct mixed orbits have distinct restrictions to \(B_0\).
\end{theorem}
\begin{proof}
Put \(C=B\otimes k\). Its degree-zero algebra \(C_0\) is Artinian and is a product of finitely many local Artin \(k\)-algebras, one for each character of \(C_0\). Their primitive idempotents also split \(C\) into open and closed parts. Fix a mixed point \(P\) in one part. Its nonzero coordinates among the chosen homogeneous generators include both signs: a homogeneous polynomial of positive (respectively negative) degree must use a generator of that sign in every monomial. Denote these coordinates by \(x_i\) and their degrees by \(d_i\), indexed by a set \(S\).

Consider any other closed point \(Q\) on the same part. It agrees on \(C_0\) with \(P\). If \(i,j\) have opposite signs, a product of positive powers of \(x_i,x_j\) has weight zero. Starting with an opposite-sign pair nonzero at \(P\), this shows that pair is nonzero at \(Q\); pairing either of them with each remaining coordinate shows that the nonzero support at \(Q\) is exactly \(S\). Write \(z_i=Q(x_i)/P(x_i)\) there. For any nonnegative integer vector \((b_i)\) with \(\sum b_i d_i=0\) we have \(\prod z_i^{b_i}=1\), by equality on \(C_0\). There exists such a zero-weight vector strictly positive in every coordinate: add, for each coordinate, a relation with a coordinate of opposite sign using positive exponents on the pair. Adding a sufficiently large multiple of this vector proves the same multiplicative relation for \emph{every} integral vector of weight zero. Consequently \((b_i)\mapsto\prod z_i^{b_i}\) factors through the degree map \(\ZZ^S\to g\ZZ\). Choose a \(g\)-th root \(t\in k^\times\) of the value at \(g\). Then \(Q(x_i)=t^{d_i}P(x_i)\) on \(S\), and all remaining coordinates vanish at both points. This identifies \(Q\) with the translate of \(P\).

The map from the coordinate ring of this part to \(k[T,T^{-1}]\) given by the orbit of \(P\) has image \(k[T^g,T^{-g}]\). Indeed the monoid generated by the supported degrees equals the group they generate: the strictly positive zero-weight relation just constructed turns every integral representation of a degree into a nonnegative one. The quotient map thus embeds this one-dimensional torus as a closed subscheme. All closed points of the part belong to it by the preceding paragraph. It remains to check reducedness, since equality on closed points alone would only detect the radical. At each closed point the local ring maps onto the local ring of this torus, so its dimension is at least one. Its embedding dimension is at most one by the tangent hypothesis; the local ring is Noetherian, and dimension is bounded by embedding dimension. It is therefore a regular local ring of dimension one, hence a discrete valuation ring. The displayed quotient has dimension one, so its localized kernel is zero, since a quotient of a discrete valuation ring by a nonzero ideal has dimension zero. The global kernel is zero as well: a nonzero module has an element with proper annihilator, surviving at a maximal ideal containing that annihilator. All maximal ideals of our finite-type \(k\)-algebra are closed points. This proves the scheme and grading assertions and the bound from the dimension of \(C_0\).
\end{proof}

Let \(B^{\rm mix}_J\) be the direct summand of \(B\otimes J\) which after extension to \(k\) is the union of the mixed parts. It is defined over \(J\): the sum of their primitive degree-zero idempotents is fixed by every automorphism of \(k/J\), since being mixed is defined by homogeneous elements and their signs. For \(j\in\ZZ\) let \(L_j\) be the image of \(B_j\) in its degree-\(j\) summand, and let \(K_j=\ker(B_j\to L_j)\). An empty union gives \(L_j=0\).

\begin{proposition}[The determinant in each weight]\label{prop:horizontal34}
Each \(L_j\) is finite free over \(\OO\). Let \(a\in B_0\) be nonzero at every mixed orbit. With one representative \(P\) per geometric mixed orbit and \(g_P\) as above, one has
\begin{equation}
 \len_{\OO}(L_j/aL_j)=\sum_{P:\ g_P\mid j}v_\ell(a(P)).
                                                               \label{eq:horiz34}
\end{equation}
There is an exact sequence
\[
 0\longrightarrow K_j/aK_j\longrightarrow B_j/aB_j
       \longrightarrow L_j/aL_j\longrightarrow0 .
\]
In particular the sum of the lengths in weights \(2\) and \(-2\) on the mixed parts is exactly twice the sum of the trace valuations when \(a\) is the weight-zero trace difference.
\end{proposition}
\begin{proof}
The component of a finitely generated graded algebra in a fixed degree is a finite module over its degree-zero algebra in this situation. In fact exponent vectors of monomials of a prescribed weight have finitely many minimal elements under coordinatewise comparison. Otherwise their infinite antichain would contain an infinite coordinatewise nondecreasing subsequence, by taking a nondecreasing subsequence one coordinate at a time (an unbounded sequence in \(\mathbb N\) has such a subsequence, as does a bounded infinite sequence). Subtracting a minimal vector lying below a given vector gives a nonnegative vector of weight zero. The finitely many minimal monomials therefore generate that component over \(B_0\). It follows that \(B_j\) is finite over \(\OO\), and that its image \(L_j\) is finite and torsion-free, hence free. Its generic fiber is the entire degree-\(j\) component of the mixed summand, since extension to \(J\) commutes with the image. By Theorem~\ref{thm:orbits34} each geometric orbit contributes one basis line if \(g_P\mid j\), and zero otherwise. On that line \(a\) acts by \(a(P)\). Its determinant valuation counts the length of its cokernel by elementary divisors; the nonzero eigenvalues are integral over \(\OO\) because \(a\in B_0\). This proves \eqref{eq:horiz34}. Finally, \(aB_j\cap K_j=aK_j\) since multiplication by \(a\) is injective on \(L_j\). This proves the exact sequence. Both \(1\) and \(2\) divide each of \(2,-2\), giving the last assertion.
\end{proof}

The trace values in this computation refer to genuine generic orbits. The term \(K_j/aK_j\) includes contributions supported away from the mixed parts. It can have positive rank. A comparison between a cotangent quotient at the augmentation and the two modules in \eqref{eq:horiz34} therefore needs control of those other contributions and of the integral map. We next establish that comparison for a class of presentations in which the opposite weight has a single coordinate.

\subsection{A graded presentation with one positive coordinate}\label{nav:IV-04}
\label{sec:slice34}

Let \(u_1,\ldots,u_n\) have degree \(-1\), let \(v\) have degree \(1\), and let \(A=\OO[u_1,\ldots,u_n,v]\) with augmentation all variables zero. Every homogeneous element \(F_i\) of degree \(-1\) can be written
\[
 F_i=\sum_{j=1}^n u_j f_{ij}(v u_1,\ldots,v u_n),\qquad
                    f_{ij}\in\OO[t_1,\ldots,t_n].
\]
Indeed each monomial has one more \(u\)-factor than \(v\)-factors. The coefficients of its linear terms are uniquely defined even though the displayed expression may have more than one choice. Put \(B=A/(F_1,\ldots,F_s)\), \(t_j=v u_j\), \(G_i(t)=\sum_j t_j f_{ij}(t)\), and
\[
       R=\OO[t_1,\ldots,t_n]/(G_1,\ldots,G_s),\qquad
       C_{-1}=\OO^n / \sum_i \OO(f_{i1}(0),\ldots,f_{in}(0)).
\]
The last group is exactly the quotient of the ambient degree-\(-1\) cotangent lattice by the differentials of the equations. Saturation of the square of the ambient augmentation ideal changes nothing, because the linear terms of a polynomial form a free \(\OO\)-module.

\begin{theorem}[An integral comparison in rank one on the positive side]\label{thm:slice34}
Suppose \(R\) is finite over \(\OO\) and its generic fiber is reduced. Write \(T\) for its \(\OO\)-torsion, \(I=\ker(R\to\OO,\ t_j\mapsto0)\), and let \(a\in B_0\) vanish at the origin but be nonzero at every mixed orbit. Then
\begin{equation}
 \len_{\OO} C_{-1}\le
 \len_{\OO}\bigl(T/(T\cap I^2)\bigr)+
                       \sum_{P\ {\rm mixed}} v_\ell(a(P)).    \label{eq:slice34}
\end{equation}
If \(R\) is torsion-free the orbit sum alone suffices. In this presentation the cotangent quotient is determined by the integral degree-zero ring, even though a first-order deformation in the individual \(u_j\) coordinates has constant values on all \(t_j\) at the origin.
\end{theorem}
\begin{proof}
The degree-zero piece of \(A\) is precisely the polynomial ring on \(t_1,\ldots,t_n\), by inspection of monomials. The degree-zero piece of the homogeneous ideal generated by the \(F_i\) is generated there by \(vF_i=G_i(t)\), since every homogeneous coefficient of degree one is \(v\) times a polynomial in the \(t_j\). Thus \(B_0=R\). In \(R\) the cotangent module at the origin is generated by the \(t_j\) modulo the linear terms of the \(G_i\): multiplying an equation by a polynomial with zero constant term contributes only to the square of the augmentation ideal. These linear terms have coefficient rows precisely \((f_{ij}(0))_j\). This gives an integral isomorphism \(I/I^2\simeq C_{-1}\).

Over \(k\) a mixed point necessarily has \(v\ne0\) and some \(u_j\ne0\). Scale uniquely to \(v=1\); its coordinates \(t_j=u_j\) solve the equations \(G_i(t)=0\) and are not all zero. Conversely any nonzero geometric point of \(R\) gives exactly this normalized mixed point. This is a bijection between mixed orbits and the geometric points of \(R\) distinct from the origin. The values of \(a\in B_0=R\) agree under it. Substitution in Theorem~\ref{thm:order31} proves \eqref{eq:slice34}. The last statement follows directly from \(t_j=v u_j\) when the basepoint has \(v=0\).
\end{proof}

For instance take \(n=s=1\) and \(F_1=u f(uv)\) for a monic polynomial \(f\in\OO[t]\) with distinct roots in characteristic zero and \(f(0)\ne0\). Then \(R=\OO[t]/(t f(t))\) is finite free and generically reduced, and \(C_{-1}=\OO/f(0)\OO\). For \(a=t\) equality holds in \eqref{eq:slice34}: the mixed orbits are parametrized by the roots \(\lambda\) of \(f\) and \(v_\ell(f(0))=\sum_\lambda v_\ell(\lambda)\). Taking \(f(t)=t-\ell^b\) realizes every positive length \(b\).

The proposed arithmetic model has several positive coordinates, including weight two, and equations of both signs and weight zero. The hypotheses recorded in \srcloc{Development/AutomorphicArithmetic/AutomorphicVarieties/CommutativeAlgebra/OrbitVanishingLocal}{103} allow that structure. The orbit conditions and the sum over representatives appear at lines 45--70 of the same file. The determinant formula here accounts for their generic mixed contributions. The special presentation in Theorem~\ref{thm:slice34} realizes an integral cotangent comparison in one case. Applying the full comparison to all arithmetic components of Proposition \ref{prop:minors33} requires a smooth chart and equations for the actual local conditions, together with automorphy and trace control for their mixed orbits.

\subsection{An upper bound from the mixed orbits}\label{nav:IV-05}
\label{sec:upper35}

Let \(\OO\) be a discrete valuation ring with uniformizer \(\pi\), fraction field \(J\) of characteristic zero, and algebraic closure \(k\) of \(J\). Use the valuation on \(k\) extending that of \(J\) with \(v(\pi)=1\), and assume \(J\) is complete. Let \(A=\bigoplus_{j\in\ZZ}A_j\) be generated over its Noetherian degree-zero ring \(D=A_0\) by finitely many elements of weights \(-2,-1,1,2\). Fix an \(\OO\)-augmentation \(x:A\to\OO\) vanishing on nonzero weights. Write \(\mathfrak b=\ker x\). Take homogeneous equations
\[
 \begin{array}{c|ccccc}
 \text{equations}&F_1,\ldots,F_r,F_*&G_1,\ldots,G_s&E_1,\ldots,E_{r'}&
                          E'_1,\ldots,E'_{s'}&Z_1,\ldots,Z_m\\
 \text{weight}&-1&-2&1&2&0 .
 \end{array}
\]
They all vanish at \(x\). Write \(I\) for their ideal, \(B=A/I\) and \(I_0=I\cap D\). Use the following hypotheses.

\begin{enumerate}\setlength{\itemsep}{1pt}
\item The invariant algebra \(B_0=D/I_0\) is finite over \(\OO\). In the cotangent space \((\mathfrak b/\mathfrak b^2)\otimes_{\OO}J\), the \(F\)'s span weight \(-1\) and the \(G\)'s span weight \(-2\).
\item For every geometric fixed point of \(B\) different from \(x\), its invariant character does not support both a purely positive and a purely negative point of \(B\). A purely positive point kills all negative weights and has some nonzero positive value. Purely negative points are defined by reversing signs. A fixed point kills all nonzero weights.
\item At each mixed geometric point of \(B\) the tangent space has dimension at most one. A mixed point has a nonzero value of each sign. Fix \(a\in D\) with \(x(a)=0\) and \(a(P)\ne0\) on all mixed points.
\end{enumerate}
All geometric points here are over \(k\). Hypothesis 1 makes \(B\otimes k\) a finite-type \(k\)-algebra. Section~\ref{sec:orbit34} gives finitely many mixed orbits, with distinct invariant characters.

\subsubsection{The three graded modules}\label{nav:IV-06}
Over \(D\), set \(U=A_{-2}\), \(V=A_2\), \(W= D F_*^2+U^2V\), and let \(D_1=UV\) denote the ideal of \(D\) generated by these products. In particular \(W\subset U\). Inside \(A[T_0,T_1,T_2]\) form the algebra
\[
 R=\bigoplus_{d\in\mathbb N^3}
 U^{d_0}V^{d_1}W^{d_2}T_0^{d_0}T_1^{d_1}T_2^{d_2}.
\]
Module products mean spans of products over \(D\) and the zeroth power means \(D\). Each \(A_j\) is a finite \(D\)-module by the monomial argument in Proposition~\ref{prop:horizontal34}. Thus \(R\) is a Noetherian standard multigraded algebra, meaning that its three unit degrees generate it over \(D\). Denote the following list of sections by \(\mathbf t\):
\begin{equation}
 F_i^2T_0\ (1\le i\le r),\quad G_jT_0,\quad E_i^2T_1,\quad
 E'_jT_1,\quad F_*^2T_2 .                                      \label{eq:sections35}
\end{equation}
Let \(H\) be any multihomogeneous ideal containing these sections and all \(Z_j\). For \(d\in\mathbb N^3\) put \(M_d=(R/H)_d\). We call a prime of \(R\) relevant if it contains none of the three ideals generated by a whole unit-degree piece. The product \(\mathfrak h\) of these three ideals is contained precisely in the irrelevant primes.

\begin{lemma}[The invariant character at a relevant prime]\label{lem:support35}
If \(\mathfrak p\supset H\) is relevant, then \(I_0\subset\mathfrak p\cap D\). For all sufficiently large \(d\), coordinate by coordinate, \(M_d\) is finite over \(\OO\). Its geometric generic support consists only of invariant characters of mixed orbits.
\end{lemma}
\begin{proof}
Choose \(u\in U,v\in V,w\in W\) whose sections in the three colors survive modulo \(\mathfrak p\). The identity \((F_*^2 T_0)(wT_2)=(wT_0)(F_*^2T_2)\) puts \(F_*^2T_0\) in \(\mathfrak p\). Every negative squared or weight-two equation \(h\in U\) in the list therefore has \(hT_0\in\mathfrak p\). If \(c\in V\), multiplication gives
\[
 (ch)(uT_0)(vT_1)=(hT_0)((cuv)T_1).
\]
The two chosen sections survive in a domain, so \(ch\in\mathfrak p\). For a negative weight-one equation \(F\) and \(b\in A_1\), apply this argument to \(h=F^2,\ c=b^2\) and take radicals. Interchanging the colors gives the same assertion for positive equations and opposite-weight coefficients. The \(Z_j\) already lie in \(\mathfrak p\). On projecting any expression in the ideal of the equations to weight zero, only these opposite-weight coefficients occur. This proves the assertion for \(I_0\).

It follows that \(I_0\mathfrak h\subset\sqrt H\). Taking a power in the Noetherian ring gives \(I_0^b\mathfrak h^b\subset H\) for some \(b\). Every degree \(d\) monomial with all \(d_i\ge b\) is divisible by a product of \(b\) generators of each color. Consequently \(I_0^b M_d=0\). The ring \(D/I_0^b\) is finite over \(\OO\). Indeed its successive factors \(I_0^j/I_0^{j+1}\) are finite modules over \(D/I_0\). This proves the finiteness assertion and leaves only finitely many possible geometric invariant characters \(\sigma:D\to k\) in the support.

Suppose one such character kills \(D_1\). We show that the fiber of \(R/H\) over \(\sigma\) has no relevant prime. Otherwise choose one, and take a field containing the fraction field of its residue domain. The following elementary identity holds for every \(n\ge0\):
\[
                      A_{2n}=V^n,\qquad A_{-2n}=U^n .
\]
In a monomial of even weight pair all the factors of odd weight. Each pair has weight \(0\) or \(\pm2\). Pair off opposite factors of weight two and absorb their products into the coefficients in \(D\). The unpaired factors have the specified sign and number.

The positive even subalgebra \(\bigoplus_{n\ge0} A_{2n}\) maps to the field just chosen by evaluating \(g\) of weight \(2n\) as \(gT_1^n\). It is nonzero on some \(v\in V\) and kills \(E_i^2,E'_j\). The full nonnegative subalgebra is integral over its even subalgebra, since the square of every odd homogeneous element has even weight and these elements generate the extension. By lying over, this character extends to the nonnegative subalgebra over an algebraically closed extension field. In particular it kills every \(E_i\). It extends further to \(A\) by setting every negative weight to zero. To verify multiplicativity, take factors of weights \(-i,j\) with \(j\ge i>0\). The square of their product belongs to \(U^iV^j=D_1^i V^{j-i}\); its positive even evaluation is zero, since \(\sigma(D_1)=0\). Products of negative total weight already go to zero. This gives a purely positive point of \(B\) with invariant character \(\sigma\). Applying the construction in the negative color, using \(F_*^2T_0\in\mathfrak p\) in the fiber as above, gives a purely negative point of \(B\). The equations, invariant character, and a required nonzero coordinate define finite-type locally closed conditions in \(\operatorname{Spec}(B\otimes_D k)\). Existence over an extension of \(k\) gives existence over \(k\) by the Nullstellensatz. Here the tensor product uses \(\sigma\). The requirements to kill all weights of one sign can equivalently be imposed by their ideal in this Noetherian fiber.

The same invariant character with all nonzero weights set to zero gives a fixed point of \(B\). Multiplicativity follows since \((fg)^2\in D_1^i\) for factors of opposite weights \(i,-i\). Hypothesis 2 now forces this fixed point to be \(x\). A purely negative solution over \(x\) is impossible. On such a solution every product of two homogeneous elements of \(\mathfrak b\) of total weight \(-1\) vanishes, as one factor has nonnegative weight. The weight-zero factors in \(\mathfrak b\) vanish by equality on \(D\). Rational cotangent spanning by the \(F\)'s then kills all of \(A_{-1}\). In total weight \(-2\) the only additional pair of negative weights is \((-1,-1)\). These products also vanish, and cotangent spanning by the \(G\)'s kills \(A_{-2}\). The algebra generators now preclude a nonzero negative value.

Thus in the finite-type fiber over each such \(\sigma\) the product of the three irrelevant ideals is nilpotent modulo \(H\). All sufficiently high pieces of that fiber vanish by standard generation. For each degree, Nakayama's lemma then excludes support over \(\sigma\) on the finite geometric generic module. There are only finitely many \(\sigma\) to exclude. At every remaining character \(D_1\) has nonzero image. Such a character of \(B_0\) extends to a point of \(B\) over \(k\): the fiber has nonzero degree-zero piece \(k\) and is of finite type. A nonzero product from \(UV\) forces both signs there. This proves the lemma.
\end{proof}

\begin{theorem}[The upper estimate for all the coordinates]\leavevmode\label{thm:upper35}
For all sufficiently large \(d\), the kernel and cokernel of multiplication by \(a\) on \(M_d\) have finite \(\OO\)-length and
\begin{equation}
 0\le \len_{\OO}(M_d/aM_d)-\len_{\OO}M_d[a]
       \le 2^{r+1+r'}\sum_{P\ {\rm mixed}}v(a(P)).              \label{eq:upper35}
\end{equation}
Use one representative per geometric orbit. The bound is uniform in the multidegree after its three coordinates pass a threshold depending on the presentation and \(H\).
\end{theorem}
\begin{proof}
At a surviving invariant character \(\sigma\), localize \(D\otimes k\) at its kernel, and call the resulting ring \(C\). Some product \(uv\), \(u\in U,v\in V\), is a unit. In the localized ambient algebra \(u\) and \(v\) are therefore units. The modules \(U,V,W\) become \(Cu,Cv,Cu\), respectively. For example any negative weight-two element divided by \(u\) has weight zero, and \(u^2v\) generates \(Cu\). Extension from \(\OO\) to \(k\) and this localization are flat, so they preserve the defining inclusions of the graded pieces. The localized ring \(R\) is a polynomial ring over \(C\) in \(uT_0,vT_1,uT_2\). Each piece of degree at least one in every coordinate, modulo \(H\), is consequently a quotient of the cyclic module over \(C\) obtained by imposing
\[
 F_i^2/u,\quad F_*^2/u,\quad G_j/u,\quad E_i^2/v,\quad E'_j/v,\quad Z_j .
\]
This module is the localized weight-zero piece of \(A\) modulo the ordinary equations with the \(r+1+r'\) weight-one equations squared. To see the equality, take weight-zero parts of the ideal. Coefficients of opposite weight two are precisely multiples of \(u\) or its inverse as appropriate.

With the original equations, every localized weight piece has dimension at most one over \(k\), by Theorem~\ref{thm:orbits34}. Replacing one homogeneous equation \(f\) of weight \(b\) by \(f^2\) at most doubles a common upper bound for these dimensions. Indeed in the algebra \(S\) with the other current equations imposed there is an exact sequence
\[
 (S/(f))_{j-b}\xrightarrow{\ f\ }(S/(f^2))_j
                         \longrightarrow (S/(f))_j\longrightarrow0 .
\]
Localization and projection to a weight are exact. Induction gives dimension at most \(2^{r+1+r'}\) for the localized piece of \(M_d\).

By Lemma~\ref{lem:support35}, \(M_d\) is finite over \(\OO\). On its torsion submodule every endomorphism has kernel and cokernel of equal length. On its free quotient the determinant of multiplication by \(a\) is nonzero, since all generic eigenvalues are the values at surviving mixed characters. After scalar extension to \(k\), the contribution to this determinant from a character \(\sigma\) is \(\sigma(a)\) to the power of its local dimension. Indeed \(a-\sigma(a)\) is nilpotent on that finite local summand. The kernel--cokernel exact sequence and elementary divisors now identify the left difference in \eqref{eq:upper35} with the determinant valuation. Each \(a(P)\) is integral since \(a\) acts through the finite invariant algebra on the original solution, so these valuations are nonnegative. The dimension bound and Lemma~\ref{lem:support35} give the assertion.
\end{proof}

The cancellation of the two lengths in \eqref{eq:upper35} accounts for torsion supported at the residue characteristic. A lower bound for this same difference in terms of the cotangent lattice requires the formal chart at \(x\).

\lead{Locations in the accompanying source} The definitions of the three modules and of the sections are in \srcloc{Development/AutomorphicArithmetic/AutomorphicVarieties/CommutativeAlgebra/GradedModulesLength/ToDegInjective}{187}. See also \srcloc{Development/AutomorphicArithmetic/AutomorphicVarieties/CommutativeAlgebra/GradedModulesLength/ExistsSecNotMem}{63} for the relevant primes and \srcloc{Development/AutomorphicArithmetic/AutomorphicVarieties/CommutativeAlgebra/GradedModulesLength/TwoSigned}{17} for the orbit estimate.

\subsection{Euler characteristics in high multidegrees}\label{nav:IV-07}
\label{sec:koszul36}

We give the integral argument that will be used at the distinguished point. Let \(R\) be a Noetherian standard \(\mathbb N^3\)-graded algebra over a Noetherian \(\OO\)-algebra \(D\), and \(L\) a finite graded \(R\)-module. A tuple \(\mathbf s=(s_1,\ldots,s_b)\) of homogeneous sections of nonnegative multidegrees is admissible on \(L\) if the pieces of \(L/(\mathbf s)L\) have finite \(\OO\)-length in all sufficiently large multidegrees. The Koszul complex of the tuple is the tensor product of the two-term complexes of multiplication \(R(-\deg s_i)\to R\). With coefficients in \(L\), write
\[
 \chi_d(\mathbf s;L)=\sum_{j=0}^{b}(-1)^j
             \len_{\OO} H_j(\mathbf s;L)_d .
\]
Here \(L(-b)_d=L_{d-b}\) for a multidegree \(b\). We use this expression only when all lengths are finite.

\begin{lemma}[Positivity and passage to a quotient]\label{lem:euler36}
For an admissible tuple, \(\chi_d(\mathbf s;L)\) exists and is nonnegative for all sufficiently large \(d\). It is additive on exact sequences when the tuple is admissible on the middle term. For an initial subtuple \(\mathbf u\) followed by \(\mathbf v\), one has eventually
\[
       \chi_d(\mathbf u,\mathbf v;L)
                         \le \chi_d(\mathbf v;L/(\mathbf u)L).
\]
A regular first section can be removed by passing to its quotient with equality.
\end{lemma}
\begin{proof}
First observe a support fact. If a finite graded module \(N\) has finite lengths in high degrees, so does any finite graded module supported on \(\operatorname{Supp} N\). Indeed, sending a scalar to its products with homogeneous generators embeds \(R/\operatorname{Ann}N\) into a finite sum of shifts of \(N\). The same high-degree property holds modulo any power of \(\operatorname{Ann}N\), by the filtration by ideal powers, whose factors are finite graded modules over \(R/\operatorname{Ann}N\). A power of this annihilator kills any finite graded module supported on \(\operatorname{Supp}N\), by the radical criterion for containment of closed sets in a Noetherian spectrum. A finite homogeneous set of generators proves the claim.

For \(N=L/(\mathbf s)L\) its support is \(\operatorname{Supp}L\cap V(\mathbf s)\) by Nakayama's lemma. Every Koszul homology is supported on this set, since a section that becomes a unit contracts its two-term complex. This proves finiteness and shows that subquotients of \(L\) inherit admissibility. The homology long exact sequence proves additivity degree by degree.

We prove positivity by induction on the number of sections. The empty tuple gives a length. A finite graded module over a Noetherian graded ring has a finite filtration with factors shifts of quotients by homogeneous primes. To recall the argument, select a nonzero homogeneous element whose annihilator is maximal among annihilators of such elements. Its annihilator is prime on homogeneous elements, as follows by applying maximality to a nonzero homogeneous multiple. It is prime on all elements by taking extreme nonzero homogeneous terms in the quotient using a translation-invariant lexicographic ordering of the degrees. Iterate in the quotient module. The ascending chain condition stops the construction. Additivity therefore reduces positivity to \(L=R/\mathfrak p\), allowing a degree shift. An irrelevant prime gives zero in high degrees by standard generation. In the domain of a relevant prime, a nonzero first section is regular. Its complex resolves the quotient, so the induction hypothesis for the remaining sections applies.

If the first section is zero in this domain, the remaining tuple \(\mathbf v\) is already admissible. For any first section of the same multidegree \(e\), taking homology first in \(\mathbf v\) shows that the Euler characteristic of the full complex is
\(\chi_d(\mathbf v;L)-\chi_{d-e}(\mathbf v;L)\). All lengths on this intermediate page are finite in high degrees and differentials preserve the alternating sum. Relevance supplies a nonzero section of degree \(e\): take a product of surviving generators of the three colors, or one if \(e=0\). Replacing the zero section by this regular section preserves the Euler characteristic and gives positivity by induction.

Finally take homology first in a single initial section \(s\) of degree \(e\). The two rows are its quotient and its kernel, the latter shifted by \(e\) and by one in homological degree. Hence
\[
 \chi_d(s,\mathbf v;L)=\chi_d(\mathbf v;L/sL)
                      -\chi_{d-e}(\mathbf v;L[s]).
\]
Both tuples on the right are admissible by the support observation. Positivity bounds the left side by its first term. Iteration proves the inequality. The same identity gives equality for a regular section.
\end{proof}

\begin{lemma}[The contribution of a scalar factor]\label{lem:peel36}
Suppose \(\pi\) acts injectively on \(L\) and every graded piece of \(L/\pi L\) has finite length. Let \(y\) be homogeneous and assume \((\pi^h y,\mathbf v)\) is admissible. Then in high multidegrees
\[
  \chi_d(\pi^h y,\mathbf v;L)=\chi_d(y,\mathbf v;L)
                                    +h\,\chi_d(\mathbf v;L/\pi L).
\]
For a function \(f(d)\) let \(\Delta_e f(d)=f(d)-f(d-e)\). If \(H(d)=\len_{\OO}(L/\pi L)_d\), then \(\chi_d(\mathbf v;L/\pi L)=\prod_{s\in\mathbf v}\Delta_{\deg s}H(d)\).
\end{lemma}
\begin{proof}
Replacing \(\pi^h y\) by \(y\) only shrinks the support of the final quotient, so admissibility follows from the proof of Lemma~\ref{lem:euler36}. The kernels of the two first sections agree because \(\pi\) is regular. Their quotients fit into
\[
 0\longrightarrow L/yL\xrightarrow{\ \pi^h\ }L/(\pi^hy)L
                                     \longrightarrow L/\pi^h L\longrightarrow0 .
\]
Regularity identifies all \(h\) successive factors of \(L/\pi^hL\) with \(L/\pi L\) without a shift in multidegree. Apply additivity and the two-row identity in the previous proof. On \(L/\pi L\) the terms of the complex themselves have finite length, so their alternating sum is the Euler characteristic. Writing the terms as direct sums of shifts gives the finite differences.
\end{proof}

\subsection{The cotangent length at the distinguished point}\label{nav:IV-08}
\label{sec:lower36}

Use the presentation and three hypotheses of Section~\ref{sec:upper35}. For the lower estimate assume \(r\ge1\) and that \(x\) has the following homogeneous formal chart. There is an \(\OO\)-algebra map from \(A\) to a formal power series ring over \(J\) whose constant term is \(x\) and which preserves weights monomial by monomial. Taking linear terms identifies \((\mathfrak b/\mathfrak b^2)\otimes J\) with the span of the variables. Their numbers are
\[
 \begin{array}{c|rrrrr}
 \text{weight}&-1&-2&1&2&0\\ \text{number}&r&s&r'+1&s'+1&m .
 \end{array}
\]
The \(m\) equations \(Z_j\) span the rational cotangent of weight zero. These are hypotheses about the actual augmented algebra, including the existence of the map on all its relations.

Let \(\mathfrak b^{(2)}=\{b\in\mathfrak b:\lambda b\in\mathfrak b^2\text{ for some }0\ne\lambda\in\OO\}\). The lattice \(C\) is the image of \(A_{-1}\) in \(\mathfrak b/\mathfrak b^{(2)}\), and \(C_F\) is its quotient by the classes of all the \(F\)'s. The ambient algebra is Noetherian, so \(\mathfrak b/\mathfrak b^2\) is a finite module over \(A/\mathfrak b=\OO\). It follows that \(C\) is finite free of rank \(r\). Rational cotangent spanning gives \(C_F\) finite length.

\begin{theorem}[Integral cotangent length and the mixed orbits]\label{thm:full36}
Under these hypotheses one has
\begin{equation}
                \len_{\OO}C_F\ \le\ 2\sum_{P\ {\rm mixed}}v(a(P)).
                                                               \label{eq:full36}
\end{equation}
There is one term for each geometric scaling orbit of the actual solution algebra \(B=A/I\).
\end{theorem}
\begin{proof}
Select a basis of \(C\) and make invertible \(\OO\)-linear combinations of the \(r+1\) negative weight-one equations. Smith normal form gives elements \(y_i\in A_{-1}\) lifting a basis of \(C\) and integers \(c_i\ge0\) with
\[
 F_i-\pi^{c_i}y_i\in\mathfrak b^{(2)}\ (1\le i\le r),\qquad
 F_*\in\mathfrak b^{(2)},\qquad \len_{\OO} C_F=\sum_{i=1}^r c_i .
\]
The ordinary equation ideal and the hypotheses of Section~\ref{sec:upper35} persist under this change. Use the new equations throughout the construction of \(R\).

Choose rational homogeneous lifts of a linear coordinate basis in every nonzero weight, choosing the \(y_i\) in weight \(-1\), and use the \(Z_j\) in weight zero. Their images are formal parameters. Indeed a substitution by power series with zero constant terms and invertible linear matrix is formally invertible: after the inverse is chosen up to degree \(n-1\), its degree-\(n\) error is canceled using the inverse linear matrix. The coefficients stabilize in every degree. Both this substitution and its inverse preserve the weights, as the recursion does. Using these parameters and then setting the zero-weight parameters to zero gives a map \(\phi\) to a power series ring in the nonzero-weight variables. It kills all \(Z_j\). Each remaining variable is the exact image of a homogeneous element of \(A\otimes J\). Elements of \(\mathfrak b^{(2)}\) have zero linear term under \(\phi\), since after multiplication by a nonzero scalar they lie in a square of the augmentation ideal.

We may suppose \(\phi(a)\ne0\). If it vanishes, choose some \(F_i\) of nonzero linear term and an integral positive weight-one element \(z\) with nonzero linear term, which exists by the chart and clearing denominators. Replace \(a\) by \(a+zF_i\). The value at \(x\), all the values on \(B\), and hence the desired bound stay the same. In the formal series domain \(\phi(zF_i)\ne0\).

Apply \(\phi\) coefficientwise to \(R\), retaining the three variables \(T_i\), and write \(K\) for its kernel. The ring \(R/K\) is a domain, \(a\) is regular there, and \(Z_j\in K\). Put \(H=K+(\mathbf t)\). Theorem~\ref{thm:upper35} makes \((\mathbf t,a)\) admissible on \(R/K\). Lemma~\ref{lem:euler36}, permutation of Koszul factors, and regularity of \(a\) give
\begin{equation*}
 \chi_d(\mathbf t;R/(K,a))=\chi_d(\mathbf t,a;R/K)
                \le\chi_d(a;R/H)
                \le 2^{r+1+r'}\sum_P v(a(P)).                 \qquad(\ast)
\end{equation*}
Here the one-section characteristic is exactly the difference of lengths in Theorem~\ref{thm:upper35}.

For a monomial in the formal coordinates, let \(n_-\) be the sum of the absolute weights of its negative factors, counted with multiplicities. Define \(n_+\) in the same way with its positive factors. With \(d=(d_0,d_1,d_2)\), set
\[
               b_-(d)=2d_0+4d_2,\qquad b_+(d)=2d_1+2d_2 .
\]
The power series coefficient of an element of \(R_d\) has \(n_-\ge b_-(d)\) in every monomial. Indeed a weight \(-2\) term has \(n_-\ge2\). A term of \(\phi(F_*)\) has \(n_-\ge2\), since in weight \(-1\) a term with \(n_-=1\) is linear. Thus \(F_*^2\) and \(U^2V\) both have \(n_-\ge4\), giving the assertion by products. The original signed weight in degree \(d\) is \(2(d_1-d_0-d_2)\). Consequently the monomials with \(n_-=b_-(d)\) have \(n_+=b_+(d)\).

In each degree retain exactly the terms with this pair of weights and carry along \(T^d\). This operation on \(R\) is an algebra homomorphism into a polynomial ring over \(J\): the minima add, so a product at the specified pair uses only terms at the respective minima. Let \(L\) be its image. In degree zero only constants remain. Thus \(D\) acts through \(x:D\to\OO\), and the map factors through \(R/(K,a)\). Each \(L_d\) is a finite torsion-free \(\OO\)-module, being the image of a finite \(D\)-module in a finite-dimensional \(J\)-space on which \(D\) acts through \(x\). It spans the entire space of the monomials with that pair of weights over \(J\). To check the last assertion, group the negative factors of such a monomial into blocks of absolute weight two. A variable of weight two is one block and two variables of weight one give one block. The total is even, so every factor is used. Make the same grouping for the positive factors. Use one negative block for each of the \(d_0\) slots, one positive for each of the \(d_1\) slots, and two negative blocks and one positive block for each of the \(d_2\) slots. The last product belongs to \(U^2V\otimes J\), and the others lie in \(U\otimes J\) or \(V\otimes J\), because the individual variables have exact rational homogeneous lifts. This realizes the prescribed monomial.

The images in \(L\) of the sections \(F_i^2T_0\) are \(\pi^{2c_i}\) times the images of \(y_i^2T_0\), since the retained terms of \(F_i\) in weight \(-1\) are its linear terms. By Lemma~\ref{lem:euler36}, the quotient onto \(L\) and its kernel give \(\chi_d(\mathbf t;R/(K,a))\ge\chi_d(\mathbf t;L)\). On \(L\) strip off the \(r\) scalar factors using Lemma~\ref{lem:peel36}, permuting the complex to put each chosen factor first. The characteristic of the tuple with all factors stripped is nonnegative by Lemma~\ref{lem:euler36}. For every stripped factor the remaining tuple has the same list of degrees. Thus, writing \(e_i\) for the three unit vectors and \(H_0(d)=\len (L/\pi L)_d\), we obtain
\[
 \chi_d(\mathbf t;L)\ \ge\ 2\Big(\sum_i c_i\Big)
      \Delta_{e_2}\Delta_{e_1}^{r'+s'}\Delta_{e_0}^{r+s-1} H_0(d).
\]

It remains to compute this integer. For \(p\ge1,q\ge0\), let \(M_{p,q}(n)\) count monomials of total weight \(2n\) in \(p\) variables of weight one and \(q\) of weight two. The finite free rank calculated above gives
\[
 H_0(d)=M_{r,s}(d_0+2d_2)\,
                   M_{r'+1,s'+1}(d_1+d_2).
\]
The generating function by total weight is
\((1-z)^{-p}(1-z^2)^{-q}=(1+z)^p(1-z^2)^{-p-q}\).
Reading its even coefficients shows that for large \(n\), \(M_{p,q}(n)\) is a polynomial of degree \(p+q-1\) with leading coefficient \(2^{p-1}/(p+q-1)!\). Indeed the even-indexed coefficients of \((1+z)^p\) sum to \(2^{p-1}\), and every resulting binomial polynomial has the displayed factorial as denominator of its leading term. Applying \(r+s-1\) differences in \(d_0\) leaves the first factor equal to \(2^{r-1}\). Applying \(r'+s'\) differences in \(d_1\) leaves the second factor linear in \(d_1+d_2\), of slope \(2^{r'}\). The final difference in \(d_2\) gives \(2^{r-1+r'}\). This also covers a constant first factor when \(r=1,s=0\). Combining the lower estimate with \((\ast)\) and dividing by \(2^{r+r'}\) yields \eqref{eq:full36}.
\end{proof}

\begin{corollary}[Strict classes represented in a graded model]\label{cor:model36}
Let \(S\) be an \(\OO\)-module of finite length. Suppose for every \(h\) its subgroup \(S[\pi^h]\) injects additively into augmentation derivations over \(\OO\) on this ambient algebra with values in \(\OO/\pi^h\). For infinite residue field assume these injections are \(\OO\)-linear. Assume the derivations are supported in weight \(-1\), kill the \(F\)'s and lift to \(\OO\)-valued augmentation derivations on the ambient algebra. Then \(\len_{\OO}S\le 2\sum_P v(a(P))\).
\end{corollary}
\begin{proof}
An integral lift kills the saturation of \(\mathfrak b^2\), by the product rule and torsion-freeness of \(\OO\). Thus the injection factors through \(\operatorname{Hom}_{\OO}(C_F,\OO/\pi^h)\). Its length is at most \(\len C_F\) by elementary divisors. For a finite residue field, take cardinalities as in Proposition~\ref{prop:sat31}. With a linear injection the same inequality follows directly from additivity of length. Take \(h\) above the exponent of \(S\) and apply Theorem~\ref{thm:full36}.
\end{proof}

For the fields of Proposition \ref{prop:minors33}, this supplies the commutative algebra implication once one constructs the arithmetic presentation, proves its chart and point hypotheses, and realizes the strict classes as such liftable derivations. The size of the orbit sum must then be bounded using the specified arithmetic trace. The separate estimate on the theta periods requires its own argument on those same intervals.

\lead{Locations in the accompanying source} The general algebraic statement occurs in \srcloc{Development/AutomorphicArithmetic/AutomorphicVarieties/CommutativeAlgebra/OrbitVanishingLocal}{111}. The filtration and the integer measuring its monomials occur in \srcloc{Development/AutomorphicArithmetic/AutomorphicVarieties/CommutativeAlgebra/GradedModulesLength/ToDegInjective}{222}. The comparison of the two bounds appears in \srcloc{Development/AutomorphicArithmetic/AutomorphicVarieties/CommutativeAlgebra/GradedModulesLength/AlgMemCoef}{92}.

\section{The ordinary locus and its tangent spaces}\label{nav:g-IV-2}

\subsection{The marked Galois representation at large primes}\label{nav:IV-09}
\label{sec:marked39}

We realize the strict classes in an integral deformation algebra on the actual power intervals. In this section \(E=E_x\), \(N=2^{l+2}Fqd\), \(t=\lfloor x^{401}\rfloor\), and
\[
       \ell>\exp(x^\eps/12),\qquad 1\le r<t .
\]
For large \(x\) this implies \(\ell>2t+3\), \(\ell>N\), and \(\ell\nmid [E:\QQ]\), by the bounds in Proposition \ref{prop:minors33}. Let \(\OO\) be the ring of integers in a finite unramified extension of \(\QQ_\ell\) containing the character values of \(\Gal(E/\QQ)\). Write \(J\) for its fraction field. Take a character \(\alpha\) with \(\alpha(c)=(-1)^r\), where \(c\) is a chosen complex conjugation, and put
\[
           \Psi=\kappa_\ell^r\alpha,\qquad
           \vartheta=\Psi^{-1},\qquad K'=\QQ(i).
\]
Here \(K'\) is fixed even as the prime-count field \(E\) varies. Let \(\eta\) be its quadratic character. Fix a decomposition subgroup \(G_\ell\subset G_\QQ\) and its inertia subgroup. Inertia acts trivially on \(K'\) and on \(\alpha\). Its mod-\(\ell\) cyclotomic image is \(\mathbf F_\ell^\times\), since \(\QQ_\ell(\mu_\ell)/\QQ_\ell\) is totally ramified of degree \(\ell-1\). Choose \(g_*\) in inertia mapping to a generator. Then \(\overline\Psi(g_*)^2\ne1\) for every \(r,\alpha\) under consideration. The same element works throughout that range at a fixed \(\ell\).

Put \(d(g)=\operatorname{diag}(\Psi(g),1,\Psi(g)^{-1})\). Use the group \(\GL_3\rtimes\{1,-1\}\) of Section~\ref{sec:deform31}, with involution \(M\mapsto J_3 M^{-\mathsf T} J_3\). The base representation has value \((d(g),\eta(g))\). A lift keeps second component \(\eta\) and has matrix \(M(g)\) reducing to \(\overline d(g)\). We impose the following conditions on the matrices themselves. We fix \(M(c)=1\) and \(M(g)=d(g)\) for every \(g\in G_\ell\). For each \(p\ne\ell\), choose inertia \(I_p\) and kill the subgroup \(I_p\cap\ker\Psi\cap\ker\eta\) in the lifted representation. The kernel condition is unchanged by conjugating the choice of inertia. These equations impose a specific basis on the local representation at \(\ell\). Independent conjugation of the global lift is not taken in defining the functor.

\begin{theorem}[A graded integral algebra for the marked lifts]\label{thm:ring39}
There are an augmented Noetherian \(\ZZ\)-graded \(\OO\)-algebra \(A\), generated in weights \(-2,-1,1,2\) over its complete Noetherian weight-zero ring, and a homogeneous ideal \(I\) of \(A\), with the following properties.
\begin{enumerate}\setlength{\itemsep}{1pt}
\item The completion of \(B=A/I\) at the maximal ideal of the residual basepoint represents the marked lifts with the above local conditions to complete Noetherian local \(\OO\)-algebras with the same residue field. Its matrices are continuous for the respective adic topologies.
\item For every \(h\ge1\), the classes in \(e_\alpha(\Sel_{\ell,r}(E)\otimes\OO)[\ell^h]\) are exactly the points of this functor over \(\OO\oplus\epsilon(\OO/\ell^h)\), \(\epsilon^2=0\), of the form
\[
               M_z(g)=(1+\epsilon z(g)(E_{21}-E_{32}))d(g).
\]
These are lifts of the \(\OO\)-valued basepoint under the square-zero projection. Their associated augmentation derivations of \(A\) are supported in weight \(-1\) and each lifts to an \(\OO\)-valued augmentation derivation of \(A\).
\end{enumerate}
\end{theorem}
\begin{proof}
All representations in question factor through the Galois group unramified outside the finite set \(S\) consisting of \(\ell\), the prime divisors of \(2N\), and infinity. Let \(G\) be this group and \(H\) the kernel of its finite residual representation \((\overline d,\eta)\). Quotient \(G\) by the kernel of the maximal pro-\(\ell\) quotient of \(H\), which is characteristic in \(H\). Call the quotient \(\Gamma\). Over an Artinian coefficient ring the principal matrix congruence kernel is a finite \(\ell\)-group, by its filtration by powers of the maximal ideal. Over a complete coefficient ring it is a pro-\(\ell\) group by passage to the inverse limit. Thus all the lifts factor through \(\Gamma\), as does the base representation. The proof of Theorem~\ref{thm:present31} applies to the finite residual field extension here. The mod-\(\ell\) characters unramified outside the primes over \(S\) have finite dimension, by Kummer theory after adjoining \(\mu_\ell\) and finiteness of the class group and the rank of the \(S\)-units. Thus the pro-\(\ell\) subgroup and then its finite extension \(\Gamma\) have finitely many topological generators.

Take generators \(g_1,\ldots,g_n\) and the free profinite group on their lifts. Set \(P=\OO[[X_{s,ij}]]\), giving \(X_{s,ij}\) weight \(k_i-k_j\) for \((k_1,k_2,k_3)=(1,0,-1)\). Assign the generators matrices \(d(g_s)+(X_{s,ij})_{ij}\) and sign \(\eta(g_s)\). The matrices are invertible since their residual determinants are units. The target over \(P\) is profinite, as \(P\) has finite quotients by powers of its maximal ideal. The universal property gives a continuous representation of the free profinite group. Impose matrix equality to one for every relator, along with the conditions at \(c,G_\ell,I_p\) on chosen lifts in the free group. The resulting closed ideal represents the required functor. Indeed any choice of generator matrices gives a unique continuous substitution on \(P\), and the equations are exactly those needed for the continuous homomorphism to descend and have the prescribed specializations.

We construct \(A\) explicitly to retain an ordinary grading. For every \(j\in\ZZ\) let \(P_j\) consist of the formal series all of whose monomials have weight \(j\), and take \(A=\bigoplus_j P_j\subset P\). The ring \(P_0\) is a quotient of a formal power series ring over \(\OO\) on finitely many zero-weight monomials. In fact the minimal nonzero exponent vectors of weight zero form a finite set by the coordinatewise monomial argument of Proposition~\ref{prop:horizontal34}. Every zero-weight monomial factors into these by induction on total degree. Choose one factorization of each monomial. Its length tends to infinity with total degree because the individual factors have bounded degree. This expresses any invariant series as a convergent series in those finitely many monomials. Hence \(P_0\) is Noetherian, complete and local, with residue field that of \(\OO\). For arbitrary \(j\), partition the monomials of weight \(j\) by a chosen minimal divisor in the same weight. Dividing by those finitely many choices leaves invariant series. This exhibits \(A\) as a finite-type algebra over \(P_0\) generated by the entry variables of nonzero weight.

The entries of the matrix of an element of the free profinite group are homogeneous of the indicated entry weights. For words this follows by conjugation with \(\operatorname{diag}(t,1,t^{-1})\), which is fixed by the involution. Equivalently, products, inverses of matrices of homogeneous entries with invariant determinant, and the reflected inverse transpose preserve the entry weights. Inverse invariant determinants here are units of \(P_0\). For general elements use continuity and the closedness of each \(P_j\) in \(P\). Subtracting the specified diagonal constants gives homogeneous equations in \(A\). Define \(I\subset A\) to be their ideal. If \(\mathfrak m\) is generated in \(A\) by \(\ell\) and all entry variables, then \(A/\mathfrak m^b=P/(\ell,X)^b\) for every \(b\). Truncations give surjectivity. For injectivity, split each pure series lying in \((\ell,X)^bP\) into terms divisible by a chosen one of the finitely many monomials \(\ell^{b-i}X^\nu\), \(|\nu|=i\le b\), and divide termwise. The resulting coefficients in each group are still of pure weight and integral. This proves the equality of the ideals on contraction. The completion is thus \(P\); quotient and completion commute for a finite ideal in a Noetherian ring. This proves the first assertion.

By character descent in Theorem \ref{thm:char}, the component in the second assertion is the strict group over \(\QQ\) with coefficients \((J/\OO)(\vartheta)\). Here \(\ell\nmid [E:\QQ]\). The element \(g_*\) acts with \(\vartheta(g_*)-1\) a unit. Thus each \(\ell^h\)-torsion class has a unique finite-coefficient cohomology class lifting it and a unique cocycle representative \(z\) with \(z(g_*)=0\). Locally at \(\ell\) the invariants also vanish. Strictness therefore means that the finite-coefficient restriction is a coboundary. Evaluating it at \(g_*\) forces its defining element to be zero. Consequently \(z|_{G_\ell}=0\). At \(c\) the character is even, so the cocycle identity gives \(2z(c)=0\), hence \(z(c)=0\).

At inertia away from \(\ell\) the image of \(\vartheta\) has order prime to \(\ell\). Its invariants on \((J/\OO)(\vartheta)\) are either zero or the entire divisible module, so a strict divisible class again restricts to a finite-coefficient coboundary. In particular \(z\) vanishes on \(I_p\cap\ker\Psi\cap\ker\eta\). Conversely vanishing there makes the restriction of \(z\) factor through the finite image of \((\Psi,\eta)\) on inertia. This image has order prime to \(\ell\), so averaging makes its cocycle a coboundary. Together with vanishing at \(G_\ell\) this is exactly the strict condition. By Proposition~\ref{prop:root31}, the matrices in the statement obey the group law precisely for a \(\vartheta\)-cocycle. They therefore give the asserted bijection with the marked root-valued points.

Their first-order changes on the generator variables occur only in weight \(-1\). A derivation at the augmentation of \(A\) obtained from these points is the first-derivative functional on formal series at \(X=0\) with the specified changes on the entries. Lift these finitely many values arbitrarily from \(\OO/\ell^h\) to \(\OO\), taking zero on all other entries. The same linear-term functional defines the promised integral lift on the ambient algebra. This lifting statement concerns \(A\); the original residue-valued derivation kills the equations in \(I\).
\end{proof}

\subsubsection{Scaling an algebraic point to an integral lift}\label{nav:IV-10}
Write \(k=\overline J\). The following assertion applies to the algebra \(B\) just constructed whenever its invariant algebra is finite, and also to its further homogeneous quotients with this property and with the basepoint retained.

\begin{proposition}[Integral representatives and their restrictions]\label{prop:point39}
If \(B_0\) is finite over \(\OO\), every \(k\)-valued point of \(B\) has a scaling translate given by a continuous marked lift over the integers of a finite extension of \(J\), reducing to the specified residual basepoint. On the mixed orbits, isomorphism of the associated three-dimensional representations of \(G_{K'}\) identifies exactly one scaling orbit. Each such representation is absolutely irreducible or is the direct sum of a line in the middle coordinate and an absolutely irreducible two-dimensional corner.
\end{proposition}
\begin{proof}
The invariant algebra is a quotient of the local ring \(P_0\) and has the same residue field via the basepoint. Finiteness makes every invariant value integral. Its maximal ideal has a power inside \(\ell B_0\) because the local algebra modulo \(\ell\) is Artinian. All invariant monomials in entry variables of positive total degree therefore have strictly positive valuation or vanish at the point. Write \(z_i\) for the nonzero values of its finitely many nonzero-weight variables, and \(b_i\) for their weights. We can choose a rational number \(s\) with \(v(z_i)+b_i s>0\) for every \(i\). Indeed, for \(b_i>0>b_j\), the invariant monomial \(z_i^{|b_j|}z_j^{b_i}\) has strictly positive valuation, which says precisely that the lower endpoint \(-v(z_i)/b_i\) is less than the upper endpoint \(v(z_j)/|b_j|\). There are finitely many endpoints, and a missing sign imposes only one-sided conditions. Scaling by an element of \(k^\times\) of valuation \(s\) puts all entries in the maximal ideal, including the weight-zero variables already controlled. All values lie in a finite extension since \(B_0\) is finite and only finitely many variables are used.

Substitution into \(P\) now converges. It agrees with the given point on \(P_0\), since invariant series there are series in invariant monomials and the character of \(B_0\) is adically continuous. It agrees on the other generators by construction. The continued map thus kills all the relations of the completed ring and gives a continuous marked lift.

For its restriction \(U(h)=M(h)\) on \(H=G_{K'}\), one has
\[
 U(c h^{-1}c^{-1})=J_3 U(h)^{\mathsf T}J_3,\qquad
 U(g_*)=\operatorname{diag}(\Psi(g_*),1,\Psi(g_*)^{-1}).
\]
The first identity follows from the semidirect product law and the marking at \(c\). The second diagonal has three distinct eigenvalues. Every invariant subspace is therefore spanned by a subset of its coordinate axes, using the spectral projections, which are polynomials in \(U(g_*)\). Form the directed graph of possible entries, with an edge from \(j\) to \(i\) if some \(U(h)_{ij}\ne0\). A coordinate subspace is stable exactly when no edge leaves its subset. The first identity pairs the edges across the antidiagonal with reversed orientation. In particular, each direction along the adjacent axes, if present, occurs on both consecutive pairs. A mixed point has both matrix weight signs: every homogeneous function is a polynomial in the entry variables over invariants, and every element outside \(H\) has matrix equal to that of its product with \(c\). With adjacent edges in both directions the graph is strongly connected. With adjacent edges in only one direction mixedness forces the reverse corner edge, again giving strong connectivity. Either case implies irreducibility. Without adjacent edges the middle line separates, and both corner directions occur. The distinct outer eigenvalues then imply irreducibility on the corner.

If two restrictions are isomorphic, an intertwiner commutes with the same regular diagonal \(U(g_*)\), so is diagonal, say \(\operatorname{diag}(d_1,d_2,d_3)\). Set \(\bar i=4-i\). Applying the intertwining equality to \(h\) and to \(ch^{-1}c^{-1}\) gives \(d_i d_{\bar i}=d_j d_{\bar j}\) whenever the \(ij\) entry is nonzero. If an adjacent edge occurs, \(d_1d_3=d_2^2\), so the intertwiner differs by the scalar \(d_2\) from a scaling matrix. In the corner case choose \(t\in k^\times\) with \(t^2=d_1/d_3\); scaling by \(\operatorname{diag}(t,1,t^{-1})\) has the same effect on all entries. In either case the matrices at \(c\) agree, so all global generator matrices agree up to this scaling. Their invariant monomial values agree. Continuity on \(B_0\) makes all invariant series values agree as well, so the two algebraic points differ by exactly a scaling action. The converse follows by conjugation.
\end{proof}

\begin{proposition}[The arithmetic tangent space modulo scaling]\label{prop:tangent39}
Under the finiteness assumption in Proposition~\ref{prop:point39}, let \(P\) be a mixed point with restriction \(U:G_{K'}\to\GL_3(k)\). The tangent space at \(P\), modulo the one-dimensional scaling direction, injects into the kernel of
\[
 H^1(G_{K'},\operatorname{End}(k^3))\longrightarrow
 \prod_{v\mid\ell}H^1(K'_v,\operatorname{End}(k^3))\ \times
 \prod_{v\nmid\ell\infty}H^1(I_v,\operatorname{End}(k^3)).
\]
The coefficient action is conjugation by \(U\). Continuous cocycles here are taken over finite coefficient extensions and then unioned in \(k\).
\end{proposition}
\begin{proof}
At an integral representative an infinitesimal change \(\dot U\) defines the usual continuous cocycle \(h\mapsto\dot U(h)U(h)^{-1}\) by differentiating multiplication. To justify continuity in this algebraic tangent space, choose a finite coefficient extension containing the values of the point and of the derivation on finitely many invariant-ring module generators and entry variables. An \(\OO\)-linear map on the finite invariant ring is bounded in this extension. Multiply the derivation by a common nonzero scalar to put its values on the generator variables in the ring of integers. Substitution of the variables plus their infinitesimal changes into the completed power series ring then converges in the ring of dual numbers over those integers, since their constant parts are in the maximal ideal. It agrees on the invariant series by continuity there, and hence on all the finite-weight functions. Undo the scalar change. The universal continuous matrices give the asserted continuous cocycle.

The cocycle vanishes on the chosen decomposition group intersected with \(G_{K'}\) by the defining equations. At the other places above \(\ell\) its restriction class also vanishes. Indeed on a conjugate decomposition group the deformation is conjugate to the fixed matrices there by its changed conjugating matrix, with the reflected inverse transpose if its sign is negative. The infinitesimal change of conjugation is a coboundary. Away from \(\ell\) the restrictions to inertia factor through the finite quotient specified by the kernel condition, so their characteristic-zero cohomology class vanishes by averaging.

If the global class vanishes, write \(\dot U=[X,U]\) for a matrix \(X\). Equality at \(g_*\) forces \(X\) to be diagonal, with entries \(x_i\). Differentiate the antidiagonal symmetry in Proposition~\ref{prop:point39}. At every nonzero entry the resulting equation is \(x_i+x_{\bar i}=x_j+x_{\bar j}\). With an adjacent edge this gives \(x_1+x_3=2x_2\), and the scalar part of \(X\) has zero commutator. In the corner case the diagonal matrix \(\tfrac12(x_1-x_3)\operatorname{diag}(1,0,-1)\) induces the same commutators. Thus in both cases the infinitesimal changes on the entries are exactly a multiple of the scaling vector. They determine the derivation on the invariant series by the continuity just proved. Every scaling vector is a coboundary, and the scaling vector is nonzero because a nonzero entry of nonzero weight at \(P\) has nonzero derivative under scaling in characteristic zero. This proves the injection.
\end{proof}

The construction realizes the strict classes for all primes in the displayed large range and gives their integral lifting property on the ambient algebra. Applying Theorem~\ref{thm:full36} additionally requires an ambient presentation with the prescribed tangent and equation counts, finiteness of its solution algebra in weight zero, and control of its fixed and mixed points. Proposition~\ref{prop:tangent39} identifies a precise sufficient arithmetic vanishing for the mixed tangent condition.

\lead{Locations in the accompanying source} For the universal representation compare \srcloc{Development/AutomorphicArithmetic/AutomorphicVarieties/UnitaryThreeGroups/UnivEx}{342}. The two shapes and the isomorphism test occur in \srcloc{Development/AutomorphicArithmetic/AutomorphicVarieties/UnitaryThreeGroups/PtShape}{85}. The fixed local basis at \(\ell\) in this section is available because the chosen marking belongs to inertia and every class being represented is strict there.

\subsection{Excluding the two-dimensional corner}\label{nav:IV-11}
\label{sec:corner40}

Keep the large-prime range and the fixed local matrices of Section~\ref{sec:marked39}. We work first with any continuous marked lift whose entries lie in the integers of a finite extension of \(J\), with residual representation \((\overline d,\eta)\). Let \(U\) denote its restriction to \(H=G_{K'}\), where \(K'=\QQ(i)\). The following exclusion also applies to points that acquire such an integral representative by scaling.

\begin{lemma}[The determinant and the middle line]\label{lem:det40}
One has \(\det U=1\). If the middle coordinate line splits off as a direct summand, its character on \(H\) is trivial.
\end{lemma}
\begin{proof}
Both characters reduce to one on \(H\). At inertia away from \(\ell\) their images are finite of order prime to \(\ell\), by the defining kernel condition. A root of unity of order prime to \(\ell\) reducing to one is one, since its order is a unit and the principal units are pro-\(\ell\). At the chosen place above \(\ell\), the prescribed diagonal matrices have determinant one and middle entry one, on the entire decomposition group in \(H\). The same statements hold at every place above \(\ell\). For conjugation within \(H\) this is immediate. Conjugation by \(c\) replaces a matrix by its reflected inverse transpose, which inverts its determinant and the character on a separated middle line. Hence the characters are everywhere unramified. Global class field theory makes every continuous everywhere unramified character of \(H\) factor through the ordinary ideal class group, since \(K'\) has no real places. That class group is trivial because the Gaussian integers are Euclidean. This proves both assertions.
\end{proof}

\begin{theorem}[There are no integral mixed corner points]\label{thm:nocorner40}
For all sufficiently large \(x\) on the common intervals, every mixed continuous marked lift in the range of Section~\ref{sec:marked39} is absolutely irreducible on \(G_{K'}\). Its restriction cannot be an irreducible representation induced from a character of an index-three subgroup. Moreover its crystalline Frobenius eigenvalues at a place of \(K'\) over \(\ell\) have three different complex absolute values.
\end{theorem}
\begin{proof}
The coordinate and graph argument of Proposition~\ref{prop:point39} applies to the given continuous representation itself. If it is reducible and mixed, write \(U=\beta\oplus\lambda\) on the outer two coordinates and the middle line, respectively. Then \(\beta\) is absolutely irreducible. Lemma~\ref{lem:det40} gives \(\lambda=1\) and \(\det\beta=1\). For \(h\in H\), with the two outer coordinates in their original order,
\[
 \beta(chc^{-1})=J_2\beta(h)^{-\mathsf T}J_2
                 =C_2\beta(h)C_2^{-1},\qquad
 J_2=\begin{pmatrix}0&1\\1&0\end{pmatrix},\quad
 C_2=\begin{pmatrix}1&0\\0&-1\end{pmatrix}.
\]
For the second equality, an inverse transpose in \(\mathrm{SL}_2\) is conjugation by the matrix of the standard alternating form. Multiplying that matrix by \(J_2\) gives \(C_2\) up to sign. Since \(c^2=1\) and \(G_\QQ=H\rtimes\langle c\rangle\), assigning \(c\mapsto C_2\) extends \(\beta\) to a continuous absolutely irreducible odd representation \(\widetilde\beta\) of \(G_\QQ\). In fact its value at \(g\) is the outer block of \(M(g)\) times \(C_2\) when \(\eta(g)=-1\), and just the outer block otherwise. This follows by writing \(g=hc\) in the negative coset and using \(M(c)=1\).

The restriction to \(G_\ell\) is exactly \(\Psi\oplus\eta\Psi^{-1}\). The finite-order factors \(\alpha,\eta\) there are unramified, so these are crystalline characters with distinct Hodge--Tate weights, differing by \(2r>0\). The global representation is unramified outside the finite set from Section~\ref{sec:marked39}. Its residual semisimplification is \(\overline\Psi\oplus\overline\eta\,\overline\Psi^{-1}\). On inertia the ratio is \(\overline\kappa_\ell^{2r}\). Since \(2\le 2r<\ell-2\), this ratio is neither trivial nor either of the first powers \(\overline\kappa_\ell^{\pm1}\). In particular the two residual local characters are distinct. Pan's modularity theorem over \(\QQ\) for odd irreducible two-dimensional de Rham representations of distinct weights applies in this residually reducible case, with \(\ell>5\). It makes \(\widetilde\beta\) a Tate twist of the representation of a classical cuspidal eigenform of weight \(2r+1\), the weight being determined by the difference of the two Hodge--Tate weights. A finite-order global twist, if present in the normalization, can be absorbed into the eigenform. \cite{VI-IV-VI-pan40}

For such a Tate twist, when the representation is crystalline at \(\ell\), its two crystalline Frobenius eigenvalues have the same absolute value under each complex embedding. To specify the classical inputs, local--global compatibility at the coefficient prime identifies the Frobenius-semisimplified Weil--Deligne parameter with the parameter of the local automorphic representation, in the common normalization. Crystallinity gives trivial inertia and zero monodromy, hence the local component is unramified. The two roots at a good prime for a holomorphic cuspidal newform of weight \(k\ge2\) have absolute value \(\ell^{(k-1)/2}\) by Deligne. A common Tate twist and a reversal of Frobenius convention preserve the equality of the two absolute values. \cite{VI-IV-VI-saito40,VI-IV-VI-deligne40}

Here the two crystalline characters have eigenvalues which, in either convention, are roots of unity times \(\ell^r\) and \(\ell^{-r}\), in some order. This follows from their explicit expressions on \(G_\ell\): the unramified factors give roots of unity, and the cyclotomic character contributes \(\ell^{-1}\) to crystalline Frobenius in the usual convention. Their absolute values differ. This contradiction excludes the corner.

For the induction assertion, an induced representation from an index-three character has three one-dimensional summands permuted by \(H\). Since \(\ell\ne2\), \(K'\) is unramified at \(\ell\). The inertia inside \(H\) contains an element \(h\) with \(\Psi(h)\) of infinite order, since its cyclotomic image contains \(1+\ell\ZZ_\ell\). Then \(U(h^6)\) fixes all three summands and is diagonal with three distinct eigenvalues. The summands are the three coordinate axes. But every \(U(g)\) in the integral representative reduces to a diagonal matrix with three unit entries. A matrix permuting the axes nontrivially has a zero diagonal entry. Thus each individual line is preserved, contradicting irreducibility. This excludes in particular all imprimitive irreducible representations in this dimension.

Finally the restriction of \(U\) to a decomposition group in \(H\) at the chosen place \(v\mid\ell\) has characters \(\Psi,1,\Psi^{-1}\). For residue degree \(f_v\), the linearized crystalline Frobenius \(\varphi^{f_v}\) therefore has absolute values \(\ell^{f_vr},1,\ell^{-f_vr}\). This calculation requires no automorphy assumption in dimension three.
\end{proof}

If the invariant algebra in Proposition~\ref{prop:point39} is finite, Theorem~\ref{thm:nocorner40} applies to every geometric mixed point of the algebra of Section~\ref{sec:marked39}. In particular the contribution of its reducible mixed points to the orbit sum is zero throughout this prime range, for every choice of separating trace. The last calculation of the theorem also excludes any irreducible point with a pure crystalline Frobenius parameter at \(v\). Establishing purity, or an alternative bound for all the remaining irreducible points, would require additional arithmetic input. The fully fixed local representation used in this section is stronger than requiring crystallinity with prescribed Hodge--Tate weights. The cotangent inequality of Section~\ref{sec:lower36} also needs the tangent and relation counts for the chosen local conditions. The identification of torsion-valued first-order points alone supplies neither those counts nor a cohomological length bound.

\lead{Locations in the accompanying source} The alternative over the imaginary quadratic field is stated in \srcloc{Development/AutomorphicArithmetic/AutomorphicVarieties/UnitaryThreeGroups/PtShape}{85}, and the passage from a corner to modular representations is treated in \srcloc{Development/AutomorphicArithmetic/AutomorphicVarieties/UnitaryThreeGroups/AdmissibleOfTypeII}{103}. The exclusion above uses the exact local matrices of Theorem~\ref{thm:ring39}.

\subsection{The algebraic group of an irreducible mixed lift}\label{nav:IV-12}
\label{sec:monodromy44}

Keep the fixed-matrix large-prime model of Section~\ref{sec:marked39}. Write \(H=G_{\QQ(i)}\), let \(c\) be its marked complex conjugation, and take an integral mixed lift as in Theorem~\ref{thm:nocorner40}. We work in its marked coordinates. Put \(T(t)=\operatorname{diag}(t,1,t^{-1})\) and let \(J_3\) be the antidiagonal matrix with entries one. Here the algebraic monodromy group means the Zariski closure of the image over the algebraic closure of the coefficient field.

\begin{theorem}[The dense case and the even two-dimensional case]\label{thm:mon44}
For such a lift \(U\) on \(H\), exactly one of the following holds.
\begin{enumerate}\setlength{\itemsep}{1pt}
\item The algebraic monodromy group of \(U\) is \(\mathrm{SL}_3\).
\item Its algebraic monodromy group is \(\mathrm{SO}(J_3)\) in the marked coordinates. There is a continuous absolutely irreducible two-dimensional representation \(\rho\) of \(G_\QQ\), over the integers of a finite coefficient extension, with
\[
 \det\rho=\Psi,\qquad \rho(c)=I_2,\qquad
 \bar\rho=\operatorname{diag}(\bar\Psi,1),\qquad
 \rho|_{G_\ell}=\operatorname{diag}(\Psi,1).
\]
It kills \(I_p\cap\ker\Psi\cap\ker\eta\) at every \(p\ne\ell\). In the explicit basis below, \(U\) is the restriction of its conjugation action on trace-zero matrices.
\end{enumerate}
Conversely, a representation with the properties in the second alternative whose trace-zero conjugation action is absolutely irreducible on \(H\) gives an integral mixed point of the fixed local model. In particular the parity in the second alternative is forced.
\end{theorem}
\begin{proof}
By Lemma~\ref{lem:det40}, the closure \(G\) is contained in \(\mathrm{SL}_3\). By Theorem~\ref{thm:nocorner40} it acts irreducibly. Its identity component \(G^0\) contains the torus of matrices \(T(t)\): the inertia image at \(\ell\) contains infinitely many such matrices, hence is Zariski dense in the one-dimensional torus. The unipotent radical of \(G^0\) is normal in \(G\). It has a nonzero fixed vector by triangularization; the fixed subspace is \(G\)-stable, hence all of the representation. Thus this radical is trivial. In characteristic zero the representation of the connected reductive group \(G^0\) is semisimple. Its isotypic components are permuted transitively by \(G\), since the sum over any orbit is invariant. More than one component would give three permuted lines in dimension three. The infinite inertia element and reduction argument in Theorem~\ref{thm:nocorner40} excludes even a permuted system of three lines. If there is just one isotypic component of one-dimensional type, \(G^0\) consists of scalars, contrary to its torus. Thus \(G^0\) acts irreducibly. Its connected center acts by scalars and is trivial in \(\mathrm{SL}_3\), so it is semisimple.

Its rank is one or two. In rank two, after a change of basis, its maximal torus is the full determinant-one diagonal torus. Each root space of the ambient Lie algebra is one-dimensional and has a distinct torus character. The Lie algebra of \(G^0\) consists of the diagonal algebra and a collection of these root spaces. The roots of a connected semisimple group occur with their negatives, so the presence of \(E_{ij}\) entails that of \(E_{ji}\). The graph of the pairs present must be connected: otherwise a sum of coordinate lines is invariant under the Lie algebra and therefore the connected group in characteristic zero. A connected graph on three vertices gives all the root spaces by taking brackets. Thus \(G^0=\mathrm{SL}_3\). In rank one the semisimple Lie algebra is \(\mathfrak{sl}_2\); its irreducible representation of dimension three is the symmetric square of the standard representation. These elementary characteristic-zero structure facts equivalently follow from the rank-one root system and the highest-weight strings for \(\mathfrak{sl}_2\). The symmetric square preserves the nondegenerate discriminant form on binary quadratics. Schur's lemma makes its invariant symmetric form unique up to scalar. The connected image is its special orthogonal group, since both have dimension three. It remains to determine the form and how the whole image acts on it.

Write \(S\) for a matrix of this form in the marked coordinates. Its line is preserved by \(G\), which normalizes \(G^0\). Each matrix \(U(h)\) multiplies the form by a scalar \(\nu(h)\) with \(\nu(h)^3=1\), by taking determinants. The torus \(T\subset G^0\) preserves \(S\), so it has only the entries \(S_{13}=S_{31}=u\) and \(S_{22}=v\), with \(uv\ne0\). After finite coefficient extension and common scaling, \(S\) is integral with at least one unit entry. Every \(U(h)\) reduces to \(T(\bar\Psi(h))\), which preserves the reduction of the form. Thus \(\nu(h)\) reduces to one. Since \(\ell>3\), the reduction of cube roots of unity is injective, and \(\nu(h)=1\). The group \(G\) is equal to \(G^0=\mathrm{SO}(S)\).

The marked polarization now gives, for \(h\in H\),
\[
  U(chc^{-1})=J_3 U(h)^{-\mathsf T}J_3
             =(J_3S)U(h)(J_3S)^{-1}.
\]
Because \(c^2=1\), the square of \(J_3S=\operatorname{diag}(u,v,u)\) centralizes an absolutely irreducible representation. Thus \(v/u=\delta\in\{1,-1\}\). Divide the form by \(u\), so its matrix is \(S_\delta=J_3\operatorname{diag}(1,\delta,1)\). Put \(T_\delta=T(\delta)=\operatorname{diag}(\delta,1,\delta)\). It differs from \(J_3S_\delta\) by a scalar, belongs to \(\mathrm{SO}(S_\delta)\), and has square one. The semidirect decomposition \(G_\QQ=H\rtimes\langle c\rangle\) therefore extends \(U\) to \(W:G_\QQ\to\mathrm{SO}(S_\delta)\) by \(W(c)=T_\delta\). If \(g\notin H\), this extension has matrix \(M(g)T_\delta\), because the original marking was \(M(c)=I_3\).

Let \(\phi=\Psi\) for \(\delta=1\) and \(\phi=\eta\Psi\) for \(\delta=-1\). We have \(\bar W(g)=T(\bar\phi(g))\) globally and \(W(g)=T(\phi(g))\) on \(G_\ell\). We lift \(W\) to a two-dimensional representation, keeping both equalities in these bases. Map the basis vectors for \(S_\delta\) to
\[
  E_{12},\qquad \operatorname{diag}(1,-1),\qquad (2/\delta)E_{21}
\]
in \(\mathfrak{sl}_2\). The trace pairing \((X,Y)\mapsto\operatorname{tr}(XY)\) in this basis is \((2/\delta)S_\delta\), and the basis is integral since \(\ell>2\). A determinant-one isometry of this pairing preserves the alternating volume form \(\operatorname{tr}([X,Y]Z)\). Nondegeneracy of the pairing shows it preserves the bracket. It extends by fixing the identity to an algebra automorphism of \(M_2\), using \(XY=\tfrac12\operatorname{tr}(XY)I+\tfrac12[X,Y]\) for trace-zero \(X,Y\). Every automorphism of the full matrix algebra over a field is conjugation by an invertible matrix, unique up to scalar. This can also be seen by transporting the two matrix-unit idempotents and then choosing generators of their images so that the transported off-diagonal units act as the usual ones.

Use a finite coefficient field containing the entries of the model and the chosen form, with integer ring \(R\). As \(W(g)\) preserves the trace-zero lattice and its inverse does too, the conjugating matrix preserves \(M_2(R)=RI\oplus\mathfrak{sl}_2(R)\) by conjugation. Its elementary divisors must have equal valuation: in a diagonal elementary-divisor basis the two off-diagonal matrix units acquire opposite valuations. Thus the conjugating matrix has a representative in \(\mathrm{GL}_2(R)\). Its residual class is that of \(\operatorname{diag}(\bar\phi(g),1)\), by the formula for \(\bar W\) and the scalar centralizer of the full matrix algebra. Scale to a representative \(A\) with exactly this reduction. The scalar \(\phi(g)/\det A\) is a principal unit, so it has a unique square root congruent to one since two is invertible in \(R\). Multiplication of \(A\) by that root gives the unique representative \(\rho(g)\) of determinant \(\phi(g)\) and the specified reduction. Uniqueness makes this assignment multiplicative. It is continuous: the compact set of integral matrices with the specified diagonal residual shape and second diagonal residue one maps continuously and injectively under \(A\mapsto(\operatorname{Ad}^0 A,\det A)\) to a Hausdorff space, hence has continuous inverse on its image. Injectivity uses that a scalar of square one and residue one is one. Consequently \(\rho|_{G_\ell}=\operatorname{diag}(\phi,1)\) exactly, \(\rho(c)=\operatorname{diag}(\delta,1)\), and \(\rho\) kills each of the inertia kernels in the statement, where both \(W\) and \(\phi\) are trivial. Its restriction to \(H\) is absolutely irreducible: a preserved line in two dimensions gives a preserved line of nilpotent trace-zero endomorphisms with image in that line and kernel containing it.

If \(\delta=-1\), this is an odd two-dimensional representation of \(G_\QQ\). It is finitely ramified and is crystalline at \(\ell\) with distinct Hodge--Tate weights differing by \(r\ge1\). The regular odd Fontaine--Mazur theorem over \(\QQ\) for \(\ell\ge5\), including the residually reducible case, makes it a Tate twist of a classical cuspidal eigenform. Pan's Theorem 1.0.4 applies here also for \(r=1\). \cite{VI-IV-VI-pan40} Local--global compatibility and purity as in Theorem~\ref{thm:nocorner40} then give equal complex absolute values for the two crystalline Frobenius eigenvalues. The explicit characters \(\phi,1\) at \(\ell\) have absolute values \(\ell^{\pm r}\) and one, since all finite-order factors of \(\phi\) are unramified there. This excludes \(\delta=-1\). For \(\delta=1\), \(S_\delta=J_3\), \(\phi=\Psi\) and \(\rho(c)=I_2\), giving the second alternative.

Conversely, in the same basis with \(\delta=1\) let \(W=\operatorname{Ad}^0\rho\). Since \(W\) takes values in \(\mathrm{SO}(J_3)\), the inverse-transpose involution with antidiagonal \(J_3\) acts trivially on its image. Thus \(g\mapsto (W(g),\eta(g))\) is a homomorphism to the marked semidirect product. It has \(M(c)=I_3\), the correct residual matrices, the fixed local diagonal matrices and the required triviality on inertia kernels. Irreducibility on \(H\) forces nonzero entries of both weight signs, since otherwise its matrices would all be upper or all lower triangular. It is therefore mixed. It cannot have monodromy \(\mathrm{SL}_3\), completing the converse.
\end{proof}

The non-dense case consequently asks for control of an even, residually split, regular two-dimensional representation with fully split local matrices. Odd modularity supplies no contradiction for this even representation. The dense case is a separate three-dimensional problem. Theorem~\ref{thm:mon44} classifies the possible characteristic-zero images of the fixed local model; the integral relations and their weighted lengths in Section~\ref{sec:lower36} still require an estimate. These assertions use the stricter fixed local matrices throughout.

\subsection{Strict classes on the ordinary component at the coefficient prime}\label{nav:IV-13}
\label{sec:original68}

Use the large primes and character summands of
Section~\ref{sec:marked39}. Thus \(\ell>\exp(x^\varepsilon/12)\),
\(1\le r<\lfloor x^{401}\rfloor\), and \(x\) lies in a sufficiently
distant common interval. Write \(\Psi=\kappa_\ell^r\alpha\) and let
\(\eta\) be the quadratic character of the auxiliary imaginary field
\(K'\) of that section. Both finite characters are unramified at
\(\ell\), their orders are prime to \(\ell\), \(\Psi(c)=1\) and
\(\eta(c)=-1\) at complex conjugation. Let \(\OO\) be the
coefficient DVR, \(J\) its fraction field, \(\varpi\) a uniformizer,
and \(D=J/\OO\). In the notation of this section \(d(g)\) denotes
\(\operatorname{diag}(\Psi(g),1,\Psi(g)^{-1})\).
Our semidirect product has involution
\(\theta(A)=J_3 A^{-\mathsf T}J_3\), where \(J_3\) reverses the
three coordinates. The base representation is \(r_0(g)=(d(g),\eta(g))\).
Choose \(g_*\) in inertia at \(\ell\) as in Section~\ref{sec:marked39}.
The three residual diagonal entries of \(d(g_*)\) are distinct.

We describe the local ideal used in the source. Index rows and columns
by \(0,1,2\), adding indices on the list of characters modulo three.
Write \(\chi_0=\Psi,\ \chi_1=1,\ \chi_2=\Psi^{-1}\).
Let \(U\) be the complete local \(\OO\)-algebra representing framed
lifts of the local residual representation with prescribed sign
\(\eta\), and \(x_0:U\to\OO\) the point \(r_0\).
The framed construction in the proof of Theorem~\ref{thm:ring39}
also gives this ring for the local group. The only finiteness
needed there is finite generation of the maximal pro-\(\ell\)
quotient of the kernel of the finite residual image. Locally its
characters of order \(\ell\) form a finite group by local class
field theory over the finite extension cut out by the residual
representation. The pro-\(\ell\) generator criterion gives the
finiteness. Matrix entries of a finite set of generators then
give a complete Noetherian power series presentation as in that proof.

Let \(M(g)\) be the matrix coordinate of the universal lift. For
inertia elements \(s\) write
\(S_a(s)=M(s)-\chi_a(s)I_3\). The ordinary ideal \(I_{\rm ord}\)
is generated by the coefficients of
\(\det(X-M(s))-\prod_a(X-\chi_a(s))\) and all entries of
\begin{equation}
       S_0(s_0')S_1(s_1')S_2(s_2'),\qquad s_i'\in I_\ell .
                                                        \label{eq:word68}
\end{equation}
In the source the inertia elements are also required to lie in the
kernels of \(\eta\) and of the finite inertial-type character.
For these large primes the latter character can be taken unramified,
and both kernels contain \(I_\ell\).

Choose \(s_0\in I_\ell\) of infinite cyclotomic image and an
element \(\gamma\) in the two kernels that induces a positive
power of arithmetic Frobenius on the residue closure. Such choices
exist because the cyclotomic image on inertia is open and the
two finite characters are unramified. Set
\begin{align*}
 c_a&=(\chi_a(s_0)-\chi_{a+1}(s_0))(\chi_a(s_0)-\chi_{a+2}(s_0)),\\
 t_a&=\operatorname{tr}\big(S_{a+1}(s_0)S_{a+2}(s_0)M(\gamma)\big),\\
 q_a&=c_{a+1}t_a-\kappa_\ell(\gamma)c_a t_{a+1},
                  \qquad \Delta=\ell q_0 q_1 .
\end{align*}
Saturate \(I_{\rm ord}\) by the powers of \(\Delta\), take the
radical, and denote the result by \(I_{\rm cr}\). The local
ideal \(I_t\) is the intersection of the minimal primes of
\(I_{\rm cr}\) contained in \(\ker x_0\). The framed point
\(x_0\) is unique with its specified matrices by representability.
These are the saturation and component choices of
\srcloc{Development/AutomorphicArithmetic/AutomorphicVarieties/PadicHodgeTheory/IsInertiaInfiniteCyclotomic}{93}
and lines 112--135 of that file.

\begin{lemma}[A family on a component through the marked local point]\label{lem:component68}
The set of primes defining \(I_t\) is nonempty. The local point
\(r_0\), its unramified diagonal variation
\(\Psi\mapsto\Psi\xi\) through the identity, and the two
lower simple-root conjugations of that variation all lie on a
common component through \(x_0\).
\end{lemma}
\begin{proof}
Let \(n(g)\in\ZZ_\ell\) be the unramified exponent given by
the projection \(\widehat\ZZ\to\ZZ_\ell\) sending arithmetic
Frobenius to one. Over \(B=\OO[[T,u,v]]\) put
\(\xi(g)=(1+T)^{n(g)}\), by the binomial expansion.
It gives a continuous character: its coefficients are integral
and continuous in \(n(g)\), and the character identity follows
from integer exponents by density. Let \(E_{ab}\) denote the
matrix units and put
\[
 X_-=E_{10}-E_{21},\qquad X_+=E_{10}+E_{21},\qquad
 H=I_3+uX_-+vX_+ .
\]
Replace \(\Psi\) by \(\Psi\xi\) in \(r_0\), then conjugate the
result in the semidirect product by \((H,1)\). The diagonal
formula defines a representation because \(\theta\) fixes
\(\operatorname{diag}(a,1,a^{-1})\). Its specialization at
\(T=u=v=0\) is \(r_0\). This determines \(U\to B\).
On inertia and at \(\gamma\) the sign is positive, so all the
ordinary equations and traces are invariant under this conjugation.
On inertia the diagonal is unchanged. Each diagonal entry of
\eqref{eq:word68} has one zero factor. Hence \(I_{\rm ord}\)
maps to zero.

For the unconjugated family let \(d_a(T)\) be its \(a\)-th
diagonal entry at \(\gamma\). The two shifted factors of
\(S(s_0)\) in the definition of \(t_a\) have product equal
to \(c_a E_{aa}\). Thus
\[
      q_a=c_a c_{a+1}
                 (d_a(T)-\kappa_\ell(\gamma)d_{a+1}(T)).
\]
Each \(c_a\ne0\) because \(s_0\) has infinite cyclotomic image.
For \(a=0,1\), the ratio of the two diagonal entries is
\(\Psi(\gamma)(1+T)^{n(\gamma)}\), a nonconstant series
since \(n(\gamma)\) is a positive integer. Thus \(\Delta\)
has nonzero image in the domain \(B\), including when its value
at \(x_0\) is zero. The prime kernel of \(U\to B\) contains the
saturated radical and is contained in \(\ker x_0\).
A minimal prime of \(I_{\rm cr}\) contained in that kernel
therefore occurs in \(I_t\). The map to \(B\) factors through
that component, proving the assertions. In particular
\(I_{\rm ord}\subseteq I_t\subseteq\ker x_0\).
\end{proof}

Let \(W\) be a finite \(\OO\)-module. Use \(\OO\oplus W\) for
the augmented ring in which the ideal \(W\) has square zero.
A framed deformation of weight minus one has matrices
\[
 M(g)=(I_3+u(g)X_-+v(g)X_+)d(g),\qquad u(g),v(g)\in W .
\]
Set \(\vartheta=\Psi^{-1}\). The involution on the Lie algebra
acts by \(-J_3 X^{\mathsf T}J_3\), with eigenvalues \(+1\)
on \(X_-\) and \(-1\) on \(X_+\). Conjugating by \(d(g)\)
multiplies these two matrix directions by \(\vartheta(g)\).
Consequently the representation condition is exactly that
\(u\) and \(v\) be continuous cocycles for \(W(\vartheta)\)
and \(W(\eta\vartheta)\), respectively. We use the convention
\(z(gh)=z(g)+g z(h)\).

\begin{theorem}[The local component detects strictness on the two root lines]\label{thm:word68}
Such a local weight-minus-one deformation factors through
\(U/I_t\) if and only if both cocycles are coboundaries over
\(G_{\QQ_\ell}\). This holds over every finite module \(W\)
at every indicated prime of the same family.
\end{theorem}
\begin{proof}
If the cocycles are local coboundaries, the representation is a
conjugate of \(r_0\) by \((I_3+uX_-+vX_+,1)\) for some
\(u,v\in W\). Indeed conjugation by \(I_3+Y\) contributes the
cocycle \(Y-gY\), and its sign can be changed by negating \(Y\).
Specialize the family of Lemma~\ref{lem:component68} to
\(T=0\) and these two values in \(W\). This proves the required
factorization integrally.

Conversely every such factorization satisfies \eqref{eq:word68}.
Write \(\lambda=\Psi(g_*)\); the differences of
\(\lambda,1,\lambda^{-1}\) are units in \(\OO\).
On both lines the action of \(g_*\) is multiplication by
\(\lambda^{-1}\). By conjugating locally in the two root
directions we can therefore make both cocycles zero at \(g_*\).
The inertia equations still hold, since on inertia conjugation
acts on all their matrix factors by the same ordinary conjugation.

In \eqref{eq:word68} take \(s_0'=s_2'=g_*\) and
\(s_1'=s\), an arbitrary inertia element. The first and last
factors are diagonal. The \((1,0)\) entry of the product is
\[
          (1-\lambda)(\lambda-\lambda^{-1})M(s)_{10}.
\]
The \((2,1)\) entry is
\((\lambda^{-1}-\lambda)(1-\lambda^{-1})M(s)_{21}\).
The scalar coefficients are units. Thus both entries of \(M(s)\)
vanish. Its diagonal characters are units and two is invertible,
so both normalized cocycles vanish on all inertia. For
\(g\in G_{\QQ_\ell}\), compare their values at
\(g_*g=g(g^{-1}g_*g)\). The conjugate inertia element on the
right has value zero. Hence \((\lambda^{-1}-1)u(g)=0\)
and likewise for \(v(g)\). The same unit proves that both
cocycles now vanish on the full local group. Undoing the local
conjugation proves that the original classes were coboundaries.
\end{proof}

\begin{corollary}[The global root cotangent quotient at the marked point]\label{cor:dual68}
Keep the inertia kernels away from \(\ell\) in
Section~\ref{sec:marked39}, use the local ideal \(I_t\) at \(\ell\)
and impose the source's two marks described below. Write
\(C_{-1}\) for the quotient of the augmented cotangent module
that represents its deformations supported in the two lower
simple-root entries. With \(S(\QQ,D(\xi))\) the group of global
classes unramified away from \(\ell\) and strict at \(\ell\),
there is an \(\OO\)-linear isomorphism
\begin{equation}
 \operatorname{Hom}_{\OO}(C_{-1},D)
       \simeq S(\QQ,D(\vartheta))\oplus
                      S(\QQ,D(\eta\vartheta)).           \label{eq:dual68}
\end{equation}
In particular the desired complementary character summand on
\(E_x\), under the large-prime descent of
Theorem~\ref{thm:ring39}, occurs with its full integral length.
\end{corollary}
\begin{proof}
The marks require \(M(c)=I_3\) and the \((0,1),(1,0)\) entries of
\(M(g_*)-M(cg_*^{-1}c)\) to vanish. These are the equations in
\srcloc{Development/AutomorphicArithmetic/AutomorphicVarieties/UnitaryThreeGroups/UnitaryThreeReflectionRight}{281}.
At \(c\) the coefficient actions are \(+1\) on \(u\) and \(-1\)
on \(v\). Thus \(u(c)=0\) already, and a unique odd-line
coboundary adjustment sets \(v(c)=0\), since two is a unit.
With this mark satisfied the second matrix at \(g_*\) in the
difference is \(J_3 M(g_*)^{\mathsf T}J_3\). The upper
entry vanishes automatically and the lower equation is
\[
              (\lambda+1)u(g_*)+(\lambda-1)v(g_*)=0 .
\]
It determines the remaining even-line coboundary adjustment
uniquely: its coefficient is, up to sign,
\((\lambda+1)(\lambda^{-1}-1)\), a unit. Thus every pair of
global cohomology classes has exactly one marked representative.

At a prime other than \(\ell\), the marked representation kills
the prescribed inertia kernel if and only if the restrictions of
both classes are coboundaries on inertia. Indeed the kernel is
the simultaneous kernel of \(\eta,\Psi\) on inertia and the
quotient has order prime to \(\ell\). The two characters act
trivially on the kernel, and averaging kills the first cohomology
of the finite quotient. The converse follows by restricting a
coboundary to the kernel. Theorem~\ref{thm:word68} supplies the
condition at \(\ell\). This proves the corresponding assertion for
every finite coefficient module \(W\).

For precision, form the global framed ring with these conditions
as in Theorem~\ref{thm:ring39}, taking the closed ideal generated
by the image of \(I_t\) from the local ring in place of the fixed
local matrix equations. Impose the two marks above. If
\(\mathfrak a\) is its augmentation ideal, maps to \(\OO\oplus W\)
over the basepoint are \(\OO\)-linear maps
\(\mathfrak a/\mathfrak a^2\to W\). Quotient by the classes of
all entries of \(M(g)d(g)^{-1}-I_3\) outside the two allowed
positions, for all \(g\). This defines \(C_{-1}\) and proves
the claimed representation of root deformations. The module
is finite type over \(\OO\), by Noetherianity.

Take \(W=D[\varpi^m]\) and pass to the union. Globally and
at \(\ell\), the invariants of \(D(\xi)\) vanish for either
\(\xi\) because \(\xi(g_*)-1\) is a unit. On inertia away
from \(\ell\) the invariants are either zero or the whole
divisible group, since the scalar character has order prime
to \(\ell\). The multiplication-by-\(\varpi^m\) cohomology
sequence therefore identifies finite-coefficient strict classes
with the \(\varpi^m\)-torsion of the divisible-coefficient
strict group. Taking these unions gives \eqref{eq:dual68}.
The character descent in Section~\ref{sec:marked39} applies
to \(\vartheta\) on the same \(E_x\) and in the same twist range.
\end{proof}

The equation of the ordinary locus is essential in the converse:
the proof uses a unit multiplying each prohibited inertia entry
over every \(\OO/\varpi^m\). A characteristic-zero tangent-space
dimension by itself would give no such conclusion. The remaining
global estimate requires an upper bound for the length of
\(C_{-1}\) as the field, twist and coefficient prime vary.

\subsection{An integral local chart at the large coefficient primes}\label{nav:IV-14}
\label{sec:chart69}

Continue with \(U,I_{\rm ord},I_t\) and the three diagonal
characters of Section~\ref{sec:original68}. The residue field
is \(k=\OO/\varpi\). On \(\mathfrak n\), the strictly upper
triangular matrices over \(k\), let \(G_{\QQ_\ell}\) act by the
adjoint action of the residual base representation, including its
inverse-transpose involution. Write \(\mathfrak n^*(1)\) for
the dual of this representation tensored with the mod-\(\ell\)
cyclotomic character.

\begin{lemma}[The ordinary inertia word and its flag]\label{lem:flag69}
Over any finite Artinian local \(\OO\)-algebra with residue
field \(k\), the equations of \(I_{\rm ord}\) are equivalent to
the existence of a full flag lifting the standard one, stable on
inertia with ordered graded characters \(\chi_0,\chi_1,\chi_2\).
This ordered flag is unique. After the unique lower unipotent
change of flag congruent to the identity, every local matrix is
upper triangular in the semidirect representation.
\end{lemma}
\begin{proof}
Let \(A\) be the Artinian ring. At \(g_*\) the characteristic
polynomial factors with the three roots \(\lambda_a=\chi_a(g_*)\)
having pairwise unit differences. Cayley--Hamilton and the
Chinese remainder theorem split \(A^3\) into the eigensummands:
their projectors are \(\prod_{b\ne a}(M(g_*)-\lambda_b I)/
\prod_{b\ne a}(\lambda_a-\lambda_b)\). The summands are free
lines by reduction and Nakayama's lemma. Choose on each a
generator reducing to its standard basis vector. In that basis
use \eqref{eq:word68} first with its last two elements equal
to \(g_*\). The last two factors multiply to a unit times
\(E_{00}\). Thus the first column of every \(M(s)\) on inertia
is \(\chi_0(s)e_0\). Next take its first two elements to be
\(g_*\). Their product is a unit times \(E_{22}\), so the
last row of \(M(s)\) equals \(\chi_2(s)e_2^{\mathsf T}\).
The trace supplies the middle diagonal entry. This proves
upper triangularity and identifies all three characters.
Conversely, for upper triangular inertia matrices with these
diagonals, the rightmost factor of the word takes the whole space
into the second flag step, the middle factor takes it into the
first, and the left factor kills the first step. The characteristic
polynomials have the required values.

Any ordered flag with those characters must have as its first
step the first eigenline of \(g_*\) and as its second the sum
of the first two eigenlines. This follows by applying the
three polynomial projectors to each invariant flag step, since
the characters on it determine the ranks of its summands.
A flag lifting the standard one is taken to it, or conversely
obtained from it, by a unique lower unipotent matrix reducing
to the identity. Indeed normalize a generator of the first step
to have first coordinate one, then select the unique vector
of the second step with its first two coordinates zero and
one, respectively. These two vectors and \(e_2\) are the
columns of the required lower unipotent matrix.

Conjugate the representation in the semidirect product to put this
flag in standard position. For \(g\) of positive sign ordinary
conjugation intertwines the inertia matrices with those at
\(gsg^{-1}\). For \(g\) of negative sign first apply \(\theta\)
to those at \(s\). This preserves upper triangularity and
reverses and inverts the three diagonal entries, which gives
the same ordered characters. On inertia those characters are
also invariant under \(s\mapsto gsg^{-1}\). Uniqueness of the
ordered flag implies in both cases that \(M(g)\) takes the
standard flag to itself. This is the assertion about all local
matrices.
\end{proof}

\begin{theorem}[Smoothness of the ordinary component over the integers]\label{thm:smooth69}
Suppose \(H^0(G_{\QQ_\ell},\mathfrak n^*(1))=0\). Then
\(I_t=I_{\rm ord}\) and \(U/I_t\) is a formal power series
ring over \(\OO\) in twelve variables. At the point \(r_0\)
one may take integral coordinates whose numbers at the five
weights are
\[
\begin{array}{c|rrrrr}
        j&-2&-1&0&1&2\\ \hline
   \text{number of coordinates}&1&2&3&4&2 .
\end{array}
\]
The weights refer to conjugation by
\(\operatorname{diag}(z,1,z^{-1})\). On the large-prime range
above, the hypothesis holds for every \(2\le r<t\).
For \(r=1\) it holds if and only if both unramified residual
characters \(\bar\alpha\) and \(\bar\eta\bar\alpha\)
are nontrivial on \(G_{\QQ_\ell}\).
\end{theorem}
\begin{proof}
Lemma~\ref{lem:flag69} identifies the ordinary framed functor
with a choice of three lower unipotent flag coordinates in the
maximal ideal and an upper triangular lift whose inertial
diagonal is fixed. Consider first its diagonal. Its ratio to
the base diagonal is a torus cocycle, with the Galois action
given by the sign involution on diagonal matrices. It vanishes
on inertia and so factors through \(\widehat\ZZ\). Its value
at arithmetic Frobenius can be any diagonal matrix congruent
to one. Indeed this value together with the sign specifies an
element of a finite semidirect product over an Artinian
coefficient ring and hence a continuous homomorphism from
\(\widehat\ZZ\) to that product. Thus the diagonal lifts freely
in three parameters, all of weight zero.

Take a small extension \(A'\to A\) of such coefficient rings
with kernel \(I\) killed by the maximal ideal of \(A'\).
Choose a lift of the given diagonal. Lift the remaining entries
of a given upper triangular representation as continuous cochains.
This is possible by finiteness of the coefficient rings, choosing
lifts on the fibers of the old entries and of the new diagonal.
Their defect of multiplication is a continuous two-cocycle with
values in \(\mathfrak n\otimes_k I\): on this square-zero
kernel conjugation is the residual adjoint action, and the
cocycle identity is associativity. Local Tate duality for finite
modules identifies the dual of \(H^2(G_{\QQ_\ell},\mathfrak n)\)
with \(H^0(G_{\QQ_\ell},\mathfrak n^*(1))\), using the
nondegenerate trace pairing of \(k/\mathbf F_\ell\) on
coefficients. By hypothesis the defect is a coboundary.
Correct the lifted cochain by an upper unipotent one-cochain
with that coboundary, with the sign chosen to cancel the defect.
The result is a representation with the chosen diagonal and
in particular with the prescribed inertia flag. The three
flag coordinates lift freely as well. This proves formal
smoothness over \(\OO\). We use the local duality and Euler
characteristic theorems for finite Galois modules over a
local field. \href{https://www.jmilne.org/math/Books/ADTnot.pdf}{Milne,
\emph{Arithmetic Duality Theorems}, Chapter I, Section 2.}

For the tangent dimension, split \(\mathfrak n=W_1\oplus W_2\)
by positive weights. Their dimensions are two and one.
With the diagonal fixed, infinitesimal upper triangular lifts
are all continuous one-cocycles in these spaces, since the
basis is part of the deformation. For \(W=W_j\), local Euler
characteristic over \(\QQ_\ell\) gives
\(
 \dim H^1=\dim W+\dim H^0+\dim H^2
\).
The exact sequence of coboundaries then gives
\(\dim Z^1=2\dim W+\dim H^2=2\dim W\).
All dimensions are over \(k\). These are four variables of
weight one and two of weight two. The unramified diagonal
variables contribute three at weight zero, and the lower
unipotent matrix contributes two at weight minus one and one
at weight minus two. This yields the table.

A formally smooth complete Noetherian local \(\OO\)-algebra with
this residue field and relative tangent dimension twelve is
\(\OO[[z_1,\ldots,z_{12}]]\). One can see the assertion directly.
Choose lifts of a relative cotangent basis to obtain a surjective
map from this power series ring by complete Nakayama. Formal
smoothness lifts an inverse on relative tangent spaces to a map
back into the power series ring: first choose a map to the
quotient by the square of its maximal ideal by the small extension
lifting property and adjust the coordinate-linear coefficients by
a derivation over \(\OO\). Then lift through the successive
powers of the maximal ideal. The composite endomorphism of the
power series ring induces an isomorphism on its full cotangent
space, since it fixes \(\OO\) and is invertible on the relative
part. Successive lifting by degree gives an inverse endomorphism.
The original surjection was consequently injective.

The matrix entries \(M(g)_{ab}-d(g)_{ab}\) span the relative
cotangent space: a map to \(k[\epsilon]/(\epsilon^2)\) that
annihilates them has the base matrices everywhere and is the
basepoint. Each entry has weight \(b-a\) under conjugation.
These weights are distinct as characters on \(k^\times\) for the
large primes in question. Choose the coordinate lifts just used
among these entries, so their values at \(x_0\) vanish and their
weights agree with the table. For any \(z\in1+\ell\ZZ_\ell\)
of infinite order the same transformation holds integrally.
Comparison of coefficients of monomials under this action shows
that a matrix entry of weight \(j\) has only monomials of total
weight \(j\) in this chart. Since the power series ring is a
domain and \(\Delta\) has nonzero image by
Lemma~\ref{lem:component68}, saturation and taking the radical
leave \(I_{\rm ord}\) fixed. Its unique minimal prime belongs to
\(\ker x_0\), giving \(I_t=I_{\rm ord}\).

Finally, on inertia the characters of \(\mathfrak n\) are
\(\bar\kappa_\ell^r\) twice and \(\bar\kappa_\ell^{2r}\)
once. For \(2\le r<t\) and \(\ell>2t+3\) neither exponent
is congruent to one modulo \(\ell-1\). Thus the dual twist has
no inertial invariants. For \(r=1\) its only possible invariants
come from weight one. On the span of \(E_{01},E_{12}\) the sign
involution is minus the swap, and its two characters in
\(\mathfrak n\) are \(\bar\Psi,\bar\eta\bar\Psi\).
After the dual twist the invariants occur precisely when one of
the stated unramified characters is trivial. This proves the
criterion.
\end{proof}

The table is the one prescribed in
\srcloc{Development/AutomorphicArithmetic/AutomorphicVarieties/UnitaryThreeGroups/ConjStable}{22}.
The local chart is used in the construction of the global
presentation in
\srcloc{Development/AutomorphicArithmetic/AutomorphicVarieties/UnitaryThreeGroups/ModelEx}{172}.
Here the integral chart has been established in the stated range
by the inertia equations and local duality. Bounding the relations
in the global representation ring is a further step.

\subsection{Numbers of generators and relations over the local chart}\label{nav:IV-15}
\label{sec:rel70}

Take \(2\le r<t\) and the large primes of
Section~\ref{sec:original68}. Let \(R_\ell=U/I_t\) be the local
ring in Theorem~\ref{thm:smooth69}. First allow all global frames
at the residual basepoint, with the prescribed sign, the away
inertia kernels and this local condition at \(\ell\). Write
\(B\) for their complete local ring. Restriction gives an
\(\OO\)-algebra map \(R_\ell\to B\). We will impose the two
global marks after counting the relative relations.

Let \(S\) consist of the finite primes where the residual data
may ramify together with \(\ell\), and write \(G_S\) for the
Galois group of \(\QQ\) unramified away from those primes and
infinity. For \(q\in S\setminus\{\ell\}\) put
\[
  N_q=I_q\cap\ker\eta\cap\ker\Psi,\qquad H_q=G_{\QQ_q}/N_q .
\]
The inertia quotient in \(H_q\) is finite of order prime to
\(\ell\). The residual adjoint action on all matrices over \(k\)
decomposes into weight spaces \(W_j\) for
\(j=-2,-1,0,1,2\), where the \((a,b)\) entry has weight \(b-a\).
These are \(G_S\)-modules, since the involution
\(X\mapsto-J_3 X^{\mathsf T} J_3\) preserves the weight.

\begin{lemma}[The relative cohomology count]\label{lem:rel70}
Let \(C^\bullet(G,W)\) denote continuous group cochains.
For \(W=W_j\) take the shifted mapping cone \(K_j^\bullet\)
of the map
\begin{align*}
 C^\bullet(G_S,W)\oplus\!
      \bigoplus_{q\in S\setminus\{\ell\}} C^\bullet(H_q,W)
 &\longrightarrow C^\bullet(G_{\QQ_\ell},W)\oplus\!
      \bigoplus_{q\in S\setminus\{\ell\}} C^\bullet(G_{\QQ_q},W),\\
 (z,(y_q))&\longmapsto
       (z|_{G_{\QQ_\ell}},(z|_{G_{\QQ_q}}-\operatorname{inf}y_q)).
\end{align*}
The shift is by minus one, so its long exact sequence has the
cohomology of the domain followed by that of the target.
Set \(a_j=\dim_k H^1(K_j^\bullet)\). Then
\begin{equation}
   H^0(K_j^\bullet)=H^3(K_j^\bullet)=0,\qquad
   \dim_k H^2(K_j^\bullet)=a_j+\mathbf 1_{|j|\le1}.
                                                      \label{eq:rel70}
\end{equation}
Moreover \(H^1(K_j^\bullet)\) is the space of framed
infinitesimal global lifts of weight \(j\) with the local
representation at \(\ell\) held fixed.
\end{lemma}
\begin{proof}
A relative one-cocycle consists of a global cocycle, local
quotient cocycles at the \(q\)'s, and degree-zero cochains
witnessing their prescribed agreements up to coboundary.
Subtract the global coboundary of the witness at \(\ell\)
to make that witness zero. Independently subtract a local
coboundary at each \(q\) to make its witness zero: degree-zero
inflation from \(H_q\) is the identity on \(W\). No residual
change of global representative remains after setting the witness
at \(\ell\) to zero. What remains is a global cocycle vanishing
on \(G_{\QQ_\ell}\), whose restriction to each \(G_{\QQ_q}\)
comes from \(H_q\). The local representatives are uniquely
determined by surjectivity onto \(H_q\). This is exactly the
framed tangent assertion, by the square-zero multiplication law.
In degree zero, the map on global invariants to the \(\ell\)-term
is injective, and inflation on the other invariants is injective.
Thus \(H^0(K_j^\bullet)=0\).

At \(q\ne\ell\) averaging over the finite inertia quotient
reduces the cohomology of \(H_q\) to that of the procyclic
unramified quotient on the inertia invariants. Its degrees zero
and one have equal dimensions and the higher groups vanish.
The local Euler characteristic of \(G_{\QQ_q}\) is also zero.
At \(\ell\) the local Euler characteristic is \(-\dim W\).
For \(G_S\) the global Euler characteristic is
\(-\dim W+\dim W^{c=1}\). Here \(\ell\) is odd,
\(S\) contains the coefficient prime and every finite ramified
place, and all the modules are finite. These hypotheses also
give vanishing above degree two for the local and global
groups. The last part of the Poitou--Tate sequence identifies
the cokernel of global-to-local restriction in degree two with
the dual of \(H^0(G_S,W^*(1))\); modified cohomology at infinity
vanishes. These are the classical Euler characteristic and
duality theorems. \href{https://www.jmilne.org/math/Books/ADTnot.pdf}{Milne,
\emph{Arithmetic Duality Theorems}, Chapter I, Sections 2, 4 and 5.}

Inertia at \(\ell\) acts on \(W_j^*(1)\) by
\(\bar\kappa_\ell^{1-rj}\). In the stated range
\(1-rj\not\equiv0\pmod{\ell-1}\) for all five weights, since
\(r\ge2\) and \(2r<\ell-2\). Thus the dual invariants vanish.
The long exact sequence of the cone gives \(H^3(K_j^\bullet)=0\).
Its Euler characteristic is the Euler characteristic of the
domain minus that of the target, hence \(\dim W_j^{c=1}\).
At \(c\) the matrix is the identity and the action on the Lie
algebra is minus anti-diagonal transposition. It fixes a line
in the three-dimensional diagonal space and a line in each
of the two-dimensional weights \(1,-1\). It acts by minus
one on each weight \(2,-2\). This proves \eqref{eq:rel70}.
\end{proof}

\begin{theorem}[A weighted presentation of the completed global ring]\label{thm:pres70}
Over the integral power series ring \(R_\ell\), the ring \(B\)
admits a presentation with \(a_j\) additional variables of each
weight \(j\). Its ideal of relations has a set of homogeneous
generators with at most \(a_j+\mathbf1_{|j|\le1}\) generators
of weight \(j\), and none of other weights. After imposing the
source's global marks, the resulting complete local ring has
a presentation over \(\OO\) with the following counts:
\[
\begin{array}{c|ccccc}
 j&-2&-1&0&1&2\\ \hline
 \text{variables}&a_{-2}+1&a_{-1}+2&a_0+3&a_1+4&a_2+2\\
 \text{relations, at most}&a_{-2}+1&a_{-1}+3&a_0+3&a_1+3&a_2+1.
\end{array}
\]
Here homogeneity refers to the integer weight, and a homogeneous
power series may have monomials of more than one ordinary degree.
\end{theorem}
\begin{proof}
The framed global ring is Noetherian by the construction of
Theorem~\ref{thm:ring39}. Its group is a quotient of \(G_S\).
For the residual-kernel finiteness used there one may also apply
the finiteness of \(H^1\) with finite coefficients over the finite
extension cut out by the residual matrices. The local condition
at \(\ell\) and the kernels away from it are imposed by closed
ideals. As in the local chart, the differences of the universal
matrix entries from their base values span the relative cotangent
space. Take a basis from these homogeneous entries, by
Lemma~\ref{lem:rel70}, and choose the corresponding relative
variables. Complete Nakayama gives a surjection
\(P=R_\ell[[X_1,\ldots,X_a]]\to B\) with kernel \(I\), where
\(a=\sum_j a_j\). In particular every relation has zero
linear part in the relative variables modulo the residual
maximal ideal. Terms linear in elements of the base ring are
permitted.

We give the relation bound. For a small extension with kernel
a \(k\)-vector space \(L\) killed by the maximal ideal, suppose
we have a prescribed lift at \(\ell\) and a global representation
over the quotient ring. Lift its global matrices as a continuous
cochain, keeping the sign. At each \(q\ne\ell\) choose a lift
on the quotient \(H_q\); this can be taken to be a homomorphism
since its second cohomology with residual adjoint coefficients
vanishes, by the procyclic computation in the lemma. The defect
of global multiplication is an adjoint-valued two-cocycle.
Comparing the restrictions of the lifted cochain with the chosen
local lifts gives local one-cochains whose coboundaries are the
restrictions of that defect, with the same sign convention throughout.
These data give an obstruction in
\[
            \bigoplus_j H^2(K_j^\bullet)\otimes_k L .
\]
Changing the global cochain, local quotient lifts or local
identifications changes the representative by a cone coboundary.
For example, two quotient lifts on \(H_q\) differ by a cocycle
there; its inflation enters as the local component of the
boundary of a degree-one cone element. If the obstruction
vanishes, subtract cochains realizing its boundary. This kills
the global defect and makes its restrictions agree with allowed
local lifts up to infinitesimal conjugation. A global
infinitesimal conjugation makes the agreement at \(\ell\)
exact. At all other primes conjugating a representation still
factoring through \(H_q\) preserves that condition. This
produces a lift over the local base. Conversely such a lift
has zero obstruction by taking its restrictions as the local data.

Apply the construction to \(P/(\mathfrak m_P I)\to B\).
Its kernel \(L=I/\mathfrak m_P I\) is a finite-dimensional
\(k\)-space. The rings here may be infinite; the continuous
cochains used above still exist. The map between the two
compact complete local rings is a surjection with finite kernel.
On a sufficiently deep maximal-ideal power it is injective, and
maps it homeomorphically onto its image, an open ideal. Choosing
lifts on the finitely many cosets gives a continuous section.
The same construction on matrices supplies the cochains, and
the obstruction still takes values in the finite discrete kernel.

Evaluation of the universal obstruction on \(L^*\) is injective.
Indeed if a nonzero functional killed it, push out the small
extension by that functional, or equivalently quotient by its
kernel in \(L\). The representation would then lift, giving an
\(R_\ell\)-algebra section from \(B\) into this pushout. Compare
the resulting map from \(P\) with the quotient map. They agree
on the base and their relative variables differ only by
elements of the square-zero residual kernel. On any \(f\in I\)
their difference is consequently the sum of these differences
multiplied by the residual linear coefficients of \(f\) in
the relative variables. All such coefficients vanish by minimality.
The map through \(B\) kills \(I\), whereas the quotient map on
\(I\) induces the chosen nonzero functional. This contradiction
proves injectivity.

We record why this bound holds weight by weight. Conjugation
by the diagonal torus acts on a universal matrix entry by its
integer weight. The matrix product, the group equations, the
away inertia kernels, and the local ordinary equations are
equivariant under this action. Thus on each finite-degree
truncation of these presentations the split torus acts
algebraically, over \(\OO\), and all maps above respect it.
Equivalently one may conjugate by
\(\operatorname{diag}(z,1,z^{-1})\) with an indeterminate
invertible \(z\) and compare Laurent coefficients modulo each
power of the maximal ideal. This construction gives integral
homogeneous projections that are compatible under truncation.
It applies to arbitrary group equations expressed by adic
limits, since each finite truncation commutes with these limits.
The kernels and their closed ideals are therefore stable under
projection to each weight.

In particular \(I/\mathfrak m_P I\) has a finite weight
decomposition. Pushing the universal obstruction by the torus
on its coefficients has the same effect as conjugating the
matrices in its adjoint factor, by naturality of the cochain
construction. Thus the component in \(H^2(K_j^\bullet)\)
pairs only with coefficients of weight \(j\). To check equality
of integer weights one may use the finite-degree torus action
just described, or pass to finite unramified coefficient
extensions of arbitrarily large residue cardinality and vary
the torus unit. Finite flat extension commutes with the matrix
presentations, their finite-dimensional relation spaces and the
continuous finite-coefficient cochains here. It follows that
\(\dim_k L_j\le\dim_k H^2(K_j^\bullet)\).
For a basis of \(L_j\), lift each basis vector in \(I\) and
take its weight-\(j\) component; closedness and the compatible
projections give a lift in \(I\) of exactly that weight.
Nakayama's lemma for the finite \(P\)-module \(I\) makes these
lifts ideal generators. This proves the relative count.

Finally \(r(c)^2=1\) and the negative sign of \(c\) imply
\(M(c)=J_3M(c)^{\mathsf T}J_3\). Thus the equation
\(M(c)=I_3\) is imposed by its two independent diagonal
differences at \((0,0),(1,1)\), its entries at \((0,1),(1,0)\)
and those at \((0,2),(2,0)\). Their weights are
\(0,0,1,-1,2,-2\). The other mark gives the two entries at
\((0,1),(1,0)\) of \(M(g_*)-M(cg_*^{-1}c)\), of weights
\(1,-1\). Lift these eight equations homogeneously from \(B\)
to \(P\) and use Theorem~\ref{thm:smooth69} for the twelve
base variables. This gives the table.
\end{proof}

The global count is the one sought in
\srcloc{Development/AutomorphicArithmetic/AutomorphicVarieties/UnitaryThreeGroups/IsLocPres}{34};
the eight marking equations are listed in
\srcloc{Development/AutomorphicArithmetic/AutomorphicVarieties/UnitaryThreeGroups/MopBase}{105}.
Theorem~\ref{thm:pres70} gives a presentation on the actual
coefficient rings. The coefficient sizes of its relations are
still needed for the length in \eqref{eq:dual68}. Applying the
intersection argument of Section~\ref{sec:lower36} also requires
control of its weight-zero algebra, characteristic-zero points
and traces. The variable and relation counts alone do not
provide the required bound on the common intervals.

\subsection{The component at the first cyclotomic twist}\label{nav:IV-16}
\label{sec:kummer72}

We treat the case \(r=1\) left open by the vanishing hypothesis in
Theorem~\ref{thm:smooth69}. We keep the same \(E_x\), the same
auxiliary field \(K'=\QQ(i)\), and the same range
\(\ell>\exp(x^\varepsilon/12)>2\lfloor x^{401}\rfloor+3\).
In particular both finite characters \(\alpha,\eta\) are unramified
at \(\ell\). Denote the cyclotomic character by \(\kappa\).

\begin{lemma}[Kummer classes with a variable unramified twist]\label{lem:kummer72}
Let \(K/\QQ_\ell\) be unramified of degree \(f\), let \(I\) be its
inertia group, and let \(\phi\) be arithmetic Frobenius on the
maximal unramified extension \(K^{\rm nr}\). Let \(C\) be a finite
local Artin \(\OO\)-algebra, \(N\) a finite \(C\)-module, and
\(\beta:G_K\to C^\times\) unramified with Frobenius value \(u\).
There is a decomposition natural in these coefficients
\begin{align}
 H^1(G_K,N(\kappa\beta)) &\simeq N[u-1]\oplus U_K(N,u),
                                                        \label{eq:kummer72}\\
 U_K(N,u)&=\ker\!\left(1-u\phi:
    \mathcal O_{K^{\rm nr}}\otimes_{\ZZ_\ell}N\longrightarrow
    \mathcal O_{K^{\rm nr}}\otimes_{\ZZ_\ell}N\right). \notag
\end{align}
Here \(N[u-1]=\{a:(u-1)a=0\}\), and \(\phi\) acts on the first
tensor factor. For a small extension \(C'\to C\) and a lift
\(u'\in C'^\times\) of \(u\), every element of
\(U_K(\mathfrak m_C,u)\) lifts to
\(U_K(\mathfrak m_{C'},u')\). On a one-dimensional residual
coefficient module \(N=k\), the dimension of \(U_K(k,\bar u)\)
over \(k\) is \(f\).
\end{lemma}
\begin{proof}
For the untwisted coefficients \(N(1)\) and a finite unramified
extension \(K_n/K\), the Kummer map identifies
\[
 H^1(G_{K_n},N(1))
       = (K_n^\times)^{\wedge}_{\ell}\otimes_{\ZZ_\ell}N .
\]
For \(N=\ZZ/\ell^b\) this is the Kummer sequence and Hilbert's
Theorem 90. For general finite \(N\), send \(a\otimes n\)
to the image of the compatible Kummer classes of \(a\) under
the coefficient map \(1\mapsto n\). The cyclic case and a
decomposition of the underlying finite abelian group prove this
natural map is an isomorphism. The uniformizer \(\ell\) splits
the valuation. The residue units have order prime to \(\ell\).
The logarithm divided by \(\ell\) identifies the principal units
of \(K_n\) with its additive ring of integers. Indeed the logarithm
and exponential series are inverse on \(1+\ell\mathcal O_{K_n}\)
and \(\ell\mathcal O_{K_n}\) because \(\ell\) is odd and
\(K_n\) is absolutely unramified. These isomorphisms commute
with Frobenius and with inclusions of unramified fields. Taking
the union over \(K_n\) gives
\[
 H^1(I,N(1))=N\oplus
       (\mathcal O_{K^{\rm nr}}\otimes_{\ZZ_\ell}N).
\]
Here we use continuity of Galois cohomology under a union of
fields with finite discrete coefficients. Equivalently, a
Kummer class over the union is represented by an element of a
finite stage, and an equality of Kummer classes also occurs
at a finite stage.

The twist \(\beta\) is trivial on \(I\) and multiplies its
Frobenius action on this decomposition by \(u\).
An inertia invariant in \(N(\kappa\beta)\) is zero, since
an inertial cyclotomic value differs from one by a unit. Inflation--restriction
therefore identifies \(H^1(G_K,N(\kappa\beta))\) with the
Frobenius invariants of the displayed direct sum. This proves
\eqref{eq:kummer72}.

For the lifting assertion write \(J=\ker(C'\to C)\), killed by
\(\mathfrak m_{C'}\). Tensoring the surjection of maximal ideals
with \(\mathcal O_{K^{\rm nr}}\) stays exact, since this ring
is torsion-free over \(\ZZ_\ell\). Lift the proposed unit term
as an element of the larger tensor product. Its error under
\(1-u'\phi\) belongs to
\(\overline{\mathbf F}_\ell\otimes_{\mathbf F_\ell}J\).
On this space \(1-\bar u\phi\) is surjective. Indeed if
\(m\) is the finite order of \(\bar u\), then
\((1-\bar u\phi)(1+\bar u\phi+\cdots+(\bar u\phi)^{m-1})
=1-\phi^m\). The last operator is surjective because the
equation \(X-X^{\ell^{fm}}=a\) has a solution in
\(\overline{\mathbf F}_\ell\) for every \(a\). Correct the
lift by a preimage of its error. Finally local Tate duality
gives \(\dim_k H^2(G_K,k(\bar\kappa\bar\beta))
=\mathbf1_{\bar u=1}\). The local Euler characteristic and
\(H^0=0\) give first cohomology of dimension
\(f+\mathbf1_{\bar u=1}\). Subtract the dimension of the
first summand of \eqref{eq:kummer72}. This proves the dimension
assertion.
\end{proof}

\begin{theorem}[The integral ordinary chart at every first twist]\label{thm:first72}
For \(r=1\), including \(\bar\alpha=1\) or
\(\bar\eta\bar\alpha=1\) on \(G_{\QQ_\ell}\), the ring \(U/I_t\)
is formally smooth over \(\OO\). The counts of its integral
coordinates in weights \(-2,-1,0,1,2\) are \(1,2,3,4,2\).
In combination with Theorem~\ref{thm:smooth69} this gives the
same chart for every \(1\le r<\lfloor x^{401}\rfloor\) on the
indicated large-prime range.
\end{theorem}
\begin{proof}
Write \(G=G_{\QQ_\ell}\) and choose an arithmetic Frobenius
lift \(F\). By Lemma~\ref{lem:flag69} a framed ordinary
deformation is a lower unipotent flag coordinate in three
variables together with an upper triangular deformation.
Its diagonal has parameter ring \(A=\OO[[t_0,t_1,t_2]]\).
The ratio of its diagonal at \(F\) to the base diagonal has
entries \(1+t_0,1+t_1,1+t_2\). This is an unramified
diagonal cocycle with the sign action, as in the proof of
Theorem~\ref{thm:smooth69}. Denote the diagonal matrix at a
positive-sign element \(g\) by \(D(g)\).

Quotient the upper unipotent group by its central entry \(02\).
The two entries \(01,12\) give an additive group in this
quotient. With the diagonal prescribed, their values are
continuous cocycles for its conjugation action, by writing
a matrix as unipotent times diagonal. For a finite coefficient
algebra \(C\) the entries of the cocycles belong to
\(N=\mathfrak m_C\). If \(\eta(F)=1\), they are separate
lines. On line \(a,a+1\), \(a=0,1\), the character is
\(\kappa\beta_a\) with \(\beta_a\) unramified and value
\[
          u_a=\alpha(F)(1+t_a)/(1+t_{a+1})
\]
at \(F\). In this case take two copies of \(K=\QQ_\ell\).
If \(\eta(F)=-1\), the sign exchanges the two lines with
a minus sign. Let \(K\) be the unramified quadratic extension.
Over \(G_K\) the first line has character \(\kappa\beta\);
its unramified factor has value
\[
                       u=\alpha(F^2)(1+t_0)/(1+t_2)
\]
at \(F^2\). Indeed the square of the diagonal at \(F\)
in the semidirect product has entries
\(D_{00}(F)/D_{22}(F),1,D_{22}(F)/D_{00}(F)\), where here
\(D_{ii}(F)\) denotes the ordinary diagonal coordinate of the
negative-sign matrix. Both adjacent ratios at \(F^2\) coincide.
The two-dimensional module is induced from the first line
over \(G_K\). Its standard vector and its image under
\(F\) form a basis since the
second vector is a unit multiple of the other standard vector.
This construction commutes with coefficient changes. Shapiro's
lemma identifies its first cohomology with that of this single
line over \(G_K\). Use one index \(\nu\) for each of the two
lines in the first case, and one index for the induced module
in the second. The sum of the associated degrees \(f_\nu\)
is always two.

Passing to cohomology removes two framed coordinates.
In fact \(g_*\) acts on the adjacent module by the scalar
\(\kappa(g_*)\) and \(\kappa(g_*)-1\) is a unit.
Subtracting a unique coboundary makes any adjacent cocycle
vanish at \(g_*\). Thus the class and two freely chosen
coboundary coordinates recover the cocycle functorially over
every \(C\). Applying Lemma~\ref{lem:kummer72} to these
classes gives valuation coordinates \(v_\nu\in\mathfrak m_C\)
with \((u_\nu-1)v_\nu=0\), and unit terms in its second summand.
Consequently the ordinary ring maps, on formal spectra, to
the base with coordinate ring
\[
      A_{\rm val}=A[[v_\nu:\nu]]/
                         ((u_\nu-1)v_\nu:\nu).
\]
The valuation assignment gives a ring homomorphism in the
opposite direction by passing through the Artinian quotients
of the universal ring. When \(u_\nu-1\) is a unit, the
corresponding variable simply vanishes.

This map of deformation functors is formally smooth.
Given a small extension with a prescribed lift of the
diagonal and the valuation coordinates, lift the unit terms
by Lemma~\ref{lem:kummer72} and the two coboundary coordinates
freely. This lifts the entire upper representation modulo its
central entry. Choose any lift of the old central entry as
a continuous function, which is possible by finiteness of
the coefficient rings. Its defect of multiplication takes
values in the residual central line tensored with the kernel
of the small extension. This is a two-cocycle by associativity.
On inertia the central line has character \(\bar\kappa^2\);
hence its dual twisted by \(1\) has no invariants. Local
Tate duality gives \(H^2=0\), so a central one-cochain
corrects the defect. The flag coordinates lift freely.
Over the residual point of \(A_{\rm val}\), the relative
tangent space has three flag coordinates, two adjacent
coboundaries, \(\sum f_\nu=2\) unit directions, and two
central cocycles, giving dimension nine. For the last number the local Euler formula
over \(\QQ_\ell\) gives \(\dim Z^1=2\dim W+\dim H^2=2\)
on the one-dimensional central line \(W\).
A formally smooth map between these complete Noetherian local
rings with the same residue field and relative tangent dimension
nine identifies the target with a formal power series ring in
nine variables over the source. The power series lifting argument
in Theorem~\ref{thm:smooth69} applies over \(A_{\rm val}\)
as well. Thus, up to a choice of coordinates,
\[
                       U/I_{\rm ord}=A_{\rm val}[[y_1,\ldots,y_9]].
\]

We now compute the saturation in the source. Work in the upper
flag basis. The elements used to define \(q_a\) have positive
sign, so their traces are unchanged by this basis choice.
Products of upper triangular matrices have diagonal equal to
the product of the diagonals. The definitions in Section
\ref{sec:original68} therefore give exactly
\[
       q_a=c_a c_{a+1}
              \big(D_{aa}(\gamma)-\kappa(\gamma)
                                  D_{a+1,a+1}(\gamma)\big).
\]
The positive-sign element \(\gamma\) induces a positive
integer power \(n_\nu\) of arithmetic Frobenius over the
associated unramified field. Its adjacent diagonal ratio
divided by \(\kappa(\gamma)\) equals \(u_\nu^{n_\nu}\).
In the induced case the two ratios are equal. It follows
that \(u_\nu-1\) divides the corresponding \(q_a\).
Thus \(\Delta v_\nu=0\) for every \(\nu\), where
\(\Delta=\ell q_0q_1\). Modulo all the \(v_\nu\), the ring
is a power series domain over \(\OO\). The image of
\(\Delta\) is nonzero there because each \(c_a\ne0\) and the
displayed formulas for the \(u_\nu\) give
\(u_\nu^{n_\nu}-1\ne0\) in the power series domain \(A\).
Hence the elements annihilated by powers of \(\Delta\) in
\(U/I_{\rm ord}\) form precisely the ideal generated by the
\(v_\nu\). Its quotient is reduced and has one minimal prime. The
basepoint factors through it because the diagonal representation
has zero adjacent classes. This identifies the quotient with \(U/I_t\).

The residual tangent directions after saturation have weights
\(-2,-1,-1\) from the flag, three zeros from the diagonal,
weight one for all four adjacent directions, and weight two
for the two central directions. The split torus acts on the
adjacent quotient by a common scalar and on its central line
by its square, so the decomposition of tangent dimensions is
compatible with the action, including in the induced case.
The saturated chart is a power series ring, and homogeneous
entry differences lift a basis of its relative cotangent space
as in Theorem~\ref{thm:smooth69}. Choose those differences as
the integral chart coordinates. This proves the stated counts.
\end{proof}

\begin{corollary}[Global relations for the first twist]\label{cor:first72}
For \(r=1\), the variable counts of Theorem~\ref{thm:pres70}
hold with the same definition of \(a_j\). So do its upper
bounds on relations, except that the weight-one bound is
\[
                 a_1+3+\mathbf1_{\alpha=\eta}.
\]
The equality of characters here is global. For the other twists
the table of Theorem~\ref{thm:pres70} applies without change.
\end{corollary}
\begin{proof}
The relative cone of Lemma~\ref{lem:rel70} still computes the
framed tangent with a fixed local lift. Global duality identifies
its third cohomology with the dual of the global invariants
of \(W_j^*(1)\). Only \(j=1\) can contribute, by inspection
on inertia at \(\ell\). At weight one the two characters in
the dual twist are \(\bar\alpha^{-1}\) and
\(\bar\eta\bar\alpha^{-1}\), since the sign involution is minus
the swap. Now \(\alpha(c)=-1\). Since the finite characters
have order prime to \(\ell\), reduction is injective on their
values. Thus the dimension of the global invariants is one
precisely when \(\alpha=\eta\), and zero otherwise.
The Euler characteristic computation in Lemma~\ref{lem:rel70}
adds this dimension to \(\dim H^2\). The obstruction argument,
the eight marking equations, and the local parameter counts
of Theorem~\ref{thm:first72} give the result.
\end{proof}

The Frobenius factors and saturation are defined in
\srcloc{Development/AutomorphicArithmetic/AutomorphicVarieties/PadicHodgeTheory/IsInertiaInfiniteCyclotomic}{99}
and lines 104--135. The outstanding global length and period
estimates are stated for these fields in
Section \ref{nav:I-28}.

\subsection{The fixed point and its nonpositive cotangent spaces}\label{nav:IV-17}
\label{sec:origin116}

Use the ordinary marked deformation ring of Sections~\ref{sec:original68}
through~\ref{sec:kummer72}. We take its input character \(\alpha\)
from the actual period field \(E=\widetilde E\) of
Section \ref{nav:I-28}.
Write \(n=[E:\QQ]\), let \(N\) be its cyclotomic modulus, and
retain \(t=\lfloor x^{401}\rfloor\), \(1\le r<t\) and
\(\ell>\exp(x^\eps/12)\). On the intervals in question
\(\ell>2t+3\), \(\ell\nmid nN\), and
\(\alpha(c)=(-1)^r\). These are the field and character
hypotheses used in the local chart and relative presentation.
Write \(H=G_{\QQ(i)}=\ker\eta\), \(\Psi=\kappa_\ell^r\alpha\),
and \(d=\operatorname{diag}(\Psi,1,\Psi^{-1})\).

We denote the completed marked ring by \(B\). Its matrix \(M\)
satisfies \(M(c)=1\). The remaining two marks set the adjacent
entries \(01,10\) of \(M(g_*)-M(cg_*^{-1}c)\) to zero.
Here \(g_*\in I_\ell\) gives three distinct residual diagonal
eigenvalues. The action by \(\operatorname{diag}(z,1,z^{-1})\)
gives the integer weight \(b-a\) to the entry \(ab\).

\begin{lemma}[Passing from the completed presentation to a graded ring]
\label{lem:algebra116}
Take the presentation \(B=P/(f_s)\) of Theorem~\ref{thm:pres70}
or Corollary~\ref{cor:first72}, where \(P=\OO[[X_1,\ldots,X_m]]\)
and each variable is chosen to vanish at the split point. Let
\(P_j\) consist of series all of whose monomials have weight \(j\).
There is a Noetherian graded ring \(A=\bigoplus_j P_j\subset P\).
Its invariant part \(P_0\) is complete local and every \(P_j\)
is finite over \(P_0\). The ring \(A/(f_s)\) embeds in \(B\),
and its degree-zero part is the full ring of invariants of \(B\).
The cotangent module of \(A\) at the split point is free over
\(\OO\) on the \(X_i\). In weight \(j\), the cotangent module
after the equations is the quotient of this free module by the
linear terms of the \(f_s\) of weight \(j\), in both presentations.
\end{lemma}
\begin{proof}
The possible weights are \(0,\pm1,\pm2\). Weight-zero monomials
are generated by variables of weight zero, opposite pairs, and
triples of weights \(1,1,-2\) or \(-1,-1,2\). Remove opposite
pairs first to see this. Choose one such factorization of every
invariant monomial. A monomial of ordinary degree \(d\) uses at
least \(d/3\) factors. Grouping a series by the chosen products
therefore expresses it as a convergent power series in finitely
many invariant monomials. Thus \(P_0\) is a quotient of a
complete Noetherian local power series ring over \(\OO\).

For fixed weight \(j\), remove invariant factors until none
remain. If all remaining weights have the same sign there are
at most \(|j|\) factors. With both signs present there can be
only weights of absolute sizes one and two with opposite signs,
and the weight of absolute size one occurs at most once.
This bounds the number of remaining factors by \(|j|+2\).
Grouping by the finitely many possible remainders proves that
\(P_j\) is finite over \(P_0\), and that \(A\) is generated
as an algebra by the original variables over \(P_0\).

Taking the weight-\(j\) part of any expression
\(\sum_s f_s b_s\) in \(P\) replaces \(b_s\) by its component
of weight \(j-\mathrm{wt}(f_s)\). This proves the asserted
injection. On each Artinian truncation torus invariance is
coefficientwise equality for an indeterminate invertible \(z\).
Project a lift of an invariant in \(B\) to weight zero. It
continues to lift that invariant on every truncation and hence
in \(B\) by separatedness.

Every homogeneous series with zero constant and linear terms is in
the square of the augmentation ideal of \(A\). In each monomial
factor off its first two variables in a fixed ordering and group
by the resulting pair. Each remaining series still has one weight.
This proves the assertion about the free cotangent module.
Taking the linear term of an arbitrary multiple of \(f_s\)
multiplies its linear term by the constant of the multiplier.
This proves the last assertion. The inclusion into
\(J[[X_1,\ldots,X_m]]\), where \(J=\operatorname{Frac}\OO\),
gives the formal chart of Section~\ref{sec:lower36} on \(A\).
\end{proof}

\begin{proposition}[The fixed subscheme is the split point]
\label{prop:fixed116}
In the graded solution ring \(A/(f_s)\), the quotient by all
elements of nonzero weight is \(\OO\), via the split augmentation.
In particular the weight-zero cotangent module of the marked
solution at this point is zero.
\end{proposition}
\begin{proof}
First work over any finite Artinian local coefficient algebra
with the prescribed residue field. A fixed lift has diagonal
matrices on \(H\), since nonzero-weight entries vanish. The
matrix on \(hc\) equals that on \(h\), so it is enough to
determine the three coordinate characters on \(H\).
On \(I_\ell\) the ordered ordinary flag of Lemma~\ref{lem:flag69}
has graded characters \(\Psi,1,\Psi^{-1}\). The coordinates
are already diagonal, and the residual eigenvalues of \(g_*\)
distinguish their order. Thus each coordinate agrees on this
inertia group with the corresponding entry of \(d\). The same
is true at the other primes over \(\ell\) by conjugation,
including the inverse transpose and reversal for \(c\).

On inertia away from \(\ell\) the quotient of a coordinate by
its base character has image of order prime to \(\ell\), by the
prescribed inertia kernel. It reduces to one and lies in the
finite \(\ell\)-group of principal units of the Artinian algebra.
It is therefore trivial. Each of the three quotient characters
on \(H\) is now everywhere unramified. Artin reciprocity and
the trivial ideal class group of \(\QQ(i)\) make it trivial
globally. The latter class number follows from Euclidean division
in \(\ZZ[i]\). Thus the only fixed lift is the split one, over
every such algebra.

Killing the nonzero-weight variables in \(P\) gives the complete
ring on the weight-zero variables, with the specialized equations.
It gives exactly the same quotient in \(A\). Indeed, in a
homogeneous series group the monomials that contain a nonzero-weight
variable by their first such variable, leaving coefficients of a
single weight in each group. Fixed lifts are represented by this
quotient because the homogeneous variables are entry differences or
parameters homogeneous in those differences, and fixedness is
equivalent to killing every nonzero weight. Apply uniqueness to
its own maximal-ideal truncations. All its augmented variables
vanish in each, so separatedness gives the isomorphism with
\(\OO\). In the cotangent module the removed variables have
nonzero weight, proving the assertion in weight zero.
\end{proof}

Let \(C_j\) be the weight-\(j\) cotangent module of the marked
solution at this augmentation, and use \(D=J/\OO\). For a
character \(\xi\) write \(S(\xi)\subset H^1(\QQ,D(\xi))\)
for the classes unramified at finite primes away from \(\ell\)
and zero on \(G_{\QQ_\ell}\).

\begin{theorem}[The two negative cotangent spaces]\label{thm:cot116}
There are identifications
\begin{align}
 \operatorname{Hom}_\OO(C_{-1},D)
       &\simeq S(\Psi^{-1})\oplus S(\eta\Psi^{-1}),\notag\\
 \operatorname{Hom}_\OO(C_{-2},D)
       &\simeq S(\eta\Psi^{-2}).                              \label{eq:roots116}
\end{align}
Both \(C_{-1}\) and \(C_{-2}\) have finite length. Hence the
negative equations span both rational negative cotangent spaces
of the ambient algebra. The weight-zero equations span its
integral weight-zero cotangent space.
\end{theorem}
\begin{proof}
Homogeneous derivations at the split point are determined by
the weighted variables in Lemma~\ref{lem:algebra116}. Their
values on matrix entries have the same weight, and those entries
generate the completed deformation. The weight-\(-1\) functor
is therefore exactly the one in Corollary~\ref{cor:dual68}.

For weight \(-2\), let \(W=\OO/\varpi^a\) and write a lift to
\(\OO\oplus W\), with square-zero ideal \(W\), as
\((1+z(g)E_{20})d(g)\) in its matrix part. The differential of
the sign involution sends \(E_{20}\) to \(-E_{20}\), and the
split diagonal acts by \(\Psi^{-2}\). Multiplication is
precisely the cocycle law for \(\xi=\eta\Psi^{-2}\).
Since \(\xi(c)=-1\), subtracting a unique coboundary makes
\(z(c)=0\). The other two marks are in adjacent entries and
are satisfied automatically.

At \(\ell\), the unit \(\xi(g_*)-1\) allows us to use a local
coboundary to set \(z(g_*)=0\). The ordinary flag must now be the
standard one, by the three distinct eigenvalues and
Lemma~\ref{lem:flag69}. The negative entry vanishes on inertia.
For \(g\in G_{\QQ_\ell}\), evaluate the cocycle at
\(g_*g=g(g^{-1}g_*g)\). This gives
\((\xi(g_*)-1)z(g)=0\), hence \(z(g)=0\). Conversely a local
coboundary gives a conjugate of the split point and belongs to
the chosen component. One may conjugate the split lift by
\(1+Y E_{20}\) over \(\OO[[Y]]\) and then specialize \(Y\)
to the square-zero value. The component is preserved by change
of frame, or directly by its defining matrix words, traces and
saturation as in Lemma~\ref{lem:component68}. Away from \(\ell\)
the inertia kernel condition gives the unramified class by
averaging on the prime-to-\(\ell\) inertia quotient.

Globally and on \(G_{\QQ_\ell}\) the divisible invariants for
\(\xi\) vanish by evaluation at \(g_*\). On each other inertia
group they are zero or divisible. Thus the multiplication sequence
identifies the finite strict group just computed with
\(S(\xi)[\varpi^a]\). Take the union over \(a\). This proves
\eqref{eq:roots116}.

All three groups on the right are character summands of strict
groups on \(E\QQ(i)\), by prime-to-\(\ell\) descent. Their twists
are \(r,r,2r\). They are finite by
Theorem \ref{thm:wild98}.
A free summand of a finitely generated \(\OO\)-module would
give an infinite summand \(D\) in its dual. Hence both cotangent
modules are torsion. Lemma~\ref{lem:algebra116} now proves the
rational spanning, and Proposition~\ref{prop:fixed116} gives
the integral assertion at weight zero.
\end{proof}

\begin{corollary}[The numerical presentation for every contributing character]
\label{cor:present116}
For all sufficiently distant intervals, every large-prime
component with \(S(\Psi^{-1})\ne0\) has the variable and equation
counts of Section~\ref{sec:lower36} and satisfies its formal
chart and rational cotangent spanning requirements. For those
counts set, in the order used in that section,
\[
  R=a_{-1}+2,\quad S=a_{-2}+1,\quad
  R'=a_1+3,\quad S'=a_2+1,\quad m_0=a_0+3 .
\]
In particular the quotient of the ambient weight-\(-1\) lattice
by the negative weight-one equations has length
\(\len_\OO S(\Psi^{-1})+\len_\OO S(\eta\Psi^{-1})\).
\end{corollary}
\begin{proof}
If \(r=1,\alpha=\eta\), the required \(S(\Psi^{-1})\) is a
summand at \(\ell\) of the first-twist strict group of the fixed
field \(\QQ(i)\). Its dual injects into the \(\ell\)-part of
the finite group \(K_2(\ZZ[i])\), by Theorem \ref{thm:wild98}.
All sufficiently large \(\ell\) therefore contribute zero,
uniformly in \(E\). This removes the only case where
Corollary~\ref{cor:first72} permits an extra equation.
In every other case the table in Theorem~\ref{thm:pres70}
gives \(R,S,R'+1,S'+1,m_0\) variables and at most
\(R+1,S,R',S',m_0\) equations in the five indicated weights.
Pad lists by zero equations if necessary. The free cotangent
module of the ambient algebra embeds into its rationalization,
so saturation of its augmentation square adds no linear terms.
The preceding theorem identifies its quotient in weight
\(-1\), including its full length by duality over the DVR.
\end{proof}

The integral intersection inequality of Theorem~\ref{thm:full36}
also requires finiteness of the weight-zero solution ring,
rigidity at the mixed points, and a nonvanishing trace difference
there. The results here establish the chart, the equations and
the distinguished-point conditions for the original ordinary
component. The comparison in the source occurs in
\srcloc{Development/AutomorphicArithmetic/AutomorphicVarieties/UnitaryThreeGroups/ModelUpper}{103}
and the strict-class injection in
\srcloc{Development/AutomorphicArithmetic/AutomorphicVarieties/UnitaryThreeGroups/VlenLe}{52}.


\subsection{Steinberg conditions at split primes and strict classes in weight one}
\label{sec:steintangent616}

Retain the fields, characters, coefficient ring \(\OO\), and marking of
Theorem~\ref{thm:ring39}. A Steinberg condition at a prime of residue
cardinality \(Q\) requires a single Jordan block for the logarithm of
tame inertia at a characteristic-zero point. Its integral closure also
contains points with smaller monodromy. We first compute its intersection
with zero monodromy. This determines which of the marked lifts can be
tested at such a prime.

\begin{proposition}[The unramified intersection]\label{prop:stintersection616}
Let \(Q>1\) be a power of a prime different from \(\ell\), with
\(\ell>3\). In the affine space of pairs of three by three matrices
\((\Phi,N)\) over \(\OO\), take the schematic closure \(\mathcal Z_Q\)
of the reduced characteristic-zero locus on which \(N\) is regular
nilpotent, \(\Phi N=QN\Phi\), and
\[
          \det(X-\Phi)=P_Q(X)=(X-Q)(X-1)(X-Q^{-1}).
\]
Its intersection with \(N=0\), as a scheme over \(\OO\), is defined by
the three coefficients of \(\det(X-\Phi)-P_Q(X)\).
\end{proposition}
\begin{proof}
Write \(A=\OO[\Phi_{ij}]\) and \(C\subset A\) for the ideal of these
three coefficients. Over an algebraic closure of the residue field the
fiber of \(A/C\) is a finite union of Jordan conjugacy classes. The
centralizer of a matrix has dimension at least three. Indeed, blocks of
sizes \(s_i,s_j\) with the same eigenvalue contribute
\(\min(s_i,s_j)\) to the dimension of the commuting algebra, and the
diagonal contributions already sum to three. Each class therefore has
dimension at most six. Krull's height theorem gives the opposite bound
for the dimension of every component of this three-equation fiber in
nine-space. Consequently \((C,\ell)\) has height four in \(A\) at each
closed point of the special fiber. In the regular local ring there the
three coefficients followed by \(\ell\) form part of a system of
parameters and hence a regular sequence. It follows that \(A/C\) has
no \(\ell\)-torsion. This assertion can be checked at those closed
points because any nonzero \(\ell\)-torsion has support in the special
fiber.

The generic fiber is reduced. In an eigenbasis the differential of
the three coefficient functions on diagonal perturbations has
Vandermonde determinant up to sign. The eigenvalues \(Q,1,Q^{-1}\)
are distinct, so this differential has rank three at each geometric
point. For any such \(\Phi\), let \(N_0\) carry its \(Q^{-1}\)-line
isomorphically to its \(1\)-line, its \(1\)-line to its \(Q\)-line,
and its \(Q\)-line to zero. The curve \((\Phi,tN_0)\) specializes a
regular block to \((\Phi,0)\).
If \(f\) vanishes on the defining locus of \(\mathcal Z_Q\), then
\(f(\Phi,0)\) vanishes on the reduced generic fiber of \(A/C\).
Torsion-freeness gives \(f(\Phi,0)\in C\). Conversely the three
coefficient differences already vanish on \(\mathcal Z_Q\).
These two ideal inclusions give the asserted scheme equality, including
its values on coefficient rings with nilpotents.
\end{proof}

\begin{proposition}[Contact with the marked weights]\label{prop:stweight616}
Let \(q\) be a rational prime different from \(\ell\) splitting completely
in \(E_x\QQ(i)\), and use arithmetic Frobenius. On the unramified
lower triangular lifts of the marked diagonal with diagonal characters
\(\Psi,1,\Psi^{-1}\), the normalized condition \(\mathcal Z_q\) cuts
out the principal scalar ideal generated by
\[
             a_{q,r}=(q^r-q)(1-q^{-r-1}).
\]
It vanishes identically in weight \(r=1\). If \(r>1\) and
\(d_q=v_\ell(q-1)\ge1\), its valuation is
\[
       2d_q+v_\ell(r-1)+v_\ell(r+1).
\]
In particular this is the largest \(m\) for which those triangular
lifts over \((\OO/\ell^m)[\epsilon]/(\epsilon^2)\) lie on
\(\mathcal Z_q\).
\end{proposition}
\begin{proof}
The characteristic polynomial of Frobenius is determined by
\(q^r,1,q^{-r}\). Its leading and constant coefficients agree with
those of \(P_q\), and the other two differences are, up to sign,
\[
       q^r+q^{-r}-q-q^{-1}=(q^r-q)(1-q^{-r-1}).
\]
Apply Proposition~\ref{prop:stintersection616}. For odd \(\ell\) and
\(d_q\ge1\), expanding successive \(\ell\)-th powers by the binomial
theorem gives \(v_\ell(q^j-1)=d_q+v_\ell(j)\) for \(j\ge1\).
Apply this to \(r-1\) and \(r+1\). Over
\(\OO\oplus\epsilon(\OO/\ell^m)\) the scalar component still takes
values in \(\OO\), so the same equation at an unramified point forces
\(r=1\) among the positive weights.
\end{proof}

There is an integral version with a variable weight when \(d_q\ge1\).
Put \(u=1+\ell\), \(a=\log_\ell(q)/\log_\ell(u)\), and assign
\(Y=(1+T)^a\) to the first Frobenius eigenvalue. Evaluating at
\(T=u^r-1\) recovers \(Y=q^r\). The same computation cuts out
\((Y-q)(Y-q^{-1})\) up to a unit. Write \(a=\ell^{d_q-1}b\), where
\(b\in\ZZ_\ell^\times\), and \(d=\ell^{d_q-1}\). In \(\OO[[T]]\)
this ideal is generated by the monic polynomial
\[
     \big((1+T)^d-u^d\big)\big((1+T)^d-u^{-d}\big).
\]
Indeed substitute \(V=((1+T)/u^{\pm1})^d\) in the binomial series
\((V^b-1)/(V-1)\). It is a unit since its residue at the closed
point is \(\bar b\). The polynomial reduces to \(T^{2d}\) in the
residue field. Its geometric generic roots are the \(2d\) distinct
values \(u^{\pm1}\zeta-1\) with \(\zeta^d=1\), since \(u^2\) has
infinite order.

\begin{theorem}[Strict tangents at ramified weight-one lifts]\label{thm:ststrict616}
On an interval with common cyclotomic modulus \(N_*\) divisible by
the conductors of the period fields, take distinct rational primes
\(q_1,q_2\equiv1\bmod 4N_*\), different from \(\ell\). Use the range
of \(\ell\) in Theorem~\ref{thm:ring39}. For every \(E_x\) on that
interval and every odd \(\alpha\) of \(\Gal(E_x/\QQ)\), there is a
marked integral family \(\rho_b\), \(b\in\ell\OO\), with the original
local conditions away from \(q_1,q_2\). At both primes its nonzero
members have regular tame monodromy and conductor exponent two.
For every \(m\ge1\), the group
\[
         e_\alpha\big(\Sel_{\ell,1}(E_x)\otimes\OO\big)[\ell^m]
\]
injects \(\OO\)-linearly into the tangent module over
\(\OO\oplus\epsilon(\OO/\ell^m)\) at the integral point
\(\rho_{\ell^m}\) satisfying both closed Steinberg conditions.
The injection persists modulo infinitesimal changes of basis.
\end{theorem}
\begin{proof}
Let \(D\) be the field cut out by \(\alpha\), let
\(\Delta=\Gal(D/\QQ)\), and put \(\Psi=\alpha\kappa_\ell\).
Consider first one of the two primes, \(q_i\). If \(h\) is the class
number of \(D\), a prime \(\mathfrak q_i\mid q_i\) satisfies
\(\mathfrak q_i^h=(a_i)\). The divisors of the elements
\(\sigma a_i\), \(\sigma\in\Delta\), are supported on distinct primes.
Their classes are multiplicatively independent modulo roots of unity.
Project \(a_i\) in \(D^\times\otimes_{\ZZ}J\) by the idempotent
\[
      e_{\alpha^{-1}}=\frac1{|\Delta|}
                   \sum_{\sigma\in\Delta}\alpha(\sigma)\sigma .
\]
Kummer theory, followed by twisting by \(\alpha\) and descent along
\(D/\QQ\), gives a class \(w_i\in H^1(\QQ,J(\Psi))\).
Here descent takes invariants, since positive-degree cohomology of the
finite group \(\Delta\) with \(J\)-coefficients is zero by averaging.
At a prime above \(q_i\) the projected valuation is nonzero. At primes
away from \(q_i,\ell\) it is zero. For a local field of residue
characteristic different from \(\ell\), the unit group has finite
pro-\(\ell\) completion, while the valuation gives a copy of
\(\ZZ_\ell\) in the completion of the multiplicative group.
The Kummer class of a uniformizer has nonzero restriction to tame
inertia, as is also seen by adjoining its successive \(\ell\)-power
roots over the maximal unramified extension. Thus \(w_i\) is ramified
at \(q_i\) and is zero on inertia over \(D\) away from \(q_i,\ell\).

Each \(a_i\) is a unit at every place over \(\ell\). Localization of
these unit Kummer classes, followed by the \(\ell\)-adic logarithm,
takes values in the \(\alpha^{-1}\)-line of \(D\otimes_{\QQ}J\).
Indeed the logarithm identifies the completed local units tensored
with \(J\) with the product of local fields tensored with \(J\) over
\(\QQ_\ell\). The normal basis theorem identifies \(D\otimes_{\QQ}J\)
with \(J[\Delta]\). Each of the two projected logarithms is nonzero.
To see this, fix an embedding in \(\overline{\QQ}_\ell\).
The logarithms of all \(\sigma a_i\) are linearly independent over
\(\QQ\), since a rational relation, after clearing denominators,
would give a product which is a local unit of logarithm zero.
Such a unit is a root of unity, and the divisor calculation forces
all exponents to vanish. Brumer's theorem on logarithms of algebraic
numbers upgrades the independence to algebraic coefficients. It
applies to these units after taking a common positive power in the
logarithm neighborhood of one. The coefficients of
\(e_{\alpha^{-1}}\) are algebraic, so its sum of logarithms is
nonzero.~\cite{VI-IV-brumer330}

There are consequently \(c_1,c_2\in J^\times\) such that
\(c_1w_1+c_2w_2\) restricts to zero on \(G_\ell\). The two distinct
ramified primes prevent cancellation on inertia at either of them.
Subtracting a coboundary makes a representing cocycle identically
zero on the specified \(G_\ell\). Its image is compact, so a common
scalar multiple gives a continuous cocycle
\(z:G_\QQ\to\OO(\Psi)\). On the old inertia kernels in
Theorem~\ref{thm:ring39} it is zero. Indeed these kernels lie in
\(G_D\) because \(\kappa_\ell\) is unramified away from \(\ell\),
and coboundaries vanish on inertia acting trivially on the
coefficients. Also \(z(c)=0\), since \(\Psi(c)=1\) and \(c^2=1\).

In one-based matrix indices set
\(U=E_{12}-E_{23}\) and \(L=E_{21}-E_{32}\).
The involution in Theorem~\ref{thm:ring39} fixes \(\exp(tU)\), and
its differential fixes \(L\). Conjugation by \(d(g)\) scales \(U\)
by \(\Psi(g)\) and \(L\) by \(\Psi(g)^{-1}\). The formula
\[
       \rho_b(g)=\big(\exp(bz(g)U)d(g),\,\eta(g)\big)
\]
therefore gives a homomorphism by the cocycle rule. Here the
exponential stops in degree two. It keeps the specified local
matrices at \(\ell\) and at \(c\). Since \(b\in\ell\OO\), it reduces
to the marked residual point. At \(q_i\), tame inertia has logarithm
a nonzero scalar multiple of \(U\) when \(b\ne0\), and Frobenius is
upper triangular with diagonal \(q_i,1,q_i^{-1}\). Thus the pair
belongs to the regular locus in characteristic zero and to its
closure over \(\OO\). Wild inertia is trivial and the inertia-fixed
space has dimension one, giving conductor exponent \(3-1=2\).

Now take \(b=\ell^m\), \(M=\OO/\ell^m\). Represent a strict class by
the normalized cocycle \(y:G_\QQ\to M(\Psi^{-1})\) of
Theorem~\ref{thm:ring39}. Multiply the matrix part of \(\rho_b(g)\)
on the left by \(1+\epsilon y(g)L\). On \(\epsilon M\) the action
of \(\rho_b\) agrees with the split action, so this is again a
homomorphism. The original marking is preserved. At \(q_i\), \(y\)
is unramified, so the logarithm of tame inertia is unchanged. Choose
a tame generator and an arithmetic Frobenius. The integral pair at
\(b=\ell^m\) reduces modulo \(\ell^m\) to \((d(\mathrm{Frob}_{q_i}),0)\).
For every polynomial equation \(f\) of \(\mathcal Z_{q_i}\), its
linear term on the square-zero ideal is its differential at this
reduction. The lower triangular change of Frobenius there leaves the
characteristic polynomial fixed. Proposition~\ref{prop:stintersection616}
therefore makes this derivative zero. The constant term already
vanishes at \(\rho_b\). This proves both Steinberg conditions for the
same tangent.

Finally \(X\mapsto(X_{21}-X_{32})/2\) retracts the inclusion
\(y\mapsto yL\) on the split adjoint module. The diagonal action
scales both entries by \(\Psi^{-1}\), and the differential of the
involution interchanges them with a minus sign. Hence the retraction
is equivariant also for the full semidirect action. It sends any
infinitesimal change of basis to a scalar coboundary, which proves
injectivity on classes and the asserted linearity.
\end{proof}

The primes can be chosen once on each sufficiently distant power
interval, for all its large coefficient primes. Linnik's theorem
applied first modulo \(4N_*\) and then modulo \(4N_*q_1\) supplies
distinct \(q_1,q_2\) bounded by a fixed power of \(4N_*\).
The interval bounds on \(N_*\) place both below
\(\exp(x^\eps/12)\) for all sufficiently large \(x\) in the interval.
Their conductor contribution on the rational base is \(q_1^2q_2^2\).
All nonzero \(\rho_b\) are conjugate over \(J\) by the matrices
\(\operatorname{diag}(b,1,b^{-1})\), which commute with the
involution in the semidirect product. They preserve the upper
triangular flag. Moreover the strict tangent has zero first variation
of every matrix trace, because \(\rho_{\ell^m}\) reduces to the
diagonal on \(M\). Converting this realization into an upper bound on
strict length would require control of integral directions over this
reducible point. Counting its rational isomorphism class gives no
such bound.

\section{Rigidity and orthogonal monodromy}\label{nav:g-IV-3}

\subsection{Ordinary deformations with a separated middle line}\label{nav:IV-18}
\label{sec:corner117}

Use the field, primes, twists and ordinary marked ring of
Section~\ref{sec:origin116}. Write \(J_2=\left(\begin{smallmatrix}
0&1\\1&0\end{smallmatrix}\right)\) and \(C=\operatorname{diag}(1,-1)\).
Let \(B^{\rm out}\) be the completed quotient obtained by setting
the four matrix entries between the middle line and the two
outer coordinates to zero, for every \(h\in H=G_{\QQ(i)}\).
The group law and \(M(c)=1\) impose the same zeros on the other
coset. Denote the invariant ring of this quotient by \(D^{\rm out}\).
The tangent directions of the full rank-three ring may include
deformations that leave this quotient.

\begin{lemma}[An odd two-dimensional representation over the integral ring]
\label{lem:outer117}
There is a continuous representation
\(W:G_\QQ\to\mathrm{GL}_2(B^{\rm out})\) of determinant \(\eta\).
For \(h\in H\), \(M(h)\) is \(W(h)\) on the outer coordinates
and \(1\) on the middle line. One has \(W(c)=C\), and the
residual representation is
\(\overline\Psi\oplus\overline\eta\,\overline\Psi^{-1}\).
On \(G_{\QQ_\ell}\) there is a stable direct-summand line
whose inertia character is \(\Psi\), with quotient inertia
character \(\Psi^{-1}\). The line splits as a module over the
coefficient ring.
\end{lemma}
\begin{proof}
Work first in a finite Artinian quotient. The determinant of
\(M|_H\) is one on inertia at \(\ell\), by the three prescribed
ordinary characters. On that inertia the middle character is
also one. In fact its line is the eigenline of \(M(g_*)\) for
the middle eigenvalue, by its distinct residue. The first
ordinary flag line belongs to the outer coordinates, and the
middle line maps isomorphically to the second graded piece.
These statements follow from the three spectral projectors
in Lemma~\ref{lem:flag69} over the Artinian ring. The ordered
characters give the assertion on all inertia. At the other
primes above \(\ell\), conjugating the flag or its reflected
dual preserves triviality of the middle and determinant
characters. On the remaining inertia groups both characters
reduce to one and factor through a finite prime-to-\(\ell\)
group. The class group argument of Proposition~\ref{prop:fixed116}
now gives determinant and middle character one on all of \(H\).

Let \(U(h)\) be the outer block. Thus \(U(h)\in\mathrm{SL}_2\),
and the global sign involution gives
\[
 U(chc)=J_2U(h)^{-\mathsf T}J_2=CU(h)C^{-1}.
\]
The second identity follows by multiplying a two-by-two matrix
of determinant one. Since \(G_\QQ=H\rtimes\langle c\rangle\),
these identities define \(W\) with the indicated value at \(c\).
Reduction at the split point gives the stated two characters.
On inertia at \(\ell\) the first ordinary line belongs to the
outer block and carries \(\Psi\). The \(\Psi(g_*)\)-eigenline is
unique. For any \(g\in G_{\QQ_\ell}\), conjugating this line by
\(W(g)\) gives another inertia line with the same character,
because \(\Psi\) is global. It must coincide with the first
line. Its quotient has the inverse character on inertia.
All the constructions are compatible on Artinian quotients.
The spectral projector at \(g_*\) exhibits the limit line as
a direct summand, proving the statements over \(B^{\rm out}\).
\end{proof}

\begin{theorem}[Finiteness of the integral trace algebra on the outer block]
\label{thm:finite117}
The ring \(D^{\rm out}\) is finite over \(\OO\) for every
large coefficient prime and every twist in the prescribed range.
In particular, on this quotient every continuous local
specialization of the invariant ring into \(k'[[s]]\) over a
finite residue extension \(k'/k\), in characteristic \(\ell\),
sends all traces to constants.
\end{theorem}
\begin{proof}
Twist \(W\) by \(\Psi^{-1}\). Its residual trace is
\(1+\bar\chi\), where its determinant is the fixed character
\(\chi=\eta\Psi^{-2}\). This determinant is odd. It is de Rham
at \(\ell\) as a cyclotomic power times an unramified finite
character. Its residual restriction at \(\ell\) is nontrivial
because on inertia it is \(\bar\kappa_\ell^{-2r}\) and
\(0<2r<\ell-1\). Choose the finite set of primes containing
\(\ell\) and all divisors of \(4N\). The representation is
unramified outside this set. Its trace on \(G_{\QQ_\ell}\) is
the sum of the two ordinary characters, and the one lifting
the trivial residual character is unramified there, by
Lemma~\ref{lem:outer117}.

Apply Pan's Theorem 5.1.2(1) for the totally real abelian base
field \(\QQ\). In the setup of his Section 5.1.1 this is the
ordinary pseudodeformation ring of the trace \(1+\bar\chi\)
with fixed determinant \(\chi\) and reducible local trace at
\(\ell\). It is finite over
\(\Lambda=\OO[[1+\ell\ZZ_\ell]]\), the completed group algebra
acting by local reciprocity through the restriction to inertia
of the local character lifting \(1\). The local residual
characters are distinct, so that character is unambiguously
labelled. These hypotheses hold here, including that \(\ell\)
is unramified in the base field.~\cite{VI-IV-VI-pan40}

The trace of \(W\Psi^{-1}\) induces a map from this universal
ring on all Artinian quotients, hence on their limit. Its weight
character is trivial on \(1+\ell\ZZ_\ell\). Thus the map factors
through the tensor product of the universal ring with \(\OO\)
along the augmentation of \(\Lambda\). The image is a finite
\(\OO\)-module. It contains the traces of \(W\), obtained by
multiplying the twisted traces by \(\Psi(g)\). A finite
\(\OO\)-module in \(B^{\rm out}\) is closed as a continuous
image of a finite product of compact copies of \(\OO\).

To identify the invariant ring, the proof of
Theorem \ref{thm:trace71}
applies in this quotient. In detail, \(Z=M(g_*)-M(cg_*^{-1}c)\)
is \(\operatorname{diag}(z,0,-z)\) for a unit \(z\), and
\(z^2=\tfrac12\operatorname{tr}Z^2\) is a polynomial in traces
on \(H\). The square root with the prescribed residue and the
inverse of \(z^2\) lie in the closed trace algebra by expansion
at their unit residues. The spectral projectors of \(Z\) are
polynomials in \(Z,1/z\). Taking traces of products with those
projectors inserted gives each diagonal entry and every opposite
pair of outer entries, also when the entries come from different
matrices. These generate the weight-zero entry monomials in this
two-coordinate block. The entry differences topologically
generate the deformation ring, so projection on finite
truncations puts every invariant in the closed trace algebra.
Conversely the traces are invariant. On \(H\) they are
\(1+\operatorname{tr}W(h)\). Hence \(D^{\rm out}\) is a
submodule of the finite module above and is finite over the
Noetherian DVR \(\OO\). Reduction of a finite local \(\OO\)-algebra
modulo \(\varpi\) is Artinian. Its image in the domain \(k'[[s]]\)
is its residue field, proving the last assertion.
\end{proof}

\begin{theorem}[The modular form and its conductor]\label{thm:form117}
Let a characteristic-zero integral point of \(B^{\rm out}\)
have some nonzero entries in both outer off-diagonal directions.
Here an integral point is a continuous local specialization into
the integers of a finite extension of \(J\). Then
\(W\otimes\kappa_\ell^r\alpha^{-1}\) is the Galois
representation of a normalized primitive cuspidal eigenform \(f\)
of weight \(2r+1\), nebentypus \(\eta\alpha^{-2}\), and level
exactly the conductor of \(\eta\alpha^{-2}\). In particular its
level divides \(4N\). If \(v\) is a rational prime split in
\(\QQ(i)\), prime to \(N\ell\), then with arithmetic Frobenius
\begin{equation}
     \operatorname{tr}M(\mathrm{Frob}_v)
              =1+v^{-r}\alpha(v)a_v(f).                 \label{eq:form117}
\end{equation}
For fixed \(\ell,r,\alpha\) the isomorphism class of this form
with its coefficient embedding determines at most one scaling
orbit of such points.
\end{theorem}
\begin{proof}
If \(W|_H\) had an invariant line over the algebraic closure,
the line would be stable under
\(U(g_*)-U(cg_*^{-1}c)=\operatorname{diag}(z,-z)\).
Its distinct eigenvalues force the line to be a coordinate line.
The two outer entry directions exclude both choices. Thus \(W\)
is absolutely irreducible and it is odd since \(W(c)=C\).

We check crystallinity. The stable line on \(G_{\QQ_\ell}\)
and its quotient are an unramified character times
\(\kappa_\ell^r\) and an unramified character times
\(\kappa_\ell^{-r}\), respectively. Their ratio \(V\) is
crystalline of the form \(\kappa_\ell^{2r}\nu\) with \(\nu\)
unramified. Over the coefficient field both \(H^0(V)\) and
\(H^0(V^*(1))\) vanish by inertia, since \(2r\ge2\).
Local duality and the local Euler characteristic give
\(\dim H^1(\QQ_\ell,V)=1\). The Bloch--Kato fundamental
exact sequence for a crystalline representation gives
\(\dim H^1_f=\dim H^0+\dim(D_{\rm dR}/\mathrm{Fil}^0)=1\)
in this positive cyclotomic twist. Hence every extension in
the indicated order is crystalline. This applies to \(W\)
and gives the two distinct Hodge--Tate weights.~\cite{VI-IV-VI-bk117}

The representation \(W\) is continuous, odd, absolutely
irreducible, finitely ramified, and de Rham with distinct weights
at \(\ell>5\). Pan's regular Fontaine--Mazur theorem over \(\QQ\)
therefore gives a cuspidal eigenform up to twist. The difference of the weights
is \(2r\). The twist displayed in the statement normalizes the
two cyclotomic exponents to \(0,2r\). Any remaining de Rham
character of weight zero over \(\QQ\), unramified almost
everywhere, has finite order. Indeed its restriction to local
units at \(\ell\) has finite image by the classification of
de Rham characters. At each other ramified prime the image of
inertia is finite by local reciprocity. Global reciprocity for
\(\QQ\) then gives a finite ray class character. Absorb such a
character into the eigenform and take its primitive
representative. The determinant equals
\(\kappa_\ell^{2r}\eta\alpha^{-2}\), giving the nebentypus.
We use arithmetic Frobenius for the usual Deligne
representation of a newform.~\cite{VI-IV-VI-pan40}

At a finite \(q\ne\ell\), inertia on \(W\) factors through the
image of \((\eta,\alpha)\) on \(I_q\), since it kills their common
kernel. This is a finite abelian group of order prime to \(\ell\).
Its two characters lift the distinct or repeated residual
characters with multiplicity, by the integral idempotents after
unramified coefficient extension. Prime-to-\(\ell\) characters
are determined by reduction. The two inertia characters are
therefore \(\alpha\) and \(\eta\alpha^{-1}\). After the indicated
twist these become \(1\) and \(\eta\alpha^{-2}\). Monodromy is
zero since inertia has finite image. Their Artin conductor is
exactly that of \(\eta\alpha^{-2}\), including at wildly
ramified primes. At \(\ell\) the crystalline representation has
unramified Weil--Deligne parameter and zero monodromy. Local--global
compatibility and the newform conductor theorem identify these
exponents with the primitive level. This gives the stated
divisibility and \eqref{eq:form117}.~\cite{VI-IV-VI-carayol117,VI-IV-VI-saito40}

An isomorphism of two such representations \(W\) commutes at
\(c\) with \(C\), hence is diagonal, say with entries \(a,b\).
On the outer block its conjugation is realized by
\(\operatorname{diag}(z,1,z^{-1})\) with \(z^2=a/b\) over the
algebraic closure. The middle line is fixed and \(M(hc)=M(h)\).
Thus the two marked systems are related by this scaling. Their
invariant-ring values agree by the closed trace description and
continuity. This gives the assertion about orbits.
\end{proof}

The character \(\alpha(v)\) in \eqref{eq:form117} records the
finite twist that gives the level \(4N\). It will enter the
integer bounding the trace differences. The ordinary
two-dimensional input in the source is specified in
\srcloc{Development/AutomorphicArithmetic/AutomorphicVarieties/AutomorphicCases/AutomorphyUnramifiedCoefficients}{19}.

\subsection{The cyclotomic restriction and the tangent space at an irreducible point}\label{nav:IV-19}
\label{sec:cyclo119}

Keep the ordinary ring \(B\) of Section~\ref{sec:origin116} and
the same \(x,E,\ell,r,\alpha\). A characteristic-zero integral
point \(P\) takes its values in the integers of a finite extension
of the coefficient field. Write \(U=M_P|_H\) for its
representation of \(H=G_{K'}\), where \(K'=\QQ(i)\). An
absolutely irreducible \(U\) means that it has no proper
nonzero invariant subspace over the algebraic closure of its
coefficient field. We also consider the integral realizations
of nonconstant trace curves in the special fiber supplied by
Theorem \ref{thm:strong118}.

\begin{lemma}[Strong irreducibility in characteristic zero]
\label{lem:strong119}
The determinant of \(M|_H\) in the entire ordinary ring is one.
For every sufficiently distant interval, an absolutely
irreducible \(U\) at a characteristic-zero integral point remains
absolutely irreducible on every open subgroup of \(H\).
\end{lemma}
\begin{proof}
Work first in a finite Artin quotient of \(B\). On inertia at
\(\ell\) the determinant on \(H\) is the product
\(\Psi\cdot1\cdot\Psi^{-1}\), also at a conjugate place.
Elsewhere it is a principal unit of order prime to \(\ell\) by
the inertia kernel condition, and thus equals one. The
everywhere unramified determinant factors through the ideal
class group of \(\QQ(i)\), which is trivial. Passing to the
inverse limit proves the assertion over \(B\).

Suppose the restriction of \(U\) became reducible over some
finite extension of \(K'\). It is then reducible on an open
normal subgroup \(H_0\) of \(H\). Over the algebraic closure
choose a simple \(H_0\)-submodule. Its translates under \(H\)
span the three-dimensional space. Their sum is semisimple and
all the simple constituents have the same dimension, which
divides three. Reducibility forces this dimension to be one.
The distinct character spaces are permuted transitively by \(H\)
and have equal dimensions. Thus \(H_0\) either acts by scalars
or has three distinct character lines.

As \(\ell\) is unramified in \(K'\), inertia at \(\ell\) lies
in \(H\) and its cyclotomic character surjects onto
\(\ZZ_\ell^\times\). The latter fact follows from total
ramification of \(\QQ_\ell(\mu_{\ell^a})\) for every \(a\).
An open subgroup of inertia therefore contains an element \(s\)
with \(\kappa_\ell(s)\) of infinite order. The ordinary
eigenvalues on \(s\) are \(\kappa_\ell(s)^r,1,\kappa_\ell(s)^{-r}\)
since \(\alpha\) is unramified there. Taking
\(s\in H_0\) excludes the scalar case.

In the other case the action on the three character lines gives
a transitive permutation quotient of the compact image of \(U\).
It is continuous, since the stabilizer of a line is closed and
there are only three lines. The kernel of reduction of this image
is pro-\(\ell\), as is seen using the successive congruence
subgroups of the integral matrix group. Its image in \(S_3\)
is trivial since \(\ell>3\). The residual matrix image is the
cyclic group of the matrices
\(\operatorname{diag}(\bar\Psi(h),1,\bar\Psi(h)^{-1})\).
A generator must act by a three-cycle. Any lift \(h\in H\) of
that generator has trace zero in the basis of the three
permuted lines. Reduction then gives \(a+1+a^{-1}=0\) for a
generator \(a\) of \(\bar\Psi(H)\), so \(a^3=1\). On inertia
the order of the image of \(\bar\Psi\) is
\((\ell-1)/\gcd(\ell-1,r)>3\), since
\(\exp(x^\varepsilon/12)>3t+1\) eventually. This is a
contradiction.
\end{proof}

\begin{theorem}[The image over the cyclotomic extension]
\label{thm:cyclo119}
Let \(U\) be as in Lemma~\ref{lem:strong119} with absolutely
irreducible image, or let \(U\) realize a nonconstant trace
curve in the special fiber. Over the algebraic closure of the
coefficient field let \(G\) denote the Zariski closure of
\(U(H)\). Its identity component \(G^0\) is semisimple and acts
irreducibly. For every finite extension \(L/K'\) the Zariski
closure of \(U(G_{L(\mu_{\ell^\infty})})\) contains \(G^0\).
In characteristic zero each of these cyclotomic images is
enormous in the sense of Newton and Thorne.
\end{theorem}
\begin{proof}
The inverse image of \(G^0\) is closed and of finite index, so
it is open. Lemma~\ref{lem:strong119}, or Theorem \ref{thm:strong118} in characteristic \(\ell\), shows that \(G^0\)
acts irreducibly. We take reduced algebraic subgroups over the
algebraic closure throughout. The unipotent radical of \(G^0\)
has a nonzero fixed vector by triangularization. Its fixed
space is \(G^0\)-stable by normality and hence is the whole
space. The matrix action is faithful, so this radical is
trivial. The connected reductive group \(G^0\) has finite
center by Schur's lemma and \(\det U=1\). It is therefore
semisimple.

The closure of \(U(G_L)\) contains \(G^0\) since finitely
many of its cosets cover \(G\). In the closure of the image of
\([G_L,G_L]\) one finds all commutators of elements of \(G^0\).
Indeed the commutator is a morphism on the product and
\(U(G_L)\times U(G_L)\) is dense in the product of its
closures. A connected semisimple group equals its algebraic
derived group. This last statement follows, for example,
from its generation by root groups and the formula
\([t,x_\beta(a)]=x_\beta((\beta(t)-1)a)\) for a root
\(\beta\) and a torus element \(t\) with \(\beta(t)\ne1\).
Thus the indicated closure contains \(G^0\). The cyclotomic
character has abelian image, so it kills \([G_L,G_L]\). This
proves the assertion for the full infinite cyclotomic extension.

In characteristic zero \(G^0\) contains an element with three
distinct eigenvalues, by applying the inertia argument of
Lemma~\ref{lem:strong119} to its open inverse image. It is
also the identity component of the closure of each cyclotomic
image above. Here are the details of the enormous-image test.
The elements with distinct eigenvalues form a nonempty Zariski
open subset of this connected component. The intersection of
the cyclotomic image with that subset is dense in \(G^0\).
Since \(G^0\) acts absolutely irreducibly, its linear span is
the whole endomorphism algebra by Burnside's theorem. For
any nonzero endomorphism \(X\), the trace pairing therefore
gives such an image element \(h\) with \(\operatorname{tr}(hX)\ne0\).
Writing \(h=\sum_\lambda\lambda e_\lambda\) in terms of its
eigenprojectors, at least one \(\operatorname{tr}(e_\lambda X)\)
is nonzero. Applying this to an element of each simple submodule
of the endomorphism algebra gives the required test. One may
enlarge the coefficient field finitely so all these eigenvalues
lie in it. Indeed an eigenvalue has degree at most three over
the initial local field, which has only finitely many extensions
of bounded degree in its algebraic closure. This is the
criterion in Newton and Thorne, Section~2.5.~\cite{VI-IV-VI-nt119}
\end{proof}

\begin{corollary}[Tangent classes and potentially automorphic points]
\label{cor:tangent119}
Suppose \(2\le r<t\) and \(P\) is a characteristic-zero integral
point for which \(U\) is absolutely irreducible. On
\(\mathcal V=\operatorname{End}(U)\) let \(H\) act by
conjugation and let \(c\) act by \(X\mapsto-J_3X^{\mathsf T}J_3\).
The tangent space of the graded solution at \(P\), modulo the
diagonal scaling direction, injects into
\(H^1_f(\QQ,\mathcal V)\), with scalars extended to the
algebraic closure. Suppose that for some CM field \(F\supset K'\)
the representation \((U\otimes\kappa_\ell^{-1})|_{G_F}\)
is attached to a regular algebraic, conjugate self-dual,
cuspidal automorphic representation of \(\mathrm{GL}_3(\mathbb A_F)\)
in the normalization of Newton and Thorne. Then this Selmer
group vanishes and the tangent space has dimension one.
\end{corollary}
\begin{proof}
An algebraic tangent vector gives continuous infinitesimal
matrices at \(P\). To verify continuity without any finiteness
assumption on the solution's invariant ring, use the ambient
graded ring in Lemma~\ref{lem:algebra116}. Its invariant part
is a quotient of a power series ring \(\OO[[Y_1,\ldots,Y_d]]\)
on finitely many invariant monomials. Enlarge the coefficient
field to contain their values \(b_j\), which are in its maximal
ideal. Extending the integer coefficients finitely commutes
with forming this power series ring by finite freeness. In
the extended ring translation by \(b=(b_j)\) is a continuous
change of variables. At \(b\) a series modulo
\((Y_1-b_1,\ldots,Y_d-b_d)^2\) is its constant plus its
ordinary linear term at \(b\). Every algebraic derivation of
this ring at \(P\) therefore acts by that linear term. The
weight-zero series and the finitely many nonzero-weight
variables generate the graded ring. After a common scalar
multiple, all their generator derivatives are integral in
one finite extension. Evaluating the variables plus their
infinitesimal changes now converges in the dual numbers over
that integer ring, and agrees with the original derivation on
every homogeneous equation. Undo the scalar multiple. This
produces continuous infinitesimal matrices over a finite
characteristic-zero field.

Writing a lift as \((1+\epsilon a(g))M_P(g)\) makes \(a\) a
one-cocycle for the displayed action, by the semidirect
product law. At primes different from \(\ell\) its restriction
to inertia factors through the prescribed finite quotient.
Averaging shows that this restriction is a coboundary.

Let \(K\) be a completion of \(K'\) at \(\ell\). On \(G_K\)
both the ordinary representation and its infinitesimal lift
have a flag with inertia exponents \(r,0,-r\), in that order.
The flag follows from the integral spectral projectors at
\(g_*\) and Lemma~\ref{lem:flag69}, applied on successive
Artin quotients. A change to the conjugate place takes a
reflected dual flag and gives the same exponents. Regard a
line over the dual numbers as a two-dimensional representation
over the coefficient field. Each graded piece is then an
unramified representation times its stated cyclotomic power,
and is crystalline. The same is true of each graded line of
the original point.

Build the representation down its flag. If \(T\) is the
crystalline part already built and \(Q\) is the next graded
piece, \(V=\operatorname{Hom}(Q,T)\) is crystalline. All its
cyclotomic exponents are at least two because \(r\ge2\).
Consequently \(H^0(K,V^*(1))=0\) and
\(\dim(D_{\rm dR}(V)/\mathrm{Fil}^0)=[K:\QQ_\ell]\dim V\),
with dimensions over the coefficient field. The first
vanishing also follows directly from the flag since all
graded exponents of \(V^*(1)\) are strictly negative. Local
duality, the Euler characteristic and the Bloch--Kato dimension
formula now give
\[
 \dim H^1(K,V)=[K:\QQ_\ell]\dim V+\dim H^0(K,V)
              =\dim H_f^1(K,V).
\]
Thus every such extension is crystalline. Induction proves
crystallinity of the specialization and of the underlying
self-extension given by its infinitesimal lift. This places
the local adjoint class in \(H^1_f(K,\mathcal V)\), by the
crystalline extension criterion.~\cite{VI-IV-VI-bk117}
It also lies in \(H^1_f(\QQ_\ell,\mathcal V)\). Indeed its
image in cohomology with \(B_{\rm cris}\)-coefficients
vanishes upon restriction to \(K\), and restriction followed
by corestriction multiplies by the nonzero scalar
\([K:\QQ_\ell]\).

If the global cohomology class vanishes, write
\(a(g)=X-gX\). The mark at \(c\) gives \(X=-J_3X^{\mathsf T}J_3\).
On \(H\), every infinitesimal matrix is the commutator of
\(X\) with its original matrix. The marked difference
\(Z=M(g_*)-M(cg_*^{-1}c)\) stays diagonal, so \([X,Z]\)
is diagonal. Since \(Z\) has distinct eigenvalues in the
marked coordinates, \(X\) is diagonal. The equation at \(c\)
then makes \(X\) a multiple of \(\operatorname{diag}(1,0,-1)\).
These multiples give precisely the infinitesimal scaling
directions on entries, hence on the graded ring by the
continuity argument. At an irreducible point some off-diagonal
entry is nonzero, so the scaling direction has dimension one.
This proves the injection.

Now assume the stated automorphy over \(F\), and write \(F^+\)
for its maximal totally real field. Then \(F=F^+K'\). The
matrix \(J_3\) is symmetric and the marked inverse-transpose
identity identifies the \(c\)-conjugate of
\(U\kappa_\ell^{-1}\) with its dual times \(\kappa_\ell^{-2}\).
This gives the polarization used in the adjoint Selmer theorem.
Its pairing is unique up to scalar by absolute irreducibility
on \(G_F\). Therefore its adjoint module on \(G_{F^+}\) is
the restriction of \(\mathcal V\), including the indicated
negative adjoint action at \(c\). The cyclotomic twist is
trivial on \(G_{F(\mu_{\ell^\infty})}\). Theorem~\ref{thm:cyclo119}
verifies the enormous-image hypothesis. Newton and Thorne's
Theorem~A thus gives \(H^1_f(F^+,\mathcal V)=0\).
Restriction preserves the local Selmer conditions, and the
restriction from \(\QQ\) is injective by corestriction in
characteristic zero. This proves the vanishing over
\(\QQ\) and the assertion about the tangent dimension.~\cite{VI-IV-VI-nt119}
\end{proof}

For the full cohomology bound on the fields of Section~\ref{sec:missing},
automorphy of all these three-dimensional points still requires
a proof. So does exclusion of the characteristic-\(\ell\)
curves in Theorem~\ref{thm:cyclo119}. In the source the
cyclotomic irreducibility argument occurs in
\srcloc{Development/AutomorphicArithmetic/AutomorphicVarieties/UnitaryThreeGroups/TypeIIrrCyc}{77}
and the adjoint application in
\srcloc{Development/AutomorphicArithmetic/AutomorphicVarieties/UnitaryThreeGroups/LiftingTheory}{55}.

\subsection{Rigidity in the full space at the two-dimensional points}\label{nav:IV-20}
\label{sec:rigid120}

Let \(P\) be a characteristic-zero integral mixed point whose
four adjacent entries vanish, in the ordinary deformation for
the fields and intervals of Section~\ref{sec:origin116}. Theorem
\ref{thm:form117} gives a modular representation \(W\) of
\(G_\QQ\) on the outer two coordinates with determinant \(\eta\).
Write \(J\) for a finite coefficient extension large enough for
the point and its normalized primitive eigenform \(f\). At a
good prime we use arithmetic Frobenius for \(V_f\), so that
\[
  k=2r+1,\qquad V_f=W\otimes\kappa_\ell^r\alpha^{-1},
  \qquad \mathrm{level}(f)=f(\eta\alpha^{-2})\mid4N .
\]
Here \(f(\xi)\) on the right denotes the conductor of a
character. We consider the full nine-dimensional adjoint
module \(\mathcal V\) from Corollary~\ref{cor:tangent119},
including the entries between the outer coordinates and the
middle coordinate. In particular this calculation tests
infinitesimal deformations in the original three-dimensional ring.

\begin{lemma}[The modular form at a mixed point has no complex multiplication]
\label{lem:nocm120}
The representation \(W|_H\) remains absolutely irreducible
on every open subgroup of \(H=G_{\QQ(i)}\) on all sufficiently
distant intervals. In particular \(f\) is a non-CM newform.
\end{lemma}
\begin{proof}
Irreducibility on \(H\) follows from the two opposite outer
entries and the marked matrix with two distinct outer eigenvalues,
as in Theorem~\ref{thm:form117}. If an open restriction is
reducible, choose an open normal subgroup \(H_0\) on which it
is reducible. The translates of a simple submodule span the
space. In dimension two the restriction is thus a sum of
characters. If the two characters are equal, an open subgroup
of inertia at \(\ell\) acts by scalars. This is impossible
since its eigenvalues on an element with cyclotomic value of
infinite order are \(\kappa_\ell(s)^r\) and
\(\kappa_\ell(s)^{-r}\).

If the characters differ, \(W(H)\) permutes their lines
transitively. The resulting continuous quotient of order two
kills the pro-\(\ell\) kernel of reduction. The residual image
of \(H\) consists of \(\operatorname{diag}(a,a^{-1})\)
for \(a\in\bar\Psi(H)\) and is cyclic. A lift of a generator
must swap the lines, so its trace is zero. Thus a generator
\(a\) satisfies \(a+a^{-1}=0\), giving order four. But on
inertia the order of the image is at least \((\ell-1)/r>4\)
throughout all sufficiently distant intervals. This proves
the assertion about restriction. A CM newform comes by
induction of a Hecke character of an imaginary quadratic
field. Its Galois representation is a sum of characters on
that field, as follows either from the induced representation
or by comparing the Frobenius polynomials of the theta series.
Restriction to its compositum with \(\QQ(i)\) would contradict
what we have just proved. A character twist does not affect
irreducibility on a subgroup.
\end{proof}

\begin{theorem}[Vanishing for the full adjoint module at a mixed outer block]
\label{thm:fullad120}
At every such \(P\) one has \(H^1_f(\QQ,\mathcal V)=0\) for
all \(1\le r<t\), including the first twist.
\end{theorem}
\begin{proof}
For \(h\in H\) the action is conjugation by the block matrix
with outer block \(W(h)\) and middle entry one. Complex
conjugation acts on endomorphisms by the negative reflected
transpose. The two scalar matrices on the outer and middle
blocks therefore give two copies of \(\eta\). On a traceless
outer matrix \(X\), the action at \(c\) is
\(-J_2X^{\mathsf T}J_2=CXC^{-1}\), with \(C=W(c)\)
as in Lemma~\ref{lem:outer117}. This gives the usual
trace-zero adjoint of \(W\). The space of maps from the
middle line to the outer plane is \(W|_H\) under \(H\).
The action of \(c\) exchanges it with the space of maps in
the reverse direction. Their direct sum is consequently
\(\operatorname{Ind}_H^{G_\QQ}(W|_H)\). Since \(W\) is already
defined on \(G_\QQ\), this induction is
\(W\otimes(1\oplus\eta)\). Altogether we obtain
\begin{equation}
 \mathcal V\simeq\operatorname{ad}^0 W
       \ \oplus\ \eta^{\oplus2}\ \oplus\ W\ \oplus\ (W\otimes\eta).
 \label{eq:ad120}
\end{equation}
Bloch--Kato local conditions commute with this direct sum.

An \(H^1_f(\QQ,\eta)\) class restricts on \(\QQ(i)\) to an
additive character unramified away from \(\ell\). At a place
above \(\ell\), the unramified classes for the trivial
coefficient field \(J\) are crystalline and have dimension
one. The Bloch--Kato dimension formula for \(J\) gives
\(\dim H_f^1=1\) there, so the character is unramified at this
place as well. It factors through the finite ideal class group,
which has no nonzero image in the additive group of \(J\).
Restriction is injective by corestriction. This disposes of
both scalar summands. For \(\operatorname{ad}^0W\) the scalar
twist relating \(W\) and \(V_f\) disappears under conjugation.
Theorem~B(1) of Newton and Thorne gives the vanishing for the
adjoint of this non-CM regular algebraic cuspidal modular
representation over \(\QQ\). Its trace-zero direct summand
therefore has zero Selmer group.~\cite{VI-IV-VI-nt119}

For the other two summands take \(L=E\QQ(i)\), which is finite
abelian over \(\QQ\). Both \(\alpha\) and \(\eta\) are trivial
on \(G_L\). Kato's representation \(V_\lambda(f)\), where
\(\lambda\mid\ell\) is the coefficient place, uses the inverse
of arithmetic Frobenius in its Euler polynomial. Thus
\(V_\lambda(f)\simeq V_f^\vee\) after extending coefficients.
Indeed the Euler polynomial of its inverse Frobenius agrees
with that of arithmetic Frobenius on \(V_f\) at every good
prime, by Kato's formula (14.10.3). Chebotarev and equality
of semisimplifications give the isomorphism since \(V_f\) is
absolutely irreducible. Consequently
\[
 V_\lambda(f)(r)
       \simeq W^\vee\otimes\alpha
       \simeq W\otimes\eta\alpha ,
\]
using the exterior pairing in dimension two. Over \(L\) this
is \(W|_{G_L}\) and also \((W\otimes\eta)|_{G_L}\).
Since \(1\le r\le k-1\) and \(2r\ne k\), Theorem 14.2(1)
of Kato gives finiteness of the Selmer group over \(L\)
with divisible torsion coefficients obtained from a stable
lattice in \(V_\lambda(f)(r)\). This is the off-center case
of that theorem.~\cite{VI-IV-VI-kato120}

We explain the passage to field coefficients. Write \(V\)
for Kato's twisted representation and \(T\) for the chosen
lattice. Fix a finite set \(S\) of places of \(L\) containing
the ramification, the archimedean places and those above \(\ell\).
Classes in \(H^1_f(L,V)\) are unramified outside \(S\), so
they lie in \(H^1(G_{L,S},V)\). Their images in
\(H^1(G_{L,S},V/T)\) lie in Kato's finite group by his local
definition, using the injection by inflation into global
cohomology. A vector space is divisible as an additive group,
and thus has zero image in a finite group. The coefficient
sequence now puts this entire vector space in the image of
\(H^1(G_{L,S},T)\).

This integral cohomology is a finitely generated module over
the coefficient DVR. In fact the degree-one groups of all
finite quotients of \(T\) are finite by finiteness of
extensions of bounded degree unramified outside \(S\). The
degree-zero groups in the same inverse system are finite as
well. Thus the integral group is their compact degree-one
inverse limit. Reduction modulo the uniformizer embeds its
quotient by that uniformizer into the first cohomology of
the finite residual module. The compact form of Nakayama's
lemma gives finite generation. The image of a finitely
generated integral module in a vector space contains no
nonzero vector subspace over the coefficient field. Hence
\(H^1_f(L,V)=0\). The conclusion persists under every finite
extension of coefficient fields by exactness of scalar
extension and the definition of the local conditions.
Finally restrict either of the last two summands in
\eqref{eq:ad120} from \(\QQ\) to \(L\). Their restrictions
are zero in Selmer cohomology by the preceding calculation,
and corestriction multiplies by \([L:\QQ]\). This proves
both vanishings over \(\QQ\). The use of \(L\) is an
injective restriction for the point associated with the
original period field \(E\).
\end{proof}

\begin{corollary}[Only scaling survives in the full ordinary tangent space]
\label{cor:corner120}
At every two-dimensional mixed point in
Theorem \ref{thm:tracebound118},
the geometric tangent space of the full graded solution has
dimension one. It is generated by diagonal scaling. This
holds uniformly over the indicated fields and intervals for
all \(1\le r<t\).
\end{corollary}
\begin{proof}
By Lemma \ref{lem:scale118} and Theorem~\ref{thm:finite117},
such a point has a continuous integral representative after
scaling. The construction of the global adjoint cocycle from
an algebraic tangent vector in Corollary~\ref{cor:tangent119}
used only integrality of the representative and the finite-weight
presentation. It applies in the full algebra \(B\), so the
cocycle may have arbitrary adjacent entries. Its inertia
classes away from \(\ell\) vanish by the same averaging.

At a place of \(K'\) over \(\ell\) write \(K\) for the local
field, of residue degree \(d\). The ordinary flag of \(P\)
has diagonal characters \(\xi_0,1,\xi_2\), of cyclotomic
exponents \(r,0,-r\). They are crystalline, as is the
representation \(W\) by Theorem~\ref{thm:form117}.
The primitive form \(f\) has level prime to \(\ell\). The two
roots of its Hecke polynomial at \(\ell\) have complex
absolute value \(\ell^r\) at every embedding by Deligne.
Crystalline local--global compatibility identifies them with
the arithmetic-normalized Weil--Deligne eigenvalues of \(V_f\)
there. Undoing the twist \(\kappa_\ell^r\alpha^{-1}\)
makes the two eigenvalues for \(W\) have complex absolute
value one. Passing to the unramified extension \(K\) raises
them to the power \(d\). The middle character has eigenvalue
one as well. The eigenvalues of the three diagonal characters
are exactly these eigenvalues, by exactness on crystalline
representations.~\cite{VI-IV-VI-deligne40,VI-IV-VI-saito40}

It follows that no ratio of two of these diagonal characters
equals \(\kappa_\ell\) on \(G_K\). Such a ratio would have a
Weil--Deligne Frobenius eigenvalue of absolute value
\(\ell^d\) in the same normalization. In an infinitesimal
ordinary deformation, each graded line, regarded over the
coefficient field, has two copies of its original character
as composition factors. Its unramified factor may acquire a
Jordan block, which is crystalline. Now repeat the flag
extension argument of Corollary~\ref{cor:tangent119}.
For an already constructed crystalline part \(T'\) and next
graded piece \(Q'\), the graded exponents of
\(V'=\operatorname{Hom}(Q',T')\) are positive. No composition
factor of \(V'\) is \(\kappa_\ell\), by the ratio calculation.
Hence \(H^0(K,(V')^*(1))=0\), also when \(r=1\). The local
Euler characteristic and the Bloch--Kato dimension formula
give \(H^1(K,V')=H^1_f(K,V')\). This proves crystallinity of
the whole self-extension and places the local adjoint class
in \(H_f^1\), descending to \(\QQ_\ell\) by corestriction.

The global class vanishes by Theorem~\ref{thm:fullad120}.
For a coboundary the mark at \(c\) makes its matrix \(X\)
satisfy \(X=-J_3X^{\mathsf T}J_3\). The diagonal mark
\(Z\) with distinct eigenvalues forces that matrix to be
diagonal, exactly as in Corollary~\ref{cor:tangent119}.
It is a multiple of \(\operatorname{diag}(1,0,-1)\).
The two opposite outer entries at a mixed point make the
corresponding scaling tangent nonzero. This proves the
assertion in the full solution.
\end{proof}

Thus on this part of the solution the integral trace algebra,
the characteristic-zero tangent dimension and the uniform
trace sum are all controlled. The intersection estimate of
Theorem~\ref{thm:full36} still requires the corresponding
arithmetic statements on the irreducible three-dimensional
part. The matrix decomposition behind the Selmer calculation
appears in
\srcloc{Development/AutomorphicArithmetic/AutomorphicVarieties/UnitaryThreeGroups/ReducibleThreeDim}{263}.

\subsection{Integral representatives and the crystalline first twist}\label{nav:IV-21}
\label{sec:crystalline121}

Let \(B\) be the ordinary marked deformation ring on the fields
and intervals in Section~\ref{sec:origin116}. Denote its
invariant ring by \(D\) and its graded algebra of finite sums
of homogeneous elements by \(B^{\rm gr}\). The coefficient
DVR is \(\OO\), with fraction field \(J\). We take geometric
characteristic-zero points with values in \(\overline J\).

\begin{lemma}[Integral representatives of all algebraic points]
\label{lem:integral121}
Every \(\OO\)-algebra homomorphism \(D\to\overline J\) is a
continuous local map to the integers of a finite extension of
\(J\). Every geometric point of \(B^{\rm gr}\) can consequently
be scaled to a continuous integral point of \(B\).
\end{lemma}
\begin{proof}
The invariant ring \(D\) is a complete local quotient of a formal
power series ring in finitely many variables over \(\OO\), with
the same residue field, by Lemma~\ref{lem:algebra116}. Let
\(\xi:D\to\overline J\), \(d\in\mathfrak m_D\), and write
\(y=\xi(d)\). If \(y\ne0\) and \(v(y)\le0\), every conjugate
of \(1/y\) over \(J\) is integral since the valuation on the
complete local field has a unique extension. Thus
\[
       f(T)=\prod_{\sigma:J(y)\hookrightarrow\overline J}
                          (1-T/\sigma(y))\in\OO[T],
       \qquad f(0)=1.
\]
Then \(f(d)\) is a unit in \(D\) with zero image under \(\xi\),
a contradiction. Every element of \(\mathfrak m_D\) therefore
has zero image or strictly positive valuation. Subtracting an
integral lift of the residue shows that \(\xi(D)\) is integral.

Let \(d_1,\ldots,d_b\) be the images of formal parameters.
Together with a uniformizer they generate \(\mathfrak m_D\).
All their images belong to one finite extension \(J'\).
The minimum valuation of their nonzero images and that of the
uniformizer is a positive number \(a\). Every element of
\(\mathfrak m_D^n\) has image of valuation at least \(na\),
since all coefficients in a finite expression in the ideal
generators have integral image. It follows that \(\xi\) is
continuous and agrees, on any power series, with the limit
of the images of the polynomial truncations. The complete
field \(J'\) contains each such limit. This proves the first
assertion. For a point of \(B^{\rm gr}\) apply
Lemma \ref{lem:scale118}
to its invariant values and the prescribed point of the fiber.
The resulting scaling gives the integral representative.
\end{proof}

\begin{theorem}[Crystallinity on the entire ordinary component]
\label{thm:crystalline121}
Let \(K\) be a completion of \(K'=\QQ(i)\) at \(\ell\).
Restrict the local ordinary component of
Sections~\ref{sec:original68}--\ref{sec:kummer72} to \(G_K\).
Every characteristic-zero specialization in a finite local
\(J\)-algebra, obtained from its framed deformation ring, is
crystalline with the labeled Hodge type of the three
cyclotomic exponents \(r,0,-r\). This holds for the whole range
\(1\le r<t\), including specializations to dual numbers at
points with an adjacent character ratio equal to \(\kappa_\ell\).
\end{theorem}
\begin{proof}
Write \(A=U_{\rm loc}/I_t\) for the ordinary local component.
It is a formal power series ring over \(\OO\) by
Theorems~\ref{thm:smooth69} and~\ref{thm:first72}. Use first
the place where its ordered flag is defined. The restricted
representation of \(G_K\) is an ordinary matrix representation
since the sign character is trivial there. At a field-valued
integral specialization it has diagonal graded characters equal
to unramified characters times \(\kappa_\ell^r,1,\kappa_\ell^{-r}\)
in that order.

Set \(a=\Delta\) as in Section~\ref{sec:kummer72} when \(r=1\),
and set \(a=1\) when \(r\ge2\). In the first case \(\Delta\ne0\)
in \(A\), by the formula for its image in the diagonal parameters
in Theorem~\ref{thm:first72}. At a field point where \(a\ne0\),
no upper graded character divided by a lower one equals
\(\kappa_\ell\) on \(G_K\). For \(r\ge2\) their exponents
already differ by at least two. For \(r=1\), the outer pair
has difference two. In an adjacent pair the factor \(q_j\)
in \(\Delta\) is a nonzero constant times
\(\xi_j(\gamma)-\kappa_\ell(\gamma)\xi_{j+1}(\gamma)\)
for the positive-sign element \(\gamma\in G_K\). Its nonzero
value excludes the asserted equality of characters.

At such points all three graded characters are crystalline.
Given a crystalline part \(T\) already constructed from the
top of the flag and the next graded character \(Q\), let
\(V=\operatorname{Hom}(Q,T)\). Its graded cyclotomic exponents
are positive and it has no composition factor equal to
\(\kappa_\ell\). Thus \(H^0(K,V^*(1))=0\) and the dimension
of \(D_{\rm dR}(V)/\mathrm{Fil}^0\) is
\([K:\QQ_\ell]\dim V\) over the coefficient field. The local
Euler characteristic and the Bloch--Kato dimension formula give
\(H^1(K,V)=H_f^1(K,V)\). Every extension at this step is
crystalline, proving the assertion at all these field points
with their indicated Hodge type.~\cite{VI-IV-VI-bk117}

Now apply the framed potentially crystalline deformation theorem
of Kisin to \(G_K\), the trivial inertial type, and the indicated
Hodge type at every embedding of \(K\). We may first enlarge
the finite coefficient field to contain the fields of definition
of the embeddings. Formation of our power series ring commutes
with this finite flat coefficient extension. Kisin gives a
quotient of the framed ring of the residual three-dimensional
representation that tests the prescribed type on every finite
local coefficient algebra in characteristic zero. Potential
crystallinity of trivial inertial type here is crystallinity.
Indeed over a finite Galois extension where a representation
is crystalline, inertia acts trivially on its crystalline
period module when the inertial type is trivial. Semilinear
descent on the maximal unramified coefficient field then gives
the full dimension of crystalline periods downstairs. For finite
local coefficients one can apply this dimension test to the
underlying vector space. A finite inertial action with trivial
type acts trivially there by semisimplicity in characteristic
zero.~\cite[Theorem 1.2.2]{VI-IV-VI-kisin121}

The restricted framed family gives a homomorphism from this
universal framed ring to \(A\). Every element in the image of
the defining ideal of the crystalline quotient vanishes at the
integral field points where \(a\ne0\). These points detect
every nonzero element of \(A\). In fact, if \(u\in A\) is
nonzero, consider the nonzero series \(ua\). A nonzero
one-variable integral series takes a nonzero value at some
sufficiently small nonzero power of the uniformizer. Its lowest
nonzero term then has strictly smaller valuation than its tail.
In several variables first make one nonzero coefficient series
nonzero by induction, then apply the one-variable assertion to
the remaining variable. All chosen coordinates have positive
valuation. Thus \(ua\) is nonzero at an integral field point.
This proves that the whole crystalline ideal maps to zero in
\(A\). Every specialization now factors through Kisin's quotient,
including a specialization with nonzero nilpotents. The simple
rank and period conditions descend by a finite extension of the
coefficient field. At a conjugate place the family is conjugate
to the original family or its reflected dual. Its Hodge type is
the same after reflection and duality, and crystallinity is
preserved. This proves the theorem at every place above \(\ell\).
\end{proof}

\begin{corollary}[The first twist at arbitrary irreducible mixed points]
\label{cor:tangent121}
For every geometric irreducible mixed point of \(B^{\rm gr}\)
the injection into \(H_f^1(\QQ,\mathcal V)\) in
Corollary~\ref{cor:tangent119} holds throughout \(1\le r<t\).
At any such point which is potentially automorphic as specified
in that corollary the geometric tangent dimension is one.
\end{corollary}
\begin{proof}
Lemma~\ref{lem:integral121} allows an integral representative.
An algebraic tangent at it yields continuous infinitesimal
matrices after a common scalar multiple, by the power series
argument of Corollary~\ref{cor:tangent119}. Theorem
\ref{thm:crystalline121} puts their restriction at every
place of \(K'\) over \(\ell\) in the crystalline self-extensions.
Therefore the class satisfies the Bloch--Kato local condition
at \(\ell\), using restriction and corestriction as in that
corollary. Its inertia classes away from \(\ell\) vanish by
finite averaging. The two marks identify the kernel with
diagonal scaling. Finally Theorem~\ref{thm:cyclo119} supplies
the cyclotomic image hypothesis over the CM field of potential
automorphy, and the adjoint Selmer vanishing descends to
\(\QQ\) by corestriction, as in
Corollary~\ref{cor:tangent119}. All these arguments now apply
at \(r=1\) as well.
\end{proof}

\subsection{Orthogonal monodromy on the ordinary component}\label{nav:IV-22}
\label{sec:ordinaryorth122}

Take \(U\) from Theorem~\ref{thm:cyclo119}. Thus \(U\) is either
an absolutely irreducible integral characteristic-zero point or
an integral realization of a nonconstant trace curve in residue
characteristic \(\ell\). Write \(F\) for its local coefficient
field, \(R\) for the valuation ring, and \(k_F\) for its finite
residue field. We may extend \(F\) finitely. In equal
characteristic the symbols \(\Psi,\alpha,\eta\) denote the
images of the original characters under the coefficient map.
On \(H=G_{\QQ(i)}\) the residual matrix in the marked lattice
is \(\operatorname{diag}(\bar\Psi,1,\bar\Psi^{-1})\).
On the intervals under consideration the order of
\(\bar\Psi(I_\ell)\) exceeds four for every sufficiently large
\(x\), uniformly for \(\ell>\exp(x^\varepsilon/12)\) and
\(1\le r<t\).

\begin{theorem}[The two possible monodromy groups]
\label{thm:orth122}
Over \(\overline F\) the Zariski closure of \(U(H)\) is
either \(\mathrm{SL}_3\) or the special orthogonal group of a
nondegenerate symmetric bilinear form \(S\). The full image
preserves \(S\) in the latter case, and its projective class
of forms is unique.
\end{theorem}
\begin{proof}
By Theorem~\ref{thm:cyclo119} the identity component \(G^0\)
is semisimple and irreducible in its faithful three-dimensional
representation, and the full group \(G\) is in \(\mathrm{SL}_3\).
The rank of \(G^0\) is one or two. Rank zero would give the
trivial group. We use the root subgroups and Weyl group of a
connected reductive group over an algebraically closed
field.~\cite{VI-IV-VI-milne122}

In rank two its maximal torus is conjugate to the full diagonal
torus of \(\mathrm{SL}_3\). A root of \(G^0\) occurs in its Lie
algebra and hence is one of the characters \(t_i/t_j\) on
the ambient off-diagonal matrix entries. Its root subgroup is
\(I+aE_{ij}\), \(a\in\overline F\). Indeed a parametrization
by the additive group is equivariant for scaling the parameter
by \(t_i/t_j\). A term \(a^m\) in an entry must have that
entry's torus character equal to \((t_i/t_j)^m\). Among
the ambient characters only \(t_i/t_j\) itself is a positive
multiple of this root in additive weight notation. Hence the
parametrization has only the indicated linear term, which is
nonzero by the inclusion of its tangent space. Roots occur in
opposite pairs. The graph of these pairs on the three coordinates
is connected, since otherwise the torus and all root groups
preserve a proper coordinate subspace and together generate
\(G^0\). A connected graph gives all three pairs using
commutators of elementary matrices. The groups for these
pairs generate \(\mathrm{SL}_3\).

In rank one the nonidentity Weyl element inverts a maximal
torus. The trace on the representation therefore equals the
trace on its dual on all conjugates of that torus, a dense
subset of \(G^0\). It agrees everywhere. These traces identify
the two irreducible representations in either coefficient
characteristic. Indeed, if two absolutely irreducible modules
are inequivalent, the image of the group algebra on their
direct sum is a product of two full matrix algebras. Its
radical annihilates both simple modules and is zero in this
faithful image, so this follows from the semisimple algebra
decomposition. A rank-one projector on just the first factor
would distinguish the traces already on a linear combination
of group elements. Thus there is a nondegenerate invariant
bilinear form, unique up to scalar by Schur's lemma. Its
transpose is a multiple by a scalar of square one. Since the
dimension is odd and the characteristic differs from two,
the form is symmetric. A connected semisimple group of rank
one has dimension three by its root decomposition. Its closed
image therefore equals \(\mathrm{SO}(S)\), also connected of
dimension three.

The group \(G\) normalizes \(G^0\) and hence preserves the
line of \(S\). Write \(U(h)^{\mathsf T}SU(h)=\nu(h)S\).
The determinant gives \(\nu^3=1\). From
\(U\simeq U^\vee\otimes\nu\), reduction in stable lattices
gives the same isomorphism on residual semisimplifications.
Here and below lattice independence follows from equality of
the characteristic polynomials after reduction and the
Brauer--Nesbitt theorem on the common finite residual quotient.
In particular \(\bar\nu\) is one of \(1,\bar\Psi,\bar\Psi^{-1}\)
on \(H\), since the residual dual has a trivial constituent.
The order on inertia excludes the last two choices. Reduction
is injective on cube roots of unity at \(\ell>3\), so
\(\nu=1\). We have \(G^0\subset G\subset\mathrm{SO}(S)=G^0\),
as asserted.
\end{proof}

\begin{lemma}[An integral lift of the projective representation]
\label{lem:spin122}
In the orthogonal case, after a finite coefficient extension
there are an integer \(e\in\{0,1\}\) and a continuous
representation \(\rho:G_\QQ\to\mathrm{GL}_2(R)\) satisfying
\begin{gather}
 \operatorname{Ad}^0\rho|_H\simeq U,\qquad
 \phi:=\det\rho=\eta^e\Psi,\qquad
 \bar\rho=\operatorname{diag}(\bar\phi,1).                 \label{eq:spin122}
\end{gather}
Here \(\operatorname{Ad}^0\) is conjugation on trace-zero
two-by-two matrices. The parity \(e\) is determined by the
marked representation \(U\). For \(e=0\) one has \(\rho(c)=I\).
For \(e=1\) its eigenvalues at \(c\) are \(1,-1\).
\end{lemma}
\begin{proof}
The symmetric form \(S\) can be defined over a finite extension.
With the marked involution \(J_3\) the relation on \(H\) becomes
\[
       U(chc)=J_3 U(h)^{-\mathsf T}J_3
             = A U(h)A^{-1},\qquad A=J_3 S .
\]
Since \(c^2=1\), irreducibility gives \(A^2=bI\) for a
nonzero scalar \(b\). Symmetry of \(S\) and \(J_3\) gives
\(A^{\mathsf T}S A=S A^2=bS\). Put
\(C=(b/\det A)A\). Taking determinants shows
\((\det A)^2=b^3\), so \(C\in\mathrm{SO}(S)\) and \(C^2=I\).
The semidirect product \(G_\QQ=H\rtimes\langle c\rangle\)
thus extends \(U\) into this orthogonal group, with \(c\)
acting by \(C\). Multiplying \(S\) by a scalar has no effect
on \(C\). It is the unique element of \(\mathrm{SO}(S)\)
implementing this conjugation, since the centralizer of \(U(H)\)
consists of scalars and a scalar in \(\mathrm{SO}_3\) is one.

Over a finite extension identify this split orthogonal group
with \(\mathrm{PGL}_2\) acting on \(\mathfrak{sl}_2\) by
conjugation. For this identification the trace pairing on
\(\mathfrak{sl}_2\) is nondegenerate since the characteristic
differs from two. Conjugation is faithful on \(\mathrm{PGL}_2\)
and identifies it with the full special orthogonal group, as
one can also see by recovering the matrix algebra from the
pairing and bracket in Theorem~\ref{thm:mon44}. Write \(w\)
for the resulting continuous projective representation of
\(G_\QQ\).

We give the lattice argument since \(S\) in the original
marked basis need not be unimodular. The local field \(F\)
has finitely many square classes by the valuation, finiteness
of its residue field and Hensel's lemma. Adjoin square roots
of their representatives. Every element of the original
\(\mathrm{PGL}_2(F)\) now has a determinant-one lift.
The preimage of the compact group \(w(G_\QQ)\) in
\(\mathrm{SL}_2\) over the larger local field is compact.
Indeed the determinant-one covering is finite and proper.
Equivalently one may cover the compact image by finitely
many compact pieces in matrix charts, normalizing one entry
to one. On each piece its determinant and the inverse
determinant are bounded. Its determinant-one representatives
are therefore bounded as well, and their union is closed.
The span over the integer ring of the orbit of a lattice under
this compact group is a stable lattice, since it contains the
original lattice and is contained in a fixed scalar multiple
of it. Choosing a basis makes \(w\) integral projectively.

Reduce in that basis. On \(H\), the adjoint of \(\bar w\)
has semisimple constituents \(\bar\Psi,1,\bar\Psi^{-1}\) by
lattice independence. In a basis adapted to a composition
series, the image of \(\ker(\bar\Psi|_H)\) therefore consists
of unipotent upper triangular matrices. Its projective image
is an \(\ell\)-group because the adjoint is faithful and the
residue field is finite. This subgroup is normal in the whole
residual projective image. If nontrivial, it fixes a unique
line of the two-dimensional projective representation. Indeed
a central element of order \(\ell\) is nontrivial projectively
unipotent and has a unique eigenline, which its centralizer
preserves. The unique common fixed line is preserved by
every normalizer. If the subgroup is trivial, \(\bar w(H)\)
is cyclic of order equal to that of \(\bar\Psi(H)\), as the
adjoint diagonal characters show. Conjugation by \(c\)
acts trivially on this cyclic quotient because \(\bar\Psi\)
is global. The centralizer of a projective semisimple generator
of order greater than two preserves both of its eigenlines
individually. Thus in both cases \(\bar w(G_\QQ)\) has an
invariant line over a finite residue extension.

Use this line to make the projective reduction upper triangular.
Integral invertible representatives then have their lower-left
entries in the maximal ideal. Conjugate by
\(\operatorname{diag}(t,1)\) with \(0<v(t)<v(\varpi_F)\)
after a finite ramified extension. Both off-diagonal entries
of every representative now belong to the maximal ideal,
and determinants stay units. The residual projective image
is diagonal and its ratio character on \(H\) is
\(\bar\Psi\) or \(\bar\Psi^{-1}\) by the same adjoint
constituents. Exchange the axes if necessary. On \(G_\QQ\)
this ratio must equal \(\bar\eta^e\bar\Psi\) for exactly
one \(e\in\{0,1\}\), since \(G_\QQ/H\) is quadratic.
The nontrivial projective involution at \(c\), if present,
stays nontrivial after reduction. The projective congruence
kernel is pro-\(\ell\), as follows from successive matrix
congruences modulo scalars. Hence \(e=0\) exactly when \(C=I\).

For each \(g\) scale an integral projective representative to
have reduction exactly \(\operatorname{diag}(\bar\phi(g),1)\).
The quotient of \(\phi(g)\) by its determinant is a principal
unit, with a unique principal square root. Scaling by that
root gives the unique matrix with the prescribed determinant
and residual matrix. Uniqueness makes the result multiplicative.
Continuity follows either from square roots on principal units
or from the continuous inverse of
\(D\mapsto([D],\det D)\) on the compact set of integral
matrices with diagonal reduction and lower-right residue one.
This map is injective there because a scalar of square and
residue one equals one. This constructs \(\rho\).
Since \(\Psi(c)=1\), its determinant at \(c\) is \((-1)^e\).
The eigenvalues of an involution and their reductions now give
the stated action. This construction applies over a Laurent
series coefficient field of characteristic \(\ell\) as well.
\end{proof}

\begin{theorem}[Local characters and crystallinity of the lift]
\label{thm:local122}
The lift \(\rho\) is absolutely irreducible on \(H\).
On \(G_{\QQ_\ell}\) it has a stable line of inertia character
\(\Psi\), and the quotient is unramified. Away from \(\ell\)
it kills \(I_p\cap\ker\eta\cap\ker\alpha\). In characteristic
zero it is crystalline at \(\ell\) with cyclotomic exponents
\(r,0\), and its prime-to-\(\ell\) Artin conductor is exactly
the conductor of \(\eta^e\alpha\).
\end{theorem}
\begin{proof}
If \(\rho|_H\) preserved a line \(L\), the endomorphisms
with image in \(L\) and vanishing on \(L\) would form an
invariant line of \(\mathfrak{sl}_2\). This contradicts the
irreducibility of \(U\).
On \(I_\ell\), the first line of the ordinary flag in \(U\)
has character \(\Psi\) and is distinguished by the simple
eigenvalues at \(g_*\). Transfer this line into
\(\mathfrak{sl}_2\). A generator \(X\) satisfies
\(\operatorname{tr}(X^2)=0\), because conjugation preserves
that trace and \(\Psi^2\) is nontrivial on inertia. Hence
\(X\) is nilpotent of rank one, since
\(\operatorname{tr}X=0\) and two is invertible. Its image
equals its kernel. The line of such matrices is stable
under \(G_{\QQ_\ell}\), since conjugating the inertia group
preserves its \(\Psi\)-character and the \(\Psi(g_*)\)-eigenspace
is one-dimensional. Its kernel therefore gives a stable line
for \(\rho\) on the whole decomposition group.

Write \(a,b\) for the characters on this line and its quotient.
Conjugation on the nilpotent direction has character \(a/b\).
Thus \(a/b=\Psi\) on \(I_\ell\) and \(ab=\phi=\Psi\) there,
so \(b^2=1\) on inertia. The two residual characters there
are \(\bar\Psi,1\). Reduction of the line and quotient
gives these characters in some order, since the intersection
of a field-valued line with a lattice is saturated. The order
of \(\bar\Psi(I_\ell)\) exceeds two, so \(\bar b=1\) there.
It follows that \(b=1\) on inertia, and \(a=\Psi\) there.

In characteristic zero restrict to \(G_K\), with
\(K=\QQ_\ell(i)\) in the chosen completion. By
Theorem~\ref{thm:crystalline121} the trace-zero adjoint is
crystalline there for every \(r\). In a flag basis write
\(\rho(g)=\left(\begin{smallmatrix}a(g)&b(g)z(g)\\0&b(g)\end{smallmatrix}\right)\).
Its action on \(E_{12}\) is by \(a/b\) and it sends
\(\operatorname{diag}(1,-1)\) to
\(\operatorname{diag}(1,-1)-2z(g)E_{12}\). The span of these
two matrices is consequently isomorphic to \(\rho\otimes b^{-1}\)
by rescaling the off-diagonal basis vector. It is a subrepresentation
of a crystalline representation, and \(b\) is unramified.
This proves crystallinity of \(\rho|_{G_K}\). As \(K/\QQ_\ell\)
is unramified of degree one or two, semilinear descent of
crystalline periods gives crystallinity over \(\QQ_\ell\).
The two exponents follow from the diagonal characters.

If \(g\in I_p\cap\ker\eta\cap\ker\alpha\) with \(p\ne\ell\),
then \(U(g)=1\) and \(\phi(g)=1\). Thus \(\rho(g)\) is a
scalar of square one reducing to \(1\), so is the identity.
The residual inertia characters at such a prime are
\(\overline{\eta^e\alpha},1\). The full inertia factors
through the finite abelian image of \((\eta,\alpha)\), of
order prime to \(\ell\). Reduction identifies its characters
with their prime-to-\(\ell\) lifts. In characteristic zero
the two characters on inertia are therefore \(\eta^e\alpha,1\),
with zero monodromy. This gives the asserted conductor,
including the possible wild conductor at \(p\).
\end{proof}

\subsection{Frobenius primes for the full invariant ring}\label{nav:IV-23}
\label{sec:frob124}

Retain the period field \(E\), \(x\in[V,V^{3/2}]\),
\(\ell>\exp(x^\varepsilon/12)\), \(1\le r<t\), and \(\alpha\)
of Section~\ref{sec:origin116}. Write \(H=G_{\QQ(i)}\),
\(\Psi=\kappa_\ell^r\alpha\), and \(M|_H\) for the universal
matrix on the ordinary component. The invariant subring
\(D\subset B\) has the augmentation \(\epsilon:D\to\OO\)
given by the diagonal representation. Set \(I=\ker\epsilon\),
write \(\mathfrak m_D,\mathfrak m_B\) for the two maximal
ideals, and let \(q=\#k\) be the cardinality of their common
residue field. For \(h\in H\) put
\[
 t_h=\operatorname{tr}M(h)-\bigl(\Psi(h)+1+\Psi(h)^{-1}\bigr)
       \in I .
\]
An oriented split prime is a rational prime split in \(\QQ(i)\)
together with a choice of one place of \(\QQ(i)\) over it.
We take arithmetic Frobenius at that place. Outside the
ramification set its \(H\)-conjugacy class determines \(t_h\).

\begin{theorem}[Frobenius generators before specialization]
\label{thm:ideal124}
Put \(d=\dim_k I/\mathfrak m_D I\). Given any finite set of
excluded primes, there exist \(d\) oriented split primes
outside it and outside \(2N\ell\) for which
\begin{equation}
                 I=(t_{\Frob_{v_1}},\ldots,t_{\Frob_{v_d}}).
                                                        \label{eq:ideal124}
\end{equation}
They can be selected from conjugacy classes in a single finite
Galois extension determined by the universal matrices modulo
\(\mathfrak m_B^4\). If \(s\) is the number of parameters in
the presentation of Lemma~\ref{lem:algebra116}, this extension
has degree at most \(2q^{9\binom{s+4}{3}}\) over \(\QQ\).
If \(s_j\) parameters have weight \(j\), the length \(d\) of
the list is bounded by
\[
 s_0+s_1s_{-1}+s_2s_{-2}
       +\binom{s_1+1}{2}s_{-2}+\binom{s_{-1}+1}{2}s_2 .
\]
\end{theorem}
\begin{proof}
The quotient of \(D\) by the ideal generated by all \(t_h\)
is topologically generated by the scalars.
Indeed \(D\) is topologically generated by the traces, by
Theorem \ref{thm:trace71}.
Every ideal of a complete Noetherian local ring is closed,
by applying Krull intersection to its quotient. The image of
\(\OO\) in this quotient is closed by compactness. Thus the
scalars fill the quotient. Augmentation splits the scalar
inclusion, proving that the ideal generated by all \(t_h\)
is \(I\). Choose \(h_1,\ldots,h_d\) whose trace differences
give a basis of \(I/\mathfrak m_DI\).

We verify what is detected by the truncated matrices. In the
ambient presentation \(P=\OO[[X_1,\ldots,X_s]]\), take each
parameter to vanish at the diagonal point as in
Lemma~\ref{lem:algebra116}. The weights belong to
\(\{0,1,-1,2,-2\}\), and the uniformizer has weight zero.
A weight-zero monomial of total degree at least four, counting
the uniformizer among its factors, splits into two nonempty
weight-zero monomials. A zero weight or a pair of opposite
weights can be taken out first. If neither occurs, the zero
sum consists of triples of weights \(1,1,-2\) or \(-1,-1,2\).
This also proves that the invariant monomials of degrees one,
two and three generate the invariant power series ring
topologically. If \(y\in D\cap\mathfrak m_B^4\), lift it into
the fourth power of the ambient maximal ideal. Taking the
weight-zero part still maps to \(y\), since the relations are
homogeneous and extraction of each weight commutes with passage
to \(B\), by Lemma~\ref{lem:algebra116}. The monomial
factorization and closedness of \(\mathfrak m_D^2\) give
\[
                    D\cap\mathfrak m_B^4\subset\mathfrak m_D^2 .
\]
The single variables of weight zero, the opposite pairs, and
the two types of triples in the parameter variables all vanish
at the augmentation. Together with \(\OO\) they generate \(D\)
topologically. Their classes span \(I/\mathfrak m_D I\), which
gives the stated count.
The augmentation extends to \(B\) at the split point and maps
\(\mathfrak m_B^4\) into \(\varpi^4\OO\), where \(\varpi\) is
a uniformizer. Also \(I\cap\mathfrak m_D^2=\mathfrak m_D I\).
Indeed \(\mathfrak m_D=\varpi\OO\oplus I\) as an
\(\OO\)-module and its square is
\(\varpi^2\OO\oplus(\varpi I+I^2)\). Hence \(M(h)\) modulo
\(\mathfrak m_B^4\) determines \(t_h\) modulo
\(\mathfrak m_DI\).

The kernel in \(H\) of this finite matrix representation is
normal in \(G_\QQ\), since conjugation by \(c\) acts on the
matrices by the marked inverse transpose. Its fixed field
\(L\) contains \(\QQ(i)\) and is unramified outside \(2N\ell\).
The maximal ideal of \(P\) has \(s+1\) generators including
\(\varpi\). There are \(\binom{s+4}{3}\) monomials of degree
at most three in them. Their graded pieces bound
\(\#(B/\mathfrak m_B^4)\) by \(q^{\binom{s+4}{3}}\).
Counting the nine entries of a matrix and the two cosets of
\(H\) gives the claimed degree bound.

For each \(h_i\) the Chebotarev theorem supplies infinitely
many rational primes with its conjugacy class in \(\Gal(L/\QQ)\).
They split in \(\QQ(i)\). Choosing a prime of \(L\) above
each one realizes the specified element of the class, fixing
the orientation. Its trace difference has the required residue
in \(I/\mathfrak m_DI\). Nakayama's lemma now gives \eqref{eq:ideal124}.
For \(d=0\) the same lemma gives \(I=0\).
\end{proof}

\begin{corollary}[Size and specialization of the prime list]
\label{cor:size124}
Let \(S\) contain \(2,3,5\), the divisors of \(N\ell\), and
any further primes to be excluded, and put \(b_S=\prod_{p\in S}p\),
\(n_*=4q^{9\binom{s+4}{3}}\). The oriented primes in \eqref{eq:ideal124}
can all be taken with
\begin{equation}
                \log v_i\ \le\ C n_*\log(b_S n_*)
                                                        \label{eq:least124}
\end{equation}
for an effective absolute constant \(C\). Two orientations may
use the same underlying rational prime in this bound.
For any integral specialization into a valuation ring \(R\)
with valuation \(v_R\), the same list computes the full trace
ideal. More precisely,
\[
   \inf_{h\in H}v_R(t_h)=\min_{1\le i\le d}v_R(t_{\Frob_{v_i}}),
\]
where the elements on both sides are evaluated in \(R\);
the value for the zero ideal or the empty list is \(+\infty\).
\end{corollary}
\begin{proof}
Adjoin \(\sqrt{b_S}\) to the field \(L\) of the theorem. The
resulting Galois field \(L'\) has degree \(n\le n_*\). Every
prime of \(S\) ramifies in the quadratic subfield since \(b_S\)
is squarefree and even. The entire extension is unramified
outside \(S\). For \(p\in S\), write \(e\) for the ramification
index in \(L'\). The local different exponent is at most
\(e-1+e\,v_p(e)\). To see this, pass to the maximal unramified
subextension of the completion. A uniformizer of the remaining
totally ramified extension generates the integer ring, and
its Eisenstein polynomial has degree \(e\). The nonzero
terms of its derivative at the uniformizer have valuations
in distinct residue classes modulo \(e\). Its valuation
is therefore at most that of the leading derivative term,
\(e-1+e\,v_p(e)\). The different is generated by this
derivative. Summing over the primes over \(p\) gives
\(v_p(|\operatorname{disc} L'|)\le n(1+v_p(n))\).
It follows that \(\log|\operatorname{disc} L'|\le
n(\log b_S+\log n)\).

Lift each required Frobenius element of \(\Gal(L/\QQ)\)
to \(\Gal(L'/\QQ)\). The unconditional least-prime theorem
of Lagarias, Montgomery and Odlyzko applied to its conjugacy
class gives an unramified rational prime
\(\ll |\operatorname{disc} L'|^{C_0}\) with an effective
absolute \(C_0\). Such a prime automatically avoids \(S\).
Orient it as in the theorem. This proves \eqref{eq:least124}. For the last
assertion \eqref{eq:ideal124} gives equality of the generated ideals after
every integral specialization. The valuation of an ideal
generated by finitely many elements in a valuation ring is
the least valuation of any generator.~\cite{VI-IV-VI-lmo124}
\end{proof}

\begin{theorem}[Simultaneous separation of irreducible points]
\label{thm:dense124}
Fix finitely many irreducible characteristic-zero points or
realizations of nonconstant trace curves from
Theorem~\ref{thm:cyclo119}, allowing different coefficient
primes, twists and period fields from the prescribed family.
Among rational primes split in \(\QQ(i)\), a set of relative
natural density one has nonzero trace difference at both
places for every member of the list. The same holds within
any nonempty finite Frobenius condition compatible with
splitting, interpreted on its split part. In \eqref{eq:ideal124} each prime
can therefore also be chosen to separate this entire finite
list, without the size assertion \eqref{eq:least124}.
\end{theorem}
\begin{proof}
Fix one representation \(U\) over its local field. Consider
the compact image \(\Gamma\) of \((U,\Psi)\) on \(H\). Over
the algebraic closure of the coefficient field let \(G\)
denote the algebraic monodromy of \(U\) and \(C_\Psi\) the
Zariski closure of \(\Psi(H)\) in the multiplicative group.
Theorem~\ref{thm:orth122} identifies \(G\) with
\(\mathrm{SL}_3\) or \(\mathrm{SO}_3\). Restrict to
\(G_{K'E(\mu_{\ell^\infty})}\). Here \(\Psi=1\) and, by
Theorem~\ref{thm:cyclo119}, the closure of the matrix image
contains \(G\). Thus the closure of \(\Gamma\) contains
\(G\times1\). Multiplying by this subgroup shows that its
closure is the full \(G\times C_\Psi\).
On each connected component the regular function
\[
                     (g,z)\longmapsto
                 \operatorname{tr}(g)-z-1-z^{-1}
\]
is nonconstant: in each of the two possible groups \(G\)
a one-parameter torus has eigenvalues \(a,1,a^{-1}\).
Consequently its zero set contains no component.

Here is a measure argument valid in both coefficient
characteristics. If a compact linear group is Zariski dense
in an algebraic group, its intersection with a closed
subvariety containing no component has Haar measure zero.
Pass first to the identity component and translate each
of the finitely many cosets back to it. In the connected
case every open subgroup of the compact group is Zariski
dense: its finitely many translates cover a dense subgroup
of an irreducible variety. The same is true of every nonempty
open subset, since here the compact groups are profinite.
If a proper closed subvariety has positive measure, choose an
irreducible one \(Z\) of least dimension with that property.
Its intersection \(A\) with the compact group is measurable.
Continuity of translations in \(L^1\) implies
\(\mu(A\cap uA)>0\) for all \(u\) in a neighborhood of
the identity. For instance approximate \(1_A\) in \(L^1\)
by functions on finite quotients. Some \(u\) in that
neighborhood satisfies \(uZ\ne Z\). Otherwise the orbit of
any point of \(Z\) under the neighborhood is dense in the
whole connected group and is contained in \(Z\). The closed
set \(Z\cap uZ\) has smaller dimension and positive measure,
contrary to the choice of \(Z\). This proves the claim.

Pulling back the zero set from \(\Gamma\) gives a closed Haar
null set in \(H\). Take the union for the finitely many
representations and for their conjugates by \(c\), to treat
both choices of place. The union is conjugacy invariant in
\(G_\QQ\), closed and null. Its images in finite Galois
quotients have inverse images of arbitrarily small measure:
the clopen neighborhoods of the closed set decrease to it
and Haar probability is regular. On each finite quotient the
Chebotarev theorem identifies the density of the corresponding
unramified primes with this measure. It follows that the
exceptional primes have upper natural density zero. Any fixed
nonempty compatible finite Frobenius condition on split primes
has positive density. Removing the exceptional primes preserves
its density, which proves all the assertions.
\end{proof}

Equations \eqref{eq:ideal124} and \eqref{eq:least124} concern the original invariant
ring even when it has positive-dimensional special fiber.
They give a finite test for congruence with the diagonal.
The size bound depends on the coefficient prime, and it bounds
the sampling primes. A uniform bound for the sum of trace
valuations on the
remaining irreducible loci is still required in the intersection
estimate of Theorem~\ref{thm:full36}.

\section{Trace congruences on the common interval}\label{nav:g-IV-4}

\subsection{The traces in the marked deformation ring}\label{nav:IV-25}
Fix \(x\) on one of the common intervals of
Theorem~\ref{thm:common16}, put \(t=\lfloor x^{401}\rfloor\), and
take a coefficient prime \(\ell>\exp(x^\eps/12)\) and a twist
\(1\le r<t\). Enlarge the lower endpoint so that \(\ell>2t+3\).
Use the character summands of the strict group for the field
\(E_x\) and the auxiliary quadratic character \(\eta\) of
\(K'=\QQ(i)\) from
Sections \ref{nav:IV-13}--\ref{nav:IV-15}.
Write \(O\) for the coefficient discrete valuation ring,
\(\varpi\) for a uniformizer and \(k\) for its finite residue field.
In particular \(\Psi=\kappa_\ell^r\alpha\), with \(\alpha\)
of order prime to \(\ell\) and unramified at \(\ell\).
For a complex conjugation \(c\), we have \(\eta(c)=-1\) and
\(\Psi(c)=1\).

Let \(B\) be the completed marked global deformation ring.
Its base matrices are
\(d(g)=\operatorname{diag}(\Psi(g),1,\Psi(g)^{-1})\).
It represents lifts with the prescribed sign \(\eta\), the
inertia kernels away from \(\ell\), and the ordinary component
through the basepoint at \(\ell\) described in Section \ref{nav:IV-13}. Write \(M(g)\) for the matrix coordinate of the
universal lift. On \(H=G_{K'}=\ker\eta\) it is a homomorphism.
We use absolute Galois groups and pull back matrices from their
specified unramified-outside-finitely-many-places quotients.
The element \(g_*\in I_\ell\subset H\) is chosen so that the
residual entries of \(d(g_*)\) are distinct. Let \(J\) reverse
the three coordinates. The sign acts by
\(A\mapsto J A^{-\mathsf T}J\).
For rows numbered \(0,1,2\), the marking at \(c\) sets the
entries \(00,11,01,10,02,20\) of \(M(c)\) equal to those of
the identity. Since \(c^2=1\), the group law gives
\(M(c)=J M(c)^{\mathsf T}J\). Hence \(M(c)=1\).
The other two marking equations set the entries \(01,10\) of
\[
                  Z=M(g_*)-M(cg_*^{-1}c)
\]
equal to zero. The split torus acts by conjugation with
\(\operatorname{diag}(a,1,a^{-1})\). Let \(D\) be the
invariant subring of \(B\). Invariance means equality of the two
Laurent expressions for an indeterminate \(a\) on every Artinian
truncation. Let \(\mathcal T\) denote the \(O\)-algebra in \(B\)
generated by \(T(h)=\operatorname{tr}M(h)\), \(h\in H\).
All closures below use the maximal-ideal topology of \(B\).

\begin{lemma}[Matrix entries around a cycle]\label{lem:cycle71}
Put \(z=M(g_*)_{00}-M(g_*)_{22}\).
Then \(z\) is a unit, \(z^2\in\mathcal T\), and every product of
matrix entries around a directed cycle of indices belongs to
\(\mathcal T[1/z^2,z]\). For this statement each edge may use
the matrix of a different group element of \(H\).
\end{lemma}
\begin{proof}
The group law and the marking at \(c\) give
\[
               M(ch^{-1}c)=J M(h)^{\mathsf T}J\qquad(h\in H).
\]
Consequently \(Z=-J Z^{\mathsf T}J\).
The two zero entries and this symmetry give
\(Z=\operatorname{diag}(z,0,-z)\) since 2 is a unit.
Its residue has \(z\ne0\), by the choice of \(g_*\).
Expanding the trace of the square of the defining difference
shows that \(z^2=\operatorname{tr}(Z^2)/2\in\mathcal T\).
The diagonal matrix units are
\[
 P_0=(Z^2+zZ)/(2z^2),\qquad P_1=1-Z^2/z^2,\qquad
 P_2=(Z^2-zZ)/(2z^2).
\]
For indices \(i_0,\ldots,i_n=i_0\), the entry product is
\[
 \prod_{j=0}^{n-1} M(h_j)_{i_j i_{j+1}}
    =\operatorname{tr}\big(P_{i_0}M(h_0)P_{i_1}M(h_1)
                            \cdots P_{i_{n-1}}M(h_{n-1})\big).
\]
Insert the three polynomial formulas. After expansion, the
trace of every matrix word is a trace of a single group element
of \(H\), or the trace of the identity, since \(g_*\) and
\(cg_*^{-1}c\) both belong to \(H\). This proves the assertion.
\end{proof}

\begin{theorem}[The invariant ring is generated topologically by traces]\label{thm:trace71}
Let \(A\) be the \(O\)-subalgebra of \(B\) generated by the
entries \(M(h)_{ij}\), \(h\in H\). Write \((A[1/z^2])_0\)
for its weight-zero part after the indicated localization. Then
\[
          (A[1/z^2])_0=\mathcal T[1/z^2,z],
          \qquad D=\overline{\mathcal T}.
\]
The algebra in the first equality is generated as a module by
\(1,z\) over \(\mathcal T[1/z^2]\). The closed algebra \(D\)
is complete Noetherian local, with residue field \(k\).
\end{theorem}
\begin{proof}
Write \(\bar i=2-i\). An entry \(M(h)_{ij}\) equals
\(M(ch^{-1}c)_{\bar j,\bar i}\) and has weight \(j-i\).
An entry of weight \(1\) may thus be used in either of the
positions \(01,12\), and a weight \(-1\) entry in either
of the positions \(10,21\). Weight \(2\) has position
\(02\) and weight \(-2\) has position \(20\).

Any list of weights in \(\{0,\pm1,\pm2\}\) with sum zero
splits into singles \((0)\), pairs \((1,-1),(2,-2)\),
and triples \((1,1,-2),(-1,-1,2)\). Remove singles and
opposite pairs. If anything remains, its positive and negative
weights must have different absolute sizes, one on each side.
The sum condition now partitions the remainder into triples.
The reflection identity in the first paragraph puts each pair
or triple around a directed cycle. Multiplication in the
coefficient ring allows us to reorder the entries. Hence
Lemma~\ref{lem:cycle71} applies to each weight-zero monomial.
Every invariant polynomial is the image of its weight-zero
part, by uniqueness of Laurent coefficients for the torus.
Localizing at the invariant \(z^2\) preserves this description.
For the opposite inclusion, traces, \(z\) and \(1/z^2\) are
invariant. The module assertion follows from \(z^2\in\mathcal T\).

Choose \(z_0\in O^\times\) lifting the residue of \(z\).
Then \(y=z^2/z_0^2-1\) is topologically nilpotent.
In \(B\), the geometric series inverts \(1+y\), and
\(z=z_0\sum_{n\ge0}\binom{1/2}{n}y^n\).
Indeed the binomial coefficients belong to \(\ZZ_\ell\), by
continuity from nonnegative integral exponents. The two sides
have the same square and their sum has unit residue \(2\bar z_0\).
Thus \(z\) and \(1/z^2\) belong to \(\overline{\mathcal T}\).

The framed presentation uses finitely many differences of entries
from their base values as topological generators. To recall why,
on the kernel of the finite residual image the congruence
subgroup over every finite Artin coefficient ring is an
\(\ell\)-group. Its universal source is the maximal pro-\(\ell\)
quotient of that kernel with the allowed ramification. This
quotient is finitely generated by the pro-\(\ell\) generator
criterion and the finiteness of the first cohomology with finite
trivial coefficients over a number field unramified outside a
fixed finite set. The latter follows from class field theory,
as in the proof of the framed presentation in Section \ref{nav:IV-09}. Add finitely many lifts of the
residual image. The entries of this list determine all the
coordinates by multiplication, inversion and limits. Inverses
of the unit determinants expand about their base values.
Elements of negative sign give the same matrices as their
products on the right with \(c\), so all entries may be taken
on \(H\). The entry differences are homogeneous, with
base constants subtracted only in weight zero.

This yields a surjection from a power series ring with finitely
many weighted variables onto \(B\). On any truncation an
invariant element of the target lifts to a polynomial, and
taking the weight-zero part still lifts that element. The
weight-zero monomials are generated by the finitely many
singles, pairs and triples of the variables just described.
The power series ring over \(O\) on those invariant monomials
maps continuously into \(B\) because each monomial vanishes
at the residual point. Its image is compact, since its
coefficients have finite residue field and its Artinian
quotients are finite. It is therefore closed in \(B\).
The polynomial truncations prove that its image is all of \(D\).
This also identifies the induced topology on \(D\) with the
quotient topology of this compact power series ring. In
particular \(D\) has the asserted local properties. The
algebraic equality already proved puts each of its polynomial
generators in \(\overline{\mathcal T}\). Traces are invariant,
and invariance is closed on every truncation. This proves the
second equality.
\end{proof}

\begin{proposition}[A test on curves over finite fields]\label{prop:curve71}
The \(O\)-module \(D\) is finite if and only if for every finite
extension \(k'/k\) and continuous local \(k\)-algebra map
\[
                   u:D/\varpi D\longrightarrow k'[[s]]
\]
all \(u(T(h))\) are constant. In these expressions a trace in
\(D\) is reduced modulo \(\varpi\) before applying \(u\).
\end{proposition}
\begin{proof}
Put \(D_k=D/\varpi D\). If \(D\) is finite, \(D_k\) is an
Artinian local \(k\)-algebra and its maximal ideal is nilpotent.
Every such map into a domain factors through \(k\). Conversely,
if \(D_k\) is Artinian, a power of the maximal ideal of \(D\)
is in \(\varpi D\). Thus the maximal-ideal and
\(\varpi\)-adic topologies on \(D\) agree. Lift a finite
\(k\)-basis of \(D_k\), expand successively modulo powers of
\(\varpi\), and use completeness of \(O\) and \(D\).
The lifts generate \(D\) over \(O\).

If \(D_k\) has positive dimension, choose a prime quotient
\(C\) of dimension one along a maximal chain of primes. It is
a complete Noetherian local domain. Its normalization is finite
because complete Noetherian local rings are Nagata.
It is local: a finite algebra over a complete local ring splits
as a product of local algebras, and the normalization is a
domain. Being one-dimensional and normal, it is a discrete
valuation ring, complete with finite residue extension \(k'/k\).
The roots of \(X^{|k'|}-X\) give its field of constants by
Hensel's lemma, and expansion in a uniformizer identifies it with
\(k'[[s]]\). The finite-normalization theorem used here is proved in the
\href{https://stacks.math.columbia.edu/tag/032E}{Stacks Project on Nagata rings}.
The composition \(D_k\twoheadrightarrow C\hookrightarrow k'[[s]]\)
is local, continuous and has a nonconstant image, since
every nonzero element of the maximal ideal of \(C\) has
strictly positive valuation. If it sent all traces to constants,
Theorem~\ref{thm:trace71} and closedness of the finite constant
field would make its entire image constant. This proves the
criterion.
\end{proof}

\begin{theorem}[Rigidity after restriction to one-dimensional characters]\label{thm:chars71}
For a map \(u\) as above, suppose there are a number field
\(L\supset K'\) finite over \(K'\), a finite extension \(F\)
of \(k'((s))\), and three continuous characters
\(\lambda_i:G_L\to F^\times\) with
\[
                 u(T(h))=\sum_{i=1}^3\lambda_i(h)
                 \qquad(h\in G_L).
\]
Then \(u(T(h))\) is constant for every \(h\in H\), and \(u\)
factors through the residue field \(k\).
\end{theorem}
\begin{proof}
Let \(P_h(X)\) be the characteristic polynomial of \(M(h)\).
Its coefficients are in \(D\). Since \(\ell>3\), Newton's
identities express these coefficients in terms of \(T(h)\),
\(T(h^2)\), and \(T(h^3)\). Thus for \(h\in G_L\),
the roots of \(u(P_h)\) are exactly the \(\lambda_i(h)\),
with multiplicities.

At a finite place of \(L\) over \(\ell\), the ordinary equations
give the three eigenvalues on inertia after reduction:
they are \(\bar\Psi(h),1,\bar\Psi(h)^{-1}\). Moving the place
by conjugation either preserves this multiset or takes its
inverse, which is the same multiset. At all other finite places
the imposed inertia kernel has finite index in inertia, since
the specified finite characters have finite image there.
Hence for every such inertia element \(h\), some \(M(h^n)\)
is the identity. Newton's identities at \(h^n\) give
\(\lambda_i(h)^n=1\) for all three characters. In particular each \(\lambda_i\) takes
roots of unity on inertia at every finite place of \(L\).

Every \(\lambda_i\) has compact image, so its valuation has
finite image in \(\ZZ\), necessarily zero. Let \(k_F\) be the
finite residue field of \(F\). Divide \(\lambda_i\) by its
reduction lifted through the constant field \(k_F\subset F\).
The quotient is a continuous character into the principal units
of \(F\). This group is torsion-free: if \(j\) is prime to
\(\ell\), then \((1+b)^j-1\) has valuation \(v(b)\) for
\(v(b)>0\), and raising to an \(\ell\)-power sends \(1+b\)
to \(1+b^{\ell}\). Hence the quotient is unramified at all
finite places. The field \(L\) is totally imaginary. By the
Hilbert class field theorem its unramified abelian Galois group
is the finite ideal class group. Thus the quotient character
is trivial, and every \(\lambda_i\) takes values in \(k_F^\times\).
The class field theorem used here is recalled in
\href{https://www.jmilne.org/math/CourseNotes/CFT.pdf}{Milne, \emph{Class Field Theory}}.

For arbitrary \(h\in H\), choose \(m>0\) with \(h^m\in G_L\);
this follows from the action on the finitely many cosets.
The roots of \(u(P_{h^m})\) are the \(m\)-th powers of the
roots of \(u(P_h)\). Indeed the coefficients of the polynomial
of the powered roots are symmetric integral polynomials in the
original roots. The resulting identity for matrices holds first
for the generic matrix with distinct eigenvalues and hence
identically over \(\ZZ\); it commutes with specialization
through \(D\). The roots of \(u(P_{h^m})\) lie in a finite
field, so the roots of \(u(P_h)\) are algebraic over a finite
field. Their sum in \(k'[[s]]\) is constant: any element of
that ring algebraic over a finite field satisfies
\(a^{q}=a\) for some prime power \(q>1\), forcing its
positive-order part to vanish. All specialized traces are
therefore constant. Theorem~\ref{thm:trace71} makes the image
of \(u\) constant, and locality makes it the image of \(k\).
\end{proof}

The local strict classes detected in
Corollary \ref{cor:dual68}
and the theta projections of Corollary~\ref{cor:minor67} thus
refer to the same \(E_x\) and the same interval in \(x\).
The preceding theorem excludes every nonconstant trace curve
whose restriction becomes a sum of characters over a finite
extension. Finiteness still requires the other trace curves to
be excluded. Even finiteness alone would give no uniform bound
on the integral length of the strict groups. The required
comparison of the nonzero full Bessel period also remains open.
For the original marking and the specialization argument compare
\srcloc{Development/AutomorphicArithmetic/AutomorphicVarieties/UnitaryThreeGroups/MopBase}{105}
and
\srcloc{Development/AutomorphicArithmetic/AutomorphicVarieties/CurveDVRSpecialization/FinCrit}{45}.

\subsection{Trace curves and the contribution from the two-dimensional block}\label{nav:IV-26}
Work on the actual field \(E=\widetilde E\) of
Proposition~\ref{prop:short107}. For each \(\ell>P=\exp(x^\eps/12)\)
and \(\alpha(c)=(-1)^r\), take the ordinary marked deformation of
\(\operatorname{diag}(\Psi,1,\Psi^{-1})\) with
\(\Psi=\kappa_\ell^r\alpha\), \(1\le r<t\). As in
Theorem~\ref{thm:trace71}, the auxiliary group is
\(H=G_{\QQ(i)}\), and the trace algebra \(D\) is the invariant
part of the completed marked deformation ring \(B\).
The sign involution is inverse transpose with the coordinates
reversed. The field \(E\) supplies the character \(\alpha\), and
the auxiliary quadratic group \(H\) is fixed for all \(x\).

\begin{lemma}[Realizing a trace specialization by integral matrices]
\label{lem:scale118}
Let \(V\) be a complete discrete valuation ring with finite
residue field and fraction field \(L\). Given a continuous local
coefficient-algebra map \(D\to V\), there is a finite extension
\(L'/L\) and a continuous local map \(B\to\mathcal O_{L'}\)
extending the given map. This holds also when \(V\) has
characteristic \(\ell\). For a prescribed point of the graded
fiber over \(\overline L\), its orbit under the diagonal scaling
contains a point giving such a map.
\end{lemma}
\begin{proof}
Use the finite-weight algebra \(A/(f_i)\) from
Lemma \ref{lem:algebra116}.
Write it \(B^{\rm gr}\). Its degree-zero part is precisely \(D\),
a direct summand as a \(D\)-module. Thus the base change to \(L\)
is nonzero. It is a finitely generated \(L\)-algebra, generated
by the finitely many nonzero-weight variables. A closed point
over \(\overline L\) exists, and the values of the variables
belong to a finite extension of \(L\), also for any prescribed
\(\overline L\)-point. Write the nonzero values of positive
weight as \(a_i\) with weight \(d_i>0\) and those of negative
weight as \(b_j\) with weight \(-e_j<0\). Normalize the
valuation by a uniformizer of \(L\).

Each \(a_i^{e_j} b_j^{d_i}\) is the image of an invariant in the
maximal ideal of \(D\), so it has positive valuation. Hence
\(-v(a_i)/d_i < v(b_j)/e_j\) for every such pair. Choose a
rational number \(u\) strictly between all lower and all upper
bounds, ignoring either list when empty. After a finite extension
there exists a scalar of valuation \(u\). Scaling by it puts
every nonzero-weight variable in the maximal ideal, including
those with value zero. The original weight-zero variables are there
by locality. Thus evaluation converges on the ambient complete
power series ring.

The two evaluations on ambient weight-zero series are continuous
and coincide on the finitely many invariant monomials of degree
at most three. They therefore agree on all such series.
For a homogeneous series of any fixed weight, grouping
by the finitely many residual monomials in the proof of Lemma \ref{lem:algebra116} expresses it as a polynomial in the variables with
weight-zero series coefficients. The two evaluations agree on
all such series and in particular on every defining relation.
This proves the claimed extension. The universal matrices give
continuous matrices over \(\mathcal O_{L'}\).
\end{proof}

\begin{theorem}[A nonconstant trace curve is strongly irreducible]
\leavevmode\label{thm:strong118}
Let \(u:D/\varpi D\to k'[[s]]\) be continuous and local,
with \(k'/k\) finite. If its image contains a nonconstant
series, each realization of Lemma~\ref{lem:scale118} has
absolutely irreducible three-dimensional restriction to
\(G_L\) for every finite extension \(L/\QQ(i)\).
Consequently the failure of \(D\) to be finite over the
coefficient ring would give a curve with precisely this
property on the large-prime ordinary component.
\end{theorem}
\begin{proof}
Let \(M\) be the specialized matrix on \(H\), now regarded over
the algebraic closure of the Laurent series field. The element
\[
 M(g_*)-M(cg_*^{-1}c)=\operatorname{diag}(z,0,-z)
\]
has distinct eigenvalues by Lemma~\ref{lem:cycle71}. Every
\(H\)-invariant subspace is therefore spanned by a subset of
the three coordinate vectors. Draw an arrow \(j\to i\) when
some \(M(h)_{ij}\ne0\), \(h\in H\), \(i\ne j\). A coordinate
subspace is invariant precisely when no arrow leaves its index
set. Since \(M(ch^{-1}c)\) is the reflected transpose of
\(M(h)\), each adjacent arrow occurs with its reflection of
the same direction along the chain \(0,1,2\).

If both adjacent directions occur, the graph is strongly
connected. If just one adjacent direction occurs, an outer arrow
in the opposite direction makes a directed cycle through all
three vertices. In each of these cases \(M\) is irreducible.
If all arrows have the same direction the representation is
triangular. Its traces are sums of three continuous characters
already on \(H\), hence constant by Theorem~\ref{thm:chars71}.
The only other reducible possibility has all adjacent entries
zero and both outer directions present. It factors through the
quotient studied in
Theorem \ref{thm:finite117}.
That integral invariant algebra is finite over the coefficient
ring, by ordinary pseudodeformation finiteness. Its image in
the characteristic-\(\ell\) valuation ring is constant.
This last assertion still holds after the finite extension
in Lemma~\ref{lem:scale118}: its residue field is finite, and
the elements algebraic over a finite field are constants.
Both reducible possibilities contradict the hypothesis.

Suppose the restriction became reducible on some finite extension.
Pass to a finite normal extension of \(\QQ(i)\) inside the
chosen algebraic closure which contains it, and let \(H_0\)
be its open normal subgroup of \(H\). Choose a simple
\(H_0\)-submodule of the representation over the algebraically
closed coefficient field. Its finitely many \(H\)-translates
are simple \(H_0\)-modules of the same dimension. Their sum
is the whole space by \(H\)-irreducibility. A sum of simple
submodules is semisimple, and one can select a direct-sum
subcollection by successively discarding a summand contained in
the previous sum. Their common dimension divides three.
Since restriction was reducible, this dimension is one. Thus
over \(H_0\) the traces are sums of three characters. A basis
for these lines is defined over a finite extension of the
coefficient field, and the characters are continuous there.
Theorem~\ref{thm:chars71} again makes all the original traces
constant. This proves strong irreducibility. The final assertion
follows by Proposition~\ref{prop:curve71}, which produces a
nonconstant one-dimensional specialization if \(D\) is not finite.
\end{proof}

Next choose a single rational prime \(v\ge7\) split in
\(\QQ(i)\) and prime to \(2F\), where \(F\) is the conductor of
\(\psi\). On sufficiently distant selected intervals the growing
auxiliary primes \(q,d_0\) avoid \(v\), so \(v\nmid N\) and
\(v<\ell\). Such a fixed prime exists by Dirichlet's theorem.
Use arithmetic Frobenius at \(v\) throughout, and write
\[
    d_{\alpha,r}=1+v^r\alpha(v)+v^{-r}\alpha(v)^{-1}.
\]
At fixed \(x,\ell,r,\alpha\), denote by \(\mathcal P_{\ell,r,\alpha}\)
the scaling orbits of geometric mixed points on the graded
solution whose four adjacent entries are zero on \(H\). Thus a
representative has some nonzero entries in each outer direction.
This definition does not assume finiteness of the invariant
algebra of the full rank-three space. Write \(b(P)\) for the
trace of its matrix at \(\mathrm{Frob}_v\).

\begin{theorem}[A uniform sum for the two-dimensional trace differences]
\label{thm:tracebound118}
With \(v_\ell(\ell)=1\), on the original period fields and
intervals one has
\begin{equation}
 \sum_{\ell>P}(\log\ell)\sum_{1\le r<t}
 \ \sum_{\alpha(c)=(-1)^r}\ \
 \sum_{P'\in\mathcal P_{\ell,r,\alpha}}
       v_\ell\!\left(b(P')-d_{\alpha,r}\right)
       \ \le\ \exp(x^{\eps/3}) .                         \label{eq:trace118}
\end{equation}
Here \(\ell\) runs over primes and \(\alpha\) over the indicated
characters of \(\operatorname{Gal}(E/\QQ)\).
\end{theorem}
\begin{proof}
The invariant algebra on this quotient is local and finite over
the coefficient DVR by Theorem \ref{thm:finite117}. At a geometric
point its values are integral and its maximal ideal maps into
the maximal ideal of a finite extension. Indeed the contraction
of that residue maximal ideal lies over the closed point of the
DVR and is maximal by finiteness. The point can thus
be scaled to a continuous integral one by
Lemma~\ref{lem:scale118} applied on the quotient. Its weight
grading, homogeneous equations and degree-zero lifting follow
from the same proof by imposing the additional homogeneous
entry equations. By Theorem \ref{thm:form117}, for fixed
\(\ell,r,\alpha\) these orbits inject into the coefficient-embedded
newforms of weight \(2r+1\), character \(\eta\alpha^{-2}\),
and primitive level \(f(\eta\alpha^{-2})\mid4N\). Their traces
are precisely of the form \(1+v^{-r}\alpha(v)a_v(f)\).

For each \(r\), take all complex characters \(\alpha\) in this
sum and all normalized primitive eigenforms \(f\) of the stated
weight, character and level, including their conjugate systems
of coefficients. Form the single number
\[
 I_r=\prod_{\alpha,f}
   \left(\alpha(v)a_v(f)-v^{2r}\alpha(v)-\alpha(v)^{-1}\right).
\]
Each factor is an algebraic integer, since the Hecke coefficient
of a normalized newform is integral. The collection is stable
under every automorphism of the algebraic numbers, so \(I_r\)
is a rational integer. Deligne's bound gives \(|a_v(f)|\le2v^r\)
at every embedding. Thus the normalized trace has absolute value
at most three, while \(|d_{\alpha,r}|\ge v^r-1-v^{-r}>3\).
Thus \(I_r\ne0\), with the empty product equal to one.
Each factor has absolute value at most \((v^r+1)^2\).
The dimension bound for cusp forms of weight \(2r+1\) and
level dividing \(4N\), proved in
Section \ref{nav:III-00},
bounds the number of factors by a fixed polynomial in \(N(r+2)\),
since the number of possible \(\alpha\)'s is at most \(n\le N\).

For each \(\ell\) choose a common coefficient embedding into
\(\overline{\QQ}_\ell\). All admissible conjugates are included.
Every orbit contributes a distinct labelled factor for fixed
\(\alpha\), even when traces agree between different forms.
Its valuation is exactly the displayed difference valuation
since \(v^r\) is an \(\ell\)-adic unit. All unselected factors
have nonnegative valuation. Summing first at each \(\ell\) and
then over primes bounds the left side by \(\sum_{r<t}\log|I_r|\).
This is at most \([N(t+2)]^A\) for a fixed \(A\) and sufficiently
large \(x\). Since \(\log N\ll_{\psi,\eps}x^{\eps/4}\) and
\(\log(t+2)=O(\log x)\), this proves \eqref{eq:trace118}.
\end{proof}

In particular all the two-dimensional orbit terms of the
intersection inequality of Theorem \ref{thm:full36}, would
together be \(o(\exp(x^\eps/2))\), even with its factor two.
Applying that inequality globally still requires finiteness on
the strongly irreducible rank-three locus, control of the
tangent spaces in the full solution and of the rank-three
trace differences. The small-prime tail in
Corollary~\ref{cor:tail114} and the nonzero comparison on
the integral theta image also remain open.

\subsection{The odd orthogonal part of the three-dimensional calculation}\label{nav:IV-27}
Keep the fields \(E\), the intervals \([V,V^{3/2}]\), the
threshold \(P=\exp(x^\eps/12)\), and the split prime \(v\) of
Theorem~\ref{thm:tracebound118}. Put \(K'=\QQ(i)\),
\(H=G_{K'}\), and let \(\eta\) be the quadratic character of
\(K'/\QQ\). At a prime \(\ell>P\), let \(\kappa_\ell\) be
the cyclotomic character and put \(\Psi=\kappa_\ell^r\alpha\)
for a character \(\alpha\) of \(\Gal(E/\QQ)\) with
\(\alpha(c)=(-1)^r\), \(1\le r<t\). In particular \(\Psi(c)=1\).
All assertions in this paragraph hold on sufficiently distant
intervals with the thresholds of Section~\ref{sec:comparison}.
The order of \(\bar\Psi\) on inertia at \(\ell\) then exceeds four.

Let \(U\) denote the three-dimensional matrix representation on
\(H\) of an absolutely irreducible point, with the marking fixed
as in Lemma~\ref{lem:scale118}. The same notation applies to a
realization of a nonconstant trace curve in
Theorem~\ref{thm:strong118}. By
Theorem \ref{thm:orth122} and Lemma \ref{lem:spin122}, if the algebraic monodromy of \(U\) preserves
a symmetric form, \(U\) is the trace-zero adjoint of a
two-dimensional representation \(\rho\) on restriction to \(H\).
The marking determines \(e\in\{0,1\}\) with
\(\det\rho=\eta^e\Psi\) and residual representation
\(\operatorname{diag}(\overline{\eta^e\Psi},1)\) on \(G_\QQ\).
We call this case \emph{odd orthogonal} if \(e=1\), that is,
if \(\rho(c)\) has determinant \(-1\). The term refers to this
two-dimensional lift. The determinant of \(U\) is always one.

\begin{proposition}[Rigidity of odd orthogonal traces in characteristic \(\ell\)]
\label{prop:oddcurve123}
There is no nonconstant trace curve of odd orthogonal type.
\end{proposition}
\begin{proof}
Suppose such a curve existed, and take the integral realization
over a complete discrete valuation ring \(R\) of characteristic
\(\ell\) and finite residue field provided by
Lemma~\ref{lem:scale118} and Lemma \ref{lem:spin122}. Enlarge the
mixed-characteristic coefficient ring \(\mathcal O\) by a finite
unramified extension to accommodate the residue field of \(R\).
The ring \(R\) then receives the coefficient map through the
finite field of constants in it. The original coefficient
map has this property since its image is a finite field.
Write \(\phi=\eta\Psi\) for the fixed mixed-characteristic
character and use its reduction for the corresponding determinant
in \(R\). It is totally odd and de Rham. Its reduction on
inertia at \(\ell\) is nontrivial.

The residual trace of \(\rho\) is \(1+\bar\phi\).
The representation is unramified outside the primes dividing
\(4N\ell\). At \(\ell\) it has two ordinary diagonal characters.
The quotient character \(b\) is unramified by
Theorem \ref{thm:local122}. Its reduction on the decomposition group is
the trivial character. Indeed the distinct residual inertia
characters distinguish the trivial summand from \(\bar\phi\).
Pan's ordinary pseudodeformation theorem over \(\QQ\) applies
to this trace with determinant \(\phi\), this set of primes and
the local reducibility condition at \(\ell>5\).
The universal ordinary trace algebra is finite over
\(\Lambda=\mathcal O[[1+\ell\ZZ_\ell]]\). The weight map
acts through the local character lifting \(1\) on principal
units under local reciprocity.~\cite[Theorem 5.1.2(1)]{VI-IV-I-pan123}

The pseudotraces modulo the powers of the maximal ideal of
\(R\) give compatible maps from this algebra, hence one map
to \(R\). Since \(b\) is unramified, the map factors through
specialization at the augmentation \(\Lambda\to\mathcal O\).
Its image in \(R\) is a domain finite over a finite field,
hence a finite field. All traces of \(\rho\) and all values
of its determinant are constant. For \(h\in H\),
\[
 \operatorname{tr}U(h)
   =\frac{(\operatorname{tr}\rho(h))^2}{\det\rho(h)}-1 .
\]
This follows by writing the conjugation action on all
endomorphisms as \(\rho\otimes\rho^\vee\), and removing its
scalar summand. Every trace of \(U\) is consequently constant.
The invariant algebra \(D\) is topologically generated by
these traces by Theorem~\ref{thm:trace71}. Its continuous
image lies in the same finite, hence closed, field of constants.
The finite extension used to realize the curve was injective on
its coefficient field, so the original image was constant as well.
\end{proof}

Let \(\mathcal Q_{\ell,r,\alpha}\) be the scaling orbits of
geometric characteristic-zero points of odd orthogonal type
with irreducible \(U\). Let \(b(Q)\) be the trace of \(U\)
at arithmetic Frobenius at \(v\). The prime \(v\) splits
in \(K'\), so this Frobenius lies in \(H\).

\begin{theorem}[Odd orthogonal points on the period fields]
\label{thm:oddtrace123}
For each such orbit there is a normalized primitive cuspidal
newform \(g\) of weight \(r+1\), character \(\eta\alpha\)
and level \(f(\eta\alpha)\mid4N\) such that
\[
     b(Q)=\frac{a_v(g)^2}{v^r\alpha(v)}-1 .
\]
The geometric tangent space of the full graded solution at \(Q\)
is generated by scaling and has dimension one. Moreover,
\begin{equation}
 \sum_{\ell>P}(\log\ell)\sum_{1\le r<t}
     \ \sum_{\alpha(c)=(-1)^r}\ 
     \sum_{Q\in\mathcal Q_{\ell,r,\alpha}}
        v_\ell\!\left(b(Q)-d_{\alpha,r}\right)
       \le \exp(x^\eps/3).                           \label{eq:oddtrace123}
\end{equation}
\end{theorem}
\begin{proof}
By Lemma \ref{lem:integral121} and Theorem \ref{thm:local122}, the point has a continuous
integral representative. The lift \(\rho\) is odd,
absolutely irreducible, finitely ramified, and crystalline at
\(\ell\), with cyclotomic exponents \(0,r\). It has split
residual semisimplification \(1\oplus\bar\phi\) with distinct
inertia characters at \(\ell\). Pan's Fontaine--Mazur theorem
over \(\QQ\) applies at \(\ell>5\) and identifies \(\rho\),
up to twist, with the representation of a cuspidal eigenform.
The weight difference gives weight \(r+1\); normalizing the
smaller cyclotomic exponent to zero leaves a de Rham character
of weight zero. Such a finitely ramified character of \(G_\QQ\)
is of finite order. Its inertia at \(\ell\) has finite image by
the local classification of de Rham characters, and at each
other ramified prime the local units have finite image in the
\(\ell\)-adic coefficient units. Global class field theory for
\(\QQ\) then makes it a ray class character. Absorb this
character into the eigenform and take its primitive newform.
With arithmetic Frobenius in the normalization of Deligne,
the resulting representation \(V_g\) equals \(\rho\) up to
isomorphism. Its determinant prescribes the character
\(\eta\alpha\). Theorem \ref{thm:local122} computes its
prime-to-\(\ell\) conductor as \(f(\eta\alpha)\). Its
Weil--Deligne representation at \(\ell\) is unramified with
zero monodromy by crystallinity. The local--global compatibility
and newform conductor theorem used in Theorem \ref{thm:form117}
give the asserted primitive level. The adjoint trace identity
of Proposition~\ref{prop:oddcurve123} now gives the displayed
formula.~\cite{VI-IV-I-pan123,VI-IV-I-compat123}

For the tangent assertion, take the unitary automorphic
normalization \(\pi_g\) of \(g\), with central character
\(\omega_g\). It has no nontrivial self-twist by a Hecke
character. The central character forces any such twist to
be quadratic. At unramified primes the self-twist would also
yield \(V_g\simeq V_g\otimes\xi\) for the corresponding
quadratic Galois character \(\xi\), by equality of Frobenius
polynomials and Chebotarev density. The intertwiner could not
be scalar, since \(\xi\ne1\). Yet \(V_g\) remains absolutely
irreducible on \(H\cap\ker\xi\). An invariant line there would
give an invariant nilpotent line in its trace-zero adjoint,
contradicting strong irreducibility of \(U\) in
Lemma \ref{lem:strong119}.
Schur's lemma gives a contradiction.

The adjoint lift of Gelbart and Jacquet,
\(\Pi=\operatorname{Sym}^2\pi_g\otimes\omega_g^{-1}\),
is therefore cuspidal. It is self-dual. Its cyclic base
change \(\Pi_{K'}\) remains cuspidal by the cuspidality
criterion of Arthur and Clozel. Indeed a failure of cuspidality
would give a self-twist by \(\eta\); in degree three this
changes the central character by the nontrivial character
\(\eta^3\). The base change is conjugate self-dual. At the
complex place its three characters are
\((z/\bar z)^r,1,(z/\bar z)^{-r}\), so it is regular
algebraic for \(r\ge1\).~\cite{VI-IV-I-gj123,VI-IV-I-ac123}

To match the twist in Corollary \ref{cor:tangent121}, let \(w\) be an
unramified finite place of \(K'\) away from \(\ell\), outside
the finitely many primes where \(g\) ramifies. For eigenvalues
\(a,b\) of arithmetic Frobenius of \(V_g\) at \(w\), the
normalized Satake multiset of \(\Pi_{K',w}\) is
\(\{a/b,1,b/a\}\) under the fixed coefficient embedding.
This multiset is invariant under inversion. In the Tate
normalization of Newton and Thorne the local parameter for
rank three uses \(\Pi_{K',w}\otimes|\det|^{-1}\), so at
geometric Frobenius these parameters acquire the factor
\(\#k(w)\). The same factor is the value of
\(\kappa_\ell^{-1}\) at that Frobenius. Thus
\(U\otimes\kappa_\ell^{-1}\) and the representation attached
to \(\Pi_{K'}\) have the same Frobenius polynomials. Their
semisimplifications are isomorphic by Chebotarev density, and
\(U\) is irreducible. This supplies the automorphic
representation required in Corollary \ref{cor:tangent121} with
\(F=K'\). That corollary gives the tangent assertion for
every \(1\le r<t\).~\cite{VI-IV-I-nt123}

We next check the multiplicity in the orbit sum. Fix
\(\ell,r,\alpha\) and an embedding of the algebraic
coefficients into \(\overline{\QQ}_\ell\). If the same
coefficient-embedded newform serves two orbits, their
representations \(U,U'\) on \(H\) are isomorphic. After
scaling to integral representatives, an intertwiner \(D_0\)
conjugates the two matrices
\[
    U(g_*)-U(cg_*^{-1}c)=\operatorname{diag}(z_*,0,-z_*),
    \qquad
    U'(g_*)-U'(cg_*^{-1}c)=\operatorname{diag}(z_*',0,-z_*').
\]
Distinct eigenvalues give \(z_*'=\pm z_*\). The prescribed
common nonzero residue at the mark distinguishes the signs
at \(\ell>2\). Thus \(z_*'=z_*\) and \(D_0\) is diagonal.
The marked conjugate-dual relation
\(U(chc)=J_3U(h)^{-\mathsf T}J_3\) for both representations
makes \(D_0^{-1}J_3D_0^{-\mathsf T}J_3\) centralize \(U(H)\).
It is scalar by absolute irreducibility. If the diagonal of
\(D_0\) is \((d_0,d_1,d_2)\), this gives \(d_0d_2=d_1^2\).
After dividing by \(d_1\), its diagonal is \((u,1,u^{-1})\).
The coset matrices are determined by the matrices on \(H\)
and the fixed value at \(c\). On \(D\), the evaluations
agree by the trace description in Theorem~\ref{thm:trace71}
and continuity for the two integral representatives. The
homogeneous parts are generated over the invariant series
by finitely many monomials in the marked entry coordinates
as in Lemma~\ref{lem:scale118}. The scalar \(u\) therefore
relates the evaluations on the entire graded algebra. The
orbits coincide. Thus at this \(\ell,r,\alpha\) a labelled
newform accounts for at most one orbit.

Put \(z=v^r\alpha(v)\). Multiplication of
\(b(Q)-d_{\alpha,r}\) by \(z\), an \(\ell\)-adic unit, gives
\(a_v(g)^2-(z+1)^2\). For each \(r\) form the product
\[
 J_r=\prod_{\alpha,g}
       \bigl(a_v(g)^2-(v^r\alpha(v)+1)^2\bigr)
\]
over every permitted complex character and all normalized
primitive newforms of the weight, character and primitive
level above, including all conjugate systems. The factors
are algebraic integers and the set of labels is Galois
stable, so \(J_r\in\ZZ\). At every complex embedding the
two Satake roots of \(g\) at \(v\) have absolute value
\(v^{r/2}\) by Deligne's theorem. Their adjoint trace has
absolute value at most three. The comparison trace has
absolute value at least \(v^r-1-v^{-r}>3\), so every factor
is nonzero. Its absolute value is at most
\(4v^r+(v^r+1)^2\). An empty product has value one.

Every selected orbit at a fixed \(\ell\) thus contributes
the valuation of a distinct labelled factor, with valuations
of all factors nonnegative. The sum in \eqref{eq:oddtrace123}
is bounded by \(\sum_{r<t}\log|J_r|\). The dimension estimate
of Section \ref{nav:III-00}
and the bound of \(N\) for the number of characters bound the
number of factors per \(r\) by a polynomial in \(N(r+2)\).
Since \(v\) is fixed, the logarithm of the bound for one
factor is \(O_v(r+1)\). For some fixed \(A\) the sum is at
most \([N(t+2)]^A\). The estimates
\(\log N\ll_{\psi,\eps}x^{\eps/4}\) and
\(\log(t+2)=O(\log x)\) give \eqref{eq:oddtrace123}.
\end{proof}

\begin{corollary}[The remaining large-prime cases]
\label{cor:remain123}
If the full weight-zero algebra fails to be finite, there is
a nonconstant residual trace curve whose algebraic monodromy
is \(\mathrm{SL}_3\) or whose monodromy is orthogonal with
an even two-dimensional lift. In characteristic zero the
two-dimensional mixed orbits and the odd orthogonal irreducible
orbits together have one-dimensional tangent spaces, and the
sum of their trace valuations with the weights in
\eqref{eq:oddtrace123} is \(o(\exp(x^\eps/2))\).
\end{corollary}
\begin{proof}
Combine Theorem~\ref{thm:strong118}, Theorem \ref{thm:orth122},
and Proposition~\ref{prop:oddcurve123} for the first assertion.
Corollary \ref{cor:corner120} gives the tangent statement in the
two-dimensional case and Theorem~\ref{thm:oddtrace123} gives
it in the odd orthogonal case. Summing \eqref{eq:trace118}
and \eqref{eq:oddtrace123} gives the bound
\(2\exp(x^\eps/3)=o(\exp(x^\eps/2))\).
\end{proof}

For the interval estimate \eqref{eq:window}, the remaining
three-dimensional cases require a proof both in the special
fiber and at characteristic-zero points. The smaller-prime
strict cohomology in Corollary~\ref{cor:tail114} and the
comparison for the nonzero projected integral theta period
on the same fields and intervals also remain open here.

\subsection{The squared characters in the rank-three argument}\label{nav:IV-28}
Write \(\Gal(E/\QQ)=P\times H\), with \(P\) the \(p\)-Sylow
subgroup. View a character \(\alpha\) of \(H\) as a Dirichlet
character by extending it trivially on \(P\), and let \(2^s\)
be the two-part of its primitive conductor. Choose
\(\alpha(c)=(-1)^r\). Let \(\eta_4\) be the quadratic character
of \(\QQ(i)\), and set \(\Psi=\kappa_p^r\alpha\) as a
character of \(G_\QQ\), where \(\kappa_p\) is the \(p\)-adic
cyclotomic character. Write
\[
 D=2\varphi(o),\qquad e_p=v_2(p^D-1),\qquad
 H_p=4(F+2)(p+1)4^{e_p},\qquad \varphi(1)=1 .
\]

\begin{proposition}[Vanishing for the two squared characters]
\label{prop:corner148}
Suppose \(2^{s-1}\ge H_p\). Let
\(\vartheta=\eta_4\Psi^2\) or \(\vartheta=(\eta_4\Psi^2)^{-1}\).
Take a finite unramified coefficient ring \(O/\ZZ_p\)
containing the values of \(\alpha\), with fraction field \(J_p\)
and residue field \(k\). Put \(\delta=\bar\vartheta\).
Every continuous cocycle with values in \(k(\delta)\) which
vanishes on a subgroup of index prime to \(p\) in
\(I_v\cap\ker\delta\) for each finite inertia group \(I_v\)
is a coboundary. Moreover,
\[
 \ker\!\left(
 H^1\big(\QQ,(J_p/O)(\vartheta)\big)
  \longrightarrow H^1\big(\QQ_p,(J_p/O)(\vartheta)\big)
       \times\prod_{v\nmid p\infty}
                 H^1\big(I_v,(J_p/O)(\vartheta)\big)\right)=0.
\]
\end{proposition}
\begin{proof}
Extend a character of \(\Gal(B_J/\QQ)\) to the product of the
Galois groups of \(K_\psi\) and \(C_J\) as in
Theorem \ref{thm:effect146}.
The independence of the two quadratic fields gives a
factorization
\[
             \alpha=\lambda\,\zeta\,\chi_q^{u}\chi_{d_0}^{v},
             \qquad u,v\in\{0,1\},
\]
where \(\lambda\) comes from \(K_\psi\) with order dividing
the prime-to-\(p\) part of \(o\), \(\zeta\) comes from \(C_J\),
and the final factors are quadratic. Squaring eliminates both
of those factors, also when \(p=q\) or \(p=d_0\).
Let \(\omega\) be the Teichm\"uller lift of
\(\bar\kappa_p\), and factor the primitive Dirichlet character
inducing \(\eta_4\omega^{2r}\alpha^2\) as \(\phi\xi\), where
\(\phi\) has odd conductor \(f\) and \(\xi\) has conductor a
power of two. This product is odd. One has \(f\mid pF\), and
the order of \(\phi\) divides
\(\operatorname{lcm}(p-1,o/p^{v_p(o)})\), hence divides
\(p^D-1\). The same divisor works for the inverse.

Our size hypothesis gives \(s\ge5\). The two-part of
\(\alpha\) takes at \(5\) a primitive \(2^{s-2}\)-th root
of unity, because \(5\) generates the units congruent to
\(1\bmod4\) at that modulus and the character is primitive
there. Squaring reduces the order to \(2^{s-3}\).
Multiplication by \(\eta_4\) preserves that value at \(5\).
On \(-1\) the value has order at most two. The conductor of
\(\xi\) is consequently exactly \(2^{s-1}\), also after
reduction since \(p\) is odd. The inequality
\(2^{s-1}\ge H_p\ge f4^{e_p}\) now gives the finite-cocycle
assertion by Theorem \ref{thm:fourier145} and Corollary \ref{cor:coc145}.

The local residual character at \(p\) is nontrivial. If
\(p\mid f\), its primitive inertia character there is
nontrivial and has prime-to-\(p\) order, so survives
reduction. Otherwise suppose it is unramified and trivial
on Frobenius at \(p\). Then the reduction of \(\xi(p)\)
lies in \(\mathbf F_{p^D}^\times\), because it is the inverse
of the reduction of \(\phi(p)\). The unit-group computation
in the proof of Theorem \ref{thm:fourier145} gives
\(p\equiv\pm1\bmod 2^{s-1-e_p}\). But
\(2^{s-1-e_p}\ge H_p/2^{e_p}>p+1\), which rules this out.
The same reasoning applies to the inverse character.

Write \(A=(J_p/O)(\vartheta)\). Nontrivial residual action
locally at \(p\) gives an element acting by a scalar whose
difference from one is a unit. In particular
\(A^{G_{\QQ_p}}=A^{G_\QQ}=0\).
The sequence \(0\to k(\delta)\to A\xrightarrow{p}A\to0\)
now identifies \(H^1(\QQ,k(\delta))\) with \(H^1(\QQ,A)[p]\);
it gives the same identification over \(\QQ_p\). A strict
class killed by \(p\) thus has a finite lift which is a local
coboundary at \(p\); its cocycle vanishes on
\(I_p\cap\ker\delta\). At \(v\nmid p\infty\), \(\kappa_p\)
is unramified and every value of the remaining inertia action
has order prime to \(p\). On \(I_v\cap\ker\delta\) the entire
action on \(A\) is consequently trivial. Local triviality in
\(H^1(I_v,A)\) forces the finite cocycle to be zero on that
subgroup, since a coboundary restricts there to zero. The
finite-cocycle assertion just proved annihilates the lift.
Finally a continuous cocycle from a profinite group to this
discrete \(p\)-primary module has finite image. Each individual
class is killed by a power of \(p\). A nonzero element in the
displayed kernel would therefore have a nonzero multiple
killed by \(p\), contrary to what we have just proved.
\end{proof}

\begin{corollary}[Highest two-power conductor on the period field]
\label{cor:corner148}
On the intervals of Theorem~\ref{thm:invar148}, take
\(s=J+2\). For every sufficiently large \(x\), every
\(\alpha\) of this conductor at two, and every indicated
twist, Proposition~\ref{prop:corner148} applies at every
odd \(p\le x^{2/3}\). For \(p\equiv3,5\bmod8\) it applies
through \(p\le x^{2+4\eps}\).
\end{corollary}
\begin{proof}
By Proposition~\ref{prop:short107}, \(2^J\ge
x^{2+8\eps}/o\). Thus \(2^{s-1}\ge2x^{2+8\eps}/o\).
equation \ref{eq:level146} and the valuation formula following it give
\(H_p\ll_\psi p^3\), uniformly in \(p\). For \(p\le x^{2/3}\)
this is \(O_\psi(x^2)\). For the indicated classes modulo
eight one has \(\max(v_2(p-1),v_2(p+1))=2\), hence
\(e_p=2+v_2(D)\) and \(H_p\ll_\psi p\). Up to
\(x^{2+4\eps}\) this is again smaller than the displayed
lower bound for \(2^{s-1}\) eventually. These inequalities
are independent of \(r,q,d_0\).
\end{proof}

The source identifies the squared character in
\srcloc{Development/AutomorphicArithmetic/AutomorphicVarieties/LevelCohomology/DirGPowOne}{23}
and its inverse on the \((3,1)\) matrix entry in
\srcloc{Development/AutomorphicArithmetic/AutomorphicVarieties/LevelCohomology/HcXppOf}{36}.
The other matrix entries used for the strict groups in
\eqref{eq:tail114} include unsquared characters. Their remaining
lengths and the contribution from \(H\ge\sqrt{g_1g_2}\) to the
prescribed, arithmetically normalized theta periods are still
required for \eqref{eq:window}.

\subsection{Trace congruences on the special fiber}\label{nav:IV-29}
Fix a large coefficient prime and a twist on the common interval as in
Theorem~\ref{thm:strong118}. Let \(H=G_{K'}\), where \(K'=\QQ(i)\),
and let \(D\) be the complete local ring of invariant functions in the
marked ordinary deformation. For \(h\in H\), put
\(t_h=T(h)-\Psi(h)-1-\Psi(h)^{-1}\). These elements vanish at the
diagonal point. On a characteristic-\(\ell\) curve through that point,
the last three terms specialize to the lifts of their residues in the
finite constant field. In particular the specialization of \(t_h\) is
the difference between a trace and its constant term.

\begin{lemma}[Integral matrix differences in the congruence subgroups]
\leavevmode\label{lem:order238}
Let \(U\) be an integral realization of a nonconstant trace curve, and
let \(V\subset H\) be any open subgroup. After a fixed finite extension
of its equal-characteristic coefficient field, write \(R\) for the
valuation ring and \(z\) for a uniformizer. For every integer \(m\ge1\),
put \(V_m=\{v\in V:U(v)\equiv1\pmod {z^m}\}\).
There is one integer \(C\ge1\), depending on \(U,V\), for which
\[
 z^{Cm}M_3(R)\ \subseteq\
     \sum_{v\in V_m}R\,(U(v)-1)\ \subseteq\ z^m M_3(R)
                   \qquad(m\ge1).
\]
\end{lemma}
\begin{proof}
The Zariski closure of \(U(H)\) over the algebraic closure is connected
and is either \(\mathrm{SL}_3\) or \(\mathrm{SO}_3\) in its irreducible
representation of degree three, by Theorem~\ref{thm:orth122}. The image
of each open subgroup is dense there, since finitely many translates of
its closure cover this irreducible group. In both groups the set of
matrices with distinct eigenvalues is nonempty: a diagonalizable torus
element has eigenvalues \(b,1,b^{-1}\) for a choice \(b^2\ne1\).
It is an open set, defined by nonvanishing of the discriminant. Choose
\(h\in V\) in the kernel of reduction with \(A=U(h)\) in that set.
Adjoin its eigenvalues to the coefficient field. In the resulting
ring \(R\) the matrix \(A-1\) still has entries in \(zR\). Choose an
eigenvalue \(\lambda\ne1\) of \(A\), and let \(P\) be the rank-one
projection onto its eigenline along the other two eigenlines.

The conjugates \(P_w=U(w)P U(w)^{-1}\), \(w\in V\), generate the
full matrix algebra. To see this over the algebraic closure, let
\(\mathcal B\) be the unital algebra they generate. Its Jacobson
radical is preserved by \(U(V)\). The radical is nilpotent, so it has
a nonzero common kernel on the natural module. That kernel is invariant
under \(U(V)\), which acts irreducibly by Theorem~\ref{thm:strong118}.
Consequently the radical acts as zero and is zero. The summands of the
module corresponding to the simple factors of the semisimple algebra
\(\mathcal B\) are permuted by \(U(V)\). Their set is finite and is
preserved by its Zariski closure. A connected group acts trivially on
this finite set. Irreducibility gives a single summand, so in a suitable
basis the algebra is a full matrix algebra acting on several copies of
its natural module. The rank of every member is a multiple of the
number of copies. The element \(P\) has rank one, forcing one copy.
This proves the assertion, which descends to the coefficient field by
linear independence. The span of the nonempty words in the \(P_w\)
is a nonzero two-sided ideal of this algebra, hence fills it. Words of
length at most nine suffice: starting at length one, the span of words
of length at most \(j\) grows strictly until it stabilizes under
multiplication by the generators, and the ambient dimension is nine.
Choose nine such words forming a basis. For some \(d_0\ge0\) their
\(R\)-span contains \(z^{d_0}M_3(R)\), by clearing the denominators of
the coordinates of the nine matrix units in this fixed basis.

Let \(S\) consist of \(1\) and the eigenvalues of \(A\) other than
\(\lambda\), counting each value once, and put
\(d_1=\sum_{\mu\in S}v_z(\lambda-\mu)\). The eigenvalues are integral.
For \(b=\ell^a\) the numbers \(\lambda^b,\mu^b\) remain distinct, and
\[
 p_b(X)=\prod_{\mu\in S}\frac{X-\mu^b}{\lambda^b-\mu^b},
 \qquad p_b(A^b)=P,\qquad p_b(1)=0.
\]
The equality \(\lambda^b-\mu^b=(\lambda-\mu)^b\) bounds the
denominators of the coefficients by \(z^{bd_1}\). Since \(p_b(1)=0\),
it is a combination of \(X^j-1\), \(1\le j\le3\), with those same
denominators. Choose \(b\) least among the powers of \(\ell\) with
\(b\ge m\). Then \(A^b-1=(A-1)^b\in z^m M_3(R)\). Thus each
\(z^{bd_1}P_w\) lies in the indicated integral span of differences
from \(V_m\), using the elements \(w h^{bj}w^{-1}\). This span is
closed under multiplication, since for group matrices
\((X-1)(Y-1)=(XY-1)-(X-1)-(Y-1)\). The nine chosen words now place
\(z^{d_0+9bd_1}M_3(R)\) inside the span. As \(b\le\ell m\), take
\(C=\max(1,d_0+9\ell d_1)\). The other inclusion is the definition of
\(V_m\).
\end{proof}

\begin{theorem}[Density of congruence to the diagonal at finite precision]
\leavevmode\label{thm:density238}
Fix a nonconstant curve as above and a nonempty finite Galois
Frobenius condition on primes of \(K'\). There is a constant
\(\delta>0\) such that for every integer \(n\ge1\) the relative
natural density within that condition of the unramified primes with
\(u(t_{\Frob_v})\equiv0\pmod{z^n}\) is at most
\((10/9)n^{-\delta}\). The same uniform exponent works on every
conjugacy class of the chosen finite quotient. The constant and the
valuation refer to the curve and that quotient.
\end{theorem}
\begin{proof}
Take \(V\) to be the kernel of the finite quotient and use
Lemma~\ref{lem:order238}. In a coset \(gV_m\) the residual trace is a
single constant \(a\). The functions \(f(v)=\operatorname{tr}(U(g)U(v))-a\)
on \(V_m\) have the following property, for each choice of \(g\).
Some value has valuation at most \(Cm\). Indeed \(U(g)\) has a unit
entry, so its trace pairing with a matrix unit is a unit. By the
lemma, the differences \(f(v)-f(1)\) cannot all be divisible by
\(z^{Cm+1}\).

Let \(d\) be the minimum valuation of \(f\) and set \(F=z^{-d}f\).
This is integral with nonzero reduction. The linear span over the local
field of its right translates has dimension at most ten: each translate
is a linear combination of the nine entries of \(U(v)\) and the
constant function. The integral-valued functions in this span form a
free \(R\)-module of rank at most ten. In fact, finitely many evaluations
embed the span in a coordinate space of its own dimension, by choosing
successive evaluations that reduce the kernel. Integral functions map
into the standard integral lattice and every submodule of a finite
free module over \(R\) is finite free. Reduction therefore puts the
span of the translates of \(F\) modulo \(z\) in dimension at most ten.
Choose a basis of that span from the translates. At every point of
\(V_m\) at least one basis function is nonzero, since a translate can
move any point to a point where \(F\) is a unit. The supports of the
basis translates all have the same Haar measure. Their union covers
\(V_m\), so \(f\) has valuation exactly \(d\) on at least one tenth of
the coset. These elements fail the congruence at precision \(Cm+1\).

Set \(m_0=1\) and \(m_{j+1}=Cm_j+1\). Matching the constant term at
precision \(m_j\) is a union of cosets of \(V_{m_j}\). Conditional on
each matching coset at least one tenth fails at precision \(m_{j+1}\).
Consequently at most \((9/10)^j\) of any initial union of \(V\)-cosets
matches at precision \(m_j\). We have \(m_j\le(C+1)^j\). Taking
\(j=\lfloor\log n/\log(C+1)\rfloor\) gives the claimed bound with
\(\delta=\log(10/9)/\log(C+1)\). For each finite precision the
condition on traces is conjugacy invariant and factors through a
finite quotient, by continuity and finiteness of \(R/z^nR\).
Chebotarev's theorem identifies this Haar probability with relative
natural density in the prescribed union of conjugacy classes, as in
Theorem~\ref{thm:dense124}.
\end{proof}

\begin{theorem}[Frobenius parameters from a set of positive upper density]
\leavevmode\label{thm:params238}
Put \(d=\dim(D/\varpi D)\). Let \(\mathcal P\) be any set of rational
primes of positive upper natural density that split in \(K'\), outside
the ramification set. There are \(d\) distinct primes in \(\mathcal P\)
and places \(v_j\) above them such that
\[
 O[[X_1,\ldots,X_d]]\longrightarrow D,\qquad
                X_j\longmapsto t_{\Frob_{v_j}}
\]
makes \(D\) a finite module. Fewer traces never suffice for such a
finite map. The places can be chosen with either orientation at each
step. For any finite list of these rings, one list of rational primes
works simultaneously, using the first \(d\) primes for a ring of
special-fiber dimension \(d\).
\end{theorem}
\begin{proof}
Work initially in \(A=D/\varpi D\). At any step consider the quotient
by the traces already chosen. For every minimal prime with a
positive-dimensional quotient, choose a further prime quotient of
dimension one on that component. Its normalization gives a continuous
map to a power series ring over a finite residue extension, nonconstant
on \(D\), by the argument of Proposition~\ref{prop:curve71}.
There are only finitely many such curves, including those from the
other rings being treated at this step. By Theorem~\ref{thm:dense124},
outside a set of rational split primes of density zero the trace
difference is nonzero on every chosen curve, at either place. Theorem~\ref{thm:density238}
also bounds the density of agreement on each curve at finite precision.
A set of positive upper density still contains infinitely many choices
after the excluded primes and the previous choices are removed.
The next trace avoids every positive-dimensional minimal prime that
was considered, because it has nonzero image on the chosen curve.
On a positive-dimensional domain of finite Krull dimension, quotient
by a nonzero element lowers dimension by at least one: prepend the
zero prime to any lifted chain of primes containing that element.
Zero-dimensional components stay zero-dimensional. This proves that
a ring of dimension \(d\) reaches dimension zero in at most \(d\)
steps. If fewer than \(d\) elements cut it to dimension zero, their
ideal would be primary to its maximal ideal, contrary to Krull's
height theorem. This also proves the lower bound for a finite map,
since its quotient by the maximal ideal of the source is a finite
algebra over the residue field.

The quotient of \(D\) by \((\varpi,t_{\Frob_{v_1}},\ldots,
t_{\Frob_{v_d}})\) is now Artinian local with finite residue field.
This ideal defines the original topology since it contains a power
of the maximal ideal. Lift a residue-field basis of the quotient to
\(D\). To express an arbitrary element, take its expansion on that
basis modulo the ideal, lift the coefficients from the residue field
to \(O\), and write the error as a sum of the ideal generators times
new elements. Repeat on these elements. At step \(j\) all unexpanded
terms lie in the \(j\)-th ideal power. The coefficient polynomials
of the chosen lifts converge in \(O[[X_1,\ldots,X_d]]\) since their
successive differences have total \((\varpi,X_1,\ldots,X_d)\)-order
tending to infinity. Their evaluations converge in the complete ring
\(D\), giving the required finite generating set. This argument
includes \(d=0\).
\end{proof}

\subsection{Ordinary cuspidal adjoints with the prescribed residual characters}\label{nav:IV-30}
Use a large coefficient prime \(\ell\), a twist \(1\le r<t(x)\)
and one of the complementary characters \(\alpha\) for \(E_x\)
in Theorem~\ref{thm:strong118}. Put \(K'=\QQ(i)\), write
\(\eta\) for its quadratic character, and set
\(\Psi=\kappa_\ell^r\alpha\), where \(\kappa_\ell\) is the
cyclotomic character. Finite characters are viewed either as
Dirichlet characters or as Galois characters by reciprocity
with arithmetic Frobenius.

\begin{theorem}[An automorphic lift of the residual diagonal]
\label{thm:adjoint240}
Use the primitive character \(\chi\) inducing \(\eta\alpha\),
its conductor \(N\), and \(k=r+1\) in Theorem~\ref{thm:eigen240}.
On every sufficiently distant common interval its hypotheses
on \(k,N,\ell\) hold. For the auxiliary characters there put
\(\xi=\xi_1\xi_2\). There is a cuspidal self-dual automorphic
representation of \(\mathrm{GL}_3/\QQ\) whose associated
Galois representation \(U\) has determinant one, is absolutely
irreducible on every open subgroup, and is ordinary crystalline
at \(\ell\) with cyclotomic exponents \(r,0,-r\).
It has a lattice whose reduction on \(G_{K'}\) is
\(\operatorname{diag}(\bar\Psi,1,\bar\Psi^{-1})\).
At every finite place away from \(\ell\) its monodromy is zero
and its inertia characters are \(\chi\xi,1,(\chi\xi)^{-1}\).
Its base change stays cuspidal along every finite solvable
Galois extension. In particular it gives a cuspidal system on
each original \(E_x\).
\end{theorem}
\begin{proof}
We have \(\chi(-1)=(-1)^{r+1}\), and \(N\) divides four times
the modulus of \(E_x\). On these intervals the logarithm of
the latter is \(o(x^\eps)\). Since \(\ell>\exp(x^\eps/12)\)
and \(r<x^{401}\), we have \(\ell>\max(k+1,(8N)^3)\) eventually.
The inertia order of \(\bar\Psi\) is
\((\ell-1)/\gcd(\ell-1,r)>4\) throughout this range, as
\(\alpha\) is unramified at \(\ell\).

Let \(f\) be the newform just constructed and \(\rho_f\) its
two-dimensional \(\ell\)-adic representation. We use the
classical theorems on Galois representations of newforms of
weight at least two. The determinant is
\(\kappa_\ell^r\chi\xi\), and the trace at any good arithmetic
Frobenius is \(a_p(f)\). The representation is absolutely
irreducible. At \(\ell\) the level is prime to \(\ell\).
Good reduction and ordinarity give a crystalline representation
with an unramified quotient line and with inertia character
\(\kappa_\ell^r\) on the subline. At a different prime its
Artin conductor exponent equals the local newform conductor
exponent. We use the normalization in which the cyclotomic
character takes arithmetic Frobenius at \(p\ne\ell\) to \(p\).
\cite{VI-IV-VI-carayol117,VI-IV-I-compat123}
\href{https://doi.org/10.1073/pnas.94.20.10520}{Skinner and
Wiles, \emph{Ordinary representations and modular forms}}.

Chebotarev and the Brauer--Nesbitt theorem give
\(\bar\rho_f^{\rm ss}=1\oplus\bar\phi\), where
\(\bar\phi=\bar\eta\,\bar\Psi\), by the determinant and the
congruences of Theorem~\ref{thm:eigen240}. One can make the
reduction split. Start with a stable lattice and an upper
triangular residual basis. If \(\pi\) is a uniformizer of
its coefficient field, pass to a field containing \(u\) with
\(u^2=\pi\) and conjugate the matrices by
\(\operatorname{diag}(u,1)\). The new upper-right entries are
divisible by \(u\), and the old lower-left entries were divisible
by \(\pi\), so the new lower-left entries are also divisible
by \(u\). Thus the residual diagonal is realized in a lattice.

The projective image of \(\rho_f\) is Zariski dense in
\(\mathrm{PGL}_2\). It is infinite, since a finite projective
image becomes scalar on a finite extension, incompatible with
the two distinct Hodge--Tate exponents there. A proper connected
positive-dimensional subgroup of \(\mathrm{PGL}_2\) is solvable.
If it has a nontrivial unipotent radical, the unique point fixed
by that radical is preserved by its normalizer, contradicting
irreducibility. The other possibility for the identity component
is a torus, in which case the full image permutes two lines.
It would then admit a nontrivial quadratic self-twist. Indeed
the operator which is \(1,-1\) on the two lines intertwines the
representation with its twist by their permutation sign.
After reduction such a twist would preserve the two characters
\(1,\bar\phi\). The quadratic character reduces nontrivially
at odd \(\ell\), so it would equal \(\bar\phi\), contrary to
the inertia order greater than four. This proves density.
Every open subgroup has the same projective Zariski closure,
since finitely many translates of a proper closed subgroup
cannot cover the connected algebraic group.

Take \(U=\operatorname{Ad}^0\rho_f\), the conjugation action on
the trace-zero matrices. The adjoint action of \(\mathrm{PGL}_2\)
is irreducible, orthogonal for the trace pairing and of
determinant one. Its automorphic realization is the adjoint lift
of \(f\) by Gelbart and Jacquet. It is cuspidal because the
quadratic self-twist characterizing the dihedral case was just
excluded. The unramified Hecke parameters match those of \(U\)
by symmetric square and determinant twist. For a cyclic
extension of prime degree, failure of cuspidality under base
change would give a nontrivial finite character self-twist of
the cuspidal source. The matching good Frobenius polynomials
and Chebotarev would then give this same self-twist of the
restricted Galois representation. On the kernel of the character
an intertwiner must be scalar by absolute irreducibility there,
which forces the character to be trivial. Induction along a
cyclic series proves the base-change assertion.
\cite{VI-IV-I-gj123,VI-IV-I-ac123}

The adjoint of the split residual lattice has characters
\(\bar\phi,1,\bar\phi^{-1}\), giving the required diagonal on
\(G_{K'}\). The ordinary flag on \(\rho_f\) gives the adjoint
flag consisting successively of the maps from quotient to subline,
the trace-zero maps preserving the subline, and all trace-zero
maps. Its three exponents are as stated, and crystallinity
is preserved by tensors and direct summands.

At \(p\mid Ns_1s_2\), \(p\ne\ell\), the conductor exponents
of \(\rho_f\) and its determinant coincide, since the
nebentype is primitive of the full newform level. For the
finite-inertia part of the Weil--Deligne parameter, the
codimension of invariants for every ramification group is at
least the codimension for the determinant. The monodromy
correction to the Artin conductor is nonnegative. Equality
therefore forces a one-dimensional inertia-fixed space when
the determinant is ramified. The finite inertia action splits
as \(1\oplus\chi\xi\). Monodromy commutes with inertia and
is nilpotent, so vanishes on both distinct character lines.
Away from \(Ns_1s_2\ell\) the representation is unramified.
Passing to the adjoint proves the claimed local description.
\end{proof}

For the full deformation ring a further issue arises in using
automorphy lifting. Theorem~1.1(iv) of Allen, Newton and Thorne
requires the distinct residual constituents to be individually
polarized by the same multiplier. On any CM extension of \(K'\),
complex conjugation fixes the restriction of \(\bar\Psi\), since
it comes from an abelian character of \(G_\QQ\). If a common
scalar twist \(\lambda\) satisfied the polarization condition,
division of the equations for \(\lambda\bar\Psi\) and
\(\lambda\) would give \(\bar\Psi^c=\bar\Psi^{-1}\).
Then \(\lambda\bar\Psi\) and \(\lambda\bar\Psi^{-1}\)
would be equal, violating distinctness. Thus the existence of
the adjoint above leaves the other components of the original
ring to be controlled.
\href{https://doi.org/10.1112/S0010437X20007484}{Allen,
Newton and Thorne, \emph{Automorphy lifting for residually
reducible \(\ell\)-adic Galois representations, II}}.

\subsection{Two extension classes from an even orthogonal point}\label{nav:III-68}
Use the large coefficient primes and marked representations
of Corollary~\ref{cor:remain123}. Suppose a characteristic-zero
irreducible point or a nonconstant residue-characteristic trace
curve is even orthogonal. By Lemma~\ref{lem:spin122} it has a
two-dimensional lift \(\rho\) over a complete discrete
valuation ring \(R\) with finite residue field, after a finite
coefficient extension. Put \(H=G_{\QQ(i)}\). Its determinant
is \(\Psi=\alpha\kappa^r\), its residual matrix is
\(\operatorname{diag}(\bar\Psi,1)\), and \(\rho(c)=1\).
In residue characteristic the character is understood through
the coefficient map. Let \(U=\operatorname{Ad}^0\rho|_H\).
Choose \(g_*\in I_\ell\subset H\) for which
\(u=\Psi(g_*)\) has residual order greater than four. Such
a value exists on the given interval. On inertia the line and
quotient characters are \(\Psi,1\), by
Theorem~\ref{thm:local122}. The projectors formed from
\(\rho(g_*)\) and the unit \(u-1\) give an integral eigenbasis
with this ordering. For \(g\in G_\QQ\) denote the matrix entries
in this basis by \(a_g,b_g,c_g,d_g\) in row order. Set
\[
 B=(b_g:g\in G_\QQ),\qquad C'=(c_g:g\in G_\QQ),\qquad
 I=(\operatorname{tr}U(h)-\Psi(h)-1-\Psi(h)^{-1}:h\in H).
\]

\begin{proposition}[The simultaneous trace congruence and the two extensions]\leavevmode
\label{prop:traceideal306}
One has \(I=BC'\), a nonzero proper ideal. In the coefficients
\(R/I\), there is a class of zero annihilator for each of
\(\Psi\) and \(\Psi^{-1}\). Both classes are unramified away
from \(\ell\), and the class for \(\Psi^{-1}\) is zero on the
whole decomposition group at \(\ell\). For the adjoint trace
ideal it suffices to use Frobenius elements at good rational
primes splitting in \(\QQ(i)\).
\end{proposition}
\begin{proof}
Both off-diagonal ideals are nonzero by irreducibility and lie
in the maximal ideal by the residual splitting. Modulo \(BC'\)
matrix multiplication makes \(a_g,d_g\) characters with
product \(\Psi(g)\). The character \(d_g\) is one on inertia
at \(\ell\), since the chosen first eigenline is the ordinary
line. Away from \(\ell\) the image of inertia in \(\rho\)
has order prime to \(\ell\), by Theorem~\ref{thm:local122}.
There \(d_g\) reduces to one and takes values in a group of
principal units of \(\ell\)-power order. Thus it is unramified
at every finite prime. The class number and the narrow class
number of \(\QQ\) are one, so \(d_g=1\) and \(a_g=\Psi(g)\)
in this quotient. Since
\(1+\operatorname{tr}U(h)=(\operatorname{tr}\rho(h))^2/\Psi(h)\),
this proves \(I\subseteq BC'\).

Conversely reduce modulo \(I\). For any \(h\in H\) use
\(g_*^j h\), for \(j=0,1,2\). At least two of these choices
have \(1+\bar u^j\bar\Psi(h)\ne0\). For each such choice the
adjoint trace identity gives
\[
 (u^j a_h+d_h)^2=(u^j\Psi(h)+1)^2\pmod I .
\]
The sum of the two expressions being squared is a unit by
residual splitting and odd residue characteristic. The two
linear congruences therefore hold with the plus sign. The
difference of the two chosen powers of \(u\) is a unit. It
follows that \(a_h=\Psi(h)\) and \(d_h=1\) modulo \(I\).
For \(g\notin H\) apply this to \(gc\), since both \(\rho(c)\)
and \(\Psi(c)\) are one. Now
\(b_g c_h=a_{gh}-a_g a_h\in I\) for every pair of group
elements. This proves the equality of ideals.

Scale the first basis vector so that the upper-right ideal
becomes \(R\) and the lower-left ideal becomes \(I\). This
preserves integrality and the diagonal and gives an upper
extension modulo \(I\). Its normalized upper entry
\(b_g/d_g\) is a cocycle for \(\Psi\). In a second lattice
scale in the opposite direction to make the lower-left ideal
\(R\) and the upper-right ideal \(I\). Its normalized lower
entry \(c_g/a_g\) is a cocycle for \(\Psi^{-1}\). Both cocycles
vanish at \(g_*\) and each takes a unit value somewhere.
A multiple which is a coboundary must vanish identically,
since evaluation at \(g_*\) forces the cobounding vector to be
zero. This proves the assertion on annihilators. Inertia away
from \(\ell\) acts through the same finite prime-to-\(\ell\)
quotient in each lattice, where averaging kills degree-one
cohomology. On the local group at \(\ell\) the lower-left
entry is identically zero.

Write \(I=(\pi^m)\). An element of \(H\) realizes the minimum
valuation \(m\), since the discrete valuation is integral.
Reduce the lift modulo \(\pi^{m+1}\), including \(\Psi\)
and the quadratic character defining \(H\). The conjugacy
class of this element in the resulting finite Galois quotient
occurs at good rational primes by Chebotarev. Those primes
split in \(\QQ(i)\), and either orientation has the same trace
difference since \(\rho(c)=1\). They realize the same minimum.
\end{proof}

\subsection{Simultaneous upper and lower extension classes}\label{nav:III-69}
For the finite family and the ring of
Lemma~\ref{lem:spinalgebra307} put \(\overline B=B/IB\)
and \(\overline L=L/IL\). These are finite-type \(O\)-modules.
Write \(D^*=\operatorname{Hom}_O(D,J/O)\) for the
\(O\)-linear dual of such a module. Denote the actual
complementary strict group over \(\QQ\) with coefficient
\((J/O)(\Psi^{-1})\) by \(S^-\). Character descent
identifies it with the prescribed component on \(E_x\).

\begin{proposition}[An integral scalar for every member of the family]\leavevmode
\label{prop:spinclasses307}
There are injections \(\overline B^*\hookrightarrow P\) and
\(\overline L^*\hookrightarrow S^-\). If
\(h=\operatorname{length}_O\overline B\), then
\[
    h\le v_\ell(L(1-r,\alpha^*)),\qquad
    (\ell^h,0,\ldots,0)\in T,\qquad
    \ell^h\overline L=0 .
\]
These statements include coefficient rings of characteristic
\(\ell\) and arbitrary congruences among the different traces.
\end{proposition}
\begin{proof}
Modulo \(I\) the diagonal entries are \(\Psi,1\).
Thus \(g\mapsto b_g\bmod IB\) is a cocycle for \(\Psi\), and
\(g\mapsto c_g/\Psi(g)\bmod IL\) is a cocycle for
\(\Psi^{-1}\). Compose them with the corresponding
\(O\)-linear functionals. These are continuous, since
the modules are finitely generated over \(O\) with their
compact quotient topology. Their values generate their
respective modules and both cocycles vanish at \(g_*\).
If the resulting class is zero, evaluation there forces
its cobounding vector to vanish since \(u-1\) is a unit.
The functional therefore vanishes. Away from \(\ell\) the
joint inertia representation is finite of order prime to
\(\ell\), so its degree-one classes vanish on inertia. The
lower entry vanishes identically on the full local group at
\(\ell\), by ordinarity.

We check the upper local condition when \(r=1\) before
quotienting by \(I\). On inertia at \(\ell\) the actual
diagonals in every tuple are \(\Psi,1\) and the lower entries
vanish. Thus \(b_g\) itself is a cocycle there with values
in \(B(\kappa)\). For any compact pro-\(\ell\) coefficient
module the Kummer valuation on inertia is defined by taking
the inverse limit of the valuations on finite abelian
quotients as in Lemma~\ref{lem:spinunit306}. The module \(B\)
is compact and a closed additive subgroup of the product of
coefficient rings. In each coordinate its Kummer valuation
is zero by the proof of that lemma, which applies in the original
integral eigenbasis as well, since the two adjacent cocycles
differ by a nonzero scalar in the fraction field. The base
coordinate is identically zero. By injectivity into the
product, the valuation in \(B\) is zero. Its image under every
functional on \(B/IB\) is zero too. This proves both injections.
The first one and Theorem~\ref{thm:primitive307} give finiteness
and the claimed bound on \(h\), since a free \(O\)-summand
would have infinite dual.

Choose lifts \(b_1,\ldots,b_s\) in \(B\) of elementary-divisor
generators of \(\overline B\), of orders
\(\ell^{a_1},\ldots,\ell^{a_s}\). They generate \(B\) over
the local ring \(T\) by Nakayama's lemma. There is a matrix
\(Z=(z_{ij})\) with entries in \(I\) such that
\(\ell^{a_i} b_i=\sum_j z_{ij}b_j\). Set
\[
             y=\det\bigl(\operatorname{diag}
                                (\ell^{a_1},\ldots,\ell^{a_s})-Z\bigr).
\]
Multiplying the relations by the adjugate shows that \(yB=0\).
Its base value is \(\ell^{\sum_i a_i}=\ell^h\). By
Lemma~\ref{lem:spinalgebra307}, \(y\) vanishes in every
irreducible coordinate and annihilates \(L\) as well. Since
\(\ell^h-y\in I\), it follows that \(\ell^h L\subset IL\).
When \(s=0\) use the empty determinant, equal to one.
\end{proof}

\subsection{The upper extension at the first twist}\label{nav:IV-57}
At \(r=1\), an unrestricted local Kummer class at \(\ell\)
can have nonzero valuation. Such a class would contribute
an extra local term precisely when \(\alpha(\ell)=1\).
We determine the valuation of the upper class in
Proposition~\ref{prop:traceideal306}, including on a curve
with coefficient field of characteristic \(\ell\).

\begin{lemma}[The upper spin extension has unit Kummer class]\leavevmode
\label{lem:spinunit306}
On the ordinary component through the marked point, the
upper class in \(H^1(\QQ,(R/I)(\alpha\kappa))\) of
Proposition~\ref{prop:traceideal306} satisfies the Kummer
unit condition at \(\ell\).
\end{lemma}
\begin{proof}
Let \(K\) be the completion of \(\QQ(i)\) at the chosen
place over \(\ell\). It is unramified over \(\QQ_\ell\),
so the two inertia groups agree. In the scaled upper lattice
from the proposition the local matrices on \(G_K\) have
diagonal characters \(a,d\), with \(d\) unramified and
\(a/d=\kappa\beta\) for an unramified \(\beta\). The upper
entry divided by \(d\) is an integral cocycle \(z\) for this
character before reduction modulo \(I\).

There is an \(R\)-linear valuation map on these cocycles
after restriction to inertia. For each finite coefficient
module \(R/\pi^j\), apply the Kummer valuation on the maximal
unramified extension, with \(\kappa\)-twisted coefficients.
Explicitly a Kummer class represented at a finite unramified
stage by \(y\) maps to its normalized valuation, extended
additively in the coefficient module. The valuations agree
under passage up the unramified tower. The cyclic case is
the Kummer sequence and the assertion for finite abelian
coefficient modules follows by additivity. These maps are
natural in the coefficients, kill coboundaries, and commute
with reduction modulo \(\pi^j\). Taking their inverse limit
defines \(v_R(z)\in R\). This is the first coordinate in the
local decomposition of Lemma~\ref{lem:kummer72}. This argument
uses finite \(\ell\)-groups at every stage for either
characteristic of \(R\).

Use the original integral rank-three lattice at this place.
The eigenvalues \(u,1,u^{-1}\) of the marked inertia element
are residually distinct, so their eigenvectors can be chosen
as an integral basis by the spectral projectors. They are
ordered along the ordinary flag. In the first two steps of
the flag the normalized adjacent cocycle has zero Kummer
valuation by Theorem~\ref{thm:first72}. Indeed its valuation
is the coordinate killed by saturation in the ordinary
chart. At a place with negative Frobenius sign the same test
is taken on \(G_K\) via the induced pair of adjacent
coordinates. Changing the adapted basis by diagonal units
rescales the class by a unit and an upper unipotent change
adds a coboundary on the adjacent quotient. Hence the
vanishing holds in the integral eigenbasis just chosen,
at all coefficient levels.

Over the fraction field of \(R\), identify the rank-three
representation with \(\operatorname{Ad}^0\rho\). In the
ordered eigenbasis
\(E_{12},\operatorname{diag}(1,-1),E_{21}\) of this adjoint,
conjugation sends the middle vector to itself minus
\(2z(g)E_{12}\). The distinct eigenvalues of \(g_*\) show
that the two ordered eigenbases correspond up to three
nonzero scalars. The first adjacent cocycle in the original
integral basis is therefore a fixed nonzero fraction-field
multiple \(s z\). It is integral and its valuation vanishes.
Choose \(h\ne0\) in \(R\) with \(hs\in R\). Linearity gives
\((hs)v_R(z)=h v_R(sz)=0\). The ring \(R\) is a domain,
so \(v_R(z)=0\). Reducing modulo \(I\) proves the unit
condition for the upper extension, including at split values
of \(\alpha(\ell)\).
\end{proof}

\subsection{The integral trace ring of several even orthogonal lifts}\label{nav:IV-58}
Fix \(\ell,r,\alpha,O,\Psi,H\) as in
Proposition~\ref{prop:traceideal306}. Consider any finite
family of its even orthogonal lifts \(\rho_i\), over complete
discrete valuation rings \(R_i\) of either coefficient
characteristic. Add the diagonal representation
\(\rho_0=\operatorname{diag}(\Psi,1)\) over \(O\). In each
original residually split lattice diagonalize the same
\(g_*\) with \(u=\Psi(g_*)\) in the first position. Denote the
tuples of entries of \(\rho_i(g)\), including position zero,
by \(a_g,b_g,c_g,d_g\). All tuples lie in
\(A=O\times\prod_i R_i\), with its product topology.
Let \(T\) be the closed \(O\)-algebra generated by the traces
of the adjoints on \(H\), and set
\[
            B=\sum_{g\in G_\QQ}T b_g,\qquad
            L=\sum_{g\in G_\QQ}T c_g,\qquad I=BL .
\]
The last product denotes the module of finite sums of
coordinatewise products.

\begin{lemma}[The augmentation and the two off-diagonal modules]\leavevmode
\label{lem:spinalgebra307}
The ring \(T\) is also the closed algebra generated by the
two-dimensional traces on \(G_\QQ\). It contains every
\(a_g,d_g\). The modules \(B,L\) are finite over \(T\), and
\(I\) is the kernel of the projection \(T\to O\) at the diagonal.
The annihilators of \(B,L\) in \(T\) coincide and consist of
the tuples zero at every irreducible member of the family.
\end{lemma}
\begin{proof}
Every trace polynomial has a common residue in the naturally
embedded copies of the residue field of \(O\). The same holds
in its closure. An element of \(T\) with nonzero common residue
is a unit in \(T\). To see this, lift its residual inverse from \(O\) and
use the convergent geometric series for the error in \(A\).
Hence \(T\) is compact and local. For \(h\in H\), the square
of \(s_j=\operatorname{tr}\rho(g_*^j h)\) is in \(T\) by the
adjoint identity, since the determinant tuple is the scalar
\(u^j\Psi(h)\). If the common residue
\(1+\bar u^j\bar\Psi(h)\) is nonzero, lift it from \(O\) and
apply Newton iteration to \(X^2-s_j^2\). The derivative is
a unit and the residual error is squared at every step
up to a unit. In every coordinate the iteration converges to
the unique square root with that residue. Closedness gives
\(s_j\in T\). Among \(j=0,1,2\) two satisfy the residue test.
Solve \(s_j=u^j a_h+d_h\) using their unit determinant to get
\(a_h,d_h\in T\). The equalities \(\rho(c)=1\) and
\(\Psi(c)=1\) give this for all \(g\). Conversely all adjoint
traces are polynomials in the two-dimensional traces divided
by the prescribed unit determinants. This proves equality of
the two closed trace rings.

The joint image of all the two-dimensional representations
is topologically finitely generated. To see this, pass to the
kernel of the common residual character in the absolute
Galois group unramified outside the finite set of allowed
primes. Its joint image lies in a product of principal
congruence kernels and is pro-\(\ell\). Its maximal pro-\(\ell\)
quotient before taking this image is finitely generated. Indeed the classes of degree \(\ell\) characters unramified outside
that finite set form a finite group. For example, adjoining
\(\ell\)-th roots of unity reduces this finiteness to the
finiteness of the class group and finite generation of
the group of units with the specified primes inverted, by
Kummer theory and prime-to-degree restriction. The generator
criterion for pro-\(\ell\) groups applies. Reinstating the
finite residual quotient gives finitely many topological
generators \(g_s\) of the joint image. Matrix multiplication
gives \(b_{gh}=a_g b_h+b_g d_h\) and its lower analogue.
Inversion divides the negative off-diagonal entries by the
known unit determinant. The sums of \(T b_{g_s}\) and of
\(T c_{g_s}\) are compact, hence closed, and contain the
entries of all words. They give \(B,L\).

The identities \(b_g c_h=a_{gh}-a_g a_h\) put \(I\) in \(T\).
It is a closed finitely generated ideal, with zero base
coordinate. Modulo \(I\) both diagonal entries are characters.
The lower one is one on inertia at \(\ell\), and unramified
away from \(\ell\) by the finite prime-to-\(\ell\) inertial
types and its trivial reduction. For this last assertion
\(T/I\) is a profinite local ring whose principal units are
pro-\(\ell\), by reduction at finite precision in the product
and passage to the quotient by the closed ideal. The ordinary
class field theorem over \(\QQ\) gives \(d_g=1\) modulo \(I\).
The determinant then gives \(a_g=\Psi(g)\) there. The trace
generators and compactness of \(O\) now give \(T/I=O\) under
base evaluation. Finally, at each irreducible coordinate some
upper entry and some lower entry are nonzero. Those coefficient
rings are domains. At position zero both modules vanish.
Coordinatewise multiplication proves the annihilator assertion.
\end{proof}

\subsection{The even orthogonal case on a common interval}\label{nav:I-74}
Take an original interval \([V,V^{3/2}]\) from
Theorem~\ref{thm:conductor232}. Let \(E_*\) be its largest
period field, \(N_*\) its cyclotomic modulus and \(g_*\) its
degree. Write \(t_*=\lfloor V^{(3/2)401}\rfloor\).
For \(2\le r<t_*\) let \(D_r(E_*)\) be the nonzero integer
of Lemma~\ref{lem:int13}, at \(n=r\), on this very field.
Define the positive integer
\[
            Z_*=|\mathrm{Cl}(E_*)|
                   \prod_{2\le r<t_*}|D_r(E_*)|.
\]
The class number here is the ordinary ideal class number. In
particular every factor is defined independently of any
deformation ring.

\begin{theorem}[Prime support and congruence depth for even orthogonal points]\leavevmode
\label{thm:evenbudget306}
At every \(x\) of a sufficiently distant such interval use the
original period field \(E_x\), a prime
\(\ell>\exp(x^\varepsilon/12)\), and a complementary character
\(\alpha\) at \(1\le r<\lfloor x^{401}\rfloor\).
An even orthogonal irreducible point in characteristic zero,
or a nonconstant even orthogonal trace curve in characteristic
\(\ell\), requires \(\ell\mid Z_*\).
Outside these primes, every nonconstant trace curve on the
marked component has monodromy \(\mathrm{SL}_3\).

For each triple \((\ell,r,[\alpha])\) let \(d_{\ell,r,\alpha}\)
be the \(O\)-length of \(C\) in Theorem~\ref{thm:positive306},
with one representative per residue Frobenius orbit.
The following inequality sums each distinct triple on the
whole interval once:
\[
 \sum_{\ell,r,[\alpha]} b_\alpha d_{\ell,r,\alpha}\log\ell
      \le \log Z_*\ll_{\psi,\varepsilon} V^{1220}.
\]
For a characteristic-zero even orthogonal point with coefficient
ring \(R\), the greatest integer \(m\) for which all its adjoint
traces agree with the diagonal traces modulo \(\pi_R^m\)
satisfies
\[
      m/e(R/O)\le \min\{j\ge0:\ell^j C=0\}
                   \le d_{\ell,r,\alpha}.
\]
Thus the same sum bounds the weighted sum of the suprema of
these normalized depths, one supremum per triple. At a single
\(x\), the construction with \(E_x,t(x)\) gives the bound
\(O_{\psi,\varepsilon}(x^{812})\).
\end{theorem}
\begin{proof}
The fields on the interval are nested. The two auxiliary
quadratic fields and \(K_\psi\) are fixed there and only the
index of the real cyclotomic layer varies. Uniformly on this
interval, \(\exp(V^\varepsilon/12)>2t_*+3\) and exceeds
\(g_*\) and every prime dividing \(N_*\) for large \(V\),
by the degree and conductor bounds of
Theorem~\ref{thm:conductor232}. Inflate each \(\alpha\) to
the largest group. The groups \(P,C\) are defined over
\(\QQ\) with the same local conditions, so this leaves them
unchanged. At \(r=1\) their dual-order description uses
character components of the class group. The sum of their
lengths at \(\ell\), weighted by \(b_\alpha\), is at most
\(v_\ell(|\mathrm{Cl}(E_*)|)\). Indeed
\(\ZZ_\ell[\operatorname{Gal}(E_*/\QQ)]\) is a product of
the unramified integer rings for the different Frobenius
orbits. An \(O\)-module of length one in the selected factor
has \(\ZZ_\ell\)-length \(b_\alpha\).

At \(r\ge2\), Theorem~\ref{thm:positive306} gives components
in the critical sign at twist \(r-1\). On each of these
unramified factors, the zeroth Fitting ideal of a finite
module is generated by \(\ell\) to its length. The injection
and Fitting inclusion of Lemma~\ref{lem:fit13} therefore
bound the weighted sum at this \(\ell,r\) by
\(v_\ell(D_r(E_*))\), as in Theorem~\ref{thm:critical13}.
Taking products proves the first inequality. At a single
field of degree \(g\) and modulus \(N\), the
conductor--discriminant formula gives
\(|\operatorname{disc}E|\le N^g\). Minkowski represents
each ideal class by an ideal of norm at most
\(2^g|\operatorname{disc}E|^{1/2}\). Counting such ideals
by \(\zeta_E(2)\le\zeta(2)^g\) gives
\(\log|\mathrm{Cl}(E)|\ll g\log(2N)\).
Lemma~\ref{lem:int13} bounds the sum of the remaining
logarithms by \(O_\psi(gt^2(g+\log(2Nt)))\). Use
\(g\ll_\psi x^{2+8\varepsilon}\),
\(\log N\ll_{\psi,\varepsilon}x^{\varepsilon/4}\),
\(t\le x^{401}\) and \(0<\varepsilon<1/52\).
This gives the claimed pointwise bound and, with
\(x=V^{3/2}\), the stated interval bound.

For an even orthogonal point let \((\pi_R^m)\) be the trace
ideal from Proposition~\ref{prop:traceideal306}. Its upper
extension class meets the conditions of \(P\), including the
unit restriction at \(r=1\) by Lemma~\ref{lem:spinunit306}.
After reduction modulo \(\pi_R\) the normalized cocycle still
takes a unit value and vanishes at the marked inertia element.
It gives a nonzero class in \(P_1\) after residue extension.
Indeed finite coefficient scalar extension commutes with
cohomology and these local conditions by using a basis of
the extension. Hence \(d_{\ell,r,\alpha}>0\) in both coefficient
characteristics by Theorem~\ref{thm:positive306}. This proves
the prime exclusion. Corollary~\ref{cor:remain123} supplies
the monodromy conclusion.

In mixed characteristic embed \(R/\pi_R^m\) as torsion in
\(\operatorname{Frac}(R)/R\). Global degree-one cohomology
injects under this map since \(\Psi(g_*)-1\) is a unit, and
all the imposed local conditions persist. The divisible
positive group over these coefficients is \(P\otimes_O R\). Indeed
\(R\) is finite free over \(O\), so the assertion follows
componentwise in a basis, also for the Kummer valuation.
The upper class has annihilator exactly \((\pi_R^m)\).
By Theorem~\ref{thm:positive306}, every element of
\(P\otimes_O R\) is killed by \(\ell^j\) for the least \(j\)
which kills \(C\). This gives the claimed depth.
\end{proof}

The estimate measures simultaneous agreement of all traces at
one point. The number of points over an exceptional prime, the
full strict group for \(\Psi^{-1}\), and the dense
three-dimensional case still require control. The interval
estimate~\eqref{eq:window} also requires the arithmetic
normalization of the theta pairing.

\subsection{One trace function on the even orthogonal locus}\label{nav:I-75}
Use the largest period field \(E_*\), its degree \(g_*\),
modulus \(N_*\), and last twist cutoff \(t_*\) from
Theorem~\ref{thm:evenbudget306}. For each primitive character
associated with a character of \(\operatorname{Gal}(E_*/\QQ)\),
let \(f_\alpha\) be its conductor. Define
\[
 M_*=\left|\prod_{r=1}^{t_*-1}
         \ \prod_{\alpha(c)=(-1)^r}
          r f_\alpha(r+1)!\,L(1-r,\alpha^*)\right|.
\]
The inner product takes every complex character of the
specified parity once.

Fix an actual marked ring on the interval at a large prime
\(\ell\) and write \(D\) for its complete invariant trace
ring with coefficient ring \(O\), as in
Theorem~\ref{thm:trace71}. Its diagonal evaluation is denoted
by \(x_0\). In the complementary strict group \(S^-\), take
the sum \(S^{\rm ev}\) of the images of \((L/IL)^*\) in
Proposition~\ref{prop:spinclasses307}, ranging over all finite
families of even orthogonal points and nonconstant even
orthogonal trace curves of this ring. This definition permits
all their integral eigenbases. It includes the lower extension
classes obtained from individual points followed by any
coefficient functional. Indeed at one point the lower-entry
lattice generates the ideal \(C'\) of
Proposition~\ref{prop:traceideal306}. Dividing the lower
entries by a generator of that ideal, then reducing modulo
the trace ideal at the point, factors through \(L/IL\).

\begin{theorem}[A common scalar and the exponents of the strict classes]\leavevmode
\label{thm:simultaneous307}
The integer \(M_*\) is positive and satisfies
\[
             \log M_*\ll g_* t_*^2\log(2N_*t_*).
\]
For each \(\ell>\exp(V^\varepsilon/12)\) on a sufficiently
distant original interval its valuation is exactly
\[
      v_\ell(M_*)=
        \sum_{r<t_*,\ r\ge1}\ \sum_{[\alpha]}
                b_\alpha\,v_\ell(L(1-r,\alpha^*)),
\]
where the inner sum uses the residue Frobenius orbits on the
same parity of characters of \(E_*\).

For every marked ring as above, put
\(d=v_\ell(L(1-r,\alpha^*))\). There is an element \(y\in D\)
with \(x_0(y)=\ell^d\) whose image is zero on every even
orthogonal irreducible characteristic-zero point and every
nonconstant even orthogonal trace curve. One may instead
prescribe the same diagonal value \(M_*\) in all these rings.
The actual strict subgroup \(S^{\rm ev}\) is killed by
\(\ell^d\). In particular, writing \(a_{\ell,r,\alpha}\)
for the smallest exponent killing it, meaning
\(\ell^{a_{\ell,r,\alpha}}S^{\rm ev}=0\), one has on the
whole interval
\[
  \sum_{\ell,r,[\alpha]} b_\alpha a_{\ell,r,\alpha}\log\ell
                           \le\log M_* .
\]
Count a repeated triple only once, taking the sum of its
detected subgroups inside the same strict group over \(\QQ\).
\end{theorem}
\begin{proof}
The Bernoulli polynomial formula in Lemma~\ref{lem:bern99}
makes each factor defining \(M_*\) an algebraic integer. Its denominator computation holds also for
\(B_1(X)=X-1/2\). In that weight all the characters are odd
and nonprincipal, and the functional equation and nonvanishing
at one give nonvanishing at zero. The higher weights use the
nonvanishing in the same lemma. Galois conjugation permutes
the factors, giving a nonzero rational integer before the
absolute value. The Bernoulli estimate there bounds the
logarithm of the absolute value of each cleared factor from
above by \(O(r\log(2N_*r))\), also at weight one. Summing
over at most \(g_*\) characters at each weight proves the
size estimate. The conductor, factorial and weight in each
clearing factor are \(\ell\)-adic units by the inequalities
in Theorem~\ref{thm:evenbudget306}. The local orbit of
\(\alpha\) has \(b_\alpha\) elements with equal valuations,
since its coefficient ring is unramified. Taking valuations
gives the identity. Theorem~\ref{thm:primitive307} identifies
each summand before weighting with the positive Selmer length.

Consider first any finite list of the indicated specializations
of \(D\) and include its base evaluation. Its image in the
product of the coefficient rings is exactly \(T\) of
Lemma~\ref{lem:spinalgebra307}. Indeed the invariant ring is
topologically generated by the traces on \(H\) by
Theorem~\ref{thm:trace71}, and it is compact. The continuous
product image therefore contains the entire closed algebra of
these trace tuples and is contained in that algebra.
Proposition~\ref{prop:spinclasses307} puts
\((\ell^h,0,\ldots,0)\) in that image, for some \(h\le d\).
Multiply by \(\ell^{d-h}\) and lift the resulting tuple to
\(D\). In the compact fiber \(x_0^{-1}(\ell^d)\) each
specialization imposes a closed vanishing condition. The
finite-list result gives the finite intersection property,
so there is a single \(y\) satisfying all of them.
Since \(M_*/\ell^d\in\ZZ\), multiplying \(y\) by this
integer gives the second specified value.

For every finite family, \(\ell^d\) kills \(L/IL\), hence
its dual and its image in the actual strict group. This
annihilation persists on their sum and under inflation of
the character from a smaller period field on the interval.
Thus \(a_{\ell,r,\alpha}\le d\), and the valuation identity
and prime factorization of \(M_*\) give the final inequality.
\end{proof}

To obtain the required bound for the order of the full
strict group one must also control the number of its generators
and its classes outside \(S^{\rm ev}\). The dense
three-dimensional deformations and the comparison with the
arithmetically normalized theta pairing remain open in the
argument for \eqref{eq:window}.

\subsection{Trace congruences and three strict character groups}\label{sec:markededges333}

Fix \(x\in[V,V^{3/2}]\), \(1\le r<\lfloor x^{401}\rfloor\), a prime
\(\ell>\exp(x^\varepsilon/12)\), and a character \(\alpha\) of the period
field with \(\alpha(c)=(-1)^r\). Put \(K'=\QQ(i)\), \(H=G_{K'}\),
\(\eta=\chi_{-4}\), \(\Psi=\alpha\kappa_\ell^r\), and
\((\chi_1,\chi_2,\chi_3)=(\Psi,1,\Psi^{-1})\). In particular
\(\chi_j(c)=1\) for each \(j\). Let \(O\) be the minimal unramified
coefficient ring of \(\alpha\). Consider an integral specialization of
the original marked ordinary representation with restriction
\(U:H\to\mathrm{GL}_3(R)\) irreducible over the fraction field of \(R\).
Here \(R\) is the integer ring of a finite extension of
\(\operatorname{Frac}(O)\), or a complete discrete valuation ring of
characteristic \(\ell\) with finite residue field extending that of \(O\).
Write \(\pi\) for its uniformizer and \(v(\pi)=1\). The split residual
basis, the marked element \(g_*\) in inertia and the polarization are
those of Lemmas~\ref{lem:flag69} and~\ref{lem:cycle71}. These hypotheses
apply to the marked integral realizations of irreducible points and
nonconstant trace curves considered above.

\begin{proposition}[The two lower extension lengths]
\label{prop:markededges333}
Let \(I_{\rm tr}\) be the ideal of \(R\) generated by
\(\operatorname{tr}U(h)-\sum_j\chi_j(h)\) for \(h\in H\). For
\(i\ne j\), let \(A_{ij}\) be the ideal generated by \(U(h)_{ij}\)
as \(h\) runs over \(H\). Put
\[
 B=A_{12}=A_{23},\qquad C=A_{21}=A_{32},\qquad
 U_0=A_{13},\qquad V_0=A_{31}.
\]
All four ideals are nonzero and proper. Write \(b,c_0,u,v_0\) for their
valuations in the displayed order. Then
\[
   I_{\rm tr}=BC+U_0V_0,\qquad m=v(I_{\rm tr})
                    =\min(b+c_0,u+v_0).
\]
The two lower adjacent entries give strict classes over \(K'\) with
character \(\Psi^{-1}\), and the lower corner gives a strict class
there with character \(\Psi^{-2}\). Their cyclic coefficient modules
have respective lengths
\[
 d_{21}=d_{32}=\min(m,v_0+b-c_0),\qquad
 d_{31}=\min(m,2c_0-v_0),\qquad d_{21}+d_{31}\ge m .
\]
Each class has zero annihilator over its cyclic coefficient ring.
\end{proposition}
\begin{proof}
Write \(\bar j=4-j\) for index reversal and let \(J\) be the reversal
matrix. The marking gives \(U(ch^{-1}c)=JU(h)^{\mathsf T}J\), which
identifies the reflected entry ideals. The orbit of each basis vector
spans the fraction-field representation by irreducibility, so each
entry ideal is nonzero. The residual matrices are diagonal.

Let \(\mathcal A\) be the \(R\)-span of \(U(H)\). This is an algebra by
the group law. The difference \(Z=U(g_*)-U(cg_*^{-1}c)\) is
\(\operatorname{diag}(z,0,-z)\) with \(z\) a unit. The same difference
for the split diagonal representation has first entry
\(z_0=\Psi(g_*)-\Psi(g_*)^{-1}\), congruent to \(z\) modulo \(\pi\).
The three residues \(z_0,0,-z_0\) are distinct. Polynomial interpolation
in \(Z\) puts each coordinate projector \(P_j\) in \(\mathcal A\).
Inserting \(P_j\) between two matrices shows
\(A_{ij}A_{jk}\subset A_{ik}\) for distinct indices.

Set \(I=BC+U_0V_0\). Taking the \(j\)-th diagonal entry modulo \(I\)
defines an algebra homomorphism \(\mathcal A\to R/I\), since every
off-diagonal term in the diagonal of a product belongs to \(I\).
We identify its restriction to \(H\). The element \(U(g_*)\) has the
three eigenvalues \(\chi_j(g_*)\) with pairwise unit differences by
ordinarity. Its spectral projectors \(\Pi_j\) are polynomials in
\(U(g_*)\), hence belong to \(\mathcal A\). The matrix
\(Q=\sum_j\Pi_j P_j\) reduces to the identity. Its inverse belongs to
\(\mathcal A\) by Cayley--Hamilton and its unit determinant. We may
scale its columns by diagonal units to give \(Q\) diagonal entries one.
The columns are ordered eigenvectors, so their successive spans give
the unique ordinary flag. In this basis the local matrices on
\(H\cap G_{\QQ_\ell}\) are upper triangular and have the prescribed
diagonal on inertia by Lemma~\ref{lem:flag69}. Passage through the
diagonal algebra homomorphism is invariant under conjugation by \(Q\).
Thus \(U(h)_{jj}\bmod I\) equals \(\chi_j(h)\) on the chosen inertia.
At a conjugate place the same assertion follows by inner conjugation
on \(H\) or by the polarized formula, which reverses the indices and
inverts the characters. Away from \(\ell\) the prescribed inertia
kernels give a finite quotient of order prime to \(\ell\). The ratio of
the two characters reduces to one and takes values in the principal
units of the finite local ring \(R/I\), an \(\ell\)-group. It is
therefore trivial on this inertia as well. The ratio is now everywhere
unramified on \(H\). The ring \(\ZZ[i]\) is Euclidean, so class field
theory makes the ratio trivial globally. This proves \(I_{\rm tr}\subset I\).

For the other inclusion, traces of all words in \(U(H)\) agree with
their split diagonal evaluations modulo \(I_{\rm tr}\), by the group
law. In particular \(z^2=\operatorname{tr}(Z^2)/2\) agrees with
\(z_0^2\), and \(z+z_0\) is a unit. Thus \(z\equiv z_0\bmod I_{\rm tr}\).
Lagrange interpolation for \(Z\) expresses \(P_i\) using the three
nodes \(z,0,-z\), all their differences units. The same expressions
for the split evaluation reduce to the same word coefficients modulo
\(I_{\rm tr}\). For \(i\ne j\) the trace
\(\operatorname{tr}(P_iU(h)P_jU(h'))\) therefore agrees modulo that
ideal with its split value zero. It equals \(U(h)_{ij}U(h')_{ji}\).
These products generate \(I\), proving equality.

For distinct \(i,j,k\) form the quotient
\[
  L_{ij}=A_{ij}/(I A_{ij}+A_{ik}A_{kj}).
\]
Choose a generator of \(A_{ij}\), the same one for each reflected
pair. Division by this generator identifies \(L_{ij}\) with
\(R/\pi^{d_{ij}}\) for some \(d_{ij}\ge0\). In this quotient the
entry \(U(h)_{ij}/\chi_j(h)\), expressed in the chosen generator, is
a cocycle for \(\chi_i/\chi_j\). Indeed the term through \(k\) in
matrix multiplication vanishes, and both diagonal factors can be
replaced by their characters modulo \(I\). If a scalar \(a\) times
the cocycle is a coboundary, then
\[
 a\,U(h)_{ij}=(\chi_i(h)-\chi_j(h))w\quad\hbox{in }L_{ij}
\]
for some \(w\in L_{ij}\), with the entry here in its original ideal.
Subtract the equations at \(g_*\) and \(cg_*^{-1}c\). The left side
vanishes because \(Z\) is diagonal. On the right the multiplier of
\(w\) is a difference of distinct elements of \(\{z_0,0,-z_0\}\),
hence a unit. Thus \(w=0\). The entries generate \(L_{ij}\), so
\(aL_{ij}=0\).

For \(i>j\) take a local \(h\in H\cap G_{\QQ_\ell}\) and put
\(U'=Q^{-1}U(h)Q\). In \(U(h)Q=QU'\) the term through \(k\) on
each side belongs to \(A_{ik}A_{kj}\), since \(Q,U'\in\mathcal A\).
We have \(U'_{ij}=0\), \(Q_{ii}=Q_{jj}=1\), and the diagonal entries
of both representations are \(\chi_j(h)\) in coordinate \(j\)
modulo \(I\). Consequently
\(U(h)_{ij}=(\chi_j(h)-\chi_i(h))Q_{ij}\) in \(L_{ij}\).
This gives a coboundary on the full decomposition group at the chosen
place. At the conjugate place, the polarized identity relates the
cocycle to the negative of its reflected lower cocycle, using the
cocycle formula for the inverse. Thus both lower adjacent classes
and the lower corner vanish locally there too. At every prime away
from \(\ell\) the inertia cocycle factors through the prescribed
finite group of order prime to \(\ell\), whose positive-degree
cohomology vanishes by averaging. These are exactly the strict
conditions on \(K'\).

Finally \(C^2\subset V_0\) and \(V_0 B\subset C\) by multiplication
of entries. The two lower quotients therefore have the asserted
nonnegative lengths. Before truncation at \(m\), the two arguments
\(v_0+b-c_0\) and \(2c_0-v_0\) sum to \(b+c_0\ge m\). If either is
at least \(m\), that term alone suffices. Otherwise truncation changes
neither. This proves their summed lower bound.
\end{proof}

For a character \(\xi\) over \(O\), let \(a(\xi)\) be the least \(a\ge0\)
such that \(\ell^a\) kills the strict group over \(\QQ\) for
\((\operatorname{Frac}(O)/O)(\xi)\). The conditions are vanishing
on the decomposition group at \(\ell\) and vanishing on inertia at
all other finite places. The groups below are finite by descent to
the period field or its compositum with \(K'\), and the finiteness in
Lemma~\ref{lem:top}. Write
\[
 a_0=a(\Psi^{-1}),\qquad a_1=a(\eta\Psi^{-1}),\qquad
 a_2=a(\eta\Psi^{-2}).
\]

\begin{theorem}[Depth at an irreducible ordinary point]
\label{thm:tracestrict333}
For each mixed-characteristic specialization above with ramification
index \(e=v(\ell)\),
\[
                  \frac{m}{e}\le \max(a_0,a_1)+a_2 .
\]
If \(a_1=a_2=0\), one lower adjacent class descends in the invariant
part to the complementary character \(\Psi^{-1}\) at full length \(m\).
In equal characteristic, under the same vanishing, this class has
zero annihilator over \(R/\pi^m\) and gives a nonzero strict class
over the residue coefficient field for \(\overline\Psi^{-1}\).
\end{theorem}
\begin{proof}
Conjugation by \(c\) sends each normalized lower cocycle to the
negative of its reflected cocycle. Indeed
\(U(chc)_{ij}=U(h^{-1})_{\bar j,\bar i}\), and division by
\(\chi_j(h)\), followed by the inverse cocycle identity, yields
the sign minus. The generators for a reflected pair agree. All
\(\chi_i(c)\) are one. Thus the two projections of the \(21\)-class
under \((1+c)/2\) and \((1-c)/2\) are strict classes for
\(\Psi^{-1}\) and \(\eta\Psi^{-1}\) over \(\QQ\), respectively.
Here restriction identifies \(H^1(\QQ,-)\) with the invariants on
\(H\), since the quotient has order two. Twisting by \(\eta\)
identifies the anti-invariants in the same way. On local inertia
and decomposition groups, the quotients also have order dividing
two, so restriction detects the required local vanishing. The corner
is anti-invariant and descends with character \(\eta\Psi^{-2}\).

The two projected adjacent classes sum on \(H\) to a class of full
length \(d_{21}\). The annihilators are ideals in a discrete valuation
ring. At least one projection has full length, since otherwise
\(\pi^{d_{21}-1}\) would kill both when \(d_{21}>0\). The corner has
full length \(d_{31}\) by Proposition~\ref{prop:markededges333}.

In mixed characteristic embed \(R/\pi^d\) into
\(\operatorname{Frac}(R)/R\) by multiplication by \(\pi^{-d}\).
The induced global degree-one map for each of these characters is
injective. In the coefficient exact sequence its possible kernel comes
from invariants of the divisible group, which vanish because
\(g_*\) acts by a scalar distinct from one modulo \(\pi\).
This also holds for the corner since \(\bar\Psi(g_*)\) has order
greater than two. The strict local restrictions persist. Extending
scalars from \(O\) to the finite free \(O\)-module \(R\) carries the
original divisible strict group to the one with coefficients in
\(\operatorname{Frac}(R)/R\), as is seen by using an \(O\)-basis in
the global cochains and their local restrictions. Consequently
\(d_{21}\le e\max(a_0,a_1)\) and \(d_{31}\le ea_2\). Their sum
bounds \(m\) from below by the proposition.

For equal characteristic the character values in \(R/\pi^d\) come
from the residue field of \(O\). A basis of this finite coefficient
space over that field separates its strict classes by coordinate
projections. A nonzero coordinate gives a strict class with finite
residue coefficients, which injects into the divisible strict group
for the corresponding characteristic-zero character by the same
invariants argument. Thus when \(a_1=a_2=0\) the corner and the
anti-invariant adjacent projection vanish, in either characteristic.
Then \(d_{31}=0\) and \(d_{21}=m\), since \(d_{21}\le m\).
The invariant projection has the full annihilator length \(m\).
As \(m>0\), at least one residue coefficient projection in equal
characteristic is nonzero. This proves all the assertions.
\end{proof}

\begin{corollary}[A single exceptional integer for the two critical groups]
\label{cor:criticaledges333}
Let \(E_*\) be the largest period field on the original interval.
Apply Lemma~\ref{lem:int13} to \(E'_*=E_*(i)\) with cyclotomic
modulus \(4N_*\), and write
\[
   B_*=\prod_{n=2}^{2t_*+1}|D_n(E'_*)|,\qquad
   t_*=\lfloor V^{(3/2)401}\rfloor .
\]
Write \(S_i\) for the strict character group defining \(a_i\), and
\(\lambda_i=\len_O S_i\). For the distinct triples
\((\ell,r,[\alpha])\) occurring anywhere
on the interval above their respective large-prime cutoffs,
\[
  \sum_{\ell,r,[\alpha]} b_\alpha\,(\lambda_1+\lambda_2)\log\ell
       \le 17\log B_*\ll_{\psi,\varepsilon} V^{1220},
       \qquad b_\alpha=[O:\ZZ_\ell].
\]
Consequently the same bound holds with \(\lambda_1+\lambda_2\) replaced by
\(\max(0,\sup m/e-a_0)\), where the supremum is taken over all the
irreducible mixed-characteristic specializations of the triple, with
value zero for an empty family. At every prime in this range with \(\ell\nmid B_*\), the
full-length conclusion of Theorem~\ref{thm:tracestrict333} applies.
\end{corollary}
\begin{proof}
All the fields \(E_x\) are subfields of \(E_*\). Inflate their
characters to this largest field and then to \(E'_*\). This last
field realizes the fixed character \(\eta\), including when
\(i\in E_*\) already. On a sufficiently distant interval the primes
considered exceed \(2t_*+3\), \(2g_*\), and every prime dividing
\(4N_*\), by Theorem~\ref{thm:evenbudget306}. Character descent
therefore identifies the two groups with components on \(E'_*\)
at weights \(s=r,2r\), respectively. Their finite characters
in the convention \(\beta^{-1}\kappa_\ell^{-s}\) are
\(\beta=\eta\alpha,\eta\alpha^2\). Both satisfy
\(\beta(c)=(-1)^{s+1}\). They lie in the critical eigenspace of
Lemma~\ref{lem:fit13}.

At a fixed \(\ell,s\), the integral group ring splits into unramified
valuation rings. On each factor, the zeroth Fitting valuation is the
module length. The Fitting inclusion and the strict dual injection
in Lemma~\ref{lem:fit13} bound the sum of the strict lengths, weighted
by the residue degrees, by \(v_\ell(D_{s+1}(E'_*))\), as in the proof
of Theorem~\ref{thm:critical13}. The proof uses only abelianity of
the CM field, its cyclotomic modulus and the two fixed good primes,
which also apply to \(E'_*\). Each exponent is at most the corresponding
length. To count multiplicities one may count each geometric character
once, since its residue orbit has size \(b_\alpha\). Unramified
extension to a larger coefficient ring preserves the exponent.
The map \(\alpha\mapsto\eta\alpha\) is injective. The fibers of
\(\alpha\mapsto\eta\alpha^2\) have at most sixteen elements. Indeed
\(\operatorname{Gal}(E_*/\QQ)\) embeds by restriction in the product
of the cyclic group of \(K_\psi\), the cyclic real cyclotomic group,
and two groups of order two. Its subgroup killed by two has at
most sixteen elements. The same holds in the character group by
the elementary divisor theorem. For a given \(s\) the adjacent
term can come only from \(r=s\), and the corner only from \(r=s/2\).
Thus the two lists together use each critical factor at most
seventeen times. Summing valuations times \(\log\ell\) proves
the first bound.

Writing \(g'_*=[E'_*:\QQ]\le2g_*\), Lemma~\ref{lem:int13} gives
\(\log B_*\ll g'_*t_*^2(g'_*+\log(8N_*t_*))\).
Theorem~\ref{thm:conductor232} bounds \(g_*\ll V^{3+12\varepsilon}\)
and \(\log(2N_*)\ll V^{3\varepsilon/8}\), while \(t_*^2\le V^{1203}\).
Together these give the displayed power of \(V\).
The assertion about suprema follows by subtracting \(a_0\) in
Theorem~\ref{thm:tracestrict333}. If \(\ell\nmid B_*\), all the
corresponding nonnegative critical exponents must vanish.
\end{proof}

\Needspace{7\baselineskip}
\begin{corollary}[Full negative cotangent lengths on the period fields]
\label{cor:midtangent333}
At each triple above, let \(C=C_{-1}\oplus C_{-2}\) be the two
negative cotangent modules at the split augmentation of the actual
ordinary ring, as in Theorem~\ref{thm:cot116}. Then
\[
       \len_O C-\lambda_0=\lambda_1+\lambda_2,\qquad
       \len_O(B_*C)=\len_O(B_* S_0).
\]
In particular, the weighted sum of the first difference, over the
distinct triples on the whole interval, is at most \(17\log B_*\).

At a given \(x\) put \(P=\exp(x^\varepsilon/12)\),
\(Y=x^{-1240}\exp(x^\varepsilon/2)\), \(B=M_*B_*\).
For every prime \(\ell>P\) and complementary triple on \(E_x\) set
\[
 C'_{\ell,r,\alpha}=
       \begin{cases} B C,&P<\ell\le Y,\\ C,&\ell>Y.\end{cases}
\]
Let \(\mathcal L_{>}(x)\) be the total logarithmic order of the
complementary strict groups for these primes and all \(1\le r<t\)
on \(E_x\). With one summand per residue character orbit, there is
a constant \(K_{\psi,\varepsilon}\) uniform on the intervals such that
\[
 -17\log B_*
 \ \le\ \mathcal L_{>}(x)-
     \sum_{\ell,r,[\alpha]} b_\alpha\log\ell\,\len_O C'_{\ell,r,\alpha}
 \ \le\ K_{\psi,\varepsilon}\,x^{-15}\exp(x^\varepsilon/2).
\]
\end{corollary}
\begin{proof}
Theorem~\ref{thm:cot116} identifies the \(O\)-linear dual of \(C\)
with \(S_0\oplus S_1\oplus S_2\), with dualizing module
\(\operatorname{Frac}(O)/O\). A cyclic \(O\)-module of length \(j\)
has dual of the same length and the same elementary divisor.
This proves the first equality. At each fixed prime each critical
component has length at most the valuation of its own determinant
\(D_{s+1}(E'_*)\) by the Fitting calculation in
Corollary~\ref{cor:criticaledges333}. Since this determinant divides
\(B_*\) up to sign, \(B_*\) kills \(S_1\) and \(S_2\) separately.
Taking images under multiplication by \(B_*\) in each cyclic factor
gives the second equality, which holds for every multiple of \(B_*\)
as well.

By descent the contribution to \(\mathcal L_{>}(x)\) from a single
triple is \(b_\alpha\log\ell\,\lambda_0\).
At \(P<\ell\le Y\), the difference of the two length terms is exactly
\(\len_O S_0[B]\), by the identity for the images and the exact
sequence of multiplication by \(B\). At \(\ell>Y\) the difference
is \(-(\lambda_1+\lambda_2)\). The positive sum is bounded by
Theorem~\ref{thm:integerlayers333} on \(E_x\) with this very \(B\).
Corollary~\ref{cor:criticaledges333} bounds the absolute value
of the negative sum. Both sums have finite support at each \(x\)
by Proposition~\ref{prop:finite15} and the nonzero determinant
product for the critical terms.
\end{proof}

The depth inequality has the direction \(m/e\le a_0\) away from
\(B_*\). Obtaining the upper bound in Section~\ref{sec:missing}
still requires control of the entire complementary group, including
its classes outside those detected by the trace congruences.

\clearpage
\section{References}\label{nav:bibliography-VI}
\begingroup\def\section*#1{}\begin{thebibliography}{99}\small
\bibitem{VI-IV-VI-wilkes} G. Wilkes, \emph{Profinite Groups}, Part III lecture notes, University of Cambridge. \href{https://www.dpmms.cam.ac.uk/~grw46/LectureNotes.pdf}{Chapters 3 and 4}.
\bibitem{VI-IV-V-bs} A. Borel and J.-P. Serre, \emph{Corners and arithmetic groups}, Comment. Math. Helv. \textbf{48} (1973), 436--491. \href{https://www.e-periodica.ch/cntmng?pid=com-001:1973:48::32}{Original paper}.
\bibitem{VI-IV-V-franke} J. Franke, \emph{Harmonic analysis in weighted \(L_2\)-spaces}, Ann. Sci. Ecole Norm. Sup. \textbf{31} (1998), 181--279. \href{https://www.numdam.org/articles/10.1016/s0012-9593(98)80015-3/}{Paper}.
\bibitem{VI-IV-V-shin} S. W. Shin, \emph{Galois representations arising from some compact Shimura varieties}, Ann. of Math. \textbf{173} (2011), 1645--1741. \href{https://annals.math.princeton.edu/2011/173-3/p09}{Theorem 1.2}.
\bibitem{VI-IV-V-newforms} F. Diamond and J. Im, \emph{Modular forms and modular curves}, in \emph{Seminar on Fermat's Last Theorem}, CMS Conf. Proc. \textbf{17} (1995), 39--133. \href{https://wstein.org/edu/Fall2003/252/references/diamond-im/Diamond-Im-Modular_forms_and_modular_curves.pdf}{Authors' text}.
\bibitem{VI-IV-V-deligne} P. Deligne, \emph{Formes modulaires et repr\'esentations \(\ell\)-adiques}, S\'eminaire Bourbaki, Exp. 355 (1969). \href{https://www.numdam.org/item/SB_1968-1969__11__139_0.pdf}{Paper}.
\bibitem{VI-IV-V-weil} P. Deligne, \emph{La conjecture de Weil I}, Publ. Math. IHES \textbf{43} (1974), 273--307. \href{https://numdam.org/articles/10.1007/BF02684373/}{Paper}.
\bibitem{VI-IV-VI-pan40} L. Pan, \emph{The Fontaine--Mazur conjecture in the residually reducible case}, J. Amer. Math. Soc. \textbf{35} (2022), 1031--1169. \href{https://arxiv.org/abs/1901.07166}{Author's paper}.
\bibitem{VI-IV-VI-saito40} T. Saito, \emph{Modular forms and \(p\)-adic Hodge theory}, Invent. Math. \textbf{129} (1997), 607--620. See also \href{https://www.ms.u-tokyo.ac.jp/~t-saito/pp/h.pdf}{\emph{Hilbert modular forms and \(p\)-adic Hodge theory}}, Compos. Math. \textbf{145} (2009), 1081--1113.
\bibitem{VI-IV-VI-deligne40} P. Deligne, \emph{La conjecture de Weil. I}, Publ. Math. IH\'ES \textbf{43} (1974), 273--307. \href{https://www.numdam.org/item/PMIHES_1974__43__273_0/}{Purity over finite fields}.
\bibitem{VI-IV-brumer330} A. Brumer, \emph{On the units of algebraic number fields}, Mathematika \textbf{14} (1967), 121--124. \href{https://doi.org/10.1112/S0025579300003703}{Original paper}.
\bibitem{VI-IV-VI-bk117} S. Bloch and K. Kato, \href{https://virtualmath1.stanford.edu/~conrad/BSDseminar/refs/BKTamagawa.pdf}{\emph{L-functions and Tamagawa numbers of motives}}, in \emph{The Grothendieck Festschrift}, Vol.~I, Progr. Math. \textbf{86} (1990), 333--400, Section~3.
\bibitem{VI-IV-VI-carayol117} H. Carayol, \href{https://www.numdam.org/articles/10.24033/asens.1512/}{\emph{Sur les repr\'esentations \(\ell\)-adiques associ\'ees aux formes modulaires de Hilbert}}, Ann. Sci. \'Ecole Norm. Sup. (4) \textbf{19} (1986), 409--468.
\bibitem{VI-IV-VI-nt119} J. Newton and J. A. Thorne, \href{https://arxiv.org/pdf/1912.11265}{\emph{Adjoint Selmer groups of automorphic Galois representations of unitary type}}, J. Eur. Math. Soc. \textbf{25} (2023), 1919--1967, Theorem~A and Section~2.5.
\bibitem{VI-IV-VI-kato120} K. Kato, \href{https://archive.numdam.org/item/AST_2004__295__117_0.pdf}{\emph{\(p\)-adic Hodge theory and values of zeta functions of modular forms}}, Ast\'erisque \textbf{295} (2004), 117--290, Theorem~14.2(1) and Section~14.10.
\bibitem{VI-IV-VI-kisin121} M. Kisin, \href{https://people.math.harvard.edu/~kisin/dvifiles/licm.pdf}{\emph{The structure of potentially semi-stable deformation rings}}, Theorem~1.2.2. See also \emph{Potentially semi-stable deformation rings}, J. Amer. Math. Soc. \textbf{21} (2008), 513--546.
\bibitem{VI-IV-VI-milne122} J. S. Milne, \href{https://www.jmilne.org/math/Books/iAG2022.pdf}{\emph{Algebraic Groups}}, Cambridge University Press, 2017, chapters on reductive groups, rank one and root data.
\bibitem{VI-IV-VI-lmo124} J. C. Lagarias, H. L. Montgomery and A. M. Odlyzko, \href{https://doi.org/10.1007/BF01390234}{\emph{A bound for the least prime ideal in the Chebotarev density theorem}}, Invent. Math. \textbf{54} (1979), 271--296.
\bibitem{VI-IV-I-pan123} L. Pan, \href{https://arxiv.org/pdf/1901.07166}{\emph{The Fontaine--Mazur conjecture in the residually reducible case}}, J. Amer. Math. Soc. \textbf{35} (2022), 1031--1169.
\bibitem{VI-IV-I-compat123} H. Carayol, \href{https://www.numdam.org/articles/10.24033/asens.1512/}{\emph{Sur les repr\'esentations \(\ell\)-adiques associ\'ees aux formes modulaires de Hilbert}}, Ann. Sci. \'Ecole Norm. Sup. (4) \textbf{19} (1986), 409--468. T. Saito, \emph{Modular forms and \(p\)-adic Hodge theory}, Invent. Math. \textbf{129} (1997), 607--620.
\bibitem{VI-IV-I-gj123} S. Gelbart and H. Jacquet, \href{https://www.numdam.org/item/ASENS_1978_4_11_4_471_0/}{\emph{A relation between automorphic representations of \(\mathrm{GL}(2)\) and \(\mathrm{GL}(3)\)}}, Ann. Sci. \'Ecole Norm. Sup. (4) \textbf{11} (1978), 471--542.
\bibitem{VI-IV-I-ac123} J. Arthur and L. Clozel, \href{https://www.claymath.org/library/cw/arthur/pdf/30.pdf}{\emph{Simple Algebras, Base Change, and the Advanced Theory of the Trace Formula}}, Annals of Mathematics Studies 120, Princeton University Press, 1989, Chapter III.
\bibitem{VI-IV-I-nt123} J. Newton and J. A. Thorne, \href{https://arxiv.org/pdf/1912.11265}{\emph{Adjoint Selmer groups of automorphic Galois representations of unitary type}}, J. Eur. Math. Soc. \textbf{25} (2023), 1919--1967, Theorem A.
\end{thebibliography}\endgroup
\end{document}
